\documentclass[12pt]{article}
\usepackage[utf8]{inputenc}
\usepackage[LGR, T1]{fontenc}
\usepackage[greek,english]{babel}
\usepackage{lmodern}
\usepackage{amsmath}
\usepackage{amssymb}
\usepackage{amsthm}
\usepackage{geometry}
\geometry{a4paper, margin=1in}

% Μορφοποίηση ενοτήτων
\usepackage{titlesec}
\titleformat{\section}
  {\normalsize\bfseries}
  {\thesection}
  {1em}
  {}
  []
\titlespacing*{\section}
  {0pt}
  {1ex plus 0.5ex minus 0.2ex}
  {0.5ex plus 0.3ex}

\titleformat{\subsection}
  {\normalsize\itshape}
  {\thesubsection}
  {1em}
  {}
  []
\titlespacing*{\subsection}
  {0pt}
  {0.8ex plus 0.3ex minus 0.1ex}
  {0.4ex plus 0.2ex}

% Κανονική εσοχή παραγράφων
\setlength{\parindent}{1.5em}

% Αφαίρεση εσοχής μόνο στην πρώτη παράγραφο μετά από τίτλους
\usepackage{etoolbox}
\makeatletter
\pretocmd{\section}{\@afterindentfalse}{}{}
\pretocmd{\subsection}{\@afterindentfalse}{}{}
\pretocmd{\subsubsection}{\@afterindentfalse}{}{}
\makeatother

\begin{document}

\selectlanguage{greek} % Θέτει την ελληνική γλώσσα ως προεπιλεγμένη

Ας αναπτύξουμε λοιπόν τη θεωρία των τανυστών σε γραμμικούς χώρους, εστιάζοντας στην ουσία της ανταλλοιότητας και συναλλοιότητας και στη γεωμετρική σημασία των αλγεβρικών πράξεων.

\section{Διανύσματα και δυϊκότητα}

Θεωρούμε έναν διανυσματικό χώρο \( V \) διάστασης \( n \) πάνω από το \( \mathbb{R} \), τον οποίο σκεφτόμαστε ως τον εφαπτόμενο χώρο σε ένα σημείο μιας πολλαπλότητας. Αυτό μας επιτρέπει να υιοθετήσουμε εξαρχής τον γεωμετρικό συμβολισμό των φυσικών βάσεων.

Μια βάση του \( V \) είναι το σύνολο των τελεστών μερικής παραγώγου \( \{ \frac{\partial}{\partial x^i} \} \). Οποιοδήποτε διάνυσμα \( \xi \) γράφεται ως
\[
\xi = \xi^i \frac{\partial}{\partial x^i}. \tag{1}
\]

Ο δυϊκός χώρος \( V^* \) αποτελείται από όλες τις γραμμικές απεικονίσεις από τον \( V \) στο \( \mathbb{R} \). Η φυσική δυϊκή βάση της \( \{ \frac{\partial}{\partial x^i} \} \) συμβολίζεται με \( \{ dx^i \} \) και ορίζεται από τη δράση
\[
dx^i \left( \frac{\partial}{\partial x^j} \right) = \delta^i_j. \tag{2}
\]
Ένα στοιχείο του \( V^* \) αναπτύσσεται ως
\[
\omega = \omega_i dx^i. \tag{3}
\]

Η καίρια στιγμή είναι η αλλαγή βάσης, την οποία γεννά μια αλλαγή συντεταγμένων \( x^i \to y^i \). Ο κανόνας της αλυσίδας υπαγορεύει πώς μετασχηματίζονται τα στοιχεία των βάσεων:

\[
\frac{\partial}{\partial y^i} = \frac{\partial x^j}{\partial y^i} \frac{\partial}{\partial x^j}, \tag{4}
\]
\[
dy^i = \frac{\partial y^i}{\partial x^j} dx^j. \tag{5}
\]

Απαιτούμε το διάνυσμα \( \xi \) και το στοιχείο \( \omega \) να είναι γεωμετρικά αντικείμενα, δηλαδή ανεξάρτητα της επιλεγμένης βάσης. Από την απαίτηση
\[
\xi = \xi^i \frac{\partial}{\partial x^i} = \tilde{\xi}^i \frac{\partial}{\partial y^i}
\]
και αντικαθιστώντας τον μετασχηματισμό της βάσης, προκύπτει ο νόμος μετασχηματισμού των ανταλλοίωτων συνιστωσών:
\[
\tilde{\xi}^i = \frac{\partial y^i}{\partial x^j} \xi^j. \tag{6}
\]
Οι δείκτες είναι πάνω και ο μετασχηματισμός γίνεται με τον αντίστροφο Ιακωβιανό πίνακα σε σχέση με τα διανύσματα βάσης. Αυτό σημαίνει πως οι συνιστώσες του διανύσματος μεταβάλλονται αντίθετα (contra) από τα διανύσματα βάσης ώστε το ίδιο το αντικείμενο να μένει αναλλοίωτο.

Αντιστοίχως, από την απαίτηση
\[
\omega = \omega_i dx^i = \tilde{\omega}_i dy^i
\]
προκύπτει ο νόμος μετασχηματισμού των συναλλοίωτων συνιστωσών:
\[
\tilde{\omega}_i = \frac{\partial x^j}{\partial y^i} \omega_j. \tag{7}
\]
Εδώ ο μετασχηματισμός γίνεται με τον ίδιο Ιακωβιανό πίνακα που μετασχηματίζει και τα διανύσματα βάσης \( \partial / \partial x^i \). Οι συνιστώσες συμμεταβάλλονται (co) με τη βάση. Η κλίση μιας βαθμωτής συνάρτησης \( f \) είναι το φυσικό παράδειγμα: οι συνιστώσες της \( \partial f / \partial x^i \) ακολουθούν ακριβώς αυτόν τον κανόνα.

\subsection{Ανταλλοιότητα και της συναλλοιότητα}

Η διάκριση ανάμεσα σε ανταλλοίωτα και συναλλοίωτα αντικείμενα δεν είναι τυπογραφική σύμβαση αλλά βαθιά γεωμετρική αναγκαιότητα. Ένας ανταλλοίωτος δείκτης δηλώνει ότι η αντίστοιχη συνιστώσα ανήκει σε ένα αντικείμενο που ζει στον εφαπτόμενο χώρο \( V \), δηλαδή σε ένα διάνυσμα. Τα διανύσματα είναι πρωταρχικά γεωμετρικά βέλη: όταν αλλάζουμε σύστημα συντεταγμένων, το βέλος παραμένει το ίδιο, αλλά οι προβολές του στους νέους άξονες αναμειγνύονται με τον αντίστροφο τρόπο από τους ίδιους τους άξονες. Ένας συναλλοίωτος δείκτης, από την άλλη, ανήκει σε ένα αντικείμενο του δυϊκού χώρου \( V^* \), μια μηχανή που τρώει διανύσματα και δίνει αριθμούς. Το αρχέτυπο είναι η κλίση μιας βαθμωτής συνάρτησης: αν η συνάρτηση μεταβάλλεται αργά κατά μήκος μιας κατεύθυνσης, η συνιστώσα της κλίσης της είναι μικρή· αυτή η σχέση διατηρείται αναλλοίωτη μόνο αν οι συνιστώσες της κλίσης μετασχηματίζονται όπως οι ίδιοι οι άξονες συντεταγμένων.

Η σύμβαση των δεικτών (πάνω για ανταλλοίωτα, κάτω για συναλλοίωτα) λειτουργεί ως αυτόματος έλεγχος συνοχής: κάθε φορά που ένας πάνω και ένας κάτω δείκτης συναντώνται σε μια άθροιση, οι αντίστροφοι Ιακωβιανοί πίνακες αλληλοεξουδετερώνονται, παράγοντας ένα μέγεθος που δεν εξαρτάται από την επιλογή συντεταγμένων. Αυτή η απλή αλγεβρική ιδιότητα κωδικοποιεί τη γεωμετρική αλήθεια ότι η δράση μιας μορφής πάνω σε ένα διάνυσμα είναι ένα βαθμωτό, ανεξάρτητο από το σύστημα αναφοράς. Κάθε φορά που βλέπουμε έναν ελεύθερο πάνω δείκτη, αναγνωρίζουμε μια ανταλλοίωτη φύση που αναμένει να συναντήσει μια συναλλοίωτη για να δώσει ένα αναλλοίωτο αποτέλεσμα, και αντιστρόφως.

\section{Αλγεβρικός ορισμός τανυστή και μετασχηματισμός τανυστικών συνιστωσών}

Ένας τανυστής τύπου \( (r, s) \) ορίζεται αλγεβρικά ως μια πολυγραμμική απεικόνιση
\[
T: \underbrace{V^* \times \dots \times V^*}_{r \text{ φορές}} \times \underbrace{V \times \dots \times V}_{s \text{ φορές}} \to \mathbb{R}. \tag{8}
\]
(Έτσι, ένα διάνυσμα είναι (1,0)-τανυστής, μια 1-μορφή είναι (0,1)-τανυστής, ένας τελεστής είναι (1,1)-τανυστής κ.ο.κ.). Η ισχύς αυτού του ορισμού είναι ότι το αντικείμενο \( T \) υπάρχει ανεξάρτητα από κάθε σύστημα συντεταγμένων.

Για να συνδεθούμε με τους υπολογισμούς, ορίζουμε τις συνιστώσες του \( T \) στη βάση \( \{ \partial / \partial x^i \} \) ως την τιμή του πάνω στα στοιχεία των βάσεων:
\[
T^{i_1 \dots i_r}_{j_1 \dots j_s} \equiv T\left( dx^{i_1}, \dots, dx^{i_r}, \frac{\partial}{\partial x^{j_1}}, \dots, \frac{\partial}{\partial x^{j_s}} \right). \tag{9}
\]

Πώς προκύπτει ο κλασικός νόμος μετασχηματισμού; Η λογική είναι άμεση: για να βρούμε τις συνιστώσες \( \tilde{T} \) στο νέο σύστημα \( y^i \), υπολογίζουμε την τιμή του ίδιου αφηρημένου τανυστή \( T \) πάνω στα νέα στοιχεία βάσης. Αντικαθιστούμε τις νέες βάσεις με τους μετασχηματισμούς τους ως προς τις παλιές, χρησιμοποιώντας τις (4) και (5), και, λόγω πολυγραμμικότητας, κάθε παράγοντας μετασχηματισμού βγαίνει έξω από τον τανυστή. Για να αναδείξουμε τη δομή της "αυτο-απαλοιφής" των δεικτών, ο θεμελιώδης νόμος γράφεται ως
\[
\tilde{T}^{i_1 \dots i_r}_{j_1 \dots j_s} = \frac{\partial y^{i_1}}{\partial x^{\mu_1}} \dots \frac{\partial y^{i_r}}{\partial x^{\mu_r}} \ T^{\mu_1 \dots \mu_r}_{\nu_1 \dots \nu_s} \ \frac{\partial x^{\nu_1}}{\partial y^{j_1}} \dots \frac{\partial x^{\nu_s}}{\partial y^{j_s}}. \tag{10}
\]

Παρατηρήστε πώς κάθε νέος ανταλλοίωτος δείκτης \( i \) "προσδένεται" σε έναν παλιό \( \mu \), και κάθε νέος συναλλοίωτος \( j \) "προσδένεται" σε έναν παλιό \( \nu \). Οι ελληνικοί δείκτες \( \mu, \nu \) είναι βωβοί και αθροίζονται, εξαφανιζόμενοι από το τελικό αποτέλεσμα. Αυτός ο μηχανισμός σύζευξης-απαλοιφής είναι η πεμπτουσία του τανυστικού λογισμού: εγγυάται ότι η πληροφορία αναδιοργανώνεται με συνεπή τρόπο, διατηρώντας το γεωμετρικό νόημα ανέπαφο.

Αυτός είναι ο "φυσικός" ορισμός: τανυστής είναι κάθε συλλογή ποσοτήτων που μετασχηματίζεται με αυτόν τον τρόπο. Η ισοδυναμία είναι πλήρης: μια πολυγραμμική απεικόνιση γεννά αριθμούς που υπακούν στον νόμο (10), και αντιστρόφως, ένα σύνολο αριθμών που τον ικανοποιεί ορίζει μια απεικόνιση ανεξάρτητη βάσης μέσω της σύμβασης άθροισης.

\section{Θεμελιώδεις πράξεις μεταξύ τανυστών }

Εστιάζουμε στις τρεις πράξεις που συνθέτουν τον πυρήνα της τανυστικής άλγεβρας.

\textbf{Τανυστικό Γινόμενο} \( \otimes \)  
Παίρνει έναν τανυστή \( T \) τύπου \( (r, s) \) και έναν \( S \) τύπου \( (p, q) \) και κατασκευάζει έναν νέο τανυστή τύπου \( (r+p, s+q) \) με συνιστώσες που ορίζονται από το απλό γινόμενο:
\[
(T \otimes S)^{i_1 \dots i_r k_1 \dots k_p}_{j_1 \dots j_s l_1 \dots l_q} = T^{i_1 \dots i_r}_{j_1 \dots j_s} S^{k_1 \dots k_p}_{l_1 \dots l_q}. \tag{11}
\]
Γεωμετρικά, η πράξη αντιστοιχεί στο να τροφοδοτήσουμε τον \( T \) με το πρώτο γκρουπ ορισμάτων και τον \( S \) με το δεύτερο και να πολλαπλασιάσουμε τα αποτελέσματα. Είναι ο μηχανισμός που χτίζει τανυστές ανώτερης τάξης από θεμελιώδη αντικείμενα, συνενώνοντας τις γραμμικές τους ιδιότητες. Δεν υπάρχει καμία σύζευξη ή απαλοιφή δεικτών· οι δείκτες απλώς παρατίθενται, επεκτείνοντας τη χωρητικότητα του νέου αντικειμένου σε ορίσματα.

\textbf{Συναίρεση}  
Είναι η πράξη που μειώνει την τάξη ενός τανυστή από \( (r, s) \) σε \( (r-1, s-1) \), εξισώνοντας έναν ανταλλοίωτο με έναν συναλλοίωτο δείκτη και αθροίζοντας. Γενικά, αν επιλέξουμε να συναίρουμε τον \( p \)-οστό ανταλλοίωτο δείκτη με τον \( q \)-οστό συναλλοίωτο, ο ορισμός είναι
\[
S^{i_1 \dots i_{p-1} i_{p+1} \dots i_r}_{j_1 \dots j_{q-1} j_{q+1} \dots j_s} = T^{i_1 \dots i_{p-1} \alpha i_{p+1} \dots i_r}_{j_1 \dots j_{q-1} \alpha j_{q+1} \dots j_s}. \tag{12}
\]
Ο ελληνικός δείκτης \( \alpha \) είναι βωβός και αθροίζεται. Για έναν τανυστή \( T \) τύπου \( (1,1) \), η μόνη δυνατή συναίρεση δίνει το ίχνος, ένα βαθμωτό \( T^\alpha_\alpha \). Γεωμετρικά, αυτή η πράξη είναι η αποτίμηση του τανυστή πάνω σε ένα ζεύγος δυϊκών βάσεων \( dx^\alpha \) και \( \partial / \partial x^\alpha \) και η άθροιση όλων των συνεισφορών. Η διαδικασία εξουδετερώνει τη συμμετοχή των δύο δεικτών στον νόμο μετασχηματισμού, αφού οι αντίστροφοι Ιακωβιανοί πίνακες συμπτύσσονται στη μονάδα, αφήνοντας ένα γνήσιο γεωμετρικό αντικείμενο με δύο λιγότερους δείκτες. Η συναίρεση συλλαμβάνει την ουσία της "αλληλεπίδρασης" ανάμεσα σε ανταλλοίωτες και συναλλοίωτες ποιότητες, παράγοντας αναλλοίωτα μεγέθη.

\textbf{Συστολή}  
Δεν είναι ανεξάρτητη πράξη αλλά σύνθεση των δύο προηγούμενων. Συστολή ενός τανυστή \( T \) με ένα διάνυσμα \( \xi \) σε μια επιλεγμένη υποδοχή σημαίνει: πρώτα παίρνουμε το τανυστικό γινόμενο \( T \otimes \xi \) και μετά εκτελούμε συναίρεση ανάμεσα στον αντίστοιχο συναλλοίωτο δείκτη του \( T \) και τον ανταλλοίωτο δείκτη του \( \xi \). Για συστολή ως προς τον \( q \)-οστό συναλλοίωτο δείκτη, ο ορισμός είναι
\[
S^{i_1 \dots i_r}_{j_1 \dots j_{q-1} j_{q+1} \dots j_s} = T^{i_1 \dots i_r}_{j_1 \dots j_{q-1} \alpha j_{q+1} \dots j_s} \, \xi^\alpha. \tag{13}
\]
Γεωμετρικά, η συστολή είναι η επιτέλεση της δράσης: "ταΐζουμε" τον τανυστή με ένα συγκεκριμένο γεωμετρικό αντικείμενο σε μία από τις υποδοχές του, μειώνοντας την τάξη του και παράγοντας ένα νέο αντικείμενο που αναμένει τα υπόλοιπα ορίσματα. Ο δείκτης \( \alpha \) που ανήκε στο όρισμα χάνεται μέσω της άθροισης, έχοντας εξαντλήσει τον ρόλο του στη γραμμική δράση.

\section{Σύνοψη: Οι Θεμελιώδεις Τύποι}

\textbf{Τανυστής τύπου} \( (r, s) \)

\[
\tilde{T}^{i_1 \dots i_r}_{j_1 \dots j_s} = \frac{\partial y^{i_1}}{\partial x^{\mu_1}} \dots \frac{\partial y^{i_r}}{\partial x^{\mu_r}} \ T^{\mu_1 \dots \mu_r}_{\nu_1 \dots \nu_s} \ \frac{\partial x^{\nu_1}}{\partial y^{j_1}} \dots \frac{\partial x^{\nu_s}}{\partial y^{j_s}}. \tag{10}
\]

Κάθε ανταλλοίωτος δείκτης μετασχηματίζεται με τον Ιακωβιανό πίνακα της αλλαγής συντεταγμένων και κάθε συναλλοίωτος με τον αντίστροφό του. Η σύζευξη παλιών και νέων δεικτών μέσω βωβής άθροισης εξασφαλίζει το γεωμετρικό νόημα.

\textbf{Τανυστικό Γινόμενο}

\[
(T \otimes S)^{i_1 \dots i_r k_1 \dots k_p}_{j_1 \dots j_s l_1 \dots l_q} = T^{i_1 \dots i_r}_{j_1 \dots j_s} S^{k_1 \dots k_p}_{l_1 \dots l_q}. \tag{11}
\]

Οι δύο τανυστές συνενώνονται χωρίς αλληλεπίδραση δεικτών. Το αποτέλεσμα δέχεται το σύνολο των ορισμάτων και των δύο, πολλαπλασιάζοντας τις εξόδους τους.

\textbf{Συναίρεση}

\[
S^{i_1 \dots i_{p-1} i_{p+1} \dots i_r}_{j_1 \dots j_{q-1} j_{q+1} \dots j_s} = T^{i_1 \dots i_{p-1} \alpha i_{p+1} \dots i_r}_{j_1 \dots j_{q-1} \alpha j_{q+1} \dots j_s}. \tag{12}
\]

Ένας πάνω και ένας κάτω δείκτης ταυτίζονται και αθροίζονται, μειώνοντας την τάξη κατά δύο. Πρόκειται για την αποτίμηση του τανυστή σε ένα ζεύγος δυϊκών βάσεων.

\textbf{Συστολή}

\[
S^{i_1 \dots i_r}_{j_1 \dots j_{q-1} j_{q+1} \dots j_s} = T^{i_1 \dots i_r}_{j_1 \dots j_{q-1} \alpha j_{q+1} \dots j_s} \, \xi^\alpha. \tag{13}
\]

Το τανυστικό γινόμενο με ένα διάνυσμα (ή 1-μορφή) ακολουθείται από συναίρεση. Η πράξη "ταΐζει" τον τανυστή με ένα συγκεκριμένο γεωμετρικό αντικείμενο, εξαντλώντας μία από τις υποδοχές του.

\section{Τανυστικά πεδία}

Όλη η προηγούμενη ανάπτυξη αφορούσε έναν μεμονωμένο διανυσματικό χώρο \( V \), σαν να κοιτάζαμε τον εφαπτόμενο χώρο σε ένα μόνο σημείο μιας πολλαπλότητας. Η μετάβαση στην έννοια του τανυστικού πεδίου είναι η φυσική γενίκευση που προκύπτει όταν σε κάθε σημείο της πολλαπλότητας αντιστοιχίζεται ένας τανυστής του αντίστοιχου εφαπτόμενου χώρου, με τρόπο ομαλό ως προς τις συντεταγμένες.

Συγκεκριμένα, αν \( M \) είναι μια λεία πολλαπλότητα διάστασης \( n \), τότε σε κάθε σημείο \( p \in M \) ορίζεται ο εφαπτόμενος χώρος \( T_pM \) και ο δυϊκός του \( T_p^*M \). Ένα τανυστικό πεδίο τύπου \( (r, s) \) είναι μια απεικόνιση που σε κάθε σημείο \( p \) αποδίδει έναν τανυστή \( T(p) \) τύπου \( (r, s) \) επί του \( T_pM \), δηλαδή μια πολυγραμμική απεικόνιση
\[
T(p): \underbrace{T_p^*M \times \dots \times T_p^*M}_{r \text{ φορές}} \times \underbrace{T_pM \times \dots \times T_pM}_{s \text{ φορές}} \to \mathbb{R}.
\]
Το καίριο σημείο που καθιστά εφικτό αυτόν τον ορισμό είναι ακριβώς ο θεμελιώδης νόμος (10): σε μια τομή δύο χαρτών με συντεταγμένες \( x^i \) και \( y^i \), οι συνιστώσες του \( T \) στα δύο συστήματα συνδέονται μέσω των Ιακωβιανών πινάκων με τρόπο που εγγυάται ότι πρόκειται για το ίδιο γεωμετρικό αντικείμενο. Χωρίς αυτή τη συνεπή συμπεριφορά υπό αλλαγή συντεταγμένων, η έννοια του πεδίου θα ήταν άνευ περιεχομένου, αφού θα εξαρτιόταν από την αυθαίρετη επιλογή χάρτη. Ο νόμος (10) είναι λοιπόν ο μηχανισμός που "συγκολλά" τους τανυστές στα επιμέρους σημεία σε ένα ενιαίο, καθολικά ορισμένο αντικείμενο.

Η απαίτηση ομαλότητας διατυπώνεται με τον απλούστερο δυνατό τρόπο: σε έναν δεδομένο χάρτη με συντεταγμένες \( x^i \), οι συναρτήσεις \( T^{i_1 \dots i_r}_{j_1 \dots j_s}(x) \) απαιτείται να είναι ομαλές (να έχουν συνεχείς μερικές παραγώγους κάθε τάξης). Λόγω του νόμου (10), αν αυτή η ιδιότητα ισχύει σε έναν χάρτη, τότε ισχύει και σε κάθε άλλον χάρτη που επικαλύπτεται με αυτόν, αφού οι Ιακωβιανοί πίνακες είναι ομαλές συναρτήσεις των συντεταγμένων. Έτσι, η διαφορισιμότητα του τανυστικού πεδίου είναι μια καλώς ορισμένη έννοια, ανεξάρτητη του χάρτη.

Αυτό που καθιστά την έννοια του τανυστικού πεδίου τόσο ισχυρή είναι ότι η άλγεβρα των δεικτών που ορίσαμε στις ενότητες 2–4 είναι απολύτως τοπική: κάθε πράξη εκτελείται σημείο προς σημείο, ανεξάρτητα από τη γεωμετρία της πολλαπλότητας. Η καμπυλότητα και η συνοχή εισέρχονται μόνο όταν θελήσουμε να συγκρίνουμε τανυστές σε γειτονικά σημεία μέσω της συναλλοίωτης παραγώγου, κάτι που ξεπερνά το πλαίσιο της παρούσας ανάπτυξης. Αυτό που έχει σημασία εδώ είναι ότι ολόκληρη η αλγεβρική υποδομή που έχουμε οικοδομήσει μεταφέρεται αυτούσια στην περίπτωση των πεδίων, απλώς οι συνιστώσες αποκτούν εξάρτηση από το σημείο.

\section{Εξειδίκευση σε Συγκεκριμένους Τύπους Τανυστών}

Η γενική θεωρία που αναπτύξαμε φιλοξενεί ως ειδικές περιπτώσεις όλα τα γνωστά αντικείμενα της γραμμικής άλγεβρας και της γεωμετρίας. Σε κάθε περίπτωση, ο γενικός συμβολισμός εξειδικεύεται φυσικά, χωρίς την εισαγωγή νέων συμβάσεων.

\subsection{Βαθμωτές ποσότητες}
Ένα βαθμωτό είναι ένας τανυστής τύπου \( (0,0) \). Δεν έχει δείκτες και δεν απαιτεί ορίσματα από τους χώρους \( V \) ή \( V^* \)· είναι απλώς μια απεικόνιση από το μονοσύνολο στο \( \mathbb{R} \), δηλαδή ένας πραγματικός αριθμός. Ο νόμος μετασχηματισμού (10) για \( r=s=0 \) δίνει \( \tilde{T} = T \), επιβεβαιώνοντας ότι ένα βαθμωτό είναι αναλλοίωτο κάτω από αλλαγές συντεταγμένων. Στον αλγεβρικό ορισμό (8), δεν υπάρχουν παράγοντες \( V^* \) ή \( V \) στο πεδίο ορισμού.

\subsection{Διανύσματα}
Ένα διάνυσμα \( \xi \) είναι τανυστής τύπου \( (1,0) \). Ως πολυγραμμική απεικόνιση, δέχεται ένα στοιχείο του \( V^* \) και επιστρέφει έναν πραγματικό αριθμό: \( \xi(\omega) = \omega(\xi) \). Οι συνιστώσες του είναι \( \xi^i \) με έναν ανταλλοίωτο δείκτη. Ο νόμος μετασχηματισμού (10) εξειδικεύεται στην
\[
\tilde{\xi}^i = \frac{\partial y^i}{\partial x^\mu} \ \xi^\mu. \tag{14}
\]
Το διάνυσμα δεν διαθέτει συναλλοίωτες υποδοχές για να "ταΐσει" κανείς· αντίθετα, το ίδιο λειτουργεί ως όρισμα σε 1-μορφές μέσω της συστολής (13), η οποία εδώ ανάγεται στην απλή σύμβαση άθροισης \( \omega_\alpha \xi^\alpha \).

\subsection{1-μορφές}
Μια 1-μορφή \( \omega \) είναι τανυστής τύπου \( (0,1) \). Ως απεικόνιση, δέχεται ένα διάνυσμα και επιστρέφει ένα βαθμωτό: \( \omega(\xi) \). Οι συνιστώσες της είναι \( \omega_i \) με έναν συναλλοίωτο δείκτη. Ο νόμος μετασχηματισμού (10) εξειδικεύεται στην
\[
\tilde{\omega}_i = \frac{\partial x^\nu}{\partial y^i} \ \omega_\nu. \tag{15}
\]
Η δράση της σε διάνυσμα είναι η συστολή (13) για \( r=0, s=1 \), που δίνει το βαθμωτό \( \omega_\alpha \xi^\alpha \), αναδεικνύοντας τη θεμελιώδη δυϊκότητα ανάμεσα σε διανύσματα και 1-μορφές.

\subsection{Γραμμικοί μετασχηματισμοί}
Ένας γραμμικός μετασχηματισμός \( a: V \to V \) είναι τανυστής τύπου \( (1,1) \). Ως πολυγραμμική απεικόνιση, δέχεται μια 1-μορφή και ένα διάνυσμα και επιστρέφει βαθμωτό: \( a(\omega, \xi) = \omega(a(\xi)) \). Οι συνιστώσες του είναι \( a^i_j \) με έναν ανταλλοίωτο και έναν συναλλοίωτο δείκτη. Ο νόμος μετασχηματισμού (10) δίνει
\[
\tilde{a}^i_j = \frac{\partial y^i}{\partial x^\mu} \ a^\mu_\nu \ \frac{\partial x^\nu}{\partial y^j}. \tag{16}
\]
Η δράση του πάνω σε ένα διάνυσμα \( \xi \) προκύπτει από τη συστολή (13) ως προς τον συναλλοίωτο δείκτη: \( (a(\xi))^i = a^i_\alpha \xi^\alpha \), που είναι ακριβώς ο πολλαπλασιασμός πίνακα με διάνυσμα. Το ίχνος του μετασχηματισμού είναι η συναίρεση (12) των δύο δεικτών του: \( a^\alpha_\alpha \).

\subsection{Τετραγωνικές μορφές}
Μια τετραγωνική μορφή \( g \) είναι ένας συμμετρικός τανυστής τύπου \( (0,2) \). Ως απεικόνιση, δέχεται δύο διανύσματα και επιστρέφει ένα βαθμωτό, γραμμικά σε κάθε όρισμα: \( g(\xi, \eta) \). Οι συνιστώσες της είναι \( g_{ij} \) με δύο συναλλοίωτους δείκτες, και η συμμετρία σημαίνει \( g_{ij} = g_{ji} \). Ο νόμος μετασχηματισμού (10) εξειδικεύεται σε
\[
\tilde{g}_{ij} = \frac{\partial x^\mu}{\partial y^i} \ g_{\mu\nu} \ \frac{\partial x^\nu}{\partial y^j}. \tag{17}
\]
Η τετραγωνική τιμή σε ένα διάνυσμα \( \xi \) είναι η διπλή συστολή \( g_{\alpha\beta} \xi^\alpha \xi^\beta \), που αντιστοιχεί στην παράσταση \( \xi^T g \xi \) της γραμμικής άλγεβρας. Το εσωτερικό γινόμενο σε έναν Ευκλείδειο χώρο είναι μια τετραγωνική μορφή, και ο νόμος μετασχηματισμού της εγγυάται ότι το μήκος ενός διανύσματος παραμένει αναλλοίωτο.

Εφόσον η τετραγωνική μορφή \( g \) είναι μη εκφυλισμένη (ως μετρική), ορίζεται ο αντίστροφος τανυστής \( g^{-1} \), ο οποίος είναι τύπου \( (2,0) \) και οι συνιστώσες του συμβολίζονται με \( g^{ij} \). Ο ορισμός του προκύπτει από την απαίτηση η συναίρεση με τον \( g \) να δίνει το ταυτοτικό στοιχείο:
\[
g^{i\alpha} g_{\alpha j} = \delta^i_j. \tag{18}
\]
Με άλλα λόγια, σε κάθε σημείο ο πίνακας \( g^{ij} \) είναι ο αντίστροφος του πίνακα \( g_{ij} \) με τη συνήθη έννοια της γραμμικής άλγεβρας. Ο νόμος μετασχηματισμού του \( g^{-1} \) είναι αυτός ενός \( (2,0) \)-τανυστή, δηλαδή
\[
\tilde{g}^{ij} = \frac{\partial y^i}{\partial x^\mu} \frac{\partial y^j}{\partial x^\nu} \ g^{\mu\nu}. \tag{19}
\]
Το γεγονός ότι οι σχέσεις (17) και (19) είναι συμβατές με την (18) οφείλεται ακριβώς στο ότι οι Ιακωβιανοί πίνακες είναι αντίστροφοι μεταξύ τους: η συναίρεση ενός πάνω και ενός κάτω δείκτη εξουδετερώνει τους μετασχηματισμούς, αφήνοντας το \( \delta^i_j \) αναλλοίωτο. Γεωμετρικά, ο \( g^{-1} \) επιτρέπει την ανύψωση δεικτών, μετατρέποντας συναλλοίωτες συνιστώσες σε ανταλλοίωτες: \( \xi^i = g^{i\alpha} \omega_\alpha \). Αντίστροφα, ο \( g \) επιτρέπει τον καταβιβασμό δεικτών: \( \omega_i = g_{i\alpha} \xi^\alpha \). Αυτή η δυϊκή σχέση ανάμεσα στον \( g \) και τον \( g^{-1} \) είναι θεμελιώδης στη διαφορική γεωμετρία, καθώς επιτρέπει την ταύτιση διανυσμάτων και 1-μορφών μέσω της μετρικής δομής, ένα βήμα που είναι απολύτως απαραίτητο για να οριστεί, για παράδειγμα, η κλίση μιας συνάρτησης ως διάνυσμα.

\subsection{Συμπλεκτική μορφή}
Μια συμπλεκτική μορφή \( \omega \) είναι ένας αντισυμμετρικός, μη εκφυλισμένος τανυστής τύπου \( (0,2) \). Ως απεικόνιση, δέχεται δύο διανύσματα και επιστρέφει ένα βαθμωτό: \( \omega(\xi, \eta) \). Οι συνιστώσες της είναι \( \omega_{ij} \) με δύο συναλλοίωτους δείκτες και ικανοποιούν \( \omega_{ij} = -\omega_{ji} \). Ο νόμος μετασχηματισμού της είναι ταυτόσημος με αυτόν κάθε \( (0,2) \)-τανυστή:
\[
\tilde{\omega}_{ij} = \frac{\partial x^\mu}{\partial y^i} \ \omega_{\mu\nu} \ \frac{\partial x^\nu}{\partial y^j}. \tag{20}
\]
Η μη εκφυλισμένη ιδιότητα σημαίνει ότι για κάθε μη μηδενικό διάνυσμα \( \xi \), η 1-μορφή \( \omega(\xi, \cdot) \) που προκύπτει από τη συστολή \( \omega_{\alpha\beta} \xi^\alpha \) είναι μη μηδενική. Η συμπλεκτική μορφή είναι ο θεμέλιος λίθος της κλασικής μηχανικής στον φασικό χώρο, και ο τανυστικός της χαρακτήρας εγγυάται ότι οι εξισώσεις Hamilton είναι ανεξάρτητες από την επιλογή κανονικών συντεταγμένων.

\subsection{Τανυστής τάσεων}
Ο τανυστής τάσεων \( \sigma \) είναι ένας συμμετρικός τανυστής τύπου \( (2,0) \) στη μηχανική των συνεχών μέσων, αν και συχνά αναπαρίσταται και ως \( (0,2) \) ή \( (1,1) \) ανάλογα με το πλαίσιο. Εδώ τον θεωρούμε τύπου \( (2,0) \) ως το φυσικό αντικείμενο που δέχεται δύο 1-μορφές (κάθετες σε επιφάνειες) και επιστρέφει το βαθμωτό της τάσης σε εκείνη την κατεύθυνση. Ως πολυγραμμική απεικόνιση, \( \sigma(\omega, \eta) \) δίνει τη συνιστώσα της δύναμης που ασκείται σε ένα επιφανειακό στοιχείο. Οι συνιστώσες του είναι \( \sigma^{ij} \) με δύο ανταλλοίωτους δείκτες και η συμμετρία σημαίνει \( \sigma^{ij} = \sigma^{ji} \). Ο νόμος μετασχηματισμού (10) εξειδικεύεται σε
\[
\tilde{\sigma}^{ij} = \frac{\partial y^i}{\partial x^\mu} \frac{\partial y^j}{\partial x^\nu} \ \sigma^{\mu\nu}. \tag{21}
\]
Η συστολή του με μια 1-μορφή \( n \) (το κάθετο διάνυσμα σε μια επιφάνεια) δίνει το διάνυσμα της τάσης \( t^i = \sigma^{i\alpha} n_\alpha \), το οποίο εκφράζει τη δύναμη ανά μονάδα επιφάνειας που ασκείται στο στοιχείο με κάθετο \( n \). Η απόκλιση του τανυστή τάσεων, \( \partial_\alpha \sigma^{\alpha i} \), εμφανίζεται στις εξισώσεις ισορροπίας και είναι ένα ανταλλοίωτο διάνυσμα, όπως απαιτεί η φυσική σημασία της δύναμης.

\subsection{Τανυστής ηλεκτρομαγνητικού πεδίου}
Ο τανυστής ηλεκτρομαγνητικού πεδίου \( F \) είναι ένας αντισυμμετρικός τανυστής τύπου \( (0,2) \). Ως απεικόνιση, δέχεται δύο διανύσματα και επιστρέφει ένα βαθμωτό που σχετίζεται με τη δύναμη Lorentz. Οι συνιστώσες του είναι \( F_{ij} \) με δύο συναλλοίωτους δείκτες και ικανοποιούν \( F_{ij} = -F_{ji} \). Σε Καρτεσιανές συντεταγμένες, οι μη μηδενικές συνιστώσες του συνδέονται με το ηλεκτρικό και το μαγνητικό πεδίο. Ο νόμος μετασχηματισμού (10) δίνει
\[
\tilde{F}_{ij} = \frac{\partial x^\mu}{\partial y^i} \ F_{\mu\nu} \ \frac{\partial x^\nu}{\partial y^j}. \tag{22}
\]
Η αντισυμμετρία διατηρείται από τον μετασχηματισμό, αφού οι Ιακωβιανοί πίνακες συστέλλονται συμμετρικά. Η δύναμη Lorentz σε ένα σωματίδιο με ταχύτητα \( \xi \) προκύπτει από τη συστολή \( F_{\alpha\beta} \xi^\beta \), η οποία δίνει μια 1-μορφή (το ηλεκτρικό πεδίο στο σύστημα ηρεμίας). Οι εξισώσεις Maxwell εκφράζονται πλήρως μέσω του \( F \) και του δυϊκού του, αναδεικνύοντας τη γεωμετρική φύση του ηλεκτρομαγνητισμού ως θεωρίας ενός τανυστικού πεδίου.

\subsection{Τανυστής καμπυλότητας Riemann}
Ο τανυστής καμπυλότητας Riemann \( R \) είναι ένας τανυστής τύπου \( (1,3) \) που περιγράφει την εσωτερική καμπυλότητα μιας πολλαπλότητας με μετρική. Ως πολυγραμμική απεικόνιση, δέχεται μια 1-μορφή και τρία διανύσματα και επιστρέφει ένα βαθμωτό που μετρά τη μη-αντιμεταθετικότητα της συναλλοίωτης παραγώγου. Οι συνιστώσες του είναι \( R^i_{jkl} \) με έναν ανταλλοίωτο και τρεις συναλλοίωτους δείκτες. Ο νόμος μετασχηματισμού (10) εξειδικεύεται σε
\[
\tilde{R}^i_{jkl} = \frac{\partial y^i}{\partial x^\mu} \ \frac{\partial x^\nu}{\partial y^j} \frac{\partial x^\rho}{\partial y^k} \frac{\partial x^\sigma}{\partial y^l} \ R^\mu_{\nu\rho\sigma}. \tag{23}
\]
Ο τανυστής Riemann ικανοποιεί μια σειρά από αλγεβρικές συμμετρίες: \( R^i_{jkl} = -R^i_{jlk} \) (αντισυμμετρία στους δύο τελευταίους δείκτες), \( R^i_{jkl} + R^i_{klj} + R^i_{ljk} = 0 \) (πρώτη ταυτότητα Bianchi), και \( R_{ijkl} = R_{klij} \) (συμμετρία ζευγών μετά από καταβιβασμό του πρώτου δείκτη με τη μετρική). Η συστολή του σε έναν ανταλλοίωτο και έναν συναλλοίωτο δείκτη δίνει τον τανυστή Ricci \( R_{jl} = R^\alpha_{j\alpha l} \), τύπου \( (0,2) \), και περαιτέρω συναίρεση με τη μετρική δίνει τη βαθμωτή καμπυλότητα \( R = g^{\alpha\beta} R_{\alpha\beta} \). Ο τανυστής Riemann είναι το κεντρικό αντικείμενο της γενικής θεωρίας της σχετικότητας, καθώς κωδικοποιεί πλήρως τη βαρυτική αλληλεπίδραση μέσω της καμπύλωσης του χωροχρόνου.

\section{Ψευδο-Ρημάννειες και Ρημάννειες Μετρικές}

Ας αναπτύξουμε τώρα τη θεωρία των ψευδο-Ρημάννειων μετρικών ως φυσική συνέχεια και εξειδίκευση της τανυστικής θεωρίας που οικοδομήσαμε, εστιάζοντας στη γεωμετρική τους σημασία και στις αλγεβρικές τους ιδιότητες.

\subsection{Ορισμός και βασικές ιδιότητες}

Στο πλαίσιο της τανυστικής θεωρίας που αναπτύξαμε, μια μετρική είναι απλώς μια τετραγωνική μορφή με πρόσημο. Συγκεκριμένα, μια \textbf{ψευδο-Ρημάννεια μετρική} \( g \) σε έναν διανυσματικό χώρο \( V \) είναι ένας συμμετρικός, μη εκφυλισμένος τανυστής τύπου \( (0,2) \), δηλαδή μια απεικόνιση
\begin{equation}
g: V \times V \to \mathbb{R}
\end{equation}
με τις ιδιότητες που ήδη συζητήσαμε στην ενότητα 6 για τις τετραγωνικές μορφές. Ο όρος "ψευδο-Ρημάννεια" αναφέρεται στο γεγονός ότι η μετρική δεν είναι απαραίτητα θετικά ορισμένη· μπορεί να έχει μικτή υπογραφή, δηλαδή να υπάρχουν μη μηδενικά διανύσματα \( \xi \) για τα οποία \( g(\xi, \xi) = 0 \) (ισοτροπικά ή φωτοειδή διανύσματα), και διανύσματα για τα οποία \( g(\xi, \xi) < 0 \). Μια \textbf{Ρημάννεια μετρική} είναι η ειδική περίπτωση όπου η μετρική είναι θετικά ορισμένη: \( g(\xi, \xi) > 0 \) για κάθε μη μηδενικό \( \xi \).

Σε μια βάση \( \{ \partial / \partial x^i \} \), οι συνιστώσες της μετρικής είναι \( g_{ij} \) και η συμμετρία εκφράζεται ως
\begin{equation}
g_{ij} = g_{ji}. \tag{24}
\end{equation}
Η μετρική επιτρέπει τον ορισμό του εσωτερικού γινομένου δύο διανυσμάτων:
\begin{equation}
\langle \xi, \eta \rangle \equiv g(\xi, \eta) = g_{ij} \xi^i \eta^j. \tag{25}
\end{equation}
Το μήκος (ή ψευδο-μήκος) ενός διανύσματος δίνεται από την τετραγωνική έκφραση
\begin{equation}
g(\xi, \xi) = g_{ij} \xi^i \xi^j, \tag{26}
\end{equation}
με την κατανόηση ότι στην ψευδο-Ρημάννεια περίπτωση αυτή η ποσότητα μπορεί να είναι αρνητική ή μηδέν. Η μη εκφυλισμένη ιδιότητα σημαίνει ότι ο πίνακας \( g_{ij} \) είναι αντιστρέψιμος σε κάθε σημείο, και ο αντίστροφός του \( g^{ij} \) είναι ο τανυστής τύπου \( (2,0) \) που ήδη ορίσαμε στην (18).

Ο νόμος μετασχηματισμού της μετρικής κάτω από αλλαγή συντεταγμένων \( x^i \to y^i \) είναι, όπως για κάθε \( (0,2) \)-τανυστή,
\begin{equation}
\tilde{g}_{ij} = \frac{\partial x^\mu}{\partial y^i} \ g_{\mu\nu} \ \frac{\partial x^\nu}{\partial y^j}. \tag{27}
\end{equation}

\subsection{Μουσικοί ισομορφισμοί}

Η ύπαρξη μιας μη εκφυλισμένης μετρικής επιτρέπει μια θεμελιώδη και κομψή ταύτιση ανάμεσα στον διανυσματικό χώρο \( V \) και τον δυϊκό του \( V^* \). Η ταύτιση αυτή δεν είναι αυθαίρετη αλλά υπαγορεύεται πλήρως από τη γεωμετρία που κωδικοποιεί η μετρική, και ονομάζεται μουσικός ισομορφισμός.

Συγκεκριμένα, σε κάθε διάνυσμα \( \xi \in V \) αντιστοιχίζεται μια 1-μορφή \( \xi^\flat \in V^* \) μέσω της πράξης του καταβιβασμού δείκτη:
\begin{equation}
(\xi^\flat)_i = g_{i\alpha} \xi^\alpha. \tag{28}
\end{equation}
Γεωμετρικά, η \( \xi^\flat \) είναι η μηχανή που δέχεται ένα άλλο διάνυσμα \( \eta \) και επιστρέφει το εσωτερικό γινόμενο \( g(\xi, \eta) \). Με άλλα λόγια, \( \xi^\flat(\eta) = g(\xi, \eta) \). Η πράξη αυτή είναι αντιστρέψιμη: σε κάθε 1-μορφή \( \omega \in V^* \) αντιστοιχίζεται ένα διάνυσμα \( \omega^\sharp \in V \) μέσω της ανύψωσης δείκτη:
\begin{equation}
(\omega^\sharp)^i = g^{i\alpha} \omega_\alpha. \tag{29}
\end{equation}
Οι δύο αυτές απεικονίσεις είναι αμοιβαία αντίστροφες λόγω της (18): \( (\xi^\flat)^\sharp = \xi \) και \( (\omega^\sharp)^\flat = \omega \). Αποτελούν λοιπόν έναν ισομορφισμό \( V \cong V^* \), γνωστό ως μουσικό ισομορφισμό (η ορολογία προέρχεται από τα σύμβολα \( \flat \) "flat" και \( \sharp \) "sharp" που χρησιμοποιούνται).

Η φυσικογεωμετρική σημασία αυτής της ταύτισης είναι τεράστια. Στη φυσική, η κλίση μιας βαθμωτής συνάρτησης \( f \) ορίζεται φυσικά ως η 1-μορφή \( df \) με συνιστώσες \( \partial f / \partial x^i \). Όμως, για να μιλήσουμε για το "διάνυσμα της κλίσης" που δείχνει στην κατεύθυνση της μέγιστης μεταβολής, χρειαζόμαστε τη μετρική: το διάνυσμα αυτό είναι ακριβώς το \( (df)^\sharp \), με συνιστώσες \( g^{i\alpha} \partial f / \partial x^\alpha \). Χωρίς μετρική, η κλίση παραμένει μια 1-μορφή και δεν μπορεί να ταυτιστεί με διάνυσμα. Αντιστρόφως, η ταχύτητα ενός σωματίδιου είναι φυσικά ένα διάνυσμα \( \xi \), αλλά για να οριστεί η ορμή του ως 1-μορφή (η οποία είναι το αντικείμενο που εμφανίζεται φυσικά στις εξισώσεις Hamilton), χρειαζόμαστε τον καταβιβασμό: η ορμή είναι η \( \xi^\flat \). Στη γενική σχετικότητα, η μουσική ισοδυναμία επιτρέπει τη μετάφραση ανάμεσα σε διανυσματικά και συναλλοίωτα αντικείμενα κατά βούληση, ανάλογα με το ποια μορφή είναι πιο βολική για έναν δεδομένο υπολογισμό, χωρίς να αλλάζει το φυσικό περιεχόμενο.

\subsection{Μετρική ως τανυστικό πεδίο}

Όπως ακριβώς συζητήσαμε στην ενότητα 5 για τα γενικά τανυστικά πεδία, μια μετρική σε μια πολλαπλότητα \( M \) είναι ένα τανυστικό πεδίο τύπου \( (0,2) \) που σε κάθε σημείο \( p \in M \) αποδίδει μια ψευδο-Ρημάννεια (ή Ρημάννεια) μετρική \( g(p) \) στον εφαπτόμενο χώρο \( T_pM \). Σε έναν χάρτη με συντεταγμένες \( x^i \), οι συνιστώσες \( g_{ij}(x) \) είναι ομαλές συναρτήσεις των συντεταγμένων. Ο νόμος μετασχηματισμού (27) εγγυάται ότι το αντικείμενο είναι καθολικά ορισμένο, ανεξάρτητα από την επιλογή χάρτη, και η ομαλότητα διατηρείται λόγω της ομαλότητας των Ιακωβιανών πινάκων. Το ζεύγος \( (M, g) \) μιας πολλαπλότητας εφοδιασμένης με μια μετρική ονομάζεται ψευδο-Ρημάννεια (ή Ρημάννεια) πολλαπλότητα. Η μετρική είναι το θεμελιώδες αντικείμενο που επιτρέπει τη μέτρηση αποστάσεων, γωνιών και όγκων στην πολλαπλότητα, και η γεωμετρία της καθορίζεται πλήρως από αυτήν.

\subsection{Το στοιχείο όγκου και οι ταυτότητες του εξωτερικού γινομένου}

Σε μια προσανατολισμένη ψευδο-Ρημάννεια πολλαπλότητα διάστασης \( n \), η μετρική επάγει με φυσικό τρόπο μια έννοια όγκου. Σε έναν χάρτη με συντεταγμένες \( x^1, \dots, x^n \), ο όγκος ενός απειροστού παραλληλεπιπέδου που ορίζεται από τα διανύσματα βάσης \( \partial / \partial x^1, \dots, \partial / \partial x^n \) δεν είναι απλώς το γινόμενο των διαφορικών \( dx^1 \dots dx^n \), αλλά πρέπει να σταθμιστεί με έναν παράγοντα που λαμβάνει υπόψη την παραμόρφωση που επιφέρει η μετρική. Ο παράγοντας αυτός είναι \( \sqrt{|\det(g_{ij})|} \), όπου \( \det(g_{ij}) \) είναι η ορίζουσα του πίνακα της μετρικής. Η απόλυτη τιμή εξασφαλίζει ότι το στοιχείο όγκου είναι θετικά ορισμένο ακόμα και στην ψευδο-Ρημάννεια περίπτωση όπου η ορίζουσα μπορεί να είναι αρνητική. Το γεγονός ότι ο παράγοντας \( \sqrt{|\det(g_{ij})|} \) μετασχηματίζεται με τον Ιακωβιανό πίνακα της αλλαγής συντεταγμένων εγγυάται ότι το ολοκλήρωμα μιας βαθμωτής συνάρτησης ως προς αυτό το στοιχείο όγκου είναι ανεξάρτητο από την επιλογή χάρτη.

Το σύμβολο Levi-Civita \( \epsilon_{i_1 \dots i_n} \) ορίζεται ως η πλήρως αντισυμμετρική ποσότητα με \( \epsilon_{12 \dots n} = 1 \). Το αντίστοιχο τανυστικό αντικείμενο που σέβεται τη μετρική είναι ο τανυστής Levi-Civita \( \varepsilon \), τύπου \( (0,n) \), με συνιστώσες
\begin{equation}
\varepsilon_{i_1 \dots i_n} = \sqrt{|\det(g_{ij})|} \ \epsilon_{i_1 \dots i_n}. \tag{30}
\end{equation}
Αυτός ο τανυστής είναι πλήρως αντισυμμετρικός και ικανοποιεί τον νόμο μετασχηματισμού ενός \( (0,n) \)-τανυστή, ακριβώς επειδή ο παράγοντας \( \sqrt{|\det(g_{ij})|} \) μετασχηματίζεται με τον Ιακωβιανό πίνακα της αλλαγής συντεταγμένων. Ο αντίστοιχος ανταλλοίωτος τανυστής \( \varepsilon^{i_1 \dots i_n} \) προκύπτει με ανύψωση των δεικτών μέσω της μετρικής και ικανοποιεί
\begin{equation}
\varepsilon^{i_1 \dots i_n} \varepsilon_{i_1 \dots i_n} = n! \ \operatorname{sgn}(\det(g_{ij})), \tag{31}
\end{equation}
όπου το πρόσημο της ορίζουσας εμφανίζεται λόγω της ανύψωσης των δεικτών. Οι χρήσιμες ταυτότητες συναίρεσης που τον διέπουν είναι οι εξής:

\begin{equation}
\varepsilon^{i_1 \dots i_n} \varepsilon_{j_1 \dots j_n} = n! \ \delta^{i_1}_{[j_1} \dots \delta^{i_n}_{j_n]}, \tag{32}
\end{equation}
\begin{equation}
\varepsilon^{i_1 \dots i_k i_{k+1} \dots i_n} \varepsilon_{i_1 \dots i_k j_{k+1} \dots j_n} = k! (n-k)! \ \delta^{i_{k+1}}_{[j_{k+1}} \dots \delta^{i_n}_{j_n]}, \tag{33}
\end{equation}
όπου οι αγκύλες \( [\dots] \) δηλώνουν πλήρη αντισυμμετρικότητα στους εσωκλειόμενους δείκτες. Οι ταυτότητες αυτές είναι θεμελιώδεις για τον χειρισμό του εξωτερικού γινομένου και των αντισυμμετρικών τανυστών.

\section{Σύνοψη: Οι Θεμελιώδεις Τύποι και Ταυτότητες των Μετρικών}

\textbf{Συμμετρία της μετρικής}

\begin{equation}
g_{ij} = g_{ji}. \tag{24}
\end{equation}

Η μετρική είναι ένας συμμετρικός τανυστής τύπου \( (0,2) \).

\textbf{Εσωτερικό γινόμενο και μήκος}

\begin{equation}
\langle \xi, \eta \rangle = g_{ij} \xi^i \eta^j, \qquad g(\xi, \xi) = g_{ij} \xi^i \xi^j. \tag{25, 26}
\end{equation}

Η μετρική ορίζει το εσωτερικό γινόμενο. Στην ψευδο-Ρημάννεια περίπτωση, το τετράγωνο του μήκους μπορεί να είναι αρνητικό ή μηδέν.

\textbf{Νόμος μετασχηματισμού της μετρικής}

\begin{equation}
\tilde{g}_{ij} = \frac{\partial x^\mu}{\partial y^i} \ g_{\mu\nu} \ \frac{\partial x^\nu}{\partial y^j}. \tag{27}
\end{equation}

Η μετρική μετασχηματίζεται ως \( (0,2) \)-τανυστής, εξασφαλίζοντας ότι το εσωτερικό γινόμενο είναι ανεξάρτητο από την επιλογή συντεταγμένων.

\textbf{Μουσική ισοδυναμία (καταβιβασμός και ανύψωση δεικτών)}

\begin{equation}
(\xi^\flat)_i = g_{i\alpha} \xi^\alpha, \qquad (\omega^\sharp)^i = g^{i\alpha} \omega_\alpha. \tag{28, 29}
\end{equation}

Η μετρική και ο αντίστροφός της παρέχουν έναν φυσικό ισομορφισμό ανάμεσα σε διανύσματα και 1-μορφές. Γεωμετρικά, η \( \xi^\flat \) είναι η μηχανή που δίνει το εσωτερικό γινόμενο με το \( \xi \), ενώ η \( \omega^\sharp \) είναι το διάνυσμα που αντιστοιχεί στην 1-μορφή \( \omega \) μέσω της μετρικής.

\textbf{Τανυστής Levi-Civita}

\begin{equation}
\varepsilon_{i_1 \dots i_n} = \sqrt{|\det(g_{ij})|} \ \epsilon_{i_1 \dots i_n}. \tag{30}
\end{equation}

Ο τανυστής Levi-Civita είναι ο πλήρως αντισυμμετρικός τανυστής που σέβεται τη μετρική και χρησιμεύει στον χειρισμό του εξωτερικού γινομένου.

\textbf{Ταυτότητες συναίρεσης του τανυστή Levi-Civita}

\begin{equation}
\varepsilon^{i_1 \dots i_n} \varepsilon_{j_1 \dots j_n} = n! \ \delta^{i_1}_{[j_1} \dots \delta^{i_n}_{j_n]}, \tag{32}
\end{equation}
\begin{equation}
\varepsilon^{i_1 \dots i_k i_{k+1} \dots i_n} \varepsilon_{i_1 \dots i_k j_{k+1} \dots j_n} = k! (n-k)! \ \delta^{i_{k+1}}_{[j_{k+1}} \dots \delta^{i_n}_{j_n]}. \tag{33}
\end{equation}

Αυτές οι ταυτότητες είναι θεμελιώδεις για τον χειρισμό του εξωτερικού γινομένου και των αντισυμμετρικών τανυστών.

\section{Παραδείγματα Μετρικών}

\textbf{Ευκλείδεια μετρική}  
Στον \( \mathbb{R}^n \) με Καρτεσιανές συντεταγμένες \( x^1, \dots, x^n \), η Ευκλείδεια μετρική είναι η απλούστερη Ρημάννεια μετρική, με συνιστώσες
\begin{equation}
g_{ij} = \delta_{ij}. \tag{34}
\end{equation}
Το εσωτερικό γινόμενο είναι το σύνηθες \( \langle \xi, \eta \rangle = \sum_{i=1}^n \xi^i \eta^i \) και ο αντίστροφος τανυστής έχει συνιστώσες \( g^{ij} = \delta^{ij} \). Ο παράγοντας όγκου είναι απλώς \( \sqrt{|\det(g_{ij})|} = 1 \).

\textbf{Μετρική Minkowski}  
Στον \( \mathbb{R}^4 \) με συντεταγμένες \( x^0, x^1, x^2, x^3 \), η μετρική Minkowski είναι μια ψευδο-Ρημάννεια μετρική υπογραφής \( (1,3) \) (ή \( (3,1) \) ανάλογα με τη σύμβαση). Υιοθετώντας τη σύμβαση της σχετικότητας, οι συνιστώσες της είναι
\begin{equation}
g_{ij} = \text{diag}(-1, 1, 1, 1), \tag{35}
\end{equation}
δηλαδή \( g_{00} = -1 \), \( g_{11} = g_{22} = g_{33} = 1 \), και όλες οι μη διαγώνιες συνιστώσες μηδέν. Το εσωτερικό γινόμενο δύο διανυσμάτων είναι
\begin{equation}
\langle \xi, \eta \rangle = -\xi^0 \eta^0 + \xi^1 \eta^1 + \xi^2 \eta^2 + \xi^3 \eta^3.
\end{equation}
Ο παράγοντας όγκου είναι \( \sqrt{|\det(g_{ij})|} = 1 \). Η μετρική Minkowski είναι η θεμελιώδης μετρική της Ειδικής Θεωρίας της Σχετικότητας, και ο νόμος μετασχηματισμού της (27) κάτω από μετασχηματισμούς Lorentz εγγυάται ότι το αναλλοίωτο διάστημα \( ds^2 = g_{ij} dx^i dx^j \) παραμένει αναλλοίωτο.

\textbf{Μετρική σε σφαιρικές συντεταγμένες}  
Στον \( \mathbb{R}^3 \) με σφαιρικές συντεταγμένες \( r, \theta, \phi \), όπου
\begin{equation}
x^1 = r, \quad x^2 = \theta, \quad x^3 = \phi,
\end{equation}
η Ευκλείδεια μετρική παίρνει τη μορφή
\begin{equation}
g_{ij} = \text{diag}(1, r^2, r^2 \sin^2\theta). \tag{36}
\end{equation}
Συγκεκριμένα, \( g_{11} = 1 \), \( g_{22} = r^2 \), \( g_{33} = r^2 \sin^2\theta \). Ο παράγοντας όγκου είναι \( \sqrt{|\det(g_{ij})|} = r^2 \sin\theta \). Αυτό το παράδειγμα αναδεικνύει πώς μια μετρική που είναι επίπεδη (έχει μηδενική καμπυλότητα) μπορεί να έχει μη σταθερές συνιστώσες σε ένα διαφορετικό σύστημα συντεταγμένων, γεγονός που οφείλεται ακριβώς στον νόμο μετασχηματισμού (27).

\section{Διαφορικές Μορφές και Σφηνοειδές Γινόμενο}

Ας αναπτύξουμε τώρα τη θεωρία των διαφορικών μορφών ως φυσική εξειδίκευση της τανυστικής θεωρίας που οικοδομήσαμε, εστιάζοντας στους αντισυμμετρικούς συναλλοίωτους τανυστές και στο σφηνοειδές γινόμενο.

\subsection{Αντισυμμετρικοί τανυστές και η έννοια της μορφής}

Στο πλαίσιο της γενικής τανυστικής θεωρίας που αναπτύξαμε, μια \textbf{διαφορική μορφή} βαθμού \( k \) (ή απλώς \( k \)-μορφή) σε έναν διανυσματικό χώρο \( V \) είναι ένας πλήρως αντισυμμετρικός τανυστής τύπου \( (0,k) \). Δηλαδή, είναι μια πολυγραμμική απεικόνιση
\[
\omega: \underbrace{V \times \dots \times V}_{k \text{ φορές}} \to \mathbb{R}
\]
που ικανοποιεί την ιδιότητα της αντισυμμετρίας: για οποιαδήποτε εναλλαγή δύο ορισμάτων, η τιμή της αλλάζει πρόσημο. Σε όρους συνιστωσών, αυτό σημαίνει ότι για κάθε μετάθεση \( \sigma \) των δεικτών,
\[
\omega_{i_{\sigma(1)} \dots i_{\sigma(k)}} = \operatorname{sgn}(\sigma) \ \omega_{i_1 \dots i_k}, \tag{37}
\]
όπου \( \operatorname{sgn}(\sigma) \) είναι το πρόσημο της μετάθεσης. Ισοδύναμα, η αντισυμμετρία μπορεί να εκφραστεί με την πλήρη αντισυμμετρικότητα ως προς οποιοδήποτε ζεύγος δεικτών:
\[
\omega_{\dots i \dots j \dots} = - \omega_{\dots j \dots i \dots}. \tag{38}
\]
Μια σημαντική συνέπεια είναι ότι αν δύο δείκτες είναι ίσοι, η αντίστοιχη συνιστώσα μηδενίζεται. Αυτό συνεπάγεται ότι σε έναν χώρο διάστασης \( n \), οι μόνες μη μηδενικές \( k \)-μορφές είναι αυτές με \( k \le n \). Για \( k > n \), η μορφή είναι ταυτοτικά μηδέν, αφού αναγκαστικά κάποιος δείκτης θα επαναληφθεί.

Οι διαφορικές μορφές συμβολίζονται συνήθως με το σύνολο των διαφορικών \( dx^i \). Μια \( k \)-μορφή γράφεται ως
\[
\omega = \frac{1}{k!} \ \omega_{i_1 \dots i_k} \ dx^{i_1} \wedge \dots \wedge dx^{i_k},
\]
όπου το σύμβολο \( \wedge \) δηλώνει το σφηνοειδές γινόμενο που θα ορίσουμε παρακάτω. Ο παράγοντας \( 1/k! \) είναι συμβατικός και εξασφαλίζει ότι η άθροιση γίνεται πάνω σε όλους τους δείκτες χωρίς υπεραρίθμηση. Σε αυτό το σημείο, αρκεί να σκεφτόμαστε την παραπάνω έκφραση ως έναν συμπαγή τρόπο γραφής ενός αντισυμμετρικού τανυστή.

Μια \textbf{απλή \( k \)-μορφή} είναι μια μορφή που μπορεί να γραφεί ως σφηνοειδές γινόμενο \( k \) 1-μορφών:
\[
\omega = \omega^1 \wedge \omega^2 \wedge \dots \wedge \omega^k,
\]
όπου κάθε \( \omega^a \) είναι μια 1-μορφή. Δεν είναι όλες οι \( k \)-μορφές απλές, αλλά κάθε \( k \)-μορφή μπορεί να γραφεί ως γραμμικός συνδυασμός απλών μορφών. Οι απλές μορφές έχουν μια σημαντική γεωμετρική ερμηνεία: μια απλή \( k \)-μορφή αντιστοιχεί σε ένα προσανατολισμένο \( k \)-διάστατο υποχώρο και, όπως θα δούμε, σχετίζεται με την έννοια του όγκου σε αυτόν τον υποχώρο.

\subsection{Το σφηνοειδές γινόμενο}

Το \textbf{σφηνοειδές γινόμενο} (ή εξωτερικό γινόμενο) είναι η φυσική πράξη που συνθέτει δύο αντισυμμετρικούς τανυστές σε έναν νέο αντισυμμετρικό τανυστή ανώτερης τάξης. Αν \( \omega \) είναι μια \( k \)-μορφή και \( \eta \) είναι μια \( l \)-μορφή, το σφηνοειδές γινόμενο τους \( \omega \wedge \eta \) είναι μια \( (k+l) \)-μορφή που ορίζεται από την πλήρη αντισυμμετρικότητα του τανυστικού γινομένου:
\[
(\omega \wedge \eta)_{i_1 \dots i_k j_1 \dots j_l} = \frac{(k+l)!}{k! \ l!} \ \omega_{[i_1 \dots i_k} \eta_{j_1 \dots j_l]}, \tag{39}
\]
όπου οι αγκύλες \( [\dots] \) δηλώνουν πλήρη αντισυμμετρικότητα στους εσωκλειόμενους δείκτες, δηλαδή άθροισμα πάνω σε όλες τις μεταθέσεις διαιρεμένο με το πλήθος τους, με πρόσημο \( \operatorname{sgn}(\sigma) \) για κάθε μετάθεση. Ο παράγοντας \( (k+l)!/(k! \ l!) \) εξασφαλίζει τη σωστή κανονικοποίηση.

Για τις 1-μορφές βάσης \( dx^i \), το σφηνοειδές γινόμενο δίνει
\[
dx^i \wedge dx^j = - dx^j \wedge dx^i, \tag{40}
\]
και ειδικότερα \( dx^i \wedge dx^i = 0 \). Γενικότερα, για μια συλλογή 1-μορφών,
\[
dx^{i_1} \wedge \dots \wedge dx^{i_k} = \sum_{\sigma \in S_k} \operatorname{sgn}(\sigma) \ dx^{i_{\sigma(1)}} \otimes \dots \otimes dx^{i_{\sigma(k)}}, \tag{41}
\]
όπου \( S_k \) είναι η ομάδα των μεταθέσεων \( k \) στοιχείων.

Το σφηνοειδές γινόμενο είναι προσεταιριστικό: \( (\omega \wedge \eta) \wedge \theta = \omega \wedge (\eta \wedge \theta) \). Είναι επίσης βαθμωτά αντιμεταθετικό: αν \( \omega \) είναι \( k \)-μορφή και \( \eta \) είναι \( l \)-μορφή, τότε
\[
\omega \wedge \eta = (-1)^{kl} \ \eta \wedge \omega. \tag{42}
\]
Αυτή η ιδιότητα είναι θεμελιώδης και αντικατοπτρίζει τη γεωμετρική φύση του προσανατολισμού: η εναλλαγή των ορισμάτων αντιστοιχεί σε αλλαγή της διάταξης των διαστάσεων.

Μια προκαταρκτική γεωμετρική νύξη είναι χρήσιμη εδώ. Το σφηνοειδές γινόμενο \( k \) 1-μορφών μετρά τον προσανατολισμένο \( k \)-διάστατο όγκο του παραλληλεπιπέδου που ορίζεται από τα αντίστοιχα διανύσματα. Αν τροφοδοτήσουμε την \( dx^1 \wedge dx^2 \) με δύο διανύσματα \( \xi, \eta \), το αποτέλεσμα \( (dx^1 \wedge dx^2)(\xi, \eta) = \xi^1 \eta^2 - \xi^2 \eta^1 \) είναι ακριβώς το εμβαδόν του παραλληλογράμμου που σχηματίζουν τα \( \xi \) και \( \eta \) στο επίπεδο \( x^1 x^2 \). Αυτή η ερμηνεία γενικεύεται σε ανώτερες διαστάσεις και αποτελεί τον πυρήνα της γεωμετρικής σημασίας των μορφών.

\subsection{Διαφορικές μορφές ως τανυστικά πεδία}

Όπως ακριβώς συζητήσαμε στην ενότητα 5 για τα γενικά τανυστικά πεδία, μια διαφορική μορφή βαθμού \( k \) σε μια πολλαπλότητα \( M \) είναι ένα τανυστικό πεδίο τύπου \( (0,k) \) που σε κάθε σημείο \( p \in M \) αποδίδει μια πλήρως αντισυμμετρική πολυγραμμική απεικόνιση στον εφαπτόμενο χώρο \( T_pM \). Σε έναν χάρτη με συντεταγμένες \( x^i \), οι συνιστώσες \( \omega_{i_1 \dots i_k}(x) \) είναι ομαλές συναρτήσεις των συντεταγμένων και είναι πλήρως αντισυμμετρικές στους δείκτες. Ο νόμος μετασχηματισμού είναι αυτός ενός \( (0,k) \)-τανυστή, και η αντισυμμετρία διατηρείται από τον μετασχηματισμό, αφού οι Ιακωβιανοί πίνακες συστέλλονται συμμετρικά. Το γεγονός ότι οι διαφορικές μορφές είναι απλώς μια ειδική κατηγορία τανυστικών πεδίων σημαίνει ότι ολόκληρη η αλγεβρική υποδομή που αναπτύξαμε (τανυστικό γινόμενο, συναίρεση, συστολή) εφαρμόζεται και εδώ, με το σφηνοειδές γινόμενο να είναι η εξειδικευμένη πράξη που σέβεται την αντισυμμετρία.

\subsection{Το στοιχείο όγκου ως διαφορική μορφή}

Σε μια προσανατολισμένη ψευδο-Ρημάννεια πολλαπλότητα διάστασης \( n \) με μετρική \( g \), το στοιχείο όγκου που συζητήσαμε στην ενότητα 7.4 μπορεί τώρα να οριστεί αυστηρά ως μια \( n \)-μορφή. Σε έναν χάρτη με συντεταγμένες \( x^1, \dots, x^n \), το στοιχείο όγκου είναι η \( n \)-μορφή
\[
dV = \sqrt{|\det(g_{ij})|} \ dx^1 \wedge dx^2 \wedge \dots \wedge dx^n. \tag{43}
\]
Οι συνιστώσες του \( dV \) είναι ακριβώς ο τανυστής Levi-Civita που ορίσαμε στην (30):
\[
(dV)_{i_1 \dots i_n} = \varepsilon_{i_1 \dots i_n} = \sqrt{|\det(g_{ij})|} \ \epsilon_{i_1 \dots i_n}.
\]
Γεωμετρικά, το στοιχείο όγκου αποδίδει σε ένα \( n \)-διάστατο παραλληλεπίπεδο που ορίζεται από \( n \) διανύσματα \( \xi_1, \dots, \xi_n \) τον προσανατολισμένο όγκο του:
\[
dV(\xi_1, \dots, \xi_n) = \sqrt{|\det(g_{ij})|} \ \det(\xi^i_a),
\]
όπου \( \xi^i_a \) είναι η \( i \)-οστή συνιστώσα του \( a \)-οστού διανύσματος. Αυτή η έκφραση είναι ανεξάρτητη από την επιλογή χάρτη ακριβώς επειδή ο παράγοντας \( \sqrt{|\det(g_{ij})|} \) μετασχηματίζεται με την ορίζουσα του Ιακωβιανού πίνακα, αντισταθμίζοντας την αλλαγή της ορίζουσας των συνιστωσών των διανυσμάτων. Το στοιχείο όγκου είναι λοιπόν η θεμελιώδης \( n \)-μορφή που επιτρέπει την ολοκλήρωση βαθμωτών συναρτήσεων με τρόπο ανεξάρτητο συντεταγμένων.

\subsection{Ο τελεστής αστερισμού Hodge}

Σε μια προσανατολισμένη ψευδο-Ρημάννεια πολλαπλότητα διάστασης \( n \) με μετρική \( g \), ο \textbf{τελεστής Hodge} \( \star \) είναι μια απεικόνιση που στέλνει \( k \)-μορφές σε \( (n-k) \)-μορφές, χρησιμοποιώντας τη μετρική και το στοιχείο όγκου. Για μια \( k \)-μορφή \( \omega \), η δυϊκή της Hodge \( \star \omega \) είναι μια \( (n-k) \)-μορφή που ορίζεται από τη σχέση
\[
(\star \omega)_{i_1 \dots i_{n-k}} = \frac{1}{k!} \ \varepsilon_{i_1 \dots i_{n-k} j_1 \dots j_k} \ \omega^{j_1 \dots j_k}, \tag{44}
\]
όπου \( \varepsilon \) είναι ο τανυστής Levi-Civita της (30) και οι δείκτες της \( \omega \) έχουν ανυψωθεί με τη μετρική:
\[
\omega^{j_1 \dots j_k} = g^{j_1 \alpha_1} \dots g^{j_k \alpha_k} \ \omega_{\alpha_1 \dots \alpha_k}.
\]
Ο παράγοντας \( 1/k! \) είναι συμβατικός και εξασφαλίζει τη σωστή κανονικοποίηση.

Ο τελεστής Hodge μπορεί να εκφραστεί και απευθείας στη βάση των διαφορικών. Για τα στοιχεία βάσης \( dx^{i_1} \wedge \dots \wedge dx^{i_k} \), η δράση του \( \star \) δίνεται από
\[
\star (dx^{i_1} \wedge \dots \wedge dx^{i_k}) = \frac{1}{(n-k)!} \ \varepsilon^{i_1 \dots i_k}_{\phantom{i_1 \dots i_k} j_1 \dots j_{n-k}} \ dx^{j_1} \wedge \dots \wedge dx^{j_{n-k}}, \tag{45}
\]
όπου οι δείκτες του τανυστή Levi-Civita έχουν ανυψωθεί με τη μετρική:
\[
\varepsilon^{i_1 \dots i_k}_{\phantom{i_1 \dots i_k} j_1 \dots j_{n-k}} = g^{i_1 \alpha_1} \dots g^{i_k \alpha_k} \ \varepsilon_{\alpha_1 \dots \alpha_k j_1 \dots j_{n-k}}.
\]
Ισοδύναμα, μπορούμε να γράψουμε
\[
\star (dx^{i_1} \wedge \dots \wedge dx^{i_k}) = \frac{\sqrt{|\det(g_{ij})|}}{(n-k)!} \ g^{i_1 \alpha_1} \dots g^{i_k \alpha_k} \ \epsilon_{\alpha_1 \dots \alpha_k j_1 \dots j_{n-k}} \ dx^{j_1} \wedge \dots \wedge dx^{j_{n-k}}.
\]
Αυτή η έκφραση αναδεικνύει πώς ο τελεστής Hodge συνδυάζει τη μετρική (μέσω της ανύψωσης δεικτών και της ορίζουσας) με την αντισυμμετρική δομή των μορφών.

Η συσχέτιση του τελεστή Hodge με το στοιχείο όγκου είναι άμεση: για τη σταθερή βαθμωτή συνάρτηση \( 1 \) (μια 0-μορφή), η δυϊκή της Hodge είναι ακριβώς το στοιχείο όγκου:
\[
\star 1 = dV = \sqrt{|\det(g_{ij})|} \ dx^1 \wedge dx^2 \wedge \dots \wedge dx^n.
\]
Αντίστροφα, η δυϊκή Hodge του στοιχείου όγκου είναι
\[
\star (dx^1 \wedge \dots \wedge dx^n) = \frac{\operatorname{sgn}(\det(g_{ij}))}{\sqrt{|\det(g_{ij})|}}.
\]
Για μια γενική \( k \)-μορφή, η διπλή εφαρμογή του τελεστή Hodge δίνει
\[
\star \star \omega = (-1)^{k(n-k) + q} \ \omega, \tag{46}
\]
όπου \( q \) είναι το πλήθος των αρνητικών ιδιοτιμών της μετρικής (ο δείκτης της ψευδο-Ρημάννειας υπογραφής). Στην Ευκλείδεια περίπτωση (\( q=0 \)), ο παράγοντας είναι απλώς \( (-1)^{k(n-k)} \). Στην περίπτωση Lorentz (\( q=1 \)), εμφανίζεται ένα επιπλέον πρόσημο μείον.

Ο τελεστής Hodge επιτρέπει τη μετάφραση ανάμεσα σε βαθμωτά και \( n \)-μορφές, και γενικότερα ανάμεσα σε αντικείμενα συμπληρωματικής διάστασης, παίζοντας κεντρικό ρόλο στη διατύπωση των εξισώσεων Maxwell και στη θεωρία ολοκλήρωσης σε πολλαπλότητες.

\subsection{Ταυτότητες του σφηνοειδούς γινομένου}

Το σφηνοειδές γινόμενο ικανοποιεί μια σειρά από χρήσιμες ταυτότητες που απορρέουν από τον ορισμό του. Για 1-μορφές \( \omega, \eta \) και διανύσματα \( \xi, \zeta \), η αποτίμηση δίνει
\[
(\omega \wedge \eta)(\xi, \zeta) = \omega(\xi) \eta(\zeta) - \omega(\zeta) \eta(\xi). \tag{47}
\]
Γενικότερα, για μια \( k \)-μορφή \( \omega \) και μια \( l \)-μορφή \( \eta \), η αποτίμηση σε \( k+l \) διανύσματα είναι
\[
(\omega \wedge \eta)(\xi_1, \dots, \xi_{k+l}) = \frac{1}{k! \ l!} \sum_{\sigma \in S_{k+l}} \operatorname{sgn}(\sigma) \ \omega(\xi_{\sigma(1)}, \dots, \xi_{\sigma(k)}) \ \eta(\xi_{\sigma(k+1)}, \dots, \xi_{\sigma(k+l)}). \tag{48}
\]
Σε όρους συνιστωσών, το σφηνοειδές γινόμενο μιας \( k \)-μορφής \( \omega \) και μιας \( l \)-μορφής \( \eta \) δίνει
\[
(\omega \wedge \eta)_{i_1 \dots i_{k+l}} = \frac{(k+l)!}{k! \ l!} \ \omega_{[i_1 \dots i_k} \eta_{i_{k+1} \dots i_{k+l}]}. \tag{49}
\]
Η προσεταιριστικότητα εκφράζεται ως
\[
(\omega \wedge \eta) \wedge \theta = \omega \wedge (\eta \wedge \theta), \tag{50}
\]
και η βαθμωτή αντιμεταθετικότητα ως
\[
\omega \wedge \eta = (-1)^{kl} \ \eta \wedge \omega. \tag{51}
\]

\section{Σύνοψη: Οι Θεμελιώδεις Τύποι και Ταυτότητες των Μορφών}

\textbf{Αντισυμμετρία μορφής}

\[
\omega_{i_{\sigma(1)} \dots i_{\sigma(k)}} = \operatorname{sgn}(\sigma) \ \omega_{i_1 \dots i_k}. \tag{37}
\]

Μια \( k \)-μορφή είναι ένας πλήρως αντισυμμετρικός τανυστής τύπου \( (0,k) \). Η αντισυμμετρία συνεπάγεται ότι η μορφή μηδενίζεται αν δύο δείκτες είναι ίσοι.

\textbf{Σφηνοειδές γινόμενο 1-μορφών βάσης}

\[
dx^i \wedge dx^j = - dx^j \wedge dx^i, \qquad dx^i \wedge dx^i = 0. \tag{40}
\]

Το σφηνοειδές γινόμενο είναι αντισυμμετρικό και μηδενίζεται για ίσα ορίσματα.

\textbf{Ορισμός σφηνοειδούς γινομένου σε συνιστώσες}

\[
(\omega \wedge \eta)_{i_1 \dots i_{k+l}} = \frac{(k+l)!}{k! \ l!} \ \omega_{[i_1 \dots i_k} \eta_{i_{k+1} \dots i_{k+l}]}. \tag{49}
\]

Το σφηνοειδές γινόμενο μιας \( k \)-μορφής και μιας \( l \)-μορφής είναι μια \( (k+l) \)-μορφή που προκύπτει από την πλήρη αντισυμμετρικότητα του τανυστικού γινομένου.

\textbf{Προσεταιριστικότητα και βαθμωτή αντιμεταθετικότητα}

\[
(\omega \wedge \eta) \wedge \theta = \omega \wedge (\eta \wedge \theta), \tag{50}
\]
\[
\omega \wedge \eta = (-1)^{kl} \ \eta \wedge \omega. \tag{51}
\]

Το σφηνοειδές γινόμενο είναι προσεταιριστικό και αντιμεταθετικό κατά βαθμό, με το πρόσημο να καθορίζεται από το γινόμενο των βαθμών.

\textbf{Αποτίμηση σφηνοειδούς γινομένου}

\[
(\omega \wedge \eta)(\xi, \zeta) = \omega(\xi) \eta(\zeta) - \omega(\zeta) \eta(\xi). \tag{47}
\]

Για 1-μορφές, η αποτίμηση του σφηνοειδούς γινομένου σε δύο διανύσματα δίνει την ορίζουσα των αποτιμήσεων, που γεωμετρικά αντιστοιχεί στο προσανατολισμένο εμβαδόν.

\textbf{Στοιχείο όγκου}

\[
dV = \sqrt{|\det(g_{ij})|} \ dx^1 \wedge dx^2 \wedge \dots \wedge dx^n. \tag{43}
\]

Το στοιχείο όγκου είναι η \( n \)-μορφή που σε κάθε σημείο αποδίδει τον προσανατολισμένο όγκο ενός \( n \)-διάστατου παραλληλεπιπέδου. Ο παράγοντας \( \sqrt{|\det(g_{ij})|} \) εξασφαλίζει την ανεξαρτησία από την επιλογή συντεταγμένων.

\textbf{Τελεστής Hodge σε συνιστώσες}

\[
(\star \omega)_{i_1 \dots i_{n-k}} = \frac{1}{k!} \ \varepsilon_{i_1 \dots i_{n-k} j_1 \dots j_k} \ \omega^{j_1 \dots j_k}. \tag{44}
\]

\textbf{Τελεστής Hodge στη βάση των διαφορικών}

\[
\star (dx^{i_1} \wedge \dots \wedge dx^{i_k}) = \frac{1}{(n-k)!} \ \varepsilon^{i_1 \dots i_k}_{\phantom{i_1 \dots i_k} j_1 \dots j_{n-k}} \ dx^{j_1} \wedge \dots \wedge dx^{j_{n-k}}. \tag{45}
\]

Ο τελεστής Hodge απεικονίζει \( k \)-μορφές σε \( (n-k) \)-μορφές χρησιμοποιώντας τη μετρική και τον τανυστή Levi-Civita. Στη βάση των διαφορικών, η δράση του εκφράζεται με ανύψωση των δεικτών του τανυστή Levi-Civita και σφηνοειδές γινόμενο των διαφορικών. Η διπλή εφαρμογή του δίνει
\[
\star \star \omega = (-1)^{k(n-k) + q} \ \omega, \tag{46}
\]
όπου \( q \) είναι το πλήθος των αρνητικών ιδιοτιμών της μετρικής. Στην Ευκλείδεια περίπτωση (\( q=0 \)), \( \star \star \omega = (-1)^{k(n-k)} \omega \). Στην Lorentz (\( q=1 \)), εμφανίζεται ένα επιπλέον πρόσημο μείον.

\section{H Παράγωγος Lie}

\subsection{Γεωμετρικό κίνητρο: ροές και παρατηρητές}

Στη γεωμετρία και τη φυσική, συχνά θέλουμε να μετρήσουμε πώς μεταβάλλεται ένα τανυστικό πεδίο όχι ως προς μια αυθαίρετη παράμετρο, αλλά κατά μήκος της φυσικής κίνησης που ορίζει ένα διανυσματικό πεδίο. Φανταστείτε έναν παρατηρητή που κινείται μέσα σε μια πολλαπλότητα, παρασυρόμενος από τη ροή ενός ρευστού. Ο παρατηρητής αυτός μεταφέρει μαζί του όργανα μέτρησης και καταγράφει τις τιμές ενός τανυστικού πεδίου $T$ κατά μήκος της τροχιάς του. Η παράγωγος Lie είναι ακριβώς ο ρυθμός μεταβολής που μετρά αυτός ο συμπαρασυρόμενος παρατηρητής: συγκρίνει την τιμή του $T$ σε ένα σημείο με την τιμή του σε ένα γειτονικό σημείο κατά μήκος της ροής, αφού πρώτα "μεταφέρει" την τιμή από το γειτονικό σημείο πίσω στο αρχικό μέσω της ροής. Αυτή η μεταφορά δεν χρησιμοποιεί καμία εξωτερική δομή όπως μια μετρική ή μια συνοχή· χρησιμοποιεί αποκλειστικά και μόνο τη ροή που ορίζει το ίδιο το διανυσματικό πεδίο.

\subsection{Ο απειροστικός ορισμός μέσω ορίου}

Έστω $\xi$ ένα λείο διανυσματικό πεδίο σε μια πολλαπλότητα $M$. Η ροή του $\xi$ είναι μια οικογένεια απεικονίσεων $\phi_t: M \to M$, ορισμένη τοπικά γύρω από το $t=0$, που ικανοποιεί
\[
\frac{d}{dt} \phi_t(p) = \xi(\phi_t(p)), \qquad \phi_0(p) = p.
\]
Η ροή $\phi_t$ μεταφέρει σημεία κατά μήκος των ολοκληρωτικών καμπυλών του $\xi$. Για ένα τανυστικό πεδίο $T$ τύπου $(r, s)$, η \textbf{παράγωγος Lie} του $T$ ως προς $\xi$ ορίζεται ως το όριο
\begin{equation}
(\mathcal{L}_\xi T)(p) = \lim_{t \to 0} \frac{(\phi_{-t})_* T(\phi_t(p)) - T(p)}{t}, \qquad (52)
\end{equation}
όπου $(\phi_{-t})_*$ είναι η εφελκυσμένη (pushforward) ή συνεφελκυσμένη (pullback) δράση της ροής, ανάλογα με τον τύπο του τανυστή. Για μια βαθμωτή συνάρτηση $f$, η εφελκυσμένη είναι απλώς η σύνθεση: $(\phi_{-t})_* f = f \circ \phi_t$. Για ένα διάνυσμα, η εφελκυσμένη είναι το διαφορικό της απεικόνισης. Για μια 1-μορφή, είναι η συνεφελκυσμένη. Η ουσία του ορισμού είναι ότι ο $\phi_t$ μεταφέρει το $T$ από το σημείο $\phi_t(p)$ πίσω στο $p$, και εμείς συγκρίνουμε με την τιμή του $T$ στο $p$. Το όριο αυτό υπάρχει λόγω της ομαλότητας της ροής και του τανυστικού πεδίου.

\subsection{Συσχέτιση με την υλική παράγωγο της ρευστομηχανικής}

Η παράγωγος Lie γενικεύει και αποσαφηνίζει μια έννοια που εμφανίζεται φυσικά στη ρευστομηχανική: την \textbf{υλική} ή \textbf{ολική παράγωγο}. Στη ρευστομηχανική, αν $\xi$ είναι το πεδίο ταχύτητας ενός ρευστού και $f$ είναι μια βαθμωτή ποσότητα (όπως η θερμοκρασία ή η πυκνότητα), η υλική παράγωγος της $f$ ορίζεται ως
\[
\frac{Df}{Dt} = \frac{\partial f}{\partial t} + \xi^\alpha \frac{\partial f}{\partial x^\alpha},
\]
όπου ο πρώτος όρος είναι η τοπική χρονική μεταβολή και ο δεύτερος η μεταφορά λόγω της κίνησης του ρευστού. Αυτή η έκφραση δεν είναι παρά η παράγωγος Lie της $f$ ως προς το διανυσματικό πεδίο $\xi$, με την προσθήκη ενός όρου μερικής παραγώγου ως προς τον χρόνο αν το πεδίο εξαρτάται ρητά από αυτόν. Σε μια γεωμετρική διατύπωση όπου ο χρόνος ενσωματώνεται ως μία από τις συντεταγμένες, η υλική παράγωγος ταυτίζεται πλήρως με την $\mathcal{L}_\xi f$.

Η σύνδεση γίνεται βαθύτερη όταν εξετάζουμε τανυστικά μεγέθη. Στη ρευστομηχανική, ο ρυθμός παραμόρφωσης ενός ρευστού περιγράφεται από τον τανυστή ρυθμού παραμόρφωσης, ο οποίος είναι ακριβώς το ήμισυ της παραγώγου Lie της μετρικής κατά μήκος του πεδίου ταχύτητας: $\frac{1}{2} \mathcal{L}_\xi g$. Η υλική παράγωγος ενός διανυσματικού πεδίου $\eta$ που είναι "παγωμένο" στο ρευστό (όπως μια γραμμή δίνης) δίνεται από την παράγωγο Lie $\mathcal{L}_\xi \eta$. Ο μηδενισμός της $\mathcal{L}_\xi \eta$ σημαίνει ότι το $\eta$ μεταφέρεται αναλλοίωτο από τη ροή, δηλαδή ότι οι ολοκληρωτικές καμπύλες του $\eta$ παρασύρονται από το ρευστό χωρίς να αλλάζουν. Αυτή η εικόνα είναι ακριβώς ο γεωμετρικός πυρήνας της έννοιας του "παγωμένου πεδίου" στη μαγνητοϋδροδυναμική.

Έτσι, η υλική παράγωγος της ρευστομηχανικής δεν είναι παρά η παράγωγος Lie εξειδικευμένη στην περίπτωση όπου η πολλαπλότητα είναι ο χώρος (ή ο χωροχρόνος) και το διανυσματικό πεδίο είναι το πεδίο ταχύτητας. Η παράγωγος Lie παρέχει το γεωμετρικό πλαίσιο που ενοποιεί τη μεταφορά βαθμωτών, διανυσμάτων και τανυστών κάτω από μια ροή, ανεξάρτητα από την ύπαρξη μετρικής ή άλλων δομών.

\subsection{Έκφραση σε τοπικές συντεταγμένες}

Για να εξαγάγουμε τον τύπο της παραγώγου Lie σε τοπικές συντεταγμένες, εργαζόμαστε απειροστικά. Σε έναν χάρτη με συντεταγμένες $x^i$, η ροή $\phi_t$ δίνεται προσεγγιστικά από
\[
\phi_t^i(p) = x^i + t \xi^i + \mathcal{O}(t^2).
\]
Το διαφορικό της $\phi_{-t}$ (που χρησιμοποιείται για την εφελκυσμένη διανυσμάτων) είναι
\[
\frac{\partial \phi_{-t}^i}{\partial x^j} = \delta^i_j - t \frac{\partial \xi^i}{\partial x^j} + \mathcal{O}(t^2),
\]
ενώ το διαφορικό της $\phi_t$ (που χρησιμοποιείται για τη συνεφελκυσμένη 1-μορφών) είναι
\[
\frac{\partial \phi_t^i}{\partial x^j} = \delta^i_j + t \frac{\partial \xi^i}{\partial x^j} + \mathcal{O}(t^2).
\]
Εφαρμόζοντας τον ορισμό (52) για διάφορους τύπους τανυστών, προκύπτουν οι ακόλουθοι τύποι σε τοπικές συντεταγμένες.

\textbf{Βαθμωτή συνάρτηση $f$:}
\begin{equation}
\mathcal{L}_\xi f = \xi^\alpha \frac{\partial f}{\partial x^\alpha}. \qquad (53)
\end{equation}
Πρόκειται απλώς για την κατευθυνόμενη παράγωγο της $f$ κατά μήκος του $\xi$. Γεωμετρικά, ο παρατηρητής μετρά τον ρυθμό μεταβολής της βαθμωτής ποσότητας καθώς παρασύρεται από τη ροή.

\textbf{Διανυσματικό πεδίο $\eta$:}
\begin{equation}
(\mathcal{L}_\xi \eta)^i = \xi^\alpha \frac{\partial \eta^i}{\partial x^\alpha} - \eta^\alpha \frac{\partial \xi^i}{\partial x^\alpha}. \qquad (54)
\end{equation}
Ο δεύτερος όρος προκύπτει από την εφελκυσμένη της ροής και αντιπροσωπεύει την παραμόρφωση του συστήματος συντεταγμένων του παρατηρητή. Η παράγωγος Lie ενός διανυσματικού πεδίου είναι ταυτόσημη με τον μεταθέτη (αγκύλη Lie) των δύο διανυσματικών πεδίων:
\[
\mathcal{L}_\xi \eta = [\xi, \eta] \equiv \xi \eta - \eta \xi.
\]

\textbf{1-μορφή $\omega$:}
\begin{equation}
(\mathcal{L}_\xi \omega)_i = \xi^\alpha \frac{\partial \omega_i}{\partial x^\alpha} + \omega_\alpha \frac{\partial \xi^\alpha}{\partial x^i}. \qquad (55)
\end{equation}
Εδώ ο δεύτερος όρος έχει πρόσημο $+$ αντί για $-$, διότι οι 1-μορφές μετασχηματίζονται με τον ανάστροφο του Ιακωβιανού πίνακα της ροής.

\textbf{Τανυστής τύπου $(r, s)$:}\\
Για ένα γενικό τανυστικό πεδίο $T$ τύπου $(r, s)$, η παράγωγος Lie δίνεται από
\begin{equation}
\begin{aligned}
(\mathcal{L}_\xi T)^{i_1 \dots i_r}_{j_1 \dots j_s} &= \xi^\alpha \frac{\partial}{\partial x^\alpha} T^{i_1 \dots i_r}_{j_1 \dots j_s} \\
&\quad - T^{\alpha i_2 \dots i_r}_{j_1 \dots j_s} \frac{\partial \xi^{i_1}}{\partial x^\alpha} - \dots - T^{i_1 \dots i_{r-1} \alpha}_{j_1 \dots j_s} \frac{\partial \xi^{i_r}}{\partial x^\alpha} \\
&\quad + T^{i_1 \dots i_r}_{\alpha j_2 \dots j_s} \frac{\partial \xi^\alpha}{\partial x^{j_1}} + \dots + T^{i_1 \dots i_r}_{j_1 \dots j_{s-1} \alpha} \frac{\partial \xi^\alpha}{\partial x^{j_s}}. \qquad (56)
\end{aligned}
\end{equation}
Ο πρώτος όρος είναι η κατευθυνόμενη παράγωγος κατά μήκος του $\xi$. Οι επόμενοι $r$ όροι (με αρνητικό πρόσημο) αντιστοιχούν στη διόρθωση για κάθε ανταλλοίωτο δείκτη, ενώ οι τελευταίοι $s$ όροι (με θετικό πρόσημο) αντιστοιχούν στη διόρθωση για κάθε συναλλοίωτο δείκτη. Η δομή αυτή είναι απολύτως γενική και αναδεικνύει την ομοιομορφία της παραγώγου Lie: κάθε δείκτης συνεισφέρει έναν όρο διόρθωσης ανάλογα με τη φύση του (ανταλλοίωτος ή συναλλοίωτος).

\subsection{Εξειδίκευση σε συγκεκριμένους τύπους τανυστών}

Εφαρμόζοντας τον γενικό τύπο (56), λαμβάνουμε τις εκφράσεις για ενδιαφέροντες ειδικούς τύπους τανυστών.

\textbf{Τετραγωνική μορφή (μετρική) $g$:}\\
Για έναν τανυστή τύπου $(0,2)$, η παράγωγος Lie είναι
\begin{equation}
(\mathcal{L}_\xi g)_{ij} = \xi^\alpha \frac{\partial g_{ij}}{\partial x^\alpha} + g_{\alpha j} \frac{\partial \xi^\alpha}{\partial x^i} + g_{i\alpha} \frac{\partial \xi^\alpha}{\partial x^j}. \qquad (57)
\end{equation}
Αυτή η έκφραση είναι θεμελιώδης στη γεωμετρία: ο μηδενισμός της $\mathcal{L}_\xi g$ σημαίνει ότι το διανυσματικό πεδίο $\xi$ είναι μια ισομετρία, δηλαδή η ροή του διατηρεί τη μετρική. Στη γενική σχετικότητα, τα διανύσματα Killing που ικανοποιούν $\mathcal{L}_\xi g = 0$ αντιστοιχούν σε συμμετρίες του χωροχρόνου.

\textbf{Συμπλεκτική μορφή $\omega$:}\\
Για μια 2-μορφή, η παράγωγος Lie είναι
\begin{equation}
(\mathcal{L}_\xi \omega)_{ij} = \xi^\alpha \frac{\partial \omega_{ij}}{\partial x^\alpha} + \omega_{\alpha j} \frac{\partial \xi^\alpha}{\partial x^i} + \omega_{i\alpha} \frac{\partial \xi^\alpha}{\partial x^j}. \qquad (58)
\end{equation}
Ο μηδενισμός της $\mathcal{L}_\xi \omega$ χαρακτηρίζει τα συμπλεκτικά διανυσματικά πεδία, που διατηρούν τη συμπλεκτική δομή. Στην κλασική μηχανική, οι Χαμιλτονιανές ροές είναι ακριβώς εκείνα τα διανυσματικά πεδία των οποίων η παράγωγος Lie της συμπλεκτικής μορφής μηδενίζεται.

\textbf{Τανυστής τάσεων $\sigma$:}\\
Για έναν τανυστή τύπου $(2,0)$, η παράγωγος Lie είναι
\begin{equation}
(\mathcal{L}_\xi \sigma)^{ij} = \xi^\alpha \frac{\partial \sigma^{ij}}{\partial x^\alpha} - \sigma^{\alpha j} \frac{\partial \xi^i}{\partial x^\alpha} - \sigma^{i\alpha} \frac{\partial \xi^j}{\partial x^\alpha}. \qquad (59)
\end{equation}
Στη μηχανική των συνεχών μέσων, ο μηδενισμός της $\mathcal{L}_\xi \sigma$ σχετίζεται με την ισορροπία του υλικού κάτω από τη ροή που ορίζει το $\xi$.

\textbf{Τανυστής ηλεκτρομαγνητικού πεδίου $F$:}\\
Για τον αντισυμμετρικό τανυστή $F$ τύπου $(0,2)$, η παράγωγος Lie είναι
\begin{equation}
(\mathcal{L}_\xi F)_{ij} = \xi^\alpha \frac{\partial F_{ij}}{\partial x^\alpha} + F_{\alpha j} \frac{\partial \xi^\alpha}{\partial x^i} + F_{i\alpha} \frac{\partial \xi^\alpha}{\partial x^j}. \qquad (60)
\end{equation}
Στην ηλεκτροδυναμική, ο μηδενισμός της $\mathcal{L}_\xi F$ σημαίνει ότι το ηλεκτρομαγνητικό πεδίο είναι αναλλοίωτο κάτω από τη ροή του $\xi$, δηλαδή ότι ο παρατηρητής που κινείται κατά μήκος του $\xi$ δεν αντιλαμβάνεται μεταβολή του πεδίου.

\subsection{Φυσικογεωμετρική ερμηνεία του μηδενισμού}

Ο μηδενισμός της παραγώγου Lie, $\mathcal{L}_\xi T = 0$, έχει μια βαθιά φυσικογεωμετρική σημασία: σημαίνει ότι το τανυστικό πεδίο $T$ είναι \textbf{αναλλοίωτο κάτω από τη ροή} που ορίζει το διανυσματικό πεδίο $\xi$. Με άλλα λόγια, αν αφήσουμε έναν παρατηρητή να παρασυρθεί από τη ροή του $\xi$, αυτός δεν θα αντιληφθεί καμία αλλαγή στο $T$. Η ροή $\phi_t$ είναι μια \textbf{συμμετρία} του $T$.

Αυτή η έννοια είναι θεμελιώδης σε όλη τη φυσική. Στη γενική σχετικότητα, τα διανύσματα Killing ($\mathcal{L}_\xi g = 0$) γεννούν διατηρούμενα μεγέθη κατά μήκος γεωδαισιακών. Στην κλασική μηχανική, ένα Χαμιλτονιανό διανυσματικό πεδίο ικανοποιεί $\mathcal{L}_{\xi_H} \omega = 0$, που είναι η γεωμετρική έκφραση του θεωρήματος Liouville για τη διατήρηση του φασικού όγκου. Στην ηλεκτροδυναμική, η απαίτηση $\mathcal{L}_\xi F = 0$ για κατάλληλο $\xi$ οδηγεί σε νόμους διατήρησης για την ενέργεια και την ορμή του πεδίου. Στη ρευστομηχανική, η συνθήκη $\mathcal{L}_\xi g = 0$ για το πεδίο ταχύτητας χαρακτηρίζει τις ισόογκες ροές. Σε όλες αυτές τις περιπτώσεις, η παράγωγος Lie είναι το μαθηματικό εργαλείο που μεταφράζει τη γεωμετρική έννοια της συμμετρίας σε αλγεβρικές συνθήκες επιλύσιμες σε συντεταγμένες.

Η παράγωγος Lie εξαρτάται αποκλειστικά από τη διαφορική δομή της πολλαπλότητας και το διανυσματικό πεδίο $\xi$. Δεν απαιτεί μετρική ή συνοχή, και αυτό την καθιστά το πιο θεμελιώδες είδος παραγώγου σε μια πολλαπλότητα. Είναι το φυσικό εργαλείο για τη μελέτη των συμμετριών και των ροών, και η άμεση σύνδεσή της με τον μεταθέτη διανυσματικών πεδίων την τοποθετεί στον πυρήνα της γεωμετρικής ανάλυσης.

\subsection{Πεδία Killing}

Ένα διανυσματικό πεδίο $\xi$ που ικανοποιεί
\begin{equation}
\mathcal{L}_\xi g = 0 \qquad (61)
\end{equation}
ονομάζεται \textbf{πεδίο Killing}. Η συνθήκη αυτή, γραμμένη σε συνιστώσες μέσω της (57), είναι η \textbf{εξίσωση Killing}:
\begin{equation}
\xi^\alpha \frac{\partial g_{ij}}{\partial x^\alpha} + g_{\alpha j} \frac{\partial \xi^\alpha}{\partial x^i} + g_{i\alpha} \frac{\partial \xi^\alpha}{\partial x^j} = 0. \qquad (62)
\end{equation}

Η φυσικογεωμετρική σημασία ενός πεδίου Killing είναι ότι η ροή του αποτελεί μια συνεχή συμμετρία της μετρικής. Με άλλα λόγια, η μετρική παραμένει αναλλοίωτη καθώς μετακινούμαστε κατά μήκος των ολοκληρωτικών καμπυλών του $\xi$. Αυτό σημαίνει ότι οι αποστάσεις, οι γωνίες και οι όγκοι που μετρά η μετρική δεν αλλάζουν κάτω από τη ροή. Στη γενική σχετικότητα, κάθε πεδίο Killing αντιστοιχεί σε έναν νόμο διατήρησης: αν ένα σωματίδιο κινείται κατά μήκος μιας γεωδαισιακής, η ποσότητα $g(\xi, \dot{\gamma})$ παραμένει σταθερή κατά μήκος της τροχιάς του, όπου $\dot{\gamma}$ είναι το εφαπτόμενο διάνυσμα της γεωδαισιακής. Αυτό γενικεύει τη διατήρηση της ορμής και της ενέργειας σε καμπύλους χωροχρόνους: τα πεδία Killing που σχετίζονται με χρονικές και χωρικές συμμετρίες οδηγούν αντίστοιχα σε διατήρηση ενέργειας και ορμής.

Σε μια ψευδο-Ρημάννεια πολλαπλότητα, το πλήθος των γραμμικά ανεξάρτητων πεδίων Killing είναι πεπερασμένο και καθορίζεται από τη γεωμετρία. Ο χωρόχρονος Minkowski, ως μεγιστικά συμμετρικός, διαθέτει δέκα πεδία Killing (τέσσερις μεταφορές, τρεις στροφές, τρία boosts). Η ύπαρξη ή η απουσία πεδίων Killing είναι ένας θεμελιώδης τρόπος ταξινόμησης των γεωμετριών: όσο περισσότερα πεδία Killing διαθέτει μια μετρική, τόσο πιο συμμετρική είναι.

\subsection{Συμμετρίες, νόμοι διατήρησης και η παράγωγος Lie}

Η παράγωγος Lie δεν είναι απλώς ένα εργαλείο για τη μέτρηση της μεταβολής τανυστικών πεδίων κατά μήκος μιας ροής· είναι το μαθηματικό θεμέλιο πάνω στο οποίο οικοδομείται η βαθύτερη σχέση ανάμεσα στις συμμετρίες και τους νόμους διατήρησης σε όλη τη φυσική. Η κεντρική ιδέα είναι απλή και κομψή: αν ένα τανυστικό πεδίο $T$ που κωδικοποιεί τη δυναμική ενός φυσικού συστήματος είναι αναλλοίωτο κάτω από μια ροή $\xi$, δηλαδή $\mathcal{L}_\xi T = 0$, τότε υπάρχει μια διατηρούμενη ποσότητα που συνδέεται με αυτή τη συμμετρία. Αυτή η αρχή, που βρίσκει την πληρέστερη έκφρασή της στο θεώρημα Noether, διατρέχει ολόκληρη τη φυσική από τη Νευτώνεια μηχανική μέχρι τη γενική σχετικότητα.

\subsubsection{Νόμοι του Νεύτωνα και γαλιλαϊκοί μετασχηματισμοί}

Στη Νευτώνεια μηχανική, ο χώρος και ο χρόνος είναι απόλυτοι και ανεξάρτητοι. Η μετρική του χώρου είναι η Ευκλείδεια $g_{ij} = \delta_{ij}$, και ο χρόνος είναι μια εξωτερική παράμετρος. Οι γαλιλαϊκοί μετασχηματισμοί ανάμεσα σε αδρανειακά συστήματα αναφοράς,
\[
x^i \to x^i + v^i t, \qquad t \to t,
\]
αφήνουν αναλλοίωτη την Ευκλείδεια μετρική. Το διανυσματικό πεδίο που γεννά μια γαλιλαϊκή ώθηση (boost) στην κατεύθυνση $i$ είναι το $\xi = t \partial / \partial x^i$. Η παράγωγος Lie της Ευκλείδειας μετρικής ως προς αυτό το πεδίο μηδενίζεται,
\[
(\mathcal{L}_\xi g)_{ij} = \xi^\alpha \frac{\partial g_{ij}}{\partial x^\alpha} + g_{\alpha j} \frac{\partial \xi^\alpha}{\partial x^i} + g_{i\alpha} \frac{\partial \xi^\alpha}{\partial x^j}.
\]
Εφόσον η Ευκλείδεια μετρική είναι σταθερή, ο πρώτος όρος μηδενίζεται. Ο δεύτερος και ο τρίτος όρος περιέχουν παραγώγους του $\xi$. Για το πεδίο της μεταφοράς $\xi = \partial / \partial x^i$, όλες οι παράγωγοι μηδενίζονται, άρα $\mathcal{L}_\xi g = 0$. Αυτό επιβεβαιώνει ότι οι χωρικές μεταφορές είναι ισομετρίες του Ευκλείδειου χώρου. Η διατηρούμενη ποσότητα που αντιστοιχεί σε αυτή τη συμμετρία είναι η ορμή. Πράγματι, για ένα σύστημα σωματιδίων με Λαγκρανζιανή
\[
L = \sum_a \frac{1}{2} m_a \delta_{ij} \dot{x}^i_a \dot{x}^j_a - V(x_a),
\]
η συνθήκη $\mathcal{L}_\xi L = 0$ για $\xi = \partial / \partial x^i$ σημαίνει ότι το δυναμικό $V$ είναι αναλλοίωτο σε χωρικές μεταφορές, δηλαδή εξαρτάται μόνο από τις σχετικές αποστάσεις των σωματιδίων. Τότε η ολική ορμή $P^i = \sum_a m_a \dot{x}^i_a$ διατηρείται.

\subsubsection{Νόμοι διατήρησης στη ρευστομηχανική}

Στη ρευστομηχανική, η παράγωγος Lie παρέχει τη φυσική γλώσσα για τη διατύπωση των νόμων διατήρησης. Η εξίσωση συνέχειας εκφράζει τη διατήρηση της μάζας. Σε έναν χώρο με Ευκλείδεια μετρική, η πυκνότητα μάζας $\rho$ και το πεδίο ταχύτητας $\xi$ ικανοποιούν
\[
\frac{\partial \rho}{\partial t} + \frac{\partial}{\partial x^\alpha} (\rho \xi^\alpha) = 0.
\]
Σε γεωμετρική γλώσσα, αυτή η εξίσωση γράφεται ως
\[
\frac{\partial \rho}{\partial t} + \mathcal{L}_\xi \rho + \rho \ \frac{\partial \xi^\alpha}{\partial x^\alpha} = 0.
\]
Ο όρος $\mathcal{L}_\xi \rho = \xi^\alpha \partial \rho / \partial x^\alpha$ είναι η υλική μεταφορά της πυκνότητας, ενώ ο όρος $\rho \ \partial \xi^\alpha / \partial x^\alpha$ εκφράζει τη διαστολή του ρευστού. Για ασυμπίεστο ρευστό, $\partial \xi^\alpha / \partial x^\alpha = 0$, και η εξίσωση ανάγεται στην $\partial \rho / \partial t + \mathcal{L}_\xi \rho = 0$. Σε μόνιμη κατάσταση, $\mathcal{L}_\xi \rho = 0$, που σημαίνει ότι η πυκνότητα είναι σταθερή κατά μήκος των τροχιών του ρευστού.

Η εξίσωση ορμής Euler για ιδανικό ρευστό σε γεωμετρική μορφή είναι
\[
\frac{\partial \xi^\flat}{\partial t} + \mathcal{L}_\xi \xi^\flat = -\frac{1}{\rho} dp,
\]
όπου $\xi^\flat$ είναι η 1-μορφή που αντιστοιχεί στο πεδίο ταχύτητας μέσω της Ευκλείδειας μετρικής, και $p$ είναι η πίεση. Ο όρος $\mathcal{L}_\xi \xi^\flat$ αναπτύσσεται, σύμφωνα με την (55), ως
\[
(\mathcal{L}_\xi \xi^\flat)_i = \xi^\alpha \frac{\partial \xi_i}{\partial x^\alpha} + \xi_\alpha \frac{\partial \xi^\alpha}{\partial x^i},
\]
όπου $\xi_i = \delta_{i\alpha} \xi^\alpha$. Ο δεύτερος όρος είναι η διόρθωση που προκύπτει από τη συναλλοίωτη φύση της 1-μορφής. Σε Καρτεσιανές συντεταγμένες, ο όρος αυτός ταυτίζεται με τον όρο μεταφοράς της ορμής στη συνήθη εξίσωση Euler. Η διατήρηση της ορμής σε ένα ιδανικό ρευστό συνδέεται με τη χωρική ομοιογένεια της Λαγκρανζιανής περιγραφής: η Λαγκρανζιανή πυκνότητα του ρευστού είναι αναλλοίωτη σε χωρικές μεταφορές, και το θεώρημα Noether δίνει τη διατήρηση της ορμής.

\subsubsection{Εξίσωση συνέχειας στην ηλεκτροδυναμική}

Στην ηλεκτροδυναμική, η διατήρηση του ηλεκτρικού φορτίου εκφράζεται από την εξίσωση συνέχειας. Σε τανυστική γλώσσα, το τετραδιάνυσμα του ρεύματος $J$ με συνιστώσες $J^i$ ικανοποιεί
\[
\frac{1}{\sqrt{|\det(g_{ij})|}} \frac{\partial}{\partial x^\alpha} \left( \sqrt{|\det(g_{ij})|} \ J^\alpha \right) = 0.
\]
Στον χωροχρόνο Minkowski, όπου $g_{ij} = \text{diag}(-1, 1, 1, 1)$ και $\sqrt{|\det(g_{ij})|} = 1$, αυτή ανάγεται στη συνήθη μορφή
\[
\frac{\partial J^\alpha}{\partial x^\alpha} = 0.
\]
Η γεωμετρική ερμηνεία αυτής της εξίσωσης μέσω της παραγώγου Lie προκύπτει ως εξής: αν ο χωροχρόνος διαθέτει ένα χρονικό πεδίο Killing $\xi = \partial / \partial x^0$, τότε η παράγωγος Lie του στοιχείου όγκου $dV$ ως προς $\xi$ ικανοποιεί $\mathcal{L}_\xi (dV) = 0$, και το ολοκλήρωμα της χρονικής συνιστώσας του $J$ σε μια χωρική υπερεπιφάνεια διατηρείται. Η ίδια η εξίσωση συνέχειας προκύπτει από τη συμμετρία βαθμίδας του ηλεκτρομαγνητικού πεδίου: η Λαγκρανζιανή πυκνότητα είναι αναλλοίωτη κάτω από μετασχηματισμούς βαθμίδας, και το θεώρημα Noether δίνει τη διατήρηση του ρεύματος.

\subsubsection{Ειδική σχετικότητα και μετασχηματισμοί Lorentz}

Στην ειδική σχετικότητα, ο χωροχρόνος είναι εφοδιασμένος με τη μετρική Minkowski $g_{ij} = \text{diag}(-1, 1, 1, 1)$. Οι μετασχηματισμοί Lorentz είναι ακριβώς εκείνοι που αφήνουν αναλλοίωτη αυτή τη μετρική, δηλαδή είναι ισομετρίες του χωροχρόνου Minkowski. Τα αντίστοιχα πεδία Killing είναι δέκα, και η εξίσωση Killing (62) είναι η συνθήκη που τα χαρακτηρίζει.

Για μια χρονική μεταφορά, το πεδίο Killing είναι $\xi = \partial / \partial x^0$. Αντικαθιστώντας στην (62), έχουμε
\[
(\mathcal{L}_\xi g)_{ij} = \xi^\alpha \frac{\partial g_{ij}}{\partial x^\alpha} + g_{\alpha j} \frac{\partial \xi^\alpha}{\partial x^i} + g_{i\alpha} \frac{\partial \xi^\alpha}{\partial x^j}.
\]
Εφόσον η μετρική Minkowski είναι σταθερή, ο πρώτος όρος μηδενίζεται. Οι παράγωγοι του $\xi$ είναι όλες μηδέν, άρα $\mathcal{L}_\xi g = 0$. Η διατηρούμενη ποσότητα που αντιστοιχεί σε αυτό το πεδίο Killing είναι η ενέργεια. Για ένα ελεύθερο σωματίδιο με τετραταχύτητα $u$, η ποσότητα $g(\xi, u) = g_{00} u^0 = -u^0$ διατηρείται, και η ενέργεια είναι $E = m u^0$.

Για μια χωρική μεταφορά, $\xi = \partial / \partial x^i$, η ίδια λογική δίνει $\mathcal{L}_\xi g = 0$ και η διατηρούμενη ποσότητα είναι η ορμή $P^i = m u^i$. Η γεωμετρική προέλευση της διατήρησης της ορμής είναι λοιπόν η χωρική ομοιογένεια του χωροχρόνου Minkowski.

\subsubsection{Διατήρηση ενέργειας και ορμής στη γενική σχετικότητα}

Στη γενική σχετικότητα, ο χωροχρόνος είναι μια καμπύλη ψευδο-Ρημάννεια πολλαπλότητα. Σε αντίθεση με την ειδική σχετικότητα, δεν υπάρχουν εκ των προτέρων πεδία Killing· η ύπαρξή τους εξαρτάται από τη συγκεκριμένη γεωμετρία. Ωστόσο, όταν υπάρχει ένα πεδίο Killing $\xi$, η εξίσωση Killing $\mathcal{L}_\xi g = 0$ εξασφαλίζει ότι η ποσότητα $g(\xi, u)$ διατηρείται κατά μήκος κάθε γεωδαισιακής με εφαπτόμενο διάνυσμα $u$. Πράγματι, η παράγωγος της ποσότητας αυτής κατά μήκος της γεωδαισιακής είναι
\[
u^\alpha \nabla_\alpha (g(\xi, u)) = u^\alpha u^\beta \nabla_\alpha \xi_\beta + g(\xi, u^\alpha \nabla_\alpha u).
\]
Ο δεύτερος όρος μηδενίζεται λόγω της εξίσωσης των γεωδαισιακών ($u^\alpha \nabla_\alpha u = 0$). Ο πρώτος όρος μηδενίζεται λόγω της εξίσωσης Killing, αφού η $\nabla_\alpha \xi_\beta + \nabla_\beta \xi_\alpha = 0$ συνεπάγεται ότι η συστολή με το συμμετρικό $u^\alpha u^\beta$ μηδενίζεται. Έτσι, σε έναν χωροχρόνο που διαθέτει χρονικό πεδίο Killing, η ενέργεια των ελεύθερα κινούμενων σωματιδίων διατηρείται. Αντίστοιχα, χωρικά πεδία Killing οδηγούν σε διατήρηση της ορμής.

\subsubsection{Διατήρηση του φασικού όγκου (Liouville)}

Στην κλασική μηχανική, ο φασικός χώρος μιας Χαμιλτονιανής συνάρτησης $H$ είναι μια συμπλεκτική πολλαπλότητα εφοδιασμένη με τη συμπλεκτική μορφή $\omega$. Το Χαμιλτονιανό διανυσματικό πεδίο $\xi_H$ ορίζεται από τη σχέση $\iota_{\xi_H} \omega = -dH$. Το θεώρημα Liouville δηλώνει ότι η Χαμιλτονιανή ροή διατηρεί τον φασικό όγκο. Σε γεωμετρική γλώσσα, αυτό εκφράζεται ως
\[
\mathcal{L}_{\xi_H} \omega = 0.
\]
Από την (58), η συνθήκη αυτή σε συνιστώσες είναι
\[
(\mathcal{L}_{\xi_H} \omega)_{ij} = \xi_H^\alpha \frac{\partial \omega_{ij}}{\partial x^\alpha} + \omega_{\alpha j} \frac{\partial \xi_H^\alpha}{\partial x^i} + \omega_{i\alpha} \frac{\partial \xi_H^\alpha}{\partial x^j} = 0.
\]
Η συνθήκη αυτή είναι ισοδύναμη με το ότι η ροή είναι συμπλεκτική. Ο φασικός όγκος είναι η $n$-μορφή που προκύπτει από το σφηνοειδές γινόμενο της συμπλεκτικής μορφής με τον εαυτό της $n$ φορές,
\[
\Omega = \underbrace{\omega \wedge \dots \wedge \omega}_{n \text{ φορές}}.
\]
Χρησιμοποιώντας την ιδιότητα της παραγώγου Lie ως προς το σφηνοειδές γινόμενο,
\[
\mathcal{L}_{\xi_H} \Omega = \mathcal{L}_{\xi_H} (\omega \wedge \dots \wedge \omega) = n (\mathcal{L}_{\xi_H} \omega) \wedge \omega^{n-1}.
\]
Εφόσον $\mathcal{L}_{\xi_H} \omega = 0$, έχουμε $\mathcal{L}_{\xi_H} \Omega = 0$. Αυτό σημαίνει ότι ο φασικός όγκος παραμένει σταθερός κατά μήκος της Χαμιλτονιανής ροής, που είναι ακριβώς το θεώρημα Liouville. Η γεωμετρική αυτή απόδειξη αναδεικνύει τη βαθιά σύνδεση ανάμεσα στη συμπλεκτική δομή και τη διατήρηση του όγκου.

\subsubsection{Θεώρημα Noether σε γεωμετρική γλώσσα}

Το θεώρημα Noether είναι η κορωνίδα της σύνδεσης ανάμεσα στις συμμετρίες και τους νόμους διατήρησης. Για να εκτιμήσουμε το βάθος του, πρέπει πρώτα να κατανοήσουμε τι είναι μια Λαγκρανζιανή και πώς σχετίζεται με τη δυναμική ενός συστήματος.

Στη φυσική, η κατάσταση ενός συστήματος περιγράφεται από μια συλλογή μεταβλητών --- αυτές μπορεί να είναι οι θέσεις σωματιδίων, οι τιμές ενός πεδίου σε κάθε σημείο του χώρου, ή ακόμα πιο αφηρημένα αντικείμενα. Η Λαγκρανζιανή είναι μια βαθμωτή συνάρτηση αυτών των μεταβλητών και των ρυθμών μεταβολής τους, που με έναν συμπυκνωμένο τρόπο κωδικοποιεί όλη τη δυναμική του συστήματος. Στη μηχανική ενός σωματιδίου, η Λαγκρανζιανή είναι η διαφορά κινητικής και δυναμικής ενέργειας. Στη θεωρία πεδίων, είναι μια πυκνότητα που εξαρτάται από τις τιμές του πεδίου και τις παραγώγους του. Το θεμελιώδες γεγονός είναι ότι η γνώση της Λαγκρανζιανής αρκεί για να εξαχθούν οι εξισώσεις κίνησης του συστήματος μέσω της αρχής της ελάχιστης δράσης.

Η αρχή αυτή είναι μια από τις βαθύτερες ιδέες της φυσικής. Ορίζουμε τη δράση $S$ ως το ολοκλήρωμα της Λαγκρανζιανής πυκνότητας πάνω σε μια χωροχρονική περιοχή. Για ένα μηχανικό σύστημα, αυτό είναι ένα απλό χρονικό ολοκλήρωμα της Λαγκρανζιανής. Για ένα πεδίο, είναι ένα ολοκλήρωμα πάνω σε μια περιοχή του χωροχρόνου. Η φύση, σύμφωνα με αυτή την αρχή, "επιλέγει" εκείνη την εξέλιξη του συστήματος για την οποία η δράση παίρνει ακρότατη τιμή --- συνήθως ελάχιστη. Από αυτή την απλή απαίτηση προκύπτουν όλες οι εξισώσεις κίνησης.

Ας δούμε ένα αδρό παράδειγμα για να γίνει αυτό πιο συγκεκριμένο. Φανταστείτε ένα απλό εκκρεμές, μια σημειακή μάζα που αιωρείται από ένα αβαρές νήμα. Η κατάστασή του περιγράφεται από μία μόνο γωνία, ας την πούμε $\theta$, που μετρά την απόκλιση από την κατακόρυφο. Η Λαγκρανζιανή του είναι η διαφορά κινητικής και δυναμικής ενέργειας, και μπορεί να γραφεί συναρτήσει της $\theta$ και του ρυθμού μεταβολής της $\dot{\theta}$. Η δράση είναι το χρονικό ολοκλήρωμα αυτής της Λαγκρανζιανής ανάμεσα σε δύο χρονικές στιγμές. Η αρχή της ελάχιστης δράσης λέει ότι η πραγματική κίνηση του εκκρεμούς, ανάμεσα σε μια αρχική και μια τελική γωνία, είναι εκείνη που κάνει αυτό το ολοκλήρωμα ελάχιστο. Από αυτή την απαίτηση προκύπτει η γνωστή εξίσωση του εκκρεμούς.

Σε αυτό το πλαίσιο, μια συμμετρία είναι ένας μετασχηματισμός των μεταβλητών του συστήματος που αφήνει τη δράση αναλλοίωτη. Για το εκκρεμές, αν το βαρυτικό πεδίο είναι ομογενές, η περιστροφή ολόκληρου του συστήματος γύρω από τον κατακόρυφο άξονα δεν αλλάζει τη δυναμική του --- αυτό είναι μια συμμετρία. Το θεώρημα Noether μας λέει ότι σε κάθε τέτοια συμμετρία αντιστοιχεί μια ποσότητα που διατηρείται κατά την κίνηση. Για την περιστροφική συμμετρία του εκκρεμούς, αυτή η ποσότητα είναι η στροφορμή γύρω από τον κατακόρυφο άξονα. Για ένα σύστημα σωματιδίων, η συμμετρία ως προς χωρικές μεταφορές --- το γεγονός ότι η δράση δεν αλλάζει αν μετατοπίσουμε ολόκληρο το σύστημα στον χώρο --- οδηγεί στη διατήρηση της ολικής ορμής. Η συμμετρία ως προς χρονικές μεταφορές --- το γεγονός ότι η Λαγκρανζιανή δεν εξαρτάται ρητά από τον χρόνο --- οδηγεί στη διατήρηση της ενέργειας.

Στη γεωμετρική γλώσσα που έχουμε αναπτύξει, η συνθήκη συμμετρίας εκφράζεται με την παράγωγο Lie: η δράση είναι αναλλοίωτη κάτω από μια ροή αν και μόνο αν η παράγωγος Lie της Λαγκρανζιανής πυκνότητας κατά μήκος του αντίστοιχου διανυσματικού πεδίου μηδενίζεται ή είναι μια ολική απόκλιση. Το θεώρημα Noether είναι λοιπόν η γεωμετρική απόδειξη ότι η αναλλοιότητα της δράσης κάτω από μια συνεχή οικογένεια μετασχηματισμών συνεπάγεται αναγκαστικά την ύπαρξη μιας διατηρούμενης ποσότητας. Κάθε φορά που η φύση μάς παρουσιάζει μια διατηρούμενη ποσότητα --- είτε είναι η ορμή, η ενέργεια, η στροφορμή ή το ηλεκτρικό φορτίο --- μπορούμε να είμαστε βέβαιοι ότι πίσω από αυτήν κρύβεται μια συμμετρία, και η παράγωγος Lie είναι το κλειδί που ξεκλειδώνει αυτή τη σύνδεση.

\section{Σύνοψη: Οι Θεμελιώδεις Τύποι της Παραγώγου Lie}

\textbf{Ορισμός μέσω ορίου}
\begin{equation}
(\mathcal{L}_\xi T)(p) = \lim_{t \to 0} \frac{(\phi_{-t})_* T(\phi_t(p)) - T(p)}{t}. \qquad (52)
\end{equation}
Η παράγωγος Lie μετρά τον ρυθμό μεταβολής ενός τανυστικού πεδίου όπως τον αντιλαμβάνεται ένας παρατηρητής που παρασύρεται από τη ροή του $\xi$.

\textbf{Βαθμωτή συνάρτηση $f$}
\begin{equation}
\mathcal{L}_\xi f = \xi^\alpha \frac{\partial f}{\partial x^\alpha}. \qquad (53)
\end{equation}
Είναι η κατευθυνόμενη παράγωγος της $f$ κατά μήκος του $\xi$. Ταυτίζεται με την υλική παράγωγο της ρευστομηχανικής όταν ο χρόνος ενσωματωθεί γεωμετρικά.

\textbf{Διανυσματικό πεδίο $\eta$}
\begin{equation}
(\mathcal{L}_\xi \eta)^i = \xi^\alpha \frac{\partial \eta^i}{\partial x^\alpha} - \eta^\alpha \frac{\partial \xi^i}{\partial x^\alpha}. \qquad (54)
\end{equation}
Ταυτίζεται με τον μεταθέτη $[\xi, \eta]$. Ο δεύτερος όρος διορθώνει για την παραμόρφωση του συστήματος συντεταγμένων.

\textbf{1-μορφή $\omega$}
\begin{equation}
(\mathcal{L}_\xi \omega)_i = \xi^\alpha \frac{\partial \omega_i}{\partial x^\alpha} + \omega_\alpha \frac{\partial \xi^\alpha}{\partial x^i}. \qquad (55)
\end{equation}
Το πρόσημο $+$ οφείλεται στον συναλλοίωτο χαρακτήρα της 1-μορφής.

\textbf{Γενικός τανυστής τύπου $(r, s)$}
\begin{equation}
\begin{aligned}
(\mathcal{L}_\xi T)^{i_1 \dots i_r}_{j_1 \dots j_s} &= \xi^\alpha \frac{\partial}{\partial x^\alpha} T^{i_1 \dots i_r}_{j_1 \dots j_s} \\
&\quad - T^{\alpha i_2 \dots i_r}_{j_1 \dots j_s} \frac{\partial \xi^{i_1}}{\partial x^\alpha} - \dots - T^{i_1 \dots i_{r-1} \alpha}_{j_1 \dots j_s} \frac{\partial \xi^{i_r}}{\partial x^\alpha} \\
&\quad + T^{i_1 \dots i_r}_{\alpha j_2 \dots j_s} \frac{\partial \xi^\alpha}{\partial x^{j_1}} + \dots + T^{i_1 \dots i_r}_{j_1 \dots j_{s-1} \alpha} \frac{\partial \xi^\alpha}{\partial x^{j_s}}. \qquad (56)
\end{aligned}
\end{equation}
Κάθε ανταλλοίωτος δείκτης συνεισφέρει έναν όρο $-T \partial \xi$, και κάθε συναλλοίωτος δείκτης έναν όρο $+T \partial \xi$. Η δομή είναι καθολική και προκύπτει από τον τρόπο που η ροή εφελκύει ή συνεφελκύει τους δείκτες.

\textbf{Μετρική $g$ και εξίσωση Killing}
\begin{equation}
(\mathcal{L}_\xi g)_{ij} = \xi^\alpha \frac{\partial g_{ij}}{\partial x^\alpha} + g_{\alpha j} \frac{\partial \xi^\alpha}{\partial x^i} + g_{i\alpha} \frac{\partial \xi^\alpha}{\partial x^j}. \qquad (57)
\end{equation}
Ο μηδενισμός της δίνει τις εξισώσεις Killing $\mathcal{L}_\xi g = 0$, που χαρακτηρίζουν τις ισομετρίες της μετρικής. Τα πεδία Killing αντιστοιχούν σε συμμετρίες του χωροχρόνου και γεννούν διατηρούμενα μεγέθη.

\textbf{Συμπλεκτική μορφή $\omega$}
\begin{equation}
(\mathcal{L}_\xi \omega)_{ij} = \xi^\alpha \frac{\partial \omega_{ij}}{\partial x^\alpha} + \omega_{\alpha j} \frac{\partial \xi^\alpha}{\partial x^i} + \omega_{i\alpha} \frac{\partial \xi^\alpha}{\partial x^j}. \qquad (58)
\end{equation}
Ο μηδενισμός της $\mathcal{L}_\xi \omega = 0$ χαρακτηρίζει τα συμπλεκτικά διανυσματικά πεδία, θεμελιώδη στην κλασική μηχανική.

\textbf{Τανυστής τάσεων $\sigma$}
\begin{equation}
(\mathcal{L}_\xi \sigma)^{ij} = \xi^\alpha \frac{\partial \sigma^{ij}}{\partial x^\alpha} - \sigma^{\alpha j} \frac{\partial \xi^i}{\partial x^\alpha} - \sigma^{i\alpha} \frac{\partial \xi^j}{\partial x^\alpha}. \qquad (59)
\end{equation}
Η έκφραση για τανυστή τύπου $(2,0)$, με δύο όρους διόρθωσης αρνητικούς.

\textbf{Τανυστής ηλεκτρομαγνητικού πεδίου $F$}
\begin{equation}
(\mathcal{L}_\xi F)_{ij} = \xi^\alpha \frac{\partial F_{ij}}{\partial x^\alpha} + F_{\alpha j} \frac{\partial \xi^\alpha}{\partial x^i} + F_{i\alpha} \frac{\partial \xi^\alpha}{\partial x^j}. \qquad (60)
\end{equation}
Η παράγωγος Lie του $F$ μετρά την αναλλοιότητα του ηλεκτρομαγνητικού πεδίου κάτω από τη ροή του $\xi$. Ο μηδενισμός της σημαίνει ότι το πεδίο είναι συμμετρικό ως προς τον παρατηρητή που ακολουθεί το $\xi$.

\section{Λογισμός των Μεταβολών}

\subsection{Εισαγωγή}

Ο λογισμός των μεταβολών είναι ο κλάδος των μαθηματικών που ασχολείται με την εύρεση ακρότατων συναρτησιακών — απεικονίσεων που στέλνουν συναρτήσεις ή καμπύλες σε πραγματικούς αριθμούς. Στη φυσική, ο λογισμός των μεταβολών είναι το θεμέλιο της Λαγκρανζιανής μηχανικής: η φυσική τροχιά ενός συστήματος είναι εκείνη που καθιστά στάσιμη τη δράση, ένα συναρτησιακό που ορίζεται ως το χρονικό ολοκλήρωμα της Λαγκρανζιανής. Σε αυτό το κεφάλαιο, αναπτύσσουμε τον λογισμό των μεταβολών στη γεωμετρική γλώσσα που έχουμε οικοδομήσει, χρησιμοποιώντας διανυσματικά πεδία, ροές και την παράγωγο Lie, και καταλήγουμε στο θεώρημα Noether, το οποίο συνδέει τις συμμετρίες με νόμους διατήρησης.

\subsection{Λαγκρανζιανή Μηχανική}

\subsubsection{Ο χώρος διαμορφώσεων και η δράση}

Θεωρούμε ένα φυσικό σύστημα του οποίου η κατάσταση σε κάθε χρονική στιγμή περιγράφεται πλήρως από ένα πεπερασμένο σύνολο γενικευμένων συντεταγμένων $q^1, \dots, q^n$. Ο χώρος όλων των δυνατών τιμών των $q^a$ είναι μια πολλαπλότητα $M$ διάστασης $n$, την οποία ονομάζουμε \textbf{χώρο διαμορφώσεων} του συστήματος. Μια τροχιά του συστήματος είναι μια ομαλή καμπύλη $\gamma: [t_1, t_2] \to M$, την οποία σε τοπικές συντεταγμένες γράφουμε ως $\gamma(t) = (q^a(t))$. Το εφαπτόμενο διάνυσμα της τροχιάς συμβολίζεται με $\dot{\gamma}(t)$ και έχει συνιστώσες $\dot{q}^a(t) = dq^a/dt$.

Η δυναμική του συστήματος καθορίζεται από μια βαθμωτή συνάρτηση $L$, τη \textbf{Λαγκρανζιανή}, η οποία είναι μια ομαλή συνάρτηση στον εφαπτόμενο χώρο $TM$ της πολλαπλότητας $M$. Σε τοπικές συντεταγμένες, η $L$ εξαρτάται από τις γενικευμένες συντεταγμένες $q^a$, τις ταχύτητες $\dot{q}^a$, και ενδεχομένως τον χρόνο $t$:
\[
L = L(q^a, \dot{q}^a, t).
\]

Η \textbf{δράση} $S$ είναι το συναρτησιακό που σε κάθε ομαλή καμπύλη $\gamma: [t_1, t_2] \to M$ αντιστοιχεί τον πραγματικό αριθμό
\begin{equation}
S[\gamma] = \int_{t_1}^{t_2} L(\gamma(t), \dot{\gamma}(t), t) \, dt. \qquad (178)
\end{equation}

\subsubsection{Η αρχή της στάσιμης δράσης και οι εξισώσεις Euler-Lagrange}

Η \textbf{αρχή της στάσιμης δράσης} (ή αρχή του Hamilton) δηλώνει ότι η φυσική τροχιά που ακολουθεί το σύστημα ανάμεσα σε δύο δεδομένες διαμορφώσεις $\gamma(t_1)$ και $\gamma(t_2)$ είναι εκείνη για την οποία η δράση είναι στάσιμη ως προς μεταβολές της καμπύλης που κρατούν σταθερά τα άκρα.

Για να διατυπώσουμε αυτή την αρχή με γεωμετρικό τρόπο, θεωρούμε ένα διανυσματικό πεδίο $\xi$ στη $M$, το οποίο παράγει μια τοπική ροή $\phi_\epsilon$. Η ροή αυτή παραμορφώνει μια δεδομένη τροχιά $\gamma$ σε μια νέα τροχιά $\gamma_\epsilon = \phi_\epsilon \circ \gamma$. Η μεταβολή της δράσης κάτω από αυτή την παραμόρφωση μετράται από την παράγωγο Lie της δράσης κατά μήκος του $\xi$:
\[
(\mathcal{L}_\xi S)[\gamma] = \left. \frac{d}{d\epsilon} \right|_{\epsilon=0} S[\phi_\epsilon \circ \gamma].
\]

Αν το διανυσματικό πεδίο $\xi$ μηδενίζεται στα άκρα της καμπύλης, $\xi(\gamma(t_1)) = \xi(\gamma(t_2)) = 0$, τότε η παραμορφωμένη καμπύλη $\gamma_\epsilon$ έχει τα ίδια άκρα με την $\gamma$. Η αρχή της στάσιμης δράσης απαιτεί
\[
\mathcal{L}_\xi S = 0,
\]
για κάθε διανυσματικό πεδίο $\xi$ που μηδενίζεται στα άκρα της τροχιάς.

Για να υπολογίσουμε την $\mathcal{L}_\xi S$, παραγωγίζουμε κάτω από το ολοκλήρωμα. Σε τοπικές συντεταγμένες, η μεταβολή της $L$ είναι
\[
\frac{d}{d\epsilon} L(\phi_\epsilon \circ \gamma, \dot{(\phi_\epsilon \circ \gamma)}, t) = \frac{\partial L}{\partial q^a} \xi^a + \frac{\partial L}{\partial \dot{q}^a} \dot{\xi}^a,
\]
όπου $\xi^a$ είναι οι συνιστώσες του $\xi$ και $\dot{\xi}^a = d\xi^a/dt$ κατά μήκος της τροχιάς. Συνεπώς,
\[
\mathcal{L}_\xi S = \int_{t_1}^{t_2} \left( \frac{\partial L}{\partial q^a} \xi^a + \frac{\partial L}{\partial \dot{q}^a} \dot{\xi}^a \right) dt.
\]

Ολοκληρώνοντας κατά μέρη τον δεύτερο όρο και χρησιμοποιώντας το γεγονός ότι $\xi^a(t_1) = \xi^a(t_2) = 0$, παίρνουμε
\begin{equation}
\mathcal{L}_\xi S = \int_{t_1}^{t_2} \left( \frac{\partial L}{\partial q^a} - \frac{d}{dt} \frac{\partial L}{\partial \dot{q}^a} \right) \xi^a \, dt. \qquad (179)
\end{equation}

Η απαίτηση $\mathcal{L}_\xi S = 0$ για κάθε $\xi$ που μηδενίζεται στα άκρα οδηγεί στις \textbf{εξισώσεις Euler-Lagrange}:
\begin{equation}
\frac{d}{dt} \left( \frac{\partial L}{\partial \dot{q}^a} \right) - \frac{\partial L}{\partial q^a} = 0, \qquad a = 1, \dots, n. \qquad (180)
\end{equation}
Αυτές είναι οι διαφορικές εξισώσεις που καθορίζουν την κίνηση του συστήματος. Ορίζουν ποιες τροχιές είναι φυσικές — ποιες τροχιές ακολουθεί πραγματικά το σύστημα.

\subsubsection{Η 1-μορφή της ορμής}

Από τη Λαγκρανζιανή αναδύεται μια θεμελιώδης γεωμετρική δομή στον χώρο διαμορφώσεων $M$: η \textbf{1-μορφή της ορμής}. Οι μερικές παράγωγοι της $L$ ως προς τις ταχύτητες,
\[
p_a = \frac{\partial L}{\partial \dot{q}^a},
\]
είναι οι συνιστώσες μιας 1-μορφής στην $M$. Για να το διαπιστώσουμε, ας εξετάσουμε πώς μετασχηματίζονται κάτω από αλλαγή συντεταγμένων $q^a \to \tilde{q}^a$ στη $M$. Οι ταχύτητες μετασχηματίζονται ανταλλοίωτα:
\[
\dot{\tilde{q}}^a = \frac{\partial \tilde{q}^a}{\partial q^b} \dot{q}^b.
\]

Η Λαγκρανζιανή στις νέες συντεταγμένες είναι $\tilde{L}(\tilde{q}, \dot{\tilde{q}}, t) = L(q(\tilde{q}), \dot{q}(\tilde{q}, \dot{\tilde{q}}), t)$. Παραγωγίζοντας ως προς $\dot{\tilde{q}}^a$ με τον κανόνα της αλυσίδας,
\[
\frac{\partial \tilde{L}}{\partial \dot{\tilde{q}}^a} = \frac{\partial L}{\partial \dot{q}^b} \frac{\partial \dot{q}^b}{\partial \dot{\tilde{q}}^a} = \frac{\partial L}{\partial \dot{q}^b} \frac{\partial q^b}{\partial \tilde{q}^a},
\]
αφού από την αντιστροφή του παραπάνω μετασχηματισμού έχουμε $\dot{q}^b = (\partial q^b / \partial \tilde{q}^a) \dot{\tilde{q}}^a$, και άρα $\partial \dot{q}^b / \partial \dot{\tilde{q}}^a = \partial q^b / \partial \tilde{q}^a$. Αυτός είναι ακριβώς ο συναλλοίωτος νόμος μετασχηματισμού μιας 1-μορφής. Συνεπώς, η ποσότητα
\[
p = p_a \, dq^a
\]
είναι μια καλώς ορισμένη 1-μορφή στην $M$, ανεξάρτητη από την επιλογή συντεταγμένων. Την ονομάζουμε \textbf{1-μορφή της ορμής}.

Η 1-μορφή της ορμής έχει μια άμεση γεωμετρική ερμηνεία: σε κάθε σημείο $q \in M$ και για κάθε εφαπτόμενο διάνυσμα $\dot{q} \in T_qM$, η τιμή $p(\dot{q}) = p_a \dot{q}^a$ είναι η παράγωγος της Λαγκρανζιανής κατά μήκος του $\dot{q}$ ως προς την κλίμακα της ταχύτητας. Με άλλα λόγια, η $p$ μετρά την ``ευαισθησία'' της Λαγκρανζιανής σε απειροστές μεταβολές της ταχύτητας. Στη φυσική, για ένα σωματίδιο με κινητική ενέργεια $T = \frac{1}{2} m \, g_{ab} \dot{q}^a \dot{q}^b$, η ορμή είναι $p_a = m \, g_{ab} \dot{q}^b$, δηλαδή η 1-μορφή που αντιστοιχεί στη συνήθη ορμή μέσω της μετρικής. Αλλά ακόμα και χωρίς μετρική, η $p$ είναι ένα καλώς ορισμένο γεωμετρικό αντικείμενο που κωδικοποιεί την απόκριση της $L$ στις ταχύτητες.

Η χρησιμότητα της $p$ έγκειται στο γεγονός ότι επιτρέπει τη γεωμετρική διατύπωση των εξισώσεων κίνησης και των νόμων διατήρησης. Οι εξισώσεις Euler-Lagrange (180) μπορούν να γραφούν ως
\[
\frac{d}{dt} p_a = \frac{\partial L}{\partial q^a},
\]
που εκφράζει τον ρυθμό μεταβολής της ορμής συναρτήσει των ``δυνάμεων'' $\partial L / \partial q^a$. Αυτή είναι η γεωμετρική γενίκευση του δεύτερου νόμου του Νεύτωνα: η χρονική παράγωγος της ορμής ισούται με τη δύναμη.

\subsection{Το Θεώρημα Noether}

\subsubsection{Συμμετρίες της Λαγκρανζιανής}

Μια \textbf{συμμετρία} της Λαγκρανζιανής είναι ένα διανυσματικό πεδίο $\xi$ στη $M$ (το οποίο μπορεί να εξαρτάται και από τον χρόνο) του οποίου η ροή αφήνει αναλλοίωτη τη δράση, modulo οριακούς όρους. Για να το διατυπώσουμε με ακρίβεια, θεωρούμε ότι η μεταβολή της Λαγκρανζιανής κάτω από τη ροή του $\xi$ είναι μια ολική χρονική παράγωγος κάποιας συνάρτησης $K$:
\begin{equation}
\frac{d}{d\epsilon} \bigg|_{\epsilon=0} L(\phi_\epsilon \circ \gamma, \dot{(\phi_\epsilon \circ \gamma)}, t) = \frac{d}{dt} K(\gamma(t), t). \qquad (181)
\end{equation}
Σε τοπικές συντεταγμένες, αυτό γράφεται
\begin{equation}
\frac{\partial L}{\partial q^a} \xi^a + \frac{\partial L}{\partial \dot{q}^a} \dot{\xi}^a = \frac{dK}{dt}. \qquad (182)
\end{equation}

Η συνθήκη αυτή σημαίνει ότι η δράση αλλάζει μόνο κατά οριακούς όρους κάτω από τη ροή του $\xi$:
\[
\mathcal{L}_\xi S = \int_{t_1}^{t_2} \frac{dK}{dt} \, dt = K(\gamma(t_2), t_2) - K(\gamma(t_1), t_1).
\]
Αν το $\xi$ μηδενίζεται στα άκρα, τότε $\mathcal{L}_\xi S = 0$ — αλλά η συνθήκη (181) είναι ισχυρότερη: απαιτεί η μεταβολή της $L$ να είναι ολική παράγωγος ακόμα και όταν το $\xi$ δεν μηδενίζεται στα άκρα. Αυτό ακριβώς χαρακτηρίζει μια συμμετρία: η δράση είναι αναλλοίωτη για κάθε τροχιά, ανεξάρτητα από τις τιμές του $\xi$ στα άκρα, αρκεί η μεταβολή της $L$ να είναι μια ολική παράγωγος.

\subsubsection{Διατύπωση και απόδειξη}

\textbf{Θεώρημα (Noether).} Έστω $L(q^a, \dot{q}^a, t)$ μια Λαγκρανζιανή και έστω $\xi$ ένα διανυσματικό πεδίο στη $M$ (με ενδεχόμενη χρονική εξάρτηση) που ικανοποιεί τη συνθήκη συμμετρίας (181) για κάποια συνάρτηση $K$. Τότε η ποσότητα
\begin{equation}
Q = \frac{\partial L}{\partial \dot{q}^a} \xi^a - K \qquad (183)
\end{equation}
διατηρείται κατά την κίνηση: για κάθε τροχιά που ικανοποιεί τις εξισώσεις Euler-Lagrange (180), ισχύει
\begin{equation}
\frac{dQ}{dt} = 0. \qquad (184)
\end{equation}

\textbf{Απόδειξη.} Υπολογίζουμε την ολική χρονική παράγωγο της $Q$ κατά μήκος μιας τροχιάς που ικανοποιεί τις εξισώσεις Euler-Lagrange. Έχουμε
\[
\frac{dQ}{dt} = \frac{d}{dt} \left( \frac{\partial L}{\partial \dot{q}^a} \xi^a \right) - \frac{dK}{dt}.
\]
Αναπτύσσοντας τον πρώτο όρο,
\[
\frac{d}{dt} \left( \frac{\partial L}{\partial \dot{q}^a} \xi^a \right) = \left( \frac{d}{dt} \frac{\partial L}{\partial \dot{q}^a} \right) \xi^a + \frac{\partial L}{\partial \dot{q}^a} \dot{\xi}^a.
\]
Χρησιμοποιώντας τις εξισώσεις Euler-Lagrange (180), αντικαθιστούμε την $d/dt (\partial L / \partial \dot{q}^a)$ με $\partial L / \partial q^a$. Έτσι,
\[
\frac{d}{dt} \left( \frac{\partial L}{\partial \dot{q}^a} \xi^a \right) = \frac{\partial L}{\partial q^a} \xi^a + \frac{\partial L}{\partial \dot{q}^a} \dot{\xi}^a.
\]
Από τη συνθήκη συμμετρίας (182), το δεξί μέλος ισούται ακριβώς με $dK/dt$. Άρα
\[
\frac{dQ}{dt} = \frac{d}{dt} \left( \frac{\partial L}{\partial \dot{q}^a} \xi^a \right) - \frac{dK}{dt} = 0.
\]
Αυτό ολοκληρώνει την απόδειξη. $\square$

\subsubsection{Γεωμετρική ερμηνεία με όρους της 1-μορφής της ορμής}

Το θεώρημα Noether αποκτά μια βαθιά γεωμετρική ερμηνεία όταν διατυπωθεί με όρους της 1-μορφής της ορμής $p = p_a \, dq^a$. Ο πρώτος όρος στο φορτίο Noether (183) είναι ακριβώς η δράση της 1-μορφής $p$ πάνω στο διανυσματικό πεδίο $\xi$:
\[
p(\xi) = p_a \xi^a = \frac{\partial L}{\partial \dot{q}^a} \xi^a.
\]
Πρόκειται για τη φυσική σύζευξη ανάμεσα σε μια 1-μορφή και ένα διανυσματικό πεδίο — την ίδια σύζευξη που βρίσκεται στον πυρήνα της δυϊκότητας $V^* \times V \to \mathbb{R}$. Το φορτίο Noether γράφεται λοιπόν ως
\begin{equation}
Q = p(\xi) - K = \iota_\xi p - K. \qquad (185)
\end{equation}

Αυτή η έκφραση είναι γεωμετρικά θεμελιώδης. Η $p$ είναι μια 1-μορφή στην $M$ — ένα αντικείμενο που σε κάθε σημείο τρώει ένα διάνυσμα και δίνει έναν αριθμό. Το $\xi$ είναι ένα διανυσματικό πεδίο που παράγει μια συμμετρία. Η τιμή $p(\xi)$ είναι λοιπόν η ``ορμή στην κατεύθυνση της συμμετρίας''. Με άλλα λόγια, μετρά πόση από την ορμή του συστήματος είναι ευθυγραμμισμένη με τον μετασχηματισμό που παράγει το $\xi$. Η συνάρτηση $K$ είναι μια διόρθωση που εμφανίζεται όταν η Λαγκρανζιανή δεν είναι απολύτως αναλλοίωτη αλλά αλλάζει κατά μια ολική παράγωγο.

Η διατήρηση της $Q$ σημαίνει ότι η τιμή της $p(\xi)$, αφού αφαιρεθεί η διόρθωση $K$, παραμένει σταθερή κατά μήκος της φυσικής τροχιάς. Με άλλα λόγια, η προβολή της ορμής στην κατεύθυνση της συμμετρίας διατηρείται. Αυτή είναι η γεωμετρική ουσία του θεωρήματος Noether: μια συμμετρία $\xi$ επιλέγει μια ``κατεύθυνση'' στον χώρο διαμορφώσεων, και η 1-μορφή της ορμής, όταν συσταλεί με αυτή την κατεύθυνση, παράγει μια διατηρούμενη ποσότητα.

\subsubsection{Φυσικό νόημα: νόμοι διατήρησης}

Το θεώρημα Noether θεμελιώνει μαθηματικά τη βαθιά φυσική αρχή ότι οι συμμετρίες γεννούν νόμους διατήρησης. Η $Q$ είναι μια σταθερή ποσότητα κατά μήκος της φυσικής τροχιάς — μια \textbf{πρώτη ολοκλήρωση} της κίνησης. Αυτό σημαίνει ότι η τιμή της $Q$ μπορεί να υπολογιστεί από τις αρχικές συνθήκες και παραμένει αναλλοίωτη για πάντα. Σε κάθε περίπτωση, η γεωμετρική δομή είναι η ίδια: ένα διανυσματικό πεδίο $\xi$ που παράγει μια συμμετρία οδηγεί, μέσω της 1-μορφής της ορμής $p$, σε μια διατηρούμενη ποσότητα $Q = \iota_\xi p - K$.

\subsection{Εξειδικεύσεις σε φυσικές θεωρίες}

Το θεώρημα Noether βρίσκει συγκεκριμένες και θεμελιώδεις εφαρμογές σε όλες τις κύριες φυσικές θεωρίες. Σε κάθε περίπτωση, η γενική δομή παραμένει η ίδια: μια συμμετρία της Λαγκρανζιανής οδηγεί σε μια διατηρούμενη ποσότητα. Αυτό που αλλάζει είναι η φύση του χώρου διαμορφώσεων, η μορφή της Λαγκρανζιανής, και η φυσική ερμηνεία του φορτίου Noether.

\subsubsection{Κλασική μηχανική: διατήρηση ενέργειας και στροφορμής}

Στην κλασική μηχανική, ο χώρος διαμορφώσεων ενός συστήματος $N$ σωματιδίων σε τρεις διαστάσεις είναι ο $\mathbb{R}^{3N}$, με συντεταγμένες τις θέσεις $\mathbf{r}_1, \dots, \mathbf{r}_N$. Η Λαγκρανζιανή έχει τη μορφή $L = T - V$, όπου
\[
T = \sum_{i=1}^N \frac{1}{2} m_i \delta_{ab} \dot{x}_i^a \dot{x}_i^b
\]
είναι η κινητική ενέργεια και $V(\mathbf{r}_1, \dots, \mathbf{r}_N)$ είναι η δυναμική ενέργεια. Οι εξισώσεις Euler-Lagrange (180) δίνουν τον δεύτερο νόμο του Νεύτωνα για κάθε σωματίδιο:
\[
m_i \ddot{\mathbf{r}}_i = -\nabla_i V.
\]
Η 1-μορφή της ορμής είναι $p = \sum_i m_i \delta_{ab} \dot{x}_i^a dx_i^b$, η συνήθης ορμή των σωματιδίων.

Αν το δυναμικό $V$ δεν εξαρτάται ρητά από τον χρόνο, η Λαγκρανζιανή είναι αναλλοίωτη σε χρονικές μεταφορές. Το διανυσματικό πεδίο που παράγει αυτή τη συμμετρία είναι απλώς η μετατόπιση στον χρόνο, $\xi = \partial / \partial t$. Εφαρμόζοντας τη συνθήκη (182), βρίσκουμε ότι η μεταβολή της $L$ είναι $dL/dt$, άρα $K = L$. Το φορτίο Noether είναι τότε
\[
Q = \sum_i m_i \delta_{ab} \dot{x}_i^a \dot{x}_i^b - L = 2T - (T - V) = T + V = E.
\]
Πρόκειται για την ολική μηχανική ενέργεια του συστήματος. Η διατήρησή της, $dE/dt = 0$, σημαίνει ότι η ενέργεια παραμένει σταθερή κατά την κίνηση. Η φυσική προέλευση αυτού του νόμου είναι η ομοιογένεια του χρόνου: οι νόμοι της φυσικής είναι οι ίδιοι ανεξάρτητα από το πότε εκτελούμε ένα πείραμα.

Αν το δυναμικό $V$ εξαρτάται μόνο από τις αποστάσεις μεταξύ των σωματιδίων, η Λαγκρανζιανή είναι αναλλοίωτη σε στροφές του συστήματος ως συνόλου. Μια απειροστή στροφή γύρω από έναν άξονα με μοναδιαίο διάνυσμα $\mathbf{n}$ περιγράφεται από το διανυσματικό πεδίο $\xi = \mathbf{n} \times \mathbf{r}$ που εφαρμόζεται σε κάθε σωματίδιο. Η συνθήκη συμμετρίας (182) δίνει $K = 0$, αφού η κινητική ενέργεια είναι αναλλοίωτη σε στροφές (το εσωτερικό γινόμενο $\delta_{ab} \dot{x}^a \dot{x}^b$ δεν αλλάζει) και το δυναμικό εξαρτάται μόνο από αποστάσεις. Το φορτίο Noether είναι
\[
Q = \sum_i m_i \dot{\mathbf{r}}_i \cdot (\mathbf{n} \times \mathbf{r}_i) = \mathbf{n} \cdot \sum_i m_i (\mathbf{r}_i \times \dot{\mathbf{r}}_i) = \mathbf{n} \cdot \mathbf{L},
\]
όπου $\mathbf{L}$ είναι η ολική στροφορμή του συστήματος. Αφού η $Q$ διατηρείται για κάθε κατεύθυνση $\mathbf{n}$, η ολική στροφορμή $\mathbf{L}$ διατηρείται. Η φυσική προέλευση αυτού του νόμου είναι η ισοτροπία του χώρου: οι νόμοι της φυσικής είναι οι ίδιοι ανεξάρτητα από τον προσανατολισμό του εργαστηρίου.

\subsubsection{Ηλεκτρομαγνητισμός: διατήρηση του ηλεκτρικού φορτίου}

Ο κλασικός ηλεκτρομαγνητισμός διατυπώνεται ως θεωρία πεδίου, όπου το θεμελιώδες δυναμικό είναι η 1-μορφή $A = A_j \, dx^j$. Ο χώρος διαμορφώσεων είναι ο απειροδιάστατος χώρος όλων των δυνατών διαμορφώσεων του πεδίου $A_j(x)$. Η Λαγκρανζιανή πυκνότητα είναι
\[
\mathcal{L} = -\frac{1}{4} F_{ij} F^{ij} + A_j J^j,
\]
όπου $F = dA$ είναι η 2-μορφή του πεδίου, με συνιστώσες
\[
F_{ij} = \frac{\partial A_j}{\partial x^i} - \frac{\partial A_i}{\partial x^j}.
\]
Ο πρώτος όρος $-\frac{1}{4} F_{ij} F^{ij}$ είναι η κινητική ενέργεια του ηλεκτρομαγνητικού πεδίου, και ο δεύτερος όρος $A_j J^j$ είναι η αλληλεπίδραση του πεδίου με την ύλη, όπου $J^j$ είναι η τετραπυκνότητα ρεύματος. Οι εξισώσεις Euler-Lagrange για το $A$ δίνουν τις εξισώσεις Maxwell:
\[
\frac{\partial}{\partial x^i} F^{ij} = J^j.
\]

Η ηλεκτροδυναμική διαθέτει μια θεμελιώδη \textbf{συμμετρία βαθμίδας}: ο μετασχηματισμός $A \to A + d\Lambda$, όπου $\Lambda$ είναι μια αυθαίρετη βαθμωτή συνάρτηση, αφήνει αναλλοίωτο το πεδίο $F$, αφού $F = dA = d(A + d\Lambda)$ λόγω της $d^2 = 0$. Αυτό σημαίνει ότι δύο δυναμικά που διαφέρουν κατά $d\Lambda$ περιγράφουν ακριβώς την ίδια φυσική κατάσταση. Πρόκειται για μια συνεχή συμμετρία, που παράγεται από το διανυσματικό πεδίο $\xi = d\Lambda$ στον χώρο των πεδίων. Σε αντίθεση με τις χωροχρονικές συμμετρίες που συζητήσαμε παραπάνω, αυτή είναι μια \textbf{εσωτερική συμμετρία} — δεν σχετίζεται με τον χωροχρόνο αλλά με τον χώρο των ίδιων των πεδίων.

Εφαρμόζοντας το θεώρημα Noether στη συμμετρία βαθμίδας, η Λαγκρανζιανή $\mathcal{L}$ είναι αναλλοίωτη κάτω από $A \to A + d\Lambda$ (ο όρος $F_{ij}F^{ij}$ είναι αναλλοίωτος, και ο όρος $A_j J^j$ αλλάζει κατά $J^j \partial \Lambda / \partial x^j$, που είναι μια ολική απόκλιση όταν το ρεύμα διατηρείται). Η συνθήκη (182) δίνει $K = 0$, και το φορτίο Noether που προκύπτει είναι το ολικό ηλεκτρικό φορτίο
\[
Q = \int_\Sigma \star J,
\]
όπου $\Sigma$ είναι μια χωρική υπερεπιφάνεια. Η διατήρηση του φορτίου εκφράζεται τοπικά από την εξίσωση συνέχειας $d \star J = 0$, την οποία έχουμε ήδη συναντήσει στη μελέτη της απόκλισης. Το θεώρημα Noether λοιπόν δείχνει ότι η συμμετρία βαθμίδας του ηλεκτρομαγνητισμού είναι υπεύθυνη για τη διατήρηση του ηλεκτρικού φορτίου — μια σύνδεση που δεν είναι καθόλου προφανής από τη συνήθη διατύπωση της θεωρίας.

\subsubsection{Ειδική σχετικότητα: διατήρηση της τετραορμής}

Στην ειδική σχετικότητα, ο χωροχρόνος είναι ο $\mathbb{R}^4$ εφοδιασμένος με τη μετρική Minkowski $\eta_{ij} = \text{diag}(-1, 1, 1, 1)$. Για ένα ελεύθερο σχετικιστικό σωματίδιο, η Λαγκρανζιανή ως προς μια αυθαίρετη παράμετρο $\tau$ της τροχιάς είναι
\[
L = -m \sqrt{-\eta_{ij} \frac{dx^i}{d\tau} \frac{dx^j}{d\tau}}.
\]
Οι εξισώσεις Euler-Lagrange δίνουν ευθύγραμμη ομαλή κίνηση: η τετραταχύτητα είναι σταθερή. Η 1-μορφή της ορμής είναι η τετραορμή
\[
p_i = \frac{m \, \eta_{ij} \frac{dx^j}{d\tau}}{\sqrt{-\eta_{kl} \frac{dx^k}{d\tau} \frac{dx^l}{d\tau}}}.
\]

Ο χωροχρόνος Minkowski είναι μεγιστικά συμμετρικός: είναι ομογενής στον χώρο και στον χρόνο, και ισότροπος. Οι χωροχρονικές μεταφορές $x^i \to x^i + \epsilon^i$ αποτελούν τέσσερις ανεξάρτητες συμμετρίες, που παράγονται από τα διανυσματικά πεδία $\xi_{(i)} = \partial / \partial x^i$ (για $i = 0, 1, 2, 3$). Για καθένα από αυτά, η Λαγκρανζιανή είναι αναλλοίωτη (δεν εξαρτάται ρητά από τις συντεταγμένες), άρα η συνθήκη (182) δίνει $K = 0$. Το φορτίο Noether για τη συμμετρία στην κατεύθυνση $i$ είναι
\[
Q_{(i)} = p(\xi_{(i)}) = p_j \xi_{(i)}^j = p_i.
\]
Έτσι, και οι τέσσερις συνιστώσες της τετραορμής διατηρούνται. Η χρονική συνιστώσα $p_0$ είναι η ενέργεια (διαιρεμένη με $c$), και οι χωρικές συνιστώσες $p_1, p_2, p_3$ είναι η ορμή. Η ομοιογένεια του χωροχρόνου Minkowski — το γεγονός ότι δεν υπάρχουν προνομιακά σημεία ή κατευθύνσεις — οδηγεί λοιπόν στη διατήρηση της τετραορμής. Αυτό γενικεύει τη Νευτώνεια διατήρηση της ενέργειας και της ορμής σε ένα πλαίσιο συμβατό με την ειδική σχετικότητα.

\subsubsection{Γενική σχετικότητα: διατήρηση ενέργειας-ορμής και πεδία Killing}

Στη γενική σχετικότητα, ο χωροχρόνος είναι μια τετραδιάστατη ψευδο-Ρημάννεια πολλαπλότητα $(M, g)$ της οποίας η μετρική είναι το δυναμικό πεδίο. Σε αντίθεση με τις προηγούμενες θεωρίες, εδώ ο ίδιος ο χωροχρόνος δεν είναι μια σταθερή σκηνή αλλά ένα δυναμικό αντικείμενο που καθορίζεται από την κατανομή ύλης και ενέργειας μέσω των εξισώσεων Einstein. Αυτό έχει μια βαθιά συνέπεια: η γενική σχετικότητα δεν διαθέτει μια καθολική έννοια ενέργειας ή ορμής που να διατηρείται με τον τρόπο που συμβαίνει στις προηγούμενες θεωρίες. Ο λόγος είναι ακριβώς ότι δεν υπάρχουν εκ των προτέρων χωροχρονικές συμμετρίες — ο χωροχρόνος μπορεί να είναι οτιδήποτε, και η έννοια της ``χρονικής μεταφοράς'' ή της ``χωρικής μεταφοράς'' δεν είναι καθολικά ορισμένη.

Ωστόσο, το θεώρημα Noether εξακολουθεί να ισχύει όταν ο χωροχρόνος διαθέτει συμμετρίες. Μια χωροχρονική συμμετρία στη γενική σχετικότητα περιγράφεται από ένα \textbf{πεδίο Killing} $\xi$, ένα διανυσματικό πεδίο που ικανοποιεί $\mathcal{L}_\xi g = 0$. Η συνθήκη αυτή σημαίνει ότι η μετρική — και συνεπώς όλη η γεωμετρία του χωροχρόνου — είναι αναλλοίωτη κάτω από τη ροή του $\xi$. Παραδείγματα χωροχρόνων με πεδία Killing είναι ο χωροχρόνος Schwarzschild (που περιγράφει μια σφαιρικά συμμετρική μαύρη τρύπα και διαθέτει χρονικό και περιστροφικό πεδίο Killing) και ο χωροχρόνος Friedmann-Lemaître-Robertson-Walker (που περιγράφει ένα ομογενές και ισότροπο σύμπαν και διαθέτει χωρικά πεδία Killing).

Για ένα ελεύθερο σωματίδιο που κινείται σε έναν τέτοιο χωροχρόνο, η τροχιά του είναι μια γεωδαισιακή. Η Λαγκρανζιανή είναι
\[
L = -m \sqrt{-g_{ij} \frac{dx^i}{d\tau} \frac{dx^j}{d\tau}},
\]
και η 1-μορφή της ορμής είναι
\[
p_i = \frac{m \, g_{ij} \frac{dx^j}{d\tau}}{\sqrt{-g_{kl} \frac{dx^k}{d\tau} \frac{dx^l}{d\tau}}}.
\]
Αν ο χωροχρόνος διαθέτει ένα πεδίο Killing $\xi$, η Λαγκρανζιανή είναι αναλλοίωτη κάτω από τη ροή του $\xi$ (αφού η μετρική είναι αναλλοίωτη). Η συνθήκη (182) δίνει $K = 0$, και το φορτίο Noether είναι
\[
Q = p(\xi) = g(\xi, \dot{\gamma}).
\]
Αυτή η ποσότητα διατηρείται κατά μήκος της γεωδαισιακής. Αν ο χωροχρόνος διαθέτει ένα χρονικό πεδίο Killing (όπως ο Schwarzschild), η $Q$ ερμηνεύεται ως ενέργεια του σωματιδίου. Αν διαθέτει χωρικά πεδία Killing (όπως ο FLRW), οι αντίστοιχες $Q$ ερμηνεύονται ως ορμή. Σε έναν γενικό χωροχρόνο χωρίς καμία συμμετρία, δεν υπάρχουν πεδία Killing, και η έννοια της ολικής διατηρούμενης ενέργειας παύει να είναι καλώς ορισμένη. Αυτό δεν είναι ένα ελάττωμα της θεωρίας αλλά μια θεμελιώδης ιδιότητά της: στη γενική σχετικότητα, η ενέργεια του βαρυτικού πεδίου δεν μπορεί να εντοπιστεί, και μόνο σε χωροχρόνους με συμμετρίες μπορούμε να ορίσουμε ολικές διατηρούμενες ποσότητες.

\section{Γεωμετρική Θεωρία των Εξισώσεων Cauchy και το Θεώρημα Cauchy-Kowalevskaya}

\subsection{Το πρόβλημα αρχικών τιμών σε γεωμετρική γλώσσα}

Στη φυσική και τη γεωμετρία, συναντάμε συχνά το ακόλουθο πρόβλημα: γνωρίζουμε την κατάσταση ενός συστήματος σε μια δεδομένη επιφάνεια και θέλουμε να προβλέψουμε την εξέλιξή του σε μια γειτονική περιοχή. Μαθηματικά, αυτό μεταφράζεται στο πρόβλημα αρχικών τιμών για ένα σύστημα μερικών διαφορικών εξισώσεων. Το πρόβλημα Cauchy είναι η ακριβής διατύπωση αυτής της ιδέας: δεδομένης μιας υπερεπιφάνειας $\Sigma$ (την "αρχική επιφάνεια") και δεδομένων των τιμών ενός συνόλου πεδίων πάνω σε αυτήν, αναζητούμε μια λύση των εξισώσεων σε μια περιοχή της $\Sigma$ που να αναπαράγει αυτές τις αρχικές τιμές.

Για να το θέσουμε σε γεωμετρική γλώσσα, θεωρούμε μια πολλαπλότητα $M$ διάστασης $n$ και μια υπερεπιφάνεια $\Sigma \subset M$ διάστασης $n-1$. Η $\Sigma$ ορίζεται τοπικά ως το σύνολο μηδενισμού μιας βαθμωτής συνάρτησης $\phi$,
\[
\Sigma = \{ p \in M \mid \phi(p) = 0 \},
\]
με την προϋπόθεση ότι η 1-μορφή $d\phi$ δεν μηδενίζεται πάνω στην $\Sigma$. Αυτή η συνθήκη εξασφαλίζει ότι η $\Sigma$ είναι μια λεία υπερεπιφάνεια και ότι η $d\phi$ ορίζει ένα κάθετο διάνυσμα στην $\Sigma$ μέσω της μουσικής ισοδυναμίας.

Θεωρούμε τώρα ένα διανυσματικό πεδίο $\xi$ στη γειτονιά της $\Sigma$. Το $\xi$ αντιπροσωπεύει την "κατεύθυνση εξέλιξης" — στη φυσική, αυτό είναι συχνά ένα χρονικό διανυσματικό πεδίο. Το πρόβλημα Cauchy πρώτης τάξης για μια βαθμωτή συνάρτηση $u$ μπορεί να διατυπωθεί ως εξής: ζητούμε μια λύση $u$ που να ικανοποιεί
\begin{equation}
\mathcal{L}_\xi u = F(x, u),
\end{equation}
σε μια περιοχή της $\Sigma$, με αρχική συνθήκη
\begin{equation}
u|_{\Sigma} = u_0,
\end{equation}
όπου $u_0$ είναι μια δεδομένη συνάρτηση στην $\Sigma$ και $F$ είναι μια γνωστή συνάρτηση. Σε τοπικές συντεταγμένες, αυτή η εξίσωση είναι η γνωστή ημιγραμμική εξίσωση πρώτης τάξης. Η παράγωγος Lie $\mathcal{L}_\xi$ είναι το φυσικό γεωμετρικό αντικείμενο που περιγράφει την εξέλιξη κατά μήκος της ροής του $\xi$.

\subsection{Τανυστική δομή των συστημάτων πρώτης τάξης}

Πριν προχωρήσουμε, είναι σημαντικό να αναγνωρίσουμε ποιοι τανυστές εμπλέκονται στη θεωρία Cauchy. Ένα γενικό σύστημα πρώτης τάξης για άγνωστες συναρτήσεις $u^\alpha$ ($\alpha = 1, \dots, m$) των μεταβλητών $x^i$ ($i = 1, \dots, n$) γράφεται ως
\begin{equation}
A^{i \alpha}_\beta(x, u) \frac{\partial u^\beta}{\partial x^i} = F^\alpha(x, u).
\end{equation}
Εδώ, οι $u^\alpha$ είναι οι $m$ άγνωστες βαθμωτές συναρτήσεις — ένα αντικείμενο με έναν εσωτερικό δείκτη $\alpha$ που ζει σε έναν αφηρημένο χώρο $\mathbb{R}^m$. Οι συντελεστές $A^{i \alpha}_\beta$ σχηματίζουν ένα τανυστικό πεδίο με τρεις δείκτες: ο $i$ είναι χωροχρονικός δείκτης (τύπου $(1,0)$ στην πολλαπλότητα $M$), ενώ οι $\alpha$ και $\beta$ είναι εσωτερικοί δείκτες (τύπου $(1,1)$ στον χώρο $\mathbb{R}^m$). Για κάθε σταθεροποιημένο $i$, ο $A^{i \alpha}_\beta$ είναι ένας $m \times m$ πίνακας. Το δεξί μέλος $F^\alpha$ είναι μια απεικόνιση από τον χώρο των αγνώστων στον εαυτό του.

Για $m=1$, ο εσωτερικός δείκτης εξαφανίζεται και ο $A^{i \alpha}_\beta$ ανάγεται σε ένα σύνηθες διανυσματικό πεδίο $A^i$ πάνω στην πολλαπλότητα. Η εξίσωση γίνεται
\begin{equation}
A^i \frac{\partial u}{\partial x^i} = F(x, u),
\end{equation}
που είναι ακριβώς η παράγωγος Lie της βαθμωτής συνάρτησης $u$ κατά μήκος του διανυσματικού πεδίου $A = A^i \partial / \partial x^i$,
\begin{equation}
\mathcal{L}_A u = F(x, u).
\end{equation}
Βλέπουμε λοιπόν ότι η βασική δομή είναι μια παράγωγος Lie που εξισώνεται με μια πηγή. Για $m>1$, η γενίκευση περιλαμβάνει πίνακες, αλλά η γεωμετρική ερμηνεία παραμένει: πρόκειται για μια παράγωγο Lie κατά μήκος ενός "γενικευμένου διανυσματικού πεδίου" που δρα σε ένα διανυσματικό χώρο αγνώστων.

\subsection{Χαρακτηριστικές επιφάνειες και η γεωμετρική τους σημασία}

Μια υπερεπιφάνεια $\Sigma$ με κάθετη 1-μορφή $n = n_i dx^i$ ονομάζεται \textbf{χαρακτηριστική} για το σύστημα
\[
A^{i \alpha}_\beta \frac{\partial u^\beta}{\partial x^i} = F^\alpha,
\]
αν ο πίνακας
\begin{equation}
A(n) \equiv A^{i} n_i
\end{equation}
είναι ιδιόμορφος, δηλαδή
\begin{equation}
\det(A^{i} n_i) = 0.
\end{equation}
Εδώ, ο $A^{i}$ είναι ο $m \times m$ πίνακας με στοιχεία $A^{i \alpha}_\beta$, και η σύμβαση άθροισης εφαρμόζεται στον δείκτη $i$.

Για να κατανοήσουμε τη γεωμετρική σημασία αυτής της συνθήκης, ας εξετάσουμε τι συμβαίνει όταν προσπαθούμε να επιλύσουμε το σύστημα ως προς τις παραγώγους στην κάθετη κατεύθυνση. Εισάγουμε ένα σύστημα συντεταγμένων προσαρμοσμένο στην $\Sigma$, όπου η $\Sigma$ δίνεται από $x^n = 0$ και η κάθετη κατεύθυνση είναι η $\partial / \partial x^n$. Τότε η κάθετη 1-μορφή είναι $n = dx^n$, δηλαδή $n_i = \delta^n_i$. Σε αυτές τις συντεταγμένες, το σύστημα γράφεται
\begin{equation}
A^{n} \frac{\partial u}{\partial x^n} + \sum_{a=1}^{n-1} A^{a} \frac{\partial u}{\partial x^a} = F.
\end{equation}
Για να μπορούμε να επιλύσουμε ως προς $\partial u / \partial x^n$, ο πίνακας $A^n$ πρέπει να είναι αντιστρέψιμος. Η συνθήκη $\det(A^{i} n_i) = 0$ σημαίνει ακριβώς ότι ο $A^n$ είναι ιδιόμορφος, άρα δεν μπορούμε να επιλύσουμε μοναδικά ως προς την κάθετη παράγωγο. Με άλλα λόγια, τα αρχικά δεδομένα πάνω σε μια χαρακτηριστική επιφάνεια δεν επαρκούν για τον προσδιορισμό της εξέλιξης έξω από αυτήν — η πληροφορία παραμένει "παγιδευμένη" κατά μήκος της επιφάνειας.

Για τη βαθμωτή περίπτωση ($m=1$), ο πίνακας $A^i$ είναι απλώς ένα διάνυσμα, και η χαρακτηριστική συνθήκη γίνεται
\begin{equation}
A^i n_i = 0.
\end{equation}
Αυτό σημαίνει ότι το διανυσματικό πεδίο $A = A^i \partial / \partial x^i$ είναι εφαπτόμενο στην $\Sigma$. Η παράγωγος Lie $\mathcal{L}_A u$ είναι τότε μια εσωτερική παράγωγος της $\Sigma$ και δεν περιέχει πληροφορία για την εξέλιξη κάθετα σε αυτήν.

\subsection{Το θεώρημα Cauchy-Kowalevskaya σε αυστηρή διατύπωση}

Το θεώρημα Cauchy-Kowalevskaya ασχολείται με συστήματα που είναι γραμμένα σε \textbf{κανονική μορφή} ως προς μια επιλεγμένη μεταβλητή, συνήθως τον χρόνο. Η διατύπωση έχει ως εξής.

\textbf{Θεώρημα (Cauchy-Kowalevskaya).} Θεωρούμε ένα σύστημα $m$ μερικών διαφορικών εξισώσεων πρώτης τάξης για $m$ άγνωστες συναρτήσεις $u^1, \dots, u^m$ των μεταβλητών $x^1, \dots, x^{n-1}, t$, γραμμένο στη μορφή
\begin{equation}
\frac{\partial u^\alpha}{\partial t} = F^\alpha \left( t, x^1, \dots, x^{n-1}, u^1, \dots, u^m, \frac{\partial u^1}{\partial x^1}, \dots, \frac{\partial u^m}{\partial x^{n-1}} \right),
\end{equation}
για $\alpha = 1, \dots, m$, με αρχικές συνθήκες
\begin{equation}
u^\alpha(0, x^1, \dots, x^{n-1}) = u_0^\alpha(x^1, \dots, x^{n-1}),
\end{equation}
όπου οι $F^\alpha$ είναι αναλυτικές συναρτήσεις όλων των ορισμάτων τους σε μια περιοχή του σημείου αναφοράς, και οι $u_0^\alpha$ είναι αναλυτικές συναρτήσεις των χωρικών μεταβλητών σε μια περιοχή του σημείου αναφοράς. Τότε υπάρχει μια μοναδική λύση $u^\alpha(t, x^1, \dots, x^{n-1})$, αναλυτική σε μια γειτονιά του σημείου $(0, x_0^1, \dots, x_0^{n-1})$, που ικανοποιεί το σύστημα και τις αρχικές συνθήκες.

Η γεωμετρική σημασία της κανονικής μορφής είναι ότι η υπερεπιφάνεια $t = 0$ (η αρχική επιφάνεια) είναι \textbf{μη χαρακτηριστική}. Πράγματι, σε αυτή τη μορφή, ο πίνακας που πολλαπλασιάζει την παράγωγο ως προς $t$ είναι ο ταυτοτικός, ο οποίος είναι προφανώς αντιστρέψιμος. Άρα το θεώρημα εγγυάται ύπαρξη και μοναδικότητα ακριβώς όταν η αρχική επιφάνεια είναι μη χαρακτηριστική.

\subsection{Σύνδεση με την παράγωγο Lie: ανάλυση και πράξεις}

Ας αναλύσουμε τώρα τη σύνδεση με την παράγωγο Lie μέσω συγκεκριμένων πράξεων. Θεωρούμε την απλούστερη μη τετριμμένη περίπτωση: μια βαθμωτή εξίσωση πρώτης τάξης
\[
\mathcal{L}_\xi u = F(x, u),
\]
όπου $\xi$ είναι ένα δεδομένο διανυσματικό πεδίο. Σε τοπικές συντεταγμένες, αυτή γράφεται
\begin{equation}
\xi^i \frac{\partial u}{\partial x^i} = F(x, u).
\end{equation}
Υποθέτουμε ότι η αρχική επιφάνεια $\Sigma$ δίνεται από $\phi(x) = 0$ και ότι το $\xi$ είναι εγκάρσιο στην $\Sigma$, δηλαδή $\xi^i \partial \phi / \partial x^i \neq 0$ πάνω στην $\Sigma$. Αυτή είναι ακριβώς η μη χαρακτηριστική συνθήκη:
\begin{equation}
\mathcal{L}_\xi \phi \neq 0.
\end{equation}

Για να επιλύσουμε το πρόβλημα Cauchy, εισάγουμε προσαρμοσμένες συντεταγμένες. Επιλέγουμε $n-1$ συναρτήσεις $y^1, \dots, y^{n-1}$ τέτοιες ώστε οι $(y^a, \phi)$ να αποτελούν σύστημα συντεταγμένων. Σε αυτές τις συντεταγμένες, η $\Sigma$ είναι $\phi = 0$ και το $\xi$ έχει συνιστώσα κατά μήκος της $\partial / \partial \phi$ μη μηδενική. Το σύστημα γίνεται
\begin{equation}
\xi^\phi \frac{\partial u}{\partial \phi} + \sum_{a=1}^{n-1} \xi^a \frac{\partial u}{\partial y^a} = F,
\end{equation}
και μπορούμε να επιλύσουμε ως προς $\partial u / \partial \phi$:
\begin{equation}
\frac{\partial u}{\partial \phi} = \frac{1}{\xi^\phi} \left( F - \sum_{a=1}^{n-1} \xi^a \frac{\partial u}{\partial y^a} \right).
\end{equation}
Αυτή είναι ακριβώς η κανονική μορφή Cauchy, και το θεώρημα Cauchy-Kowalevskaya εφαρμόζεται.

Η παράγωγος Lie εμφανίζεται εδώ με δύο ρόλους. Πρώτον, η ίδια η εξίσωση είναι $\mathcal{L}_\xi u = F$. Δεύτερον, η μη χαρακτηριστική συνθήκη εκφράζεται ως $\mathcal{L}_\xi \phi \neq 0$, που σημαίνει ότι η ροή του $\xi$ διασχίζει την $\Sigma$. Για να δούμε πώς η ροή κατασκευάζει τη λύση, θεωρούμε την απειροστική μεταβολή της $u$ κατά μήκος της ροής του $\xi$. Από τον ορισμό της παραγώγου Lie,
\begin{equation}
\frac{d}{dt} u(\phi_t(p)) = \mathcal{L}_\xi u (\phi_t(p)).
\end{equation}
Αν $\mathcal{L}_\xi u = F$, τότε η μεταβολή της $u$ κατά μήκος της ροής είναι πλήρως καθορισμένη από τα δεδομένα. Ολοκληρώνοντας κατά μήκος των ολοκληρωτικών καμπυλών του $\xi$ που ξεκινούν από την $\Sigma$, κατασκευάζουμε τη λύση σε μια ολόκληρη γειτονιά. Αυτή ακριβώς είναι η γεωμετρική διαδικασία πίσω από την απόδειξη του θεωρήματος: η ροή του $\xi$ "σαρώνει" μια περιοχή της $\Sigma$ και μεταφέρει τα αρχικά δεδομένα.

Για ένα σύστημα $m$ εξισώσεων, η γενίκευση περιλαμβάνει τον πίνακα $A^i$. Αν ορίσουμε τον "πίνακα-διανυσματικό τελεστή"
\begin{equation}
\mathcal{L}_A = A^i \frac{\partial}{\partial x^i},
\end{equation}
που δρα στο διάνυσμα στήλη $u = (u^\alpha)$, το σύστημα γράφεται
\begin{equation}
\mathcal{L}_A u = F.
\end{equation}
Αυτός ο τελεστής είναι μια γενίκευση της παραγώγου Lie για συστήματα. Η μη χαρακτηριστική συνθήκη $\det(A^i n_i) \neq 0$ εξασφαλίζει ότι ο τελεστής $\mathcal{L}_A$ περιέχει μια εγκάρσια παράγωγο που μπορεί να αντιστραφεί.

\subsection{Χαρακτηριστικές επιφάνειες και διάδοση πληροφορίας}

Η έννοια της χαρακτηριστικής επιφάνειας έχει μια βαθιά φυσική ερμηνεία: είναι οι επιφάνειες κατά μήκος των οποίων διαδίδεται η πληροφορία. Ας το δούμε αυτό μέσω της παραγώγου Lie. Σε μια χαρακτηριστική επιφάνεια, το διανυσματικό πεδίο $A = A^i \partial / \partial x^i$ είναι εφαπτόμενο στην $\Sigma$. Αυτό σημαίνει ότι η ροή του $A$ παραμένει πάνω στην $\Sigma$ — δεν μπορεί να τη διασχίσει. Κατά συνέπεια, η εξίσωση $\mathcal{L}_A u = F$ δεν μπορεί να μεταφέρει πληροφορία από την $\Sigma$ στο εξωτερικό της· όλη η δυναμική περιορίζεται στην ίδια την επιφάνεια.

Στην κυματική εξίσωση, για παράδειγμα, οι χαρακτηριστικές επιφάνειες είναι οι κώνοι φωτός — οι επιφάνειες που ορίζονται από
\begin{equation}
g^{ij} n_i n_j = 0
\end{equation}
στη μετρική Lorentz. Πάνω σε αυτές τις επιφάνειες, οι εξισώσεις δεν μπορούν να προσδιορίσουν μοναδικά την εξέλιξη, και αυτό ακριβώς επιτρέπει τη διάδοση ασυνεχειών (κυμάτων) κατά μήκος των χαρακτηριστικών. Η παράγωγος Lie κατά μήκος ενός φωτοειδούς διανύσματος Killing είναι ο τελεστής που περιγράφει αυτή τη διάδοση.

\subsection{Φυσικογεωμετρική ερμηνεία και συμμετρίες}

Η σύνδεση με τις συμμετρίες είναι άμεση. Αν $\mathcal{L}_\xi u = 0$, τότε το $\xi$ είναι μια συμμετρία της λύσης: η $u$ είναι σταθερή κατά μήκος των ολοκληρωτικών καμπυλών του $\xi$. Το πρόβλημα Cauchy σε αυτή την περίπτωση είναι τετριμμένο: η λύση είναι απλώς η μεταφορά των αρχικών δεδομένων κατά μήκος της ροής.

Αν $\mathcal{L}_\xi u = F \neq 0$, η $u$ δεν είναι πλέον σταθερή κατά μήκος της ροής, αλλά ο ρυθμός μεταβολής της είναι γνωστός. Το θεώρημα Cauchy-Kowalevskaya εγγυάται ότι αυτή η γνώση, μαζί με τα αρχικά δεδομένα, αρκεί για τον πλήρη προσδιορισμό της $u$. Η γεωμετρική εικόνα είναι ότι ο παρατηρητής που κινείται κατά μήκος της ροής του $\xi$ μπορεί να προβλέψει την εξέλιξη του συστήματος, δεδομένων των αρχικών μετρήσεων στην $\Sigma$ και της γνώσης του ρυθμού μεταβολής $F$.

Συνοψίζοντας, οι θεμελιώδεις τανυστές που εμπλέκονται στη θεωρία Cauchy είναι το διανυσματικό πεδίο $\xi$ τύπου $(1,0)$ που ορίζει την κατεύθυνση εξέλιξης, η 1-μορφή $d\phi$ τύπου $(0,1)$ που ορίζει την αρχική επιφάνεια, ο τανυστής $A^{i \alpha}_\beta$ τύπου $(1,1)$ στον χώρο των αγνώστων και $(1,0)$ στον χωροχρόνο για συστήματα, και η βαθμωτή συνάρτηση $u$ τύπου $(0,0)$ ή το διανυσματικό πεδίο $u^\alpha$ με εσωτερικό δείκτη ως άγνωστη. Η παράγωγος Lie $\mathcal{L}_\xi$, που δρα σε βαθμωτές, διανύσματα ή γενικευμένα διανύσματα, είναι το κεντρικό αναλυτικό εργαλείο. Το θεώρημα Cauchy-Kowalevskaya, διατυπωμένο γεωμετρικά, είναι η απόδειξη ότι η ροή ενός μη χαρακτηριστικού διανυσματικού πεδίου μπορεί να "σαρώσει" μια ολόκληρη γειτονιά μιας αρχικής επιφάνειας και να κατασκευάσει μοναδική λύση, δεδομένων αναλυτικών αρχικών δεδομένων. Η παράγωγος Lie είναι ο απειροστικός γεννήτορας αυτής της ροής, και η μελέτη των χαρακτηριστικών επιφανειών είναι η μελέτη των συνόλων όπου αυτή η διαδικασία αποτυγχάνει να είναι μοναδική — ακριβώς επειδή η ροή εφάπτεται στην επιφάνεια και δεν μπορεί να τη διασχίσει.

\subsection{Αρχετυπικά παραδείγματα από τη μαθηματική φυσική}

Η θεωρία Cauchy βρίσκει τις πιο εντυπωσιακές της εφαρμογές στις θεμελιώδεις εξισώσεις της μαθηματικής φυσικής. Ας εξετάσουμε τρία αρχετυπικά παραδείγματα, αναλύοντας το καθένα με τα γεωμετρικά εργαλεία που αναπτύξαμε.

\subsubsection{Η εξίσωση της θερμότητας}

Η εξίσωση της θερμότητας σε μία χωρική διάσταση είναι η απλούστερη παραβολική εξίσωση,
\begin{equation}
\frac{\partial u}{\partial t} = \frac{\partial^2 u}{\partial x^2}.
\end{equation}
Αυτή είναι μια εξίσωση δεύτερης τάξης, αλλά μπορεί να αναχθεί σε σύστημα πρώτης τάξης εισάγοντας μια βοηθητική μεταβλητή. Θέτουμε $v = \partial u / \partial x$, οπότε η εξίσωση γίνεται
\[
\frac{\partial u}{\partial t} = \frac{\partial v}{\partial x}, \qquad \frac{\partial u}{\partial x} = v.
\]
Σε μορφή πίνακα, αυτό γράφεται ως
\begin{equation}
\begin{pmatrix} 1 & 0 \\ 0 & 0 \end{pmatrix} \frac{\partial}{\partial t} \begin{pmatrix} u \\ v \end{pmatrix} + \begin{pmatrix} 0 & -1 \\ -1 & 0 \end{pmatrix} \frac{\partial}{\partial x} \begin{pmatrix} u \\ v \end{pmatrix} = \begin{pmatrix} 0 \\ -v \end{pmatrix}.
\end{equation}
Ο πίνακας που πολλαπλασιάζει την παράγωγο ως προς $t$ είναι ιδιόμορφος ($\det = 0$). Αυτό σημαίνει ότι η επιφάνεια $t = \text{σταθερά}$ είναι χαρακτηριστική. Πράγματι, η εξίσωση της θερμότητας δεν έχει πεπερασμένη ταχύτητα διάδοσης — η πληροφορία διαδίδεται ακαριαία σε όλο τον χώρο. Γεωμετρικά, η απουσία μιας μη χαρακτηριστικής χρονικής κατεύθυνσης σημαίνει ότι δεν μπορούμε να εφαρμόσουμε το θεώρημα Cauchy-Kowalevskaya με αρχική επιφάνεια $t = 0$ στην παραπάνω μορφή. Αντίθετα, αν θεωρήσουμε το πρόβλημα αρχικών τιμών με δεδομένο το $u$ και τις χωρικές του παραγώγους στην $t=0$, η εξίσωση γίνεται μη χαρακτηριστική ως προς τη χωρική κατεύθυνση αν την αντιστρέψουμε. Αυτό αντικατοπτρίζει τη θεμελιώδη ασυμμετρία ανάμεσα στον χρόνο και τον χώρο στην παραβολική διάδοση.

\subsubsection{Η κυματική εξίσωση}

Η κυματική εξίσωση σε μία χωρική διάσταση,
\begin{equation}
\frac{\partial^2 u}{\partial t^2} - \frac{\partial^2 u}{\partial x^2} = 0,
\end{equation}
ανάγεται σε σύστημα πρώτης τάξης θέτοντας $v = \partial u / \partial t$ και $w = \partial u / \partial x$. Το σύστημα γίνεται
\begin{equation}
\frac{\partial v}{\partial t} - \frac{\partial w}{\partial x} = 0, \qquad \frac{\partial w}{\partial t} - \frac{\partial v}{\partial x} = 0,
\end{equation}
με τη συμβατότητα $\partial u / \partial t = v$, $\partial u / \partial x = w$. Σε μορφή πίνακα,
\begin{equation}
\frac{\partial}{\partial t} \begin{pmatrix} v \\ w \end{pmatrix} + \begin{pmatrix} 0 & -1 \\ -1 & 0 \end{pmatrix} \frac{\partial}{\partial x} \begin{pmatrix} v \\ w \end{pmatrix} = 0.
\end{equation}
Εδώ ο πίνακας που πολλαπλασιάζει τη χρονική παράγωγο είναι ο ταυτοτικός, άρα η επιφάνεια $t = 0$ είναι μη χαρακτηριστική. Το θεώρημα Cauchy-Kowalevskaya εφαρμόζεται: δεδομένων αναλυτικών αρχικών συνθηκών $u(0,x)$ και $\partial u / \partial t (0,x)$, υπάρχει μοναδική αναλυτική λύση. Οι χαρακτηριστικές επιφάνειες εδώ είναι οι γραμμές $t \pm x = \text{σταθερά}$, που αντιστοιχούν σε διάδοση με ταχύτητα φωτός. Αυτές είναι ακριβώς οι επιφάνειες κατά μήκος των οποίων διαδίδονται οι ασυνέχειες. Στη γεωμετρική γλώσσα της παραγώγου Lie, η κυματική εξίσωση μπορεί να γραφεί ως
\begin{equation}
(\mathcal{L}_{\partial_t}^2 - \mathcal{L}_{\partial_x}^2) u = 0,
\end{equation}
αναδεικνύοντας ότι η εξέλιξη καθορίζεται από δύο ανταγωνιστικές ροές κατά μήκος των χαρακτηριστικών κατευθύνσεων.

\subsubsection{Οι εξισώσεις Maxwell}

Οι εξισώσεις Maxwell στο κενό, σε μορφή πρώτης τάξης, είναι
\begin{equation}
\frac{\partial E^i}{\partial t} = \epsilon^{ijk} \frac{\partial B^k}{\partial x^j}, \qquad \frac{\partial B^i}{\partial t} = -\epsilon^{ijk} \frac{\partial E^k}{\partial x^j},
\end{equation}
μαζί με τις συνθήκες απόκλισης
\begin{equation}
\frac{\partial E^i}{\partial x^i} = 0, \qquad \frac{\partial B^i}{\partial x^i} = 0.
\end{equation}
Εδώ, $E^i$ και $B^i$ είναι οι συνιστώσες του ηλεκτρικού και του μαγνητικού πεδίου, και $\epsilon^{ijk}$ είναι το σύμβολο Levi-Civita. Το σύστημα αυτό είναι έξι εξισώσεων πρώτης τάξης για έξι αγνώστους. Σε μορφή πίνακα,
\begin{equation}
\frac{\partial}{\partial t} \begin{pmatrix} E \\ B \end{pmatrix} + \begin{pmatrix} 0 & -\epsilon^{ijk} \frac{\partial}{\partial x^j} \\ \epsilon^{ijk} \frac{\partial}{\partial x^j} & 0 \end{pmatrix} \begin{pmatrix} E \\ B \end{pmatrix} = 0.
\end{equation}
Ο πίνακας που πολλαπλασιάζει τη χρονική παράγωγο είναι ο ταυτοτικός, άρα η επιφάνεια $t = 0$ είναι μη χαρακτηριστική. Το θεώρημα Cauchy-Kowalevskaya εγγυάται ύπαρξη και μοναδικότητα της λύσης για αναλυτικά αρχικά δεδομένα.

Οι χαρακτηριστικές επιφάνειες των εξισώσεων Maxwell είναι ακριβώς οι κώνοι φωτός. Πράγματι, αναζητώντας επιφάνειες $\phi(t, x) = 0$ για τις οποίες ο πίνακας των συντελεστών της κάθετης παραγώγου γίνεται ιδιόμορφος, βρίσκουμε τη συνθήκη
\begin{equation}
\left( \frac{\partial \phi}{\partial t} \right)^2 - \sum_{i=1}^3 \left( \frac{\partial \phi}{\partial x^i} \right)^2 = 0.
\end{equation}
Αυτή είναι ακριβώς η εξίσωση του κώνου φωτός. Γεωμετρικά, αυτό σημαίνει ότι οι ασυνέχειες του ηλεκτρομαγνητικού πεδίου διαδίδονται κατά μήκος φωτοειδών κατευθύνσεων. Σε τανυστική γλώσσα, ο τανυστής ηλεκτρομαγνητικού πεδίου $F_{ij}$ ικανοποιεί την εξίσωση κύματος $\square F_{ij} = 0$, και οι χαρακτηριστικές επιφάνειες είναι εκείνες για τις οποίες η παράγωγος Lie κατά μήκος ενός φωτοειδούς διανύσματος Killing δεν μπορεί να αντιστραφεί. Αυτή η γεωμετρική εικόνα ενοποιεί τη διάδοση του φωτός με την έννοια των χαρακτηριστικών επιφανειών: το φως διαδίδεται ακριβώς κατά μήκος των επιφανειών όπου οι εξισώσεις Maxwell χάνουν την προβλεπτική τους ισχύ, επιτρέποντας έτσι τη μεταφορά πληροφορίας χωρίς προκαθορισμό από τα αρχικά δεδομένα.

\section{Το Γινόμενο Lie}

\subsection{Ορισμός ως ειδική περίπτωση της παραγώγου Lie}

Στην ενότητα 12 ορίσαμε την παράγωγο Lie ενός τανυστικού πεδίου $T$ κατά μήκος ενός διανυσματικού πεδίου $\xi$. Μια ιδιαίτερα σημαντική ειδική περίπτωση προκύπτει όταν το ίδιο το $T$ είναι ένα άλλο διανυσματικό πεδίο, έστω $\eta$. Τότε η παράγωγος Lie του $\eta$ κατά μήκος του $\xi$ ονομάζεται \textbf{γινόμενο Lie} (ή αγκύλη Lie) των $\xi$ και $\eta$, και συμβολίζεται με $[\xi, \eta]$. Από την εξίσωση (54) έχουμε αμέσως την έκφραση σε τοπικές συντεταγμένες:
\begin{equation}
[\xi, \eta]^i = (\mathcal{L}_\xi \eta)^i = \xi^\alpha \frac{\partial \eta^i}{\partial x^\alpha} - \eta^\alpha \frac{\partial \xi^i}{\partial x^\alpha}.
\end{equation}
Το γινόμενο Lie είναι λοιπόν ένα νέο διανυσματικό πεδίο, που προκύπτει από τα $\xi$ και $\eta$ μέσω μιας διαδικασίας που περιλαμβάνει τόσο τις τιμές των πεδίων όσο και τις παραγώγους τους. Η έκφραση αυτή αναδεικνύει την αντισυμμετρία του γινομένου:
\begin{equation}
[\xi, \eta] = -[\eta, \xi],
\end{equation}
η οποία είναι άμεση συνέπεια του ορισμού. Το γινόμενο Lie ικανοποιεί επίσης την ταυτότητα Jacobi, που είναι η θεμελιώδης αλγεβρική ιδιότητα που το καθιστά μια άλγεβρα Lie στον χώρο των διανυσματικών πεδίων:
\begin{equation}
[\xi, [\eta, \zeta]] + [\eta, [\zeta, \xi]] + [\zeta, [\xi, \eta]] = 0.
\end{equation}
Η ταυτότητα αυτή δεν είναι προφανής από τον ορισμό, αλλά αποδεικνύεται με άμεσο υπολογισμό σε τοπικές συντεταγμένες. Γεωμετρικά, η ταυτότητα Jacobi εκφράζει τη συμβατότητα των ροών τριών διανυσματικών πεδίων: η διαδοχική εφαρμογή των ροών με διαφορετική σειρά δεν οδηγεί σε αντιφάσεις.

\subsection{Βασικές ταυτότητες του γινομένου Lie}

Πέρα από την αντισυμμετρία και την ταυτότητα Jacobi, το γινόμενο Lie ικανοποιεί μια σειρά από χρήσιμες ταυτότητες που το συνδέουν με άλλες γεωμετρικές πράξεις. Για κάθε βαθμωτή συνάρτηση $f$, ισχύει ο κανόνας Leibniz:
\begin{equation}
[\xi, f \eta] = (\mathcal{L}_\xi f) \eta + f [\xi, \eta] = (\xi^\alpha \frac{\partial f}{\partial x^\alpha}) \eta + f [\xi, \eta].
\end{equation}
Η ταυτότητα αυτή δείχνει πώς το γινόμενο Lie αλληλεπιδρά με τον πολλαπλασιασμό διανυσματικών πεδίων με βαθμωτές συναρτήσεις. Είναι άμεση συνέπεια του ορισμού και του κανόνα Leibniz για τις μερικές παραγώγους.

Μια άλλη σημαντική ταυτότητα συνδέει το γινόμενο Lie με την παράγωγο Lie μιας 1-μορφής. Για μια 1-μορφή $\omega$ και διανυσματικά πεδία $\xi, \eta$, ισχύει:
\begin{equation}
(\mathcal{L}_\xi \omega)(\eta) = \mathcal{L}_\xi (\omega(\eta)) - \omega([\xi, \eta]).
\end{equation}
Η ταυτότητα αυτή είναι ουσιαστικά ο ορισμός της παραγώγου Lie για 1-μορφές μέσω της δράσης της σε διανύσματα.

Για το γινόμενο Lie δύο διανυσματικών πεδίων και μιας βαθμωτής συνάρτησης, έχουμε την ταυτότητα:
\begin{equation}
[\xi, \eta](f) = \xi(\eta(f)) - \eta(\xi(f)),
\end{equation}
όπου $\xi(f) = \mathcal{L}_\xi f$ είναι η κατευθυνόμενη παράγωγος. Η ταυτότητα αυτή είναι ακριβώς η σύνδεση του γινομένου Lie με τον μεταθέτη των τελεστών διαφόρισης: το γινόμενο Lie μετρά την αποτυχία των δύο διανυσματικών πεδίων να αντιμετατίθενται ως διαφορικοί τελεστές.

\subsection{Γεωμετρική ερμηνεία: η μη-αντιμεταθετικότητα των ροών}

Η βαθύτερη γεωμετρική σημασία του γινομένου Lie αποκαλύπτεται όταν το συνδέσουμε με τις ροές των διανυσματικών πεδίων. Θεωρούμε δύο διανυσματικά πεδία $\xi$ και $\eta$ σε μια πολλαπλότητα $M$, με αντίστοιχες ροές $\phi_t$ και $\psi_s$. Το ερώτημα που τίθεται φυσικά είναι: αν ξεκινήσουμε από ένα σημείο $p$, μετακινηθούμε κατά μήκος της ροής του $\xi$ για χρόνο $t$ και μετά κατά μήκος της ροής του $\eta$ για χρόνο $s$, θα φτάσουμε στο ίδιο σημείο που θα φτάναμε αν αντιστρέφαμε τη σειρά των κινήσεων; Με άλλα λόγια, αντιμετατίθενται οι ροές $\phi_t$ και $\psi_s$;

Για να απαντήσουμε, υπολογίζουμε τη διαφορά ανάμεσα στις δύο διαδρομές. Ξεκινώντας από το $p$, η πρώτη διαδρομή μας πηγαίνει στο $\psi_s(\phi_t(p))$, ενώ η δεύτερη στο $\phi_t(\psi_s(p))$. Αναπτύσσουμε τις ροές σε σειρά Taylor γύρω από το $p$. Για τη ροή $\phi_t$, έχουμε:
\begin{equation}
\phi_t^i(p) = x^i + t \xi^i(p) + \frac{t^2}{2} \xi^\alpha(p) \frac{\partial \xi^i}{\partial x^\alpha}(p) + \mathcal{O}(t^3).
\end{equation}
Το ανάπτυγμα αυτό προκύπτει από την εξίσωση της ροής: η πρώτη παράγωγος είναι το ίδιο το διανυσματικό πεδίο, και η δεύτερη παράγωγος προκύπτει από την παραγώγιση της $\xi^i(\phi_t(p))$ ως προς $t$. Αντίστοιχα, για τη ροή $\psi_s$:
\begin{equation}
\psi_s^i(q) = y^i + s \eta^i(q) + \frac{s^2}{2} \eta^\alpha(q) \frac{\partial \eta^i}{\partial x^\alpha}(q) + \mathcal{O}(s^3).
\end{equation}

Τώρα υπολογίζουμε τη σύνθεση $\psi_s \circ \phi_t$ στο σημείο $p$. Αντικαθιστώντας το $q = \phi_t(p)$ και αναπτύσσοντας μέχρι όρους τάξης $t^2$, $s^2$ και $ts$, λαμβάνουμε:
\begin{equation}
\psi_s^i(\phi_t(p)) = x^i + t \xi^i + s \eta^i + \frac{t^2}{2} \xi^\alpha \frac{\partial \xi^i}{\partial x^\alpha} + \frac{s^2}{2} \eta^\alpha \frac{\partial \eta^i}{\partial x^\alpha} + ts \, \xi^\alpha \frac{\partial \eta^i}{\partial x^\alpha} + \cdots
\end{equation}
Αντιστρέφοντας τη σειρά, υπολογίζουμε τη σύνθεση $\phi_t \circ \psi_s$. Το αποτέλεσμα είναι ανάλογο, με εναλλαγή των ρόλων των $\xi$ και $\eta$ και των $t$ και $s$:
\begin{equation}
\phi_t^i(\psi_s(p)) = x^i + t \xi^i + s \eta^i + \frac{t^2}{2} \xi^\alpha \frac{\partial \xi^i}{\partial x^\alpha} + \frac{s^2}{2} \eta^\alpha \frac{\partial \eta^i}{\partial x^\alpha} + ts \, \eta^\alpha \frac{\partial \xi^i}{\partial x^\alpha} + \cdots
\end{equation}

Η διαφορά των δύο εκφράσεων δίνει:
\begin{equation}
\psi_s^i(\phi_t(p)) - \phi_t^i(\psi_s(p)) = ts \left( \xi^\alpha \frac{\partial \eta^i}{\partial x^\alpha} - \eta^\alpha \frac{\partial \xi^i}{\partial x^\alpha} \right) + \mathcal{O}(t^3, s^3, t^2 s, t s^2).
\end{equation}
Ο όρος που πολλαπλασιάζει το $ts$ είναι ακριβώς η $i$-οστή συνιστώσα του γινομένου Lie $[\xi, \eta]$. Έτσι, καταλήγουμε στη θεμελιώδη γεωμετρική σχέση:
\begin{equation}
\psi_s(\phi_t(p)) - \phi_t(\psi_s(p)) = ts \, [\xi, \eta](p) + \text{ανώτεροι όροι}.
\end{equation}

Η σχέση αυτή αποκαλύπτει ότι το γινόμενο Lie $[\xi, \eta]$ είναι το απειροστικό μέτρο της μη-αντιμεταθετικότητας των ροών των $\xi$ και $\eta$. Γεωμετρικά, το παραλληλόγραμμο που σχηματίζεται από τις δύο ροές δεν κλείνει ακριβώς· το "κενό" μετράται από το $[\xi, \eta]$. Η συνθήκη $[\xi, \eta] = 0$ σημαίνει ότι οι ροές αντιμετατίθενται τοπικά: το παραλληλόγραμμο κλείνει, και μπορούμε να ορίσουμε ένα σύστημα συντεταγμένων στο οποίο τα $\xi$ και $\eta$ είναι τα διανυσματικά πεδία των αξόνων. Αυτή είναι η γεωμετρική βάση του θεωρήματος Frobenius.

\subsection{Θεωρία ολοκληρωσιμότητας και το θεώρημα Frobenius}

Το γινόμενο Lie δεν είναι απλώς ένα μέτρο της μη-αντιμεταθετικότητας δύο μεμονωμένων ροών· είναι το κλειδί για την κατανόηση του πότε μια ολόκληρη οικογένεια διανυσματικών πεδίων μπορεί να "υφάνει" μια υποπολλαπλότητα. Αυτή είναι η ουσία της θεωρίας ολοκληρωσιμότητας, και το κεντρικό της αποτέλεσμα είναι το θεώρημα Frobenius.

\subsubsection{Τι σημαίνει ολοκληρωσιμότητα}

Φανταστείτε ότι σε μια πολλαπλότητα $M$ διάστασης $n$ έχουμε $k$ διανυσματικά πεδία $\xi_1, \dots, \xi_k$, γραμμικά ανεξάρτητα σε κάθε σημείο μιας περιοχής. Τα πεδία αυτά ορίζουν σε κάθε σημείο $p$ έναν $k$-διάστατο υπόχωρο του εφαπτόμενου χώρου $T_pM$, το γραμμικό τους άνοιγμα. Μια τέτοια αντιστοίχιση ονομάζεται \textbf{κατανομή} $\mathcal{D}$ διάστασης $k$.

Το ερώτημα της ολοκληρωσιμότητας είναι το εξής: υπάρχουν υποπολλαπλότητες $N$ διάστασης $k$ που είναι παντού εφαπτόμενες στην κατανομή; Δηλαδή, για κάθε σημείο $q \in N$, ισχύει $T_q N = \mathcal{D}_q$; Αν ναι, η κατανομή ονομάζεται \textbf{ολοκληρώσιμη}, και οι υποπολλαπλότητες $N$ ονομάζονται \textbf{ολοκληρωμένες υποπολλαπλότητες} ή \textbf{φύλλα} της κατανομής.

Για να καταλάβουμε τη γεωμετρική σημασία αυτής της έννοιας, ας θυμηθούμε τι συμβαίνει όταν $k=1$. Μια κατανομή διάστασης $1$ είναι απλώς ένα διανυσματικό πεδίο. Το θεώρημα ύπαρξης και μοναδικότητας των συνήθων διαφορικών εξισώσεων μας λέει ότι από κάθε σημείο διέρχεται ακριβώς μία ολοκληρωμένη καμπύλη. Άρα κάθε κατανομή διάστασης $1$ είναι ολοκληρώσιμη.

Όταν $k > 1$, η κατάσταση είναι δραματικά διαφορετική. Δεν είναι καθόλου προφανές ότι οι ολοκληρωτικές καμπύλες των $\xi_a$ μπορούν να συντεθούν σε μια $k$-διάστατη επιφάνεια. Το εμπόδιο είναι ακριβώς η μη-αντιμεταθετικότητα των ροών που αναλύσαμε προηγουμένως. Αν οι ροές δύο πεδίων $\xi_a$ και $\xi_b$ δεν αντιμετατίθενται, τότε το "παραλληλόγραμμο" που σχηματίζουν δεν κλείνει, και η απόπειρα να χτίσουμε μια επιφάνεια από τις ολοκληρωτικές καμπύλες οδηγεί σε αντιφάσεις. Αντίθετα, αν όλα τα γινόμενα Lie $[\xi_a, \xi_b]$ ανήκουν στην κατανομή, τότε τα "κενά" που δημιουργούν οι μη-αντιμεταθετικές ροές είναι εφαπτόμενα στην κατανομή και μπορούμε να τα "γεμίσουμε" χωρίς να βγούμε από αυτήν.

\subsubsection{Το θεώρημα Frobenius: διατύπωση και γεωμετρική απόδειξη}

\textbf{Θεώρημα (Frobenius).} Έστω $\mathcal{D}$ μια ομαλή κατανομή διάστασης $k$ σε μια πολλαπλότητα $M$ διάστασης $n$. Υποθέτουμε ότι σε μια περιοχή κάθε σημείου, η $\mathcal{D}$ παράγεται από $k$ ομαλά διανυσματικά πεδία $\xi_1, \dots, \xi_k$, γραμμικά ανεξάρτητα σε κάθε σημείο. Τότε η $\mathcal{D}$ είναι ολοκληρώσιμη αν και μόνο αν είναι κλειστή ως προς το γινόμενο Lie: για κάθε $a, b = 1, \dots, k$, υπάρχουν ομαλές συναρτήσεις $C_{ab}^c$ τέτοιες ώστε
\begin{equation}
[\xi_a, \xi_b] = C_{ab}^c \xi_c.
\end{equation}
Αν αυτή η συνθήκη ικανοποιείται, τότε για κάθε σημείο $p$ υπάρχει ένα σύστημα τοπικών συντεταγμένων $(y^1, \dots, y^k, y^{k+1}, \dots, y^n)$ τέτοιο ώστε οι ολοκληρωμένες υποπολλαπλότητες να δίνονται από τις εξισώσεις $y^{k+1} = \text{σταθερά}, \dots, y^n = \text{σταθερά}$, και τα διανυσματικά πεδία $\xi_a$ να είναι εφαπτόμενα σε αυτές.

Ας αναπτύξουμε ένα πλήρες γεωμετρικό επιχείρημα για την απόδειξη της κεντρικής κατεύθυνσης: ότι η συνθήκη συνεπάγεται την ολοκληρωσιμότητα. Το επιχείρημα βασίζεται στην επαναληπτική κατασκευή ενός συστήματος συντεταγμένων, χρησιμοποιώντας διαδοχικά τις ροές των διανυσματικών πεδίων.

Βήμα 1: Επιλέγουμε ένα σημείο αναφοράς $p$ και θεωρούμε το πρώτο διανυσματικό πεδίο $\xi_1$. Η ροή του $\phi^1_{t_1}$ ορίζει μια οικογένεια ολοκληρωτικών καμπυλών. Εισάγουμε την πρώτη συντεταγμένη $y^1 = t_1$ κατά μήκος αυτών των καμπυλών, έτσι ώστε $\xi_1 = \partial / \partial y^1$.

Βήμα 2: Σε κάθε σημείο της ροής του $\xi_1$, θεωρούμε το δεύτερο διανυσματικό πεδίο $\xi_2$. Θα θέλαμε να εισάγουμε τη δεύτερη συντεταγμένη $y^2$ κατά μήκος της ροής του $\xi_2$, έτσι ώστε $\xi_2 = \partial / \partial y^2$. Ωστόσο, για να είναι αυτό συνεπές, οι ροές των $\xi_1$ και $\xi_2$ πρέπει να αντιμετατίθενται όταν κινούμαστε κατά μήκος των αξόνων τους. Αν $[\xi_1, \xi_2] \neq 0$, αυτό δεν ισχύει γενικά. Αλλά εδώ ακριβώς επεμβαίνει η συνθήκη Frobenius: μας επιτρέπει να αντικαταστήσουμε τα αρχικά διανυσματικά πεδία με νέα, που παράγουν την ίδια κατανομή αλλά έχουν μηδενικό γινόμενο Lie — δηλαδή αντιμετατίθενται.

Η κατασκευή των νέων πεδίων γίνεται ως εξής. Από τη συνθήκη, έχουμε
\[
[\xi_1, \xi_2] = C_{12}^1 \xi_1 + C_{12}^2 \xi_2 + \sum_{c=3}^k C_{12}^c \xi_c.
\]
Θεωρούμε τον γραμμικό συνδυασμό $\tilde{\xi}_2 = \xi_2 + f \xi_1$, όπου $f$ είναι μια συνάρτηση που θα προσδιοριστεί. Το γινόμενο Lie με το $\xi_1$ είναι
\[
[\xi_1, \tilde{\xi}_2] = [\xi_1, \xi_2] + [\xi_1, f \xi_1] = [\xi_1, \xi_2] + (\mathcal{L}_{\xi_1} f) \xi_1 + f [\xi_1, \xi_1].
\]
Ο τελευταίος όρος μηδενίζεται ($[\xi_1, \xi_1] = 0$). Αντικαθιστώντας την έκφραση για το $[\xi_1, \xi_2]$, έχουμε
\[
[\xi_1, \tilde{\xi}_2] = (C_{12}^1 + \mathcal{L}_{\xi_1} f) \xi_1 + C_{12}^2 \xi_2 + \sum_{c=3}^k C_{12}^c \xi_c.
\]
Εκφράζοντας το $\xi_2$ ως $\tilde{\xi}_2 - f \xi_1$, βλέπουμε ότι ο όρος $C_{12}^2 \xi_2$ συνεισφέρει έναν όρο ανάλογο του $\tilde{\xi}_2$ και έναν ανάλογο του $\xi_1$. Επιλέγοντας τη συνάρτηση $f$ έτσι ώστε να ικανοποιεί τη διαφορική εξίσωση
\[
\mathcal{L}_{\xi_1} f + C_{12}^1 - f C_{12}^2 = 0,
\]
μπορούμε να εξαλείψουμε τον όρο ανάλογο του $\xi_1$ από το $[\xi_1, \tilde{\xi}_2]$. Το αποτέλεσμα είναι ότι το $[\xi_1, \tilde{\xi}_2]$ περιέχει μόνο όρους ανάλογους των $\tilde{\xi}_2, \xi_3, \dots, \xi_k$. Δεν έχει συνιστώσα κατά μήκος του $\xi_1$.

Συνεχίζοντας επαγωγικά αυτή τη διαδικασία, μπορούμε να κατασκευάσουμε ένα νέο σύνολο διανυσματικών πεδίων $\tilde{\xi}_1, \dots, \tilde{\xi}_k$ που παράγουν την ίδια κατανομή $\mathcal{D}$ και ικανοποιούν
\[
[\tilde{\xi}_a, \tilde{\xi}_b] = 0 \quad \text{για κάθε } a, b = 1, \dots, k.
\]
Με άλλα λόγια, η συνθήκη Frobenius μας επιτρέπει να "ευθυγραμμίσουμε" τα διανυσματικά πεδία ώστε να αντιμετατίθενται όλα μεταξύ τους.

Βήμα 3: Αφού έχουμε $k$ αντιμετατιθέμενα διανυσματικά πεδία $\tilde{\xi}_1, \dots, \tilde{\xi}_k$, μπορούμε να ορίσουμε τις συντεταγμένες $y^1, \dots, y^k$ ως τις παραμέτρους των αντίστοιχων ροών. Επειδή τα πεδία αντιμετατίθενται, η σειρά με την οποία ακολουθούμε τις ροές δεν έχει σημασία — το παραλληλόγραμμο κλείνει. Οι ροές ορίζουν μια $k$-διάστατη επιφάνεια, και οι συντεταγμένες $(y^1, \dots, y^k)$ είναι ένα σύστημα συντεταγμένων πάνω σε αυτήν.

Βήμα 4: Συμπληρώνουμε το σύστημα συντεταγμένων με $n-k$ επιπλέον συναρτήσεις $y^{k+1}, \dots, y^n$ που είναι σταθερές κατά μήκος των $\tilde{\xi}_a$ (δηλαδή $\mathcal{L}_{\tilde{\xi}_a} y^{b} = 0$ για $b > k$). Αυτές ορίζουν τις εγκάρσιες συντεταγμένες. Σε αυτό το σύστημα, οι ολοκληρωμένες υποπολλαπλότητες είναι οι επιφάνειες $y^{k+1} = c^{k+1}, \dots, y^n = c^n$, και τα $\tilde{\xi}_a$ (άρα και τα αρχικά $\xi_a$) είναι εφαπτόμενα σε αυτές.

Αυτή η κατασκευή αποδεικνύει το θεώρημα Frobenius και αποκαλύπτει τον βαθύ γεωμετρικό του πυρήνα: η συνθήκη κλειστότητας ως προς το γινόμενο Lie είναι ακριβώς αυτό που χρειαζόμαστε για να μπορούμε να "ισιώσουμε" την κατανομή και να την εκφράσουμε ως το εφαπτόμενο επίπεδο ενός τοπικού συστήματος συντεταγμένων.

\subsection{Σύνδεση ολοκληρωσιμότητας Frobenius και θεωρίας Cauchy-Kowalevskaya}

Η θεωρία ολοκληρωσιμότητας Frobenius και η θεωρία Cauchy-Kowalevskaya είναι δύο όψεις του ίδιου νομίσματος, διατυπωμένες σε διαφορετικές γλώσσες: η μία στη γλώσσα των διανυσματικών πεδίων και των ροών, η άλλη στη γλώσσα των μερικών διαφορικών εξισώσεων. Η σύνδεσή τους αναδεικνύει μια από τις πιο κομψές ενοποιήσεις της γεωμετρίας και της ανάλυσης.

Για να κατανοήσουμε αυτή τη σχέση, ας επιστρέψουμε στο πρόβλημα Cauchy πρώτης τάξης. Θεωρούμε μια βαθμωτή εξίσωση $\mathcal{L}_\xi u = F(x, u)$ με αρχική συνθήκη $u|_{\Sigma} = u_0$ σε μια υπερεπιφάνεια $\Sigma$. Η μη χαρακτηριστική συνθήκη ήταν $\mathcal{L}_\xi \phi \neq 0$, όπου $\phi = 0$ ορίζει την $\Sigma$. Αυτή η συνθήκη εξασφαλίζει ότι το διανυσματικό πεδίο $\xi$ είναι εγκάρσιο στην $\Sigma$, και το θεώρημα Cauchy-Kowalevskaya εγγυάται την ύπαρξη και μοναδικότητα λύσης.

Από τη σκοπιά της ολοκληρωσιμότητας, εισάγουμε μια βοηθητική άγνωστη συνάρτηση $u^0$ και θεωρούμε το σύστημα
\[
\mathcal{L}_\xi u^0 = 1, \qquad \mathcal{L}_\xi u = F,
\]
με αρχικές συνθήκες $u^0|_{\Sigma} = 0$ και $u|_{\Sigma} = u_0$. Η $u^0$ είναι ουσιαστικά η παράμετρος κατά μήκος των ολοκληρωτικών καμπυλών του $\xi$. Αυτό το σύστημα μπορεί να θεωρηθεί ως ένα σύστημα δύο εξισώσεων πρώτης τάξης, και η μη χαρακτηριστική συνθήκη είναι ότι το $\xi$ δεν εφάπτεται στην $\Sigma$.

Η βαθύτερη σύνδεση εμφανίζεται όταν περάσουμε σε συστήματα περισσότερων εξισώσεων. Ας υποθέσουμε ότι έχουμε $m$ διανυσματικά πεδία $\xi_1, \dots, \xi_m$ σε μια πολλαπλότητα $M$, και θέλουμε να βρούμε $m$ συναρτήσεις $u^1, \dots, u^m$ τέτοιες ώστε
\begin{equation}
\mathcal{L}_{\xi_a} u^\alpha = \delta_a^\alpha, \qquad a, \alpha = 1, \dots, m.
\end{equation}
Αυτό σημαίνει ότι τα $\xi_a$ είναι τα διανυσματικά πεδία των αξόνων ενός συστήματος συντεταγμένων. Το ερώτημα είναι πότε ένα τέτοιο σύστημα υπάρχει. Η απάντηση είναι ακριβώς η συνθήκη Frobenius: τα $\xi_a$ πρέπει να αντιμετατίθενται,
\begin{equation}
[\xi_a, \xi_b] = 0,
\end{equation}
ώστε να μπορούν να αποτελέσουν τους άξονες ενός συστήματος συντεταγμένων. Αν $[\xi_a, \xi_b] \neq 0$, τότε η σειρά με την οποία ολοκληρώνουμε κατά μήκος των $\xi_a$ και $\xi_b$ επηρεάζει το αποτέλεσμα, και δεν μπορούμε να ορίσουμε συνεπείς συντεταγμένες. Αυτό είναι το ανάλογο της συνθήκης ολοκληρωσιμότητας για μεικτές μερικές παραγώγους.

Για ένα γενικότερο σύστημα $\mathcal{L}_{\xi_a} u^\alpha = F_a^\alpha(x, u)$ με αρχικές συνθήκες σε μια υπερεπιφάνεια $\Sigma$ διάστασης $n-m$, το πρόβλημα Cauchy έχει μοναδική λύση αν και μόνο αν η $\Sigma$ είναι εγκάρσια στα $\xi_a$ (μη χαρακτηριστική) και τα $\xi_a$ ικανοποιούν τη συνθήκη Frobenius. Η πρώτη συνθήκη εξασφαλίζει ότι μπορούμε να ξεκινήσουμε την ολοκλήρωση από την $\Sigma$. Η δεύτερη συνθήκη εξασφαλίζει ότι η ολοκλήρωση κατά μήκος των διαφορετικών $\xi_a$ είναι συμβατή — ότι το αποτέλεσμα δεν εξαρτάται από τη σειρά με την οποία ολοκληρώνουμε.

Γεωμετρικά, η εικόνα είναι η εξής: η $\Sigma$ είναι μια "αρχική επιφάνεια" εγκάρσια στην κατανομή που ορίζουν τα $\xi_a$. Η συνθήκη Frobenius εγγυάται ότι η κατανομή είναι ολοκληρώσιμη, δηλαδή ότι υπάρχουν $m$-διάστατες ολοκληρωμένες υποπολλαπλότητες εφαπτόμενες στα $\xi_a$. Αυτές οι υποπολλαπλότητες τέμνουν την $\Sigma$ σε ένα μοναδικό σημείο η καθεμία. Το πρόβλημα Cauchy συνίσταται στην εύρεση αυτών των υποπολλαπλοτήτων και των συναρτήσεων $u^\alpha$ που είναι σταθερές κατά μήκος τους. Η μη χαρακτηριστική συνθήκη εξασφαλίζει ότι κάθε τέτοια υποπολλαπλότητα τέμνει την $\Sigma$ ακριβώς μία φορά, και η συνθήκη Frobenius εξασφαλίζει ότι αυτές οι υποπολλαπλότητες υπάρχουν και είναι καλά ορισμένες.

Συνοψίζοντας, η θεωρία Cauchy-Kowalevskaya και το θεώρημα Frobenius είναι συμπληρωματικές όψεις του ίδιου γεωμετρικού προβλήματος. Και οι δύο ανάγονται στην ίδια θεμελιώδη απαίτηση: την αντιμεταθετικότητα των αντίστοιχων ροών, εκφρασμένη μέσω του μηδενισμού του γινομένου Lie.

\subsection{Φυσικογεωμετρικές εξειδικεύσεις υπό το πρίσμα της ολοκληρωσιμότητας}

Το θεώρημα Frobenius, η θεωρία Cauchy-Kowalevskaya και η έννοια της ολοκληρωσιμότητας βρίσκουν εντυπωσιακές εφαρμογές στις θεμελιώδεις φυσικές θεωρίες.

\subsubsection{Κλασική μηχανική: η αγκύλη Poisson και οι ολοκληρωμένες επιφάνειες}

Στην κλασική μηχανική, ο φασικός χώρος είναι μια συμπλεκτική πολλαπλότητα εφοδιασμένη με τη συμπλεκτική μορφή $\omega$. Για κάθε παρατηρήσιμη συνάρτηση $f$, το Χαμιλτονιανό διανυσματικό πεδίο $\xi_f$ ορίζεται από τη σχέση $\iota_{\xi_f} \omega = -df$, που σε συνιστώσες γράφεται
\begin{equation}
\xi_f^i = \omega^{i\alpha} \frac{\partial f}{\partial x^\alpha}.
\end{equation}
Το γινόμενο Lie δύο Χαμιλτονιανών διανυσματικών πεδίων συνδέεται με την αγκύλη Poisson μέσω της θεμελιώδους ταυτότητας
\begin{equation}
[\xi_f, \xi_g] = -\xi_{\{f, g\}}, \qquad \{f, g\} = \omega^{i\alpha} \frac{\partial f}{\partial x^i} \frac{\partial g}{\partial x^\alpha}.
\end{equation}

Από την οπτική της ολοκληρωσιμότητας, ας θεωρήσουμε $k$ συναρτήσεις $f_1, \dots, f_k$ των οποίων οι αγκύλες Poisson ικανοποιούν $\{f_a, f_b\} = C_{ab}^c f_c$ για κάποιες σταθερές $C_{ab}^c$. Αυτό σημαίνει ότι τα αντίστοιχα Χαμιλτονιανά διανυσματικά πεδία ικανοποιούν $[\xi_{f_a}, \xi_{f_b}] = -C_{ab}^c \xi_{f_c}$, που είναι ακριβώς η συνθήκη Frobenius. Συνεπώς, τα $\xi_{f_1}, \dots, \xi_{f_k}$ ορίζουν μια ολοκληρώσιμη κατανομή, και ο φασικός χώρος φυλλώνεται από $k$-διάστατες υποπολλαπλότητες. Αυτές είναι οι επιφάνειες στάθμης των συναρτήσεων $f_1, \dots, f_k$ (εφόσον είναι σε ενέλιξη), και η ολοκληρωσιμότητά τους είναι η γεωμετρική βάση του θεωρήματος Liouville-Arnold για τα πλήρως ολοκληρώσιμα συστήματα. Στην ειδική περίπτωση που $k = n$ (ο φασικός χώρος έχει διάσταση $2n$), οι επιφάνειες αυτές είναι οι τόροι Liouville, και η ύπαρξή τους επιτρέπει την πλήρη επίλυση των εξισώσεων κίνησης μέσω τετραγωνισμών. Η σύνδεση με τη θεωρία Cauchy-Kowalevskaya είναι άμεση: οι εξισώσεις κίνησης $\dot{f}_a = \{f_a, H\}$ μπορούν να θεωρηθούν ως ένα σύστημα Cauchy, και η ολοκληρωσιμότητα κατά Frobenius των αντίστοιχων Χαμιλτονιανών ροών εξασφαλίζει ότι το σύστημα είναι συμβατό και επιλύσιμο.

\subsubsection{Κβαντομηχανική: ο μεταθέτης και η ταυτόχρονη διαγωνοποίηση}

Στην κβαντομηχανική, η μετάβαση από την αγκύλη Poisson στον μεταθέτη,
\begin{equation}
\{f, g\} \longleftrightarrow \frac{i}{\hbar} [\hat{f}, \hat{g}],
\end{equation}
μεταφράζει τη γεωμετρική συνθήκη ολοκληρωσιμότητας σε αλγεβρική συνθήκη ταυτόχρονης διαγωνοποίησης. Ένα σύνολο τελεστών $\hat{f}_1, \dots, \hat{f}_k$ που αντιμετατίθενται μεταξύ τους, $[\hat{f}_a, \hat{f}_b] = 0$, ικανοποιεί το κβαντικό ανάλογο της συνθήκης Frobenius. Αυτό σημαίνει ότι μπορούν να διαγωνοποιηθούν ταυτόχρονα: υπάρχει μια βάση του χώρου Hilbert στην οποία όλοι οι $\hat{f}_a$ είναι διαγώνιοι. Οι ιδιοτιμές τους παρέχουν ένα πλήρες σύνολο κβαντικών αριθμών που χαρακτηρίζουν την κατάσταση του συστήματος. Αυτή είναι η κβαντική εκδοχή της ολοκληρωσιμότητας: η δυνατότητα ταυτόχρονης μέτρησης των $f_a$ με αυθαίρετη ακρίβεια αντιστοιχεί στη γεωμετρική δυνατότητα φύλλωσης του φασικού χώρου από τις επιφάνειες στάθμης τους. Από τη σκοπιά της θεωρίας Cauchy, η ταυτόχρονη διαγωνοποίηση αντιστοιχεί στην επίλυση ενός συστήματος ιδιοτιμικών εξισώσεων με συμβατές αρχικές συνθήκες — η συνθήκη αντιμεταθετικότητας είναι ακριβώς η συνθήκη συμβατότητας που εξασφαλίζει ότι οι εξισώσεις μπορούν να λυθούν ταυτόχρονα.

\subsubsection{Γενική σχετικότητα: πεδία Killing και ολοκληρωσιμότητα του χωροχρόνου}

Στη γενική σχετικότητα, η άλγεβρα Lie των πεδίων Killing παίζει κεντρικό ρόλο στην ταξινόμηση των χωροχρόνων. Αν ένας χωροχρόνος διαθέτει $k$ πεδία Killing $\xi_1, \dots, \xi_k$ που ικανοποιούν $[\xi_a, \xi_b] = C_{ab}^c \xi_c$, τότε τα πεδία αυτά ορίζουν μια ολοκληρώσιμη κατανομή, και ο χωροχρόνος φυλλώνεται από $k$-διάστατες επιφάνειες που είναι οι τροχιές της ομάδας συμμετρίας. Στην περίπτωση του χωροχρόνου Minkowski, τα δέκα πεδία Killing της άλγεβρας Poincaré ικανοποιούν αυτή τη συνθήκη με τις δομικές σταθερές της ομάδας Lorentz.

Πιο γενικά, η ολοκληρωσιμότητα οικογενειών χρονικών ή φωτοειδών γεωδαισιακών διέπεται από το θεώρημα Frobenius. Μια οικογένεια χρονικών καμπυλών μπορεί να θεωρηθεί ως οι ολοκληρωτικές καμπύλες ενός μοναδιαίου χρονικού διανυσματικού πεδίου $\xi$. Το ερώτημα αν αυτές οι καμπύλες είναι ορθογώνιες σε μια οικογένεια χωρικών υπερεπιφανειών (δηλαδή αν ορίζουν έναν παρατηρητή "χωρίς περιστροφή") ισοδυναμεί με την ολοκληρωσιμότητα της κατανομής που είναι ορθογώνια στο $\xi$. Η συνθήκη Frobenius για αυτή την κατανομή είναι ακριβώς ο μηδενισμός της στροβιλότητας του $\xi$, που εκφράζεται μέσω της παραγώγου Lie της μετρικής κατά μήκος του $\xi$. Οι χωροχρόνοι που επιδέχονται έναν τέτοιο παρατηρητή χωρίς περιστροφή ονομάζονται στατικοί ή στάσιμοι, και η ύπαρξή τους είναι θεμελιώδης για τη διατύπωση της θερμοδυναμικής των μελανών οπών και τη μελέτη των βαρυτικών καταρρεύσεων. Η σύνδεση με τη θεωρία Cauchy-Kowalevskaya είναι εδώ ιδιαίτερα βαθιά: οι εξισώσεις Einstein, όταν διατυπωθούν ως πρόβλημα αρχικών τιμών, απαιτούν την επιλογή μιας χωρικής υπερεπιφάνειας $\Sigma$ και μιας χρονικής κατεύθυνσης εξέλιξης. Η συνθήκη Frobenius για την κατανομή των κάθετων διανυσμάτων στην $\Sigma$ είναι ακριβώς η συνθήκη που εξασφαλίζει ότι η $\Sigma$ μπορεί να ενσωματωθεί σε μια οικογένεια χωρικών υπερεπιφανειών που καλύπτουν τον χωροχρόνο. Χωρίς αυτήν, η διατύπωση του προβλήματος Cauchy για τις εξισώσεις Einstein δεν θα ήταν δυνατή — δεν θα μπορούσαμε καν να ορίσουμε τι σημαίνει "αρχική επιφάνεια" με συνεπή τρόπο. Το γεγονός ότι η γενική σχετικότητα επιδέχεται μια καλώς ορισμένη διατύπωση αρχικών τιμών είναι, σε τελική ανάλυση, μια συνέπεια του θεωρήματος Frobenius και της δομής της άλγεβρας Lie των διανυσματικών πεδίων που εμπλέκονται.

\section{Γεωμετρική Ερμηνεία των Απλών $k$-Μορφών}

\subsection{Από τα διανύσματα στις ορίζουσες: δράση σε $k$-παραλληλεπίπεδα}

Θεωρούμε έναν διανυσματικό χώρο $V$ διάστασης $n$ πάνω από το $\mathbb{R}$, τον οποίο σκεφτόμαστε ως τον εφαπτόμενο χώρο $T_pM$ σε ένα σημείο μιας πολλαπλότητας. Μια \textbf{απλή $k$-μορφή} είναι ένας πλήρως αντισυμμετρικός συναλλοίωτος τανυστής τύπου $(0,k)$ που μπορεί να γραφεί ως το σφηνοειδές γινόμενο $k$ 1-μορφών:
\[
\omega = \omega^1 \wedge \omega^2 \wedge \dots \wedge \omega^k,
\]
όπου κάθε $\omega^a \in V^*$ είναι μια 1-μορφή. Σε ένα σύστημα συντεταγμένων $x^1, \dots, x^n$, οι 1-μορφές αναπαρίστανται ως
\[
\omega^a = \omega^a_j \, dx^j, \qquad a = 1, \dots, k,
\]
όπου $\omega^a_j$ είναι οι συνιστώσες της $\omega^a$ στη βάση $\{dx^j\}$. Η δράση της $\omega$ σε $k$ διανύσματα $\xi_1, \dots, \xi_k \in V$ — δηλαδή η δράση της σε ένα $k$-παραλληλεπίπεδο που ορίζεται από αυτά τα διανύσματα — δίνεται από την ορίζουσα
\begin{equation}
\omega(\xi_1, \dots, \xi_k) = \det \left( \omega^a(\xi_i) \right)_{a,i=1}^k. \qquad (111)
\end{equation}
Αυτή η έκφραση είναι η πεμπτουσία της γεωμετρικής ερμηνείας: η τιμή $\omega(\xi_1, \dots, \xi_k)$ είναι η \textbf{ορίζουσα} ενός $k \times k$ πίνακα, του οποίου τα στοιχεία στη γραμμή $a$ και στήλη $i$ είναι η τιμή της 1-μορφής $\omega^a$ πάνω στο διάνυσμα $\xi_i$. Για να καταλάβουμε τι σημαίνει αυτή η ορίζουσα, πρέπει να προσδιορίσουμε γεωμετρικά τον υπόχωρο στον οποίο προβάλλει η $\omega$ μέσω των διανυσμάτων που είναι δυϊκά στις $\omega^a$.

Εφόσον οι $\omega^1, \dots, \omega^k$ είναι γραμμικά ανεξάρτητες, υπάρχουν $k$ γραμμικά ανεξάρτητα διανύσματα στον $V$, τα οποία συμβολίζουμε με $\partial / \partial \omega^1, \dots, \partial / \partial \omega^k$, τέτοια ώστε
\begin{equation}
\omega^a \left( \frac{\partial}{\partial \omega^b} \right) = \delta^a_b, \qquad a, b = 1, \dots, k. \qquad (114)
\end{equation}
Ο συμβολισμός αυτός είναι το φυσικό ανάλογο των διανυσμάτων βάσης $\partial / \partial x^i$ που είναι δυϊκά στα $dx^j$: όπως $dx^j(\partial / \partial x^i) = \delta^j_i$, έτσι και εδώ $\omega^b(\partial / \partial \omega^a) = \delta^b_a$. Τα διανύσματα $\partial / \partial \omega^a$ είναι τα μοναδικά διανύσματα που είναι δυϊκά στις $\omega^a$ και αποτελούν τη δυϊκή βάση των $\omega^a$ στον υπόχωρο που παράγουν. Ορίζουμε τον $k$-διάστατο υπόχωρο
\begin{equation}
\pi(\omega) = \text{span} \left\{ \frac{\partial}{\partial \omega^1}, \dots, \frac{\partial}{\partial \omega^k} \right\}. \qquad (115)
\end{equation}
Ο $\pi(\omega)$ είναι ακριβώς ο υπόχωρος στον οποίο προβάλλει η $\omega$. Από την ίδια τη δομή της ορίζουσας στην (111), είναι σαφές ότι κάθε γραμμή του πίνακα αντιστοιχεί στην τιμή της $\omega^a$ πάνω στα $\xi_i$. Οι τιμές αυτές δεν είναι παρά οι συντεταγμένες των προβολών των $\xi_i$ πάνω στον $\pi(\omega)$, εκφρασμένες στη βάση $\partial / \partial \omega^1, \dots, \partial / \partial \omega^k$.

Για να κατανοήσουμε ποια διανύσματα "αγνοεί" η προβολή αυτή, ορίζουμε τον \textbf{ριζικό υπόχωρο} (radical subspace) της $\omega$, τον οποίο συμβολίζουμε με $\rho(\omega)$, ως την τομή των πυρήνων των 1-μορφών που την απαρτίζουν:
\begin{equation}
\rho(\omega) = \bigcap_{a=1}^k \ker \omega^a. \qquad (116)
\end{equation}
Ο $\rho(\omega)$ είναι ο μέγιστος υπόχωρος του $V$ με την ιδιότητα ότι η $\omega$ μηδενίζεται όταν οποιοδήποτε από τα ορίσματά της ανήκει σε αυτόν. Ισοδύναμα, είναι το σύνολο των διανυσμάτων $\xi \in V$ για τα οποία η συστολή $\iota_\xi \omega$ είναι ταυτοτικά μηδέν:
\begin{equation}
\rho(\omega) = \{ \xi \in V \mid \iota_\xi \omega = 0 \}. \qquad (117)
\end{equation}
Ο $\rho(\omega)$ είναι ακριβώς ο πυρήνας της προβολής στον $\pi(\omega)$. Πράγματι, αν $\xi \in \rho(\omega)$, τότε $\omega^a(\xi) = 0$ για κάθε $a$, άρα η προβολή του $\xi$ στον $\pi(\omega)$ είναι μηδέν. Αντίστροφα, αν η προβολή του $\xi$ στον $\pi(\omega)$ είναι μηδέν, τότε $\omega^a(\xi) = 0$ για κάθε $a$, άρα $\xi \in \rho(\omega)$. Η ορίζουσα στην (111) είναι αναλλοίωτη υπό την προσθήκη σε οποιαδήποτε στήλη ενός διανύσματος που ανήκει στον $\rho(\omega)$ (αφού τότε η αντίστοιχη στήλη παραμένει αμετάβλητη), και μηδενίζεται αν μια στήλη ανήκει εξ ολοκλήρου στον $\rho(\omega)$. Αυτό σημαίνει ότι η ορίζουσα εξαρτάται μόνο από τις κλάσεις ισοδυναμίας των $\xi_i$ modulo $\rho(\omega)$, δηλαδή από τις προβολές τους στον $\pi(\omega)$.

Συνεπώς, η ορίζουσα στην (111) μετρά τον \textbf{προσανατολισμένο $k$-διάστατο όγκο} του παραλληλεπιπέδου που σχηματίζουν οι προβολές των $\xi_1, \dots, \xi_k$ πάνω στον $\pi(\omega)$, χρησιμοποιώντας ως μονάδα όγκου το παραλληλεπίπεδο που ορίζουν τα ίδια τα διανύσματα βάσης $\partial / \partial \omega^1, \dots, \partial / \partial \omega^k$.

Μια απλή $k$-μορφή, ως αντισυμμετρικός τανυστής, έχει την ιδιότητα να μηδενίζεται όταν τα ορίσματά της είναι γραμμικά εξαρτημένα. Αυτό σημαίνει ότι αν τα διανύσματα $\xi_1, \dots, \xi_k$ δεν παράγουν έναν πλήρη $k$-διάστατο όγκο — δηλαδή αν το $k$-παραλληλεπίπεδο που ορίζουν είναι εκφυλισμένο — τότε η ορίζουσα στην (111) μηδενίζεται. Αλλά η αντισυμμετρία επιτρέπει κάτι ισχυρότερο: η $\omega$ μπορεί να μηδενίζεται ακόμα και για μη εκφυλισμένα $k$-παραλληλεπίπεδα, αν αυτά βρίσκονται εξ ολοκλήρου μέσα στον $\rho(\omega)$ ή τον τέμνουν μη τετριμμένα. Αυτό μας οδηγεί στην έννοια του μηδενοχώρου της μορφής.

\subsection{Ο ριζικός υπόχωρος σε συντεταγμένες και ο μηδενοχώρος}

Ο $\rho(\omega)$ μπορεί να περιγραφεί ρητά μέσω ενός συστήματος γραμμικών εξισώσεων. Ένα διάνυσμα $\xi = \xi^j \, \partial / \partial x^j$ ανήκει στον $\ker \omega^a$ αν και μόνο αν $\omega^a(\xi) = 0$. Χρησιμοποιώντας τη δράση της $\omega^a$ στη βάση $\{\partial / \partial x^j\}$,
\[
\omega^a(\xi) = (\omega^a_j \, dx^j) \left( \xi^i \, \frac{\partial}{\partial x^i} \right) = \omega^a_j \, \xi^i \, dx^j \left( \frac{\partial}{\partial x^i} \right) = \omega^a_j \, \xi^i \, \delta^j_i = \omega^a_j \, \xi^j.
\]
Συνεπώς, η συνθήκη $\omega^a(\xi) = 0$ γράφεται ως $\omega^a_j \, \xi^j = 0$. Για να ανήκει το $\xi$ στην τομή όλων των πυρήνων, πρέπει να ικανοποιεί αυτή τη συνθήκη για κάθε $a = 1, \dots, k$. Συμβολίζοντας τις συνιστώσες του διανύσματος με $x^j$ αντί για $\xi^j$, ο ριζικός υπόχωρος $\rho(\omega)$ περιγράφεται από το ομογενές γραμμικό σύστημα
\begin{equation}
\omega^a_j \, x^j = 0, \qquad a = 1, \dots, k. \qquad (118)
\end{equation}
Αυτό είναι ένα σύστημα $k$ γραμμικών εξισώσεων με $n$ αγνώστους $x^1, \dots, x^n$. Εφόσον οι $\omega^a$ είναι γραμμικά ανεξάρτητες, ο $k \times n$ πίνακας $(\omega^a_j)$ έχει βαθμό $k$, και ο χώρος λύσεων $\rho(\omega)$ έχει διάσταση $n-k$. Τα διανύσματα που ανήκουν στον $\rho(\omega)$ μηδενίζονται από όλες τις $\omega^a$, και συνεπώς η $\omega$ μηδενίζεται αν οποιοδήποτε από τα ορίσματά της ανήκει στον $\rho(\omega)$.

Ορίζουμε τώρα τον \textbf{μηδενοχώρο} της $\omega$, και τον συμβολίζουμε με $\text{null}(\omega)$, ως το σύνολο όλων των $k$-διάστατων υποχώρων $U \subseteq V$ τέτοιων ώστε η $\omega$ να μηδενίζεται σε κάθε $k$-παραλληλεπίπεδο που περιέχεται στον $U$. Ισοδύναμα, ένας $k$-διάστατος υπόχωρος $U$ ανήκει στον $\text{null}(\omega)$ αν και μόνο αν για οποιαδήποτε βάση $\eta_1, \dots, \eta_k$ του $U$, ισχύει $\omega(\eta_1, \dots, \eta_k) = 0$.

Λόγω της πολυγραμμικότητας και της αντισυμμετρίας, ο $\text{null}(\omega)$ καθορίζεται πλήρως από τον $\rho(\omega)$. Συγκεκριμένα, ένας $k$-διάστατος υπόχωρος $U$ ανήκει στον $\text{null}(\omega)$ αν και μόνο αν η τομή του με τον $\rho(\omega)$ είναι μη τετριμμένη:
\begin{equation}
U \in \text{null}(\omega) \iff U \cap \rho(\omega) \neq \{0\}. \qquad (119)
\end{equation}
Αυτή η συνθήκη επιδέχεται μια ιδιαίτερα διαφωτιστική γεωμετρική αναδιατύπωση σε όρους προβολών. Η δράση της $\omega$ σε ένα $k$-παραλληλεπίπεδο που ορίζεται από τα $\xi_1, \dots, \xi_k$ είναι η ορίζουσα των προβολών τους στον $\pi(\omega)$. Αυτή η ορίζουσα μηδενίζεται ακριβώς όταν οι προβολές είναι γραμμικά εξαρτημένες, δηλαδή όταν ο $k$-διάστατος υπόχωρος $U$ που παράγουν τα $\xi_1, \dots, \xi_k$ έχει προβολή στον $\pi(\omega)$ διάστασης μικρότερης του $k$. Ισοδύναμα, η συνθήκη $\dim(\text{proj}_{\pi(\omega)}(U)) < k$ σημαίνει ότι ο πυρήνας της προβολής — που είναι ακριβώς ο $\rho(\omega)$ — έχει μη τετριμμένη τομή με τον $U$, δηλαδή $U \cap \rho(\omega) \neq \{0\}$. Έτσι, ο μηδενοχώρος $\text{null}(\omega)$ μπορεί να χαρακτηριστεί ισοδύναμα ως εξής:
\begin{equation}
\text{null}(\omega) = \{ U \subseteq V \mid \dim U = k \text{ και } \dim(\text{proj}_{\pi(\omega)}(U)) < k \}. \qquad (120)
\end{equation}
Με άλλα λόγια, ο $\text{null}(\omega)$ αποτελείται από εκείνους τους $k$-διάστατους υποχώρους που, όταν προβληθούν στον $k$-διάστατο υπόχωρο $\pi(\omega)$ που ορίζουν τα δυϊκά διανύσματα των $\omega^a$, "καταρρέουν" σε μικρότερη διάσταση, χάνοντας έτσι τον $k$-διάστατο όγκο τους — και συνεπώς μηδενίζονται από την $\omega$.

Γεωμετρικά, αυτό σημαίνει ότι η $\omega$ "αδιαφορεί" για οποιαδήποτε συνιστώσα των διανυσμάτων κατά μήκος του $\rho(\omega)$. Η $\omega$ "βλέπει" μόνο τις προβολές των διανυσμάτων στον $\pi(\omega)$, ο οποίος έχει διάσταση $k$. Στον $\pi(\omega)$, τα διανύσματα $\partial / \partial \omega^1, \dots, \partial / \partial \omega^k$ ορίζουν ένα σύστημα συντεταγμένων, και η ορίζουσα μετρά τον \textbf{προσανατολισμένο $k$-διάστατο όγκο} του παραλληλεπιπέδου που ορίζουν οι προβολές των $\xi_1, \dots, \xi_k$ μέσα στον $\pi(\omega)$, εκφρασμένο στις συντεταγμένες που ορίζουν τα ίδια τα $\partial / \partial \omega^1, \dots, \partial / \partial \omega^k$.

Με άλλα λόγια, μια απλή $k$-μορφή είναι μια μηχανή που τρώει $k$ διανύσματα, τα προβάλλει στον $k$-διάστατο υπόχωρο $\pi(\omega)$ (αγνοώντας τις συνιστώσες τους κατά μήκος του $\rho(\omega)$), και υπολογίζει τον προσανατολισμένο όγκο του παραλληλεπιπέδου που σχηματίζουν εκεί, χρησιμοποιώντας ως μονάδα μέτρησης τον όγκο του παραλληλεπιπέδου που ορίζεται από τα $\partial / \partial \omega^1, \dots, \partial / \partial \omega^k$. Ο $\rho(\omega)$ είναι ο πυρήνας της προβολής στον $\pi(\omega)$, και ο $\text{null}(\omega)$ αποτελείται ακριβώς από εκείνους τους $k$-διάστατους υποχώρους που δεν είναι πλήρως εγκάρσιοι στον $\rho(\omega)$, ή ισοδύναμα, από εκείνους των οποίων η προβολή στον $\pi(\omega)$ έχει διάσταση μικρότερη του $k$.

\subsection{Μορφές που προβάλλουν στο ίδιο επίπεδο}

Μια θεμελιώδης ιδιότητα των απλών $k$-μορφών είναι ότι ο υπόχωρος $\pi(\omega)$ στον οποίο προβάλλουν, ή ισοδύναμα ο ριζικός υπόχωρος $\rho(\omega)$, τις καθορίζει με μοναδικό τρόπο, modulo βαθμωτό πολλαπλασιασμό. Ας θεωρήσουμε δύο απλές $k$-μορφές
\[
\omega = \omega^1 \wedge \dots \wedge \omega^k, \qquad \psi = \psi^1 \wedge \dots \wedge \psi^k,
\]
και ας υποθέσουμε ότι προβάλλουν στον ίδιο $k$-διάστατο υπόχωρο:
\[
\pi(\omega) = \pi(\psi).
\]
Αυτό είναι ισοδύναμο με το ότι έχουν τον ίδιο ριζικό υπόχωρο:
\[
\rho(\omega) = \rho(\psi),
\]
και συνεπώς, λόγω του χαρακτηρισμού (119), τον ίδιο μηδενοχώρο:
\[
\text{null}(\omega) = \text{null}(\psi).
\]
Γεωμετρικά, και οι δύο μορφές μετρούν προσανατολισμένο $k$-διάστατο όγκο πάνω στο ίδιο επίπεδο $\pi = \pi(\omega) = \pi(\psi)$, απλώς χρησιμοποιώντας διαφορετικές μονάδες μέτρησης.

Αφού οι $\omega^1, \dots, \omega^k$ και οι $\psi^1, \dots, \psi^k$ είναι βάσεις του ίδιου $k$-διάστατου χώρου 1-μορφών που μηδενίζονται στον $\rho(\omega)$, κάθε $\psi^a$ μπορεί να εκφραστεί ως γραμμικός συνδυασμός των $\omega^b$. Συνεπώς, υπάρχουν βαθμωτοί συντελεστές $\gamma^a_b$ τέτοιοι ώστε
\[
\psi^a = \gamma^a_b \, \omega^b, \qquad a = 1, \dots, k,
\]
όπου η άθροιση εννοείται στον δείκτη $b$. Τότε, λόγω της πολυγραμμικότητας και της αντισυμμετρίας του σφηνοειδούς γινομένου, έχουμε
\[
\psi = \psi^1 \wedge \dots \wedge \psi^k = (\gamma^1_{b_1} \omega^{b_1}) \wedge \dots \wedge (\gamma^k_{b_k} \omega^{b_k}) = \gamma^1_{b_1} \dots \gamma^k_{b_k} \, \omega^{b_1} \wedge \dots \wedge \omega^{b_k}.
\]
Από την αντισυμμετρία, το $\omega^{b_1} \wedge \dots \wedge \omega^{b_k}$ είναι μη μηδενικό μόνο αν οι δείκτες $b_1, \dots, b_k$ είναι μια μετάθεση των $1, \dots, k$, και τότε ισούται με $\operatorname{sgn}(\sigma) \, \omega^1 \wedge \dots \wedge \omega^k$. Το άθροισμα πάνω σε όλες τις μεταθέσεις δίνει ακριβώς την ορίζουσα του πίνακα $(\gamma^a_b)$. Έτσι,
\begin{equation}
\psi = \det(\gamma^a_b) \ \omega. \qquad (121)
\end{equation}
Με άλλα λόγια, δύο απλές $k$-μορφές που προβάλλουν στον ίδιο υπόχωρο $\pi$ είναι αναγκαστικά ανάλογες: διαφέρουν μόνο κατά έναν βαθμωτό παράγοντα, ο οποίος είναι η ορίζουσα του πίνακα αλλαγής βάσης από τις $\omega^a$ στις $\psi^a$. Γεωμετρικά, αυτός ο παράγοντας είναι ακριβώς ο λόγος των μονάδων όγκου που χρησιμοποιούν οι δύο μορφές για τη μέτρηση του προσανατολισμένου $k$-διάστατου όγκου πάνω στο $\pi$.

Το αποτέλεσμα αυτό αναδεικνύει μια θεμελιώδη αναλλοιότητα: ο ριζικός υπόχωρος $\rho(\omega)$ και ο μηδενοχώρος $\text{null}(\omega)$ είναι πλήρως καθορισμένοι από την κλάση ισοδυναμίας της $\omega$ modulo βαθμωτό πολλαπλασιασμό. Με άλλα λόγια, η γεωμετρική ουσία μιας απλής $k$-μορφής — το επίπεδο προβολής και το σύνολο των υποχώρων που μηδενίζονται — είναι ανεξάρτητη από την επιλογή μονάδων μέτρησης. Μια απλή $k$-μορφή είναι, ουσιαστικά, ένα προσανατολισμένο $k$-διάστατο επίπεδο $\pi(\omega)$ εφοδιασμένο με μια μονάδα μέτρησης όγκου, και η αλλαγή μονάδων αντιστοιχεί σε πολλαπλασιασμό με βαθμωτό.

\subsection{Ο νόμος μετασχηματισμού και η γεωμετρική του σημασία}

Ο νόμος μετασχηματισμού μιας $k$-μορφής κάτω από αλλαγή συντεταγμένων $x^i \to y^i$ είναι, όπως για κάθε τανυστή τύπου $(0,k)$,
\begin{equation}
\tilde{\omega}_{i_1 \dots i_k} = \frac{\partial x^{\mu_1}}{\partial y^{i_1}} \dots \frac{\partial x^{\mu_k}}{\partial y^{i_k}} \ \omega_{\mu_1 \dots \mu_k}. \qquad (112)
\end{equation}
Για μια απλή $k$-μορφή $\omega = \omega^1 \wedge \dots \wedge \omega^k$, αυτό σημαίνει ότι οι συνιστώσες της μετασχηματίζονται με την ορίζουσα του Ιακωβιανού πίνακα που αντιστοιχεί στις $k$ κατευθύνσεις που "ανιχνεύει" η μορφή. Γεωμετρικά, αυτό είναι απολύτως φυσικό: ο προσανατολισμένος όγκος ενός παραλληλεπιπέδου αλλάζει ακριβώς κατά την ορίζουσα του μετασχηματισμού των συντεταγμένων. Αν οι συντεταγμένες αλλάζουν με τρόπο που "τεντώνει" τον χώρο, ο όγκος πολλαπλασιάζεται με τον παράγοντα της ορίζουσας. Η απλή $k$-μορφή, μέσω του νόμου μετασχηματισμού της, ενσωματώνει ακριβώς αυτή τη γεωμετρική πληροφορία. Ο ριζικός υπόχωρος $\rho(\omega)$ και ο μηδενοχώρος $\text{null}(\omega)$, ως γεωμετρικά αντικείμενα, μετασχηματίζονται συναλλοίωτα κάτω από αλλαγή συντεταγμένων, διατηρώντας τη γεωμετρική τους σημασία ανέπαφη.

\subsection{Η παρουσία Ρημάννειας μετρικής}

Όταν ο διανυσματικός χώρος $V$ είναι εφοδιασμένος με μια Ρημάννεια (ή ψευδο-Ρημάννεια) μετρική $g$, η γεωμετρική δομή των απλών $k$-μορφών εμπλουτίζεται δραματικά. Η μετρική παρέχει μια κανονική ταύτιση ανάμεσα σε διανύσματα και 1-μορφές μέσω των μουσικών ισομορφισμών $\flat$ και $\sharp$, και αυτό επιτρέπει νέους χαρακτηρισμούς του ριζικού υπόχωρου και του μηδενοχώρου.

\subsubsection{Μετρικός χαρακτηρισμός του ριζικού υπόχωρου}

Χωρίς μετρική, ο $\rho(\omega)$ ορίζεται ως η τομή των πυρήνων των $\omega^a$ στην (116). Με την παρουσία της μετρικής, μπορούμε να δώσουμε έναν ισοδύναμο, καθαρά γεωμετρικό χαρακτηρισμό: ο $\rho(\omega)$ είναι το ορθογώνιο συμπλήρωμα του $\pi(\omega)$ ως προς τη μετρική $g$. Για να το αποδείξουμε αυτό, θεωρούμε ένα διάνυσμα $\xi \in V$. Χρησιμοποιώντας τους μουσικούς ισομορφισμούς, η 1-μορφή $\omega^a$ μπορεί να γραφεί ως $\omega^a = (\xi_a)^\flat$ για κάποιο μοναδικό διάνυσμα $\xi_a = (\omega^a)^\sharp$. Τότε, για κάθε $\eta \in V$, έχουμε
\[
\omega^a(\eta) = (\xi_a)^\flat(\eta) = g(\xi_a, \eta).
\]
Συνεπώς, η συνθήκη $\eta \in \ker \omega^a$ είναι ισοδύναμη με $g(\xi_a, \eta) = 0$, δηλαδή το $\eta$ είναι ορθογώνιο στο $\xi_a$. Ο $\rho(\omega)$ είναι η τομή αυτών των ορθογώνιων συμπληρωμάτων για $a = 1, \dots, k$, άρα
\[
\rho(\omega) = \{ \xi_a \mid a = 1, \dots, k \}^\perp.
\]
Αλλά τα $\xi_a = (\omega^a)^\sharp$ είναι ακριβώς τα διανύσματα που παράγουν τον $\pi(\omega)$. Πράγματι, από την (114) και τον ορισμό του $\sharp$, έχουμε
\[
\delta^a_b = \omega^a \left( \frac{\partial}{\partial \omega^b} \right) = g \left( (\omega^a)^\sharp, \frac{\partial}{\partial \omega^b} \right) = g(\xi_a, \partial / \partial \omega^b),
\]
που δείχνει ότι τα $\xi_a$ και $\partial / \partial \omega^a$ είναι δυϊκές βάσεις ως προς τη μετρική στον $\pi(\omega)$. Συνεπώς, $\text{span}\{\xi_1, \dots, \xi_k\} = \pi(\omega)$. Άρα,
\begin{equation}
\rho(\omega) = \pi(\omega)^\perp = \{ \eta \in V \mid g(\eta, \zeta) = 0 \ \forall \zeta \in \pi(\omega) \}. \qquad (122)
\end{equation}
Αυτός ο χαρακτηρισμός είναι θεμελιώδης: ο ριζικός υπόχωρος είναι ακριβώς ο υπόχωρος των διανυσμάτων που είναι κάθετα σε όλα τα διανύσματα του $\pi(\omega)$. Χωρίς μετρική, η έννοια της καθετότητας δεν υπάρχει, και ο $\rho(\omega)$ ορίζεται μόνο μέσω των $\omega^a$. Με τη μετρική, αποκτά ένα ανεξάρτητο γεωμετρικό νόημα.

\subsubsection{Μετρικός χαρακτηρισμός του μηδενοχώρου}

Ο χαρακτηρισμός (122) επιτρέπει μια νέα διατύπωση του μηδενοχώρου. Από την (119), ένας $k$-διάστατος υπόχωρος $U$ ανήκει στον $\text{null}(\omega)$ αν και μόνο αν $U \cap \rho(\omega) \neq \{0\}$. Χρησιμοποιώντας την (122), αυτό γράφεται ισοδύναμα ως
\begin{equation}
U \in \text{null}(\omega) \iff U \cap \pi(\omega)^\perp \neq \{0\}. \qquad (123)
\end{equation}
Με άλλα λόγια, ένας $k$-διάστατος υπόχωρος μηδενίζεται από την $\omega$ αν και μόνο αν περιέχει ένα μη μηδενικό διάνυσμα κάθετο στον $\pi(\omega)$. Αυτή η διατύπωση είναι ιδιαίτερα διαφωτιστική στη φυσική: αν ο $\pi(\omega)$ είναι ο χώρος που "βλέπει" η μορφή, τότε ο $\text{null}(\omega)$ αποτελείται από εκείνους τους $k$-διάστατους υποχώρους που δεν είναι πλήρως "εγκάρσιοι" στον $\pi(\omega)$, αλλά αντίθετα περιέχουν μια κάθετη συνιστώσα.

Ισοδύναμα, η συνθήκη μπορεί να εκφραστεί σε όρους της προβολής. Από την (120), έχουμε $\dim(\text{proj}_{\pi(\omega)}(U)) < k$. Με τη μετρική, η προβολή στον $\pi(\omega)$ είναι η ορθογώνια προβολή, και η συνθήκη αυτή σημαίνει ότι ο $U$ "χάνει" τουλάχιστον μία διάσταση όταν προβληθεί ορθογώνια στον $\pi(\omega)$. Ο πυρήνας αυτής της ορθογώνιας προβολής είναι ακριβώς ο $\rho(\omega) = \pi(\omega)^\perp$, και η απώλεια διάστασης οφείλεται στο γεγονός ότι ο $U$ περιέχει ένα μη μηδενικό διάνυσμα σε αυτόν τον πυρήνα.

\subsubsection{Το στοιχείο όγκου και ο τελεστής Hodge}

Η μετρική επιλέγει ένα κανονικό στοιχείο όγκου, την $n$-μορφή
\[
dV = \sqrt{|\det(g_{ij})|} \ dx^1 \wedge \dots \wedge dx^n,
\]
η οποία έχει την ιδιότητα ότι ο όγκος του παραλληλεπιπέδου που ορίζουν τα ορθοκανονικά διανύσματα βάσης είναι ακριβώς $1$. Αυτό παρέχει μια φυσική μονάδα μέτρησης όγκων, ανεξάρτητη από την επιλογή βάσης.

Για μια απλή $k$-μορφή $\omega$, η μετρική επιτρέπει τον ορισμό της δυϊκής Hodge $\star \omega$, μιας $(n-k)$-μορφής. Ο τελεστής Hodge εναλλάσσει τους ρόλους του $\pi(\omega)$ και του $\rho(\omega)$. Συγκεκριμένα, αν η $\omega$ είναι μια απλή $k$-μορφή, τότε η $\star \omega$ είναι μια απλή $(n-k)$-μορφή για την οποία ισχύει
\begin{equation}
\pi(\star \omega) = \rho(\omega), \qquad \rho(\star \omega) = \pi(\omega). \qquad (124)
\end{equation}
Αυτό αποδεικνύεται ως εξής: από τον ορισμό του τελεστή Hodge, η $\star \omega$ δρα σε $(n-k)$ διανύσματα συστέλλοντας την $\omega$ (με ανυψωμένους δείκτες) με τον τανυστή Levi-Civita. Τα διανύσματα που ανήκουν στον $\pi(\omega)$ δίνουν μη μηδενική συνεισφορά μόνο αν είναι πλήρως εγκάρσια στον $\rho(\omega)$, και το αποτέλεσμα είναι ότι η $\star \omega$ μηδενίζεται ακριβώς σε εκείνα τα $(n-k)$-παραλληλεπίπεδα που δεν είναι πλήρως εγκάρσια στον $\pi(\omega)$. Αυτό σημαίνει ότι $\rho(\star \omega) = \pi(\omega)$, και αφού ο $\pi(\star \omega)$ είναι το ορθογώνιο συμπλήρωμα του $\rho(\star \omega)$, έχουμε $\pi(\star \omega) = \pi(\omega)^\perp = \rho(\omega)$.

Αυτή η εναλλαγή ρόλων είναι μια από τις πιο κομψές συνέπειες της παρουσίας της μετρικής. Χωρίς μετρική, η έννοια του δυϊσμού Hodge δεν υπάρχει καν. Με τη μετρική, ο Hodge παρέχει μια πλήρη συμμετρία ανάμεσα στις $k$-μορφές και τις $(n-k)$-μορφές, ανταλλάσσοντας τον υπόχωρο προβολής με τον ριζικό υπόχωρο και αντιστρόφως.

\subsection{Οριακές περιπτώσεις υπό το πρίσμα της μετρικής}

Η γενική ερμηνεία που αναπτύξαμε, εμπλουτισμένη τώρα με τη μετρική δομή, εξειδικεύεται με διαφωτιστικό τρόπο στις οριακές τιμές του $k$.

\textbf{Περίπτωση $k=1$.} Μια 1-μορφή $\omega$ αντιστοιχεί σε $k=1$, οπότε η ορίζουσα είναι απλώς ο αριθμός $\omega(\xi_1)$. Ο υπόχωρος $\pi(\omega)$ είναι η μονοδιάστατη ευθεία που παράγεται από το δυϊκό διάνυσμα $\partial / \partial \omega$, και ο $\rho(\omega) = \ker \omega$ είναι ο υπερεπίπεδος χώρος διάστασης $n-1$ που περιγράφεται από τη γραμμική εξίσωση $\omega_j x^j = 0$. Με τη μετρική, ο $\rho(\omega)$ είναι το ορθογώνιο συμπλήρωμα της ευθείας $\pi(\omega)$, δηλαδή το υπερεπίπεδο που είναι κάθετο στο διάνυσμα $(\omega)^\sharp$. Η $\omega$ μετρά την ορθογώνια προβολή του $\xi_1$ πάνω στην ευθεία $\pi(\omega)$ — δηλαδή, μετρά το μήκος του $\xi_1$ κατά μήκος της κατεύθυνσης $\partial / \partial \omega$, σε μονάδες που ορίζονται από το ίδιο το $\partial / \partial \omega$. Αν κανονικοποιήσουμε την $\omega$ ώστε να έχει μοναδιαίο μήκος ως προς τη μετρική ($\|\omega\| = 1$), τότε η $\omega(\xi_1)$ δίνει ακριβώς το μήκος της ορθογώνιας προβολής του $\xi_1$ στην κατεύθυνση $(\omega)^\sharp$. Ο μηδενοχώρος $\text{null}(\omega)$ αποτελείται από όλους τους $1$-διάστατους υποχώρους που περιέχονται στον $\rho(\omega)$, δηλαδή από όλες τις ευθείες που είναι κάθετες στην κατεύθυνση $(\omega)^\sharp$. Ισοδύναμα, είναι οι $1$-διάστατοι υποχώροι των οποίων η ορθογώνια προβολή στον $\pi(\omega)$ έχει διάσταση $0 < 1$. Ο τελεστής Hodge απεικονίζει την $\omega$ σε μια $(n-1)$-μορφή $\star \omega$, της οποίας ο ριζικός υπόχωρος είναι η ευθεία $\pi(\omega)$ και της οποίας ο υπόχωρος προβολής είναι το υπερεπίπεδο $\rho(\omega)$. Έτσι, η $\star \omega$ μετρά τον προσανατολισμένο όγκο της προβολής ενός $(n-1)$-παραλληλεπιπέδου στο υπερεπίπεδο που είναι κάθετο στην αρχική κατεύθυνση.

\textbf{Περίπτωση $k=n-1$.} Μια $(n-1)$-μορφή $\omega$ σε έναν $n$-διάστατο χώρο αντιστοιχεί σε $k=n-1$. Ο υπόχωρος $\pi(\omega)$ έχει διάσταση $n-1$ και παράγεται από τα δυϊκά διανύσματα $\partial / \partial \omega^1, \dots, \partial / \partial \omega^{n-1}$, ενώ ο $\rho(\omega)$ έχει διάσταση $1$ και περιγράφεται από το γραμμικό σύστημα $\omega^a_j x^j = 0$, $a = 1, \dots, n-1$. Με τη μετρική, ο $\rho(\omega)$ είναι η μονοδιάστατη ευθεία που είναι κάθετη στο υπερεπίπεδο $\pi(\omega)$. Η $\omega$ τρώει $n-1$ διανύσματα και επιστρέφει τον προσανατολισμένο όγκο του παραλληλεπιπέδου που σχηματίζουν οι ορθογώνιες προβολές τους στον $\pi(\omega)$, αγνοώντας τη συνιστώσα τους κατά μήκος της κάθετης κατεύθυνσης $\rho(\omega)$. Ισοδύναμα, η $\omega$ μετρά τον όγκο της ορθογώνιας προβολής του παραλληλεπιπέδου στο υπερεπίπεδο $\pi(\omega)$. Ο μηδενοχώρος $\text{null}(\omega)$ αποτελείται από όλους τους $(n-1)$-διάστατους υποχώρους που τέμνουν μη τετριμμένα τον $\rho(\omega)$, δηλαδή από εκείνους που δεν είναι πλήρως εγκάρσιοι στην κάθετη κατεύθυνση. Ισοδύναμα, είναι οι $(n-1)$-διάστατοι υποχώροι των οποίων η ορθογώνια προβολή στον $\pi(\omega)$ έχει διάσταση μικρότερη του $n-1$, δηλαδή ακριβώς εκείνοι που περιέχουν την κάθετη ευθεία $\rho(\omega)$. Ο τελεστής Hodge απεικονίζει την $\omega$ σε μια 1-μορφή $\star \omega$, της οποίας ο ριζικός υπόχωρος είναι το υπερεπίπεδο $\pi(\omega)$ και της οποίας ο υπόχωρος προβολής είναι η ευθεία $\rho(\omega)$. Αυτή η 1-μορφή μετρά τη ροή μέσα από το υπερεπίπεδο: το διάνυσμα $(\star \omega)^\sharp$ είναι ακριβώς το κάθετο διάνυσμα στο υπερεπίπεδο $\pi(\omega)$, και το μήκος του μετρά την "πυκνότητα" της ροής.

\textbf{Περίπτωση $k=n$.} Μια $n$-μορφή $\omega$ σε έναν $n$-διάστατο χώρο είναι η μέγιστη δυνατή τάξη. Εδώ, $k=n$, οπότε ο $\pi(\omega)$ ταυτίζεται με ολόκληρο τον $V$ (τα $n$ δυϊκά διανύσματα $\partial / \partial \omega^1, \dots, \partial / \partial \omega^n$ αποτελούν βάση του $V$), και ο $\rho(\omega)$ έχει διάσταση $0$ (αποτελείται μόνο από το μηδενικό διάνυσμα), αφού το ομογενές σύστημα $\omega^a_j x^j = 0$, $a = 1, \dots, n$, έχει μόνο τη μηδενική λύση. Με τη μετρική, ο $\rho(\omega) = \{0\}$ είναι τετριμμένος, και ο $\pi(\omega) = V$ είναι ολόκληρος ο χώρος. Ο μηδενοχώρος $\text{null}(\omega)$ είναι κενός ως προς τους $n$-διάστατους υποχώρους: δεν υπάρχει κανένας $n$-διάστατος υπόχωρος του $V$ (εκτός από τον ίδιο τον $V$) που να τέμνει μη τετριμμένα τον $\rho(\omega) = \{0\}$, και η προβολή οποιουδήποτε $n$-διάστατου υποχώρου στον $V$ έχει διάσταση $n$, όχι μικρότερη. Συνεπώς, κανένας μη εκφυλισμένος $n$-διάστατος υπόχωρος δεν ανήκει στον $\text{null}(\omega)$. Η $\omega$ τρώει $n$ διανύσματα και επιστρέφει τον προσανατολισμένο $n$-διάστατο όγκο του παραλληλεπιπέδου που σχηματίζουν, εκφρασμένο στις μονάδες που ορίζουν τα $\partial / \partial \omega^1, \dots, \partial / \partial \omega^n$. Αυτή είναι ακριβώς η έννοια ενός στοιχείου όγκου: μια $n$-μορφή είναι μια μηχανή μέτρησης όγκων στον $V$. Με τη μετρική, το κανονικό στοιχείο όγκου $dV$ είναι η μοναδική $n$-μορφή που αποδίδει όγκο $1$ στο παραλληλεπίπεδο που ορίζουν τα ορθοκανονικά διανύσματα βάσης. Σε αυτή την περίπτωση, ο νόμος μετασχηματισμού γίνεται
\begin{equation}
\tilde{\omega}_{1 \dots n} = \det \left( \frac{\partial x^\mu}{\partial y^i} \right) \ \omega_{1 \dots n}, \qquad (113)
\end{equation}
που είναι ακριβώς ο μετασχηματισμός ενός όγκου κάτω από αλλαγή συντεταγμένων: ο όγκος πολλαπλασιάζεται με την ορίζουσα του Ιακωβιανού πίνακα. Ο τελεστής Hodge απεικονίζει την $n$-μορφή $\omega$ στη βαθμωτή συνάρτηση $\star \omega$, και αντιστρόφως, μια βαθμωτή συνάρτηση $f$ απεικονίζεται στην $n$-μορφή $f \, dV$. Αυτό επιβεβαιώνει ότι η $n$-μορφή είναι το γεωμετρικό αντικείμενο που κωδικοποιεί την έννοια του όγκου, και η μετρική παρέχει τη φυσική μονάδα μέτρησης.

\subsection{Σύνοψη}

Μια απλή $k$-μορφή $\omega = \omega^1 \wedge \dots \wedge \omega^k$ είναι, σε σημειακό επίπεδο, ένας τανυστής που αντιστοιχεί σε ένα $k$-διάστατο προσανατολισμένο στοιχείο όγκου στον υπόχωρο $\pi(\omega) = \text{span}\{\partial / \partial \omega^1, \dots, \partial / \partial \omega^k\}$, όπου τα $\partial / \partial \omega^a$ είναι τα δυϊκά διανύσματα των $\omega^a$, που ικανοποιούν την (114). Η δράση της σε $k$ διανύσματα είναι η ορίζουσα (111) των προβολών των διανυσμάτων πάνω στον $\pi(\omega)$. Ο ριζικός υπόχωρος $\rho(\omega)$, που ορίζεται ρητά στην (116) ως η τομή των πυρήνων των $\omega^a$, και χαρακτηρίζεται ισοδύναμα από την (117) ως ο μέγιστος υπόχωρος που μηδενίζεται από την $\omega$ και από το γραμμικό σύστημα (118), έχει διάσταση $n-k$ και είναι ο πυρήνας της προβολής στον $\pi(\omega)$. Ο μηδενοχώρος $\text{null}(\omega)$, που χαρακτηρίζεται από τις (119) και (120), αποτελείται από τους $k$-διάστατους υποχώρους που τέμνουν μη τετριμμένα τον $\rho(\omega)$. Δύο απλές $k$-μορφές προβάλλουν στον ίδιο υπόχωρο αν και μόνο αν διαφέρουν κατά βαθμωτό παράγοντα, όπως εκφράζεται από την (121), γεγονός που αναδεικνύει την αναλλοιότητα του $\rho(\omega)$ και του $\text{null}(\omega)$ υπό αλλαγή μονάδων μέτρησης.

Με την παρουσία μιας Ρημάννειας μετρικής $g$, ο ριζικός υπόχωρος αποκτά τον γεωμετρικό χαρακτηρισμό (122) ως το ορθογώνιο συμπλήρωμα του $\pi(\omega)$: $\rho(\omega) = \pi(\omega)^\perp$. Ο μηδενοχώρος χαρακτηρίζεται ισοδύναμα από την (123): $U \in \text{null}(\omega)$ αν και μόνο αν $U \cap \pi(\omega)^\perp \neq \{0\}$. Ο τελεστής Hodge εναλλάσσει τους ρόλους του $\pi(\omega)$ και του $\rho(\omega)$ μέσω της (124), παρέχοντας μια πλήρη συμμετρία ανάμεσα στις $k$-μορφές και τις $(n-k)$-μορφές. Ο νόμος μετασχηματισμού (112) αντανακλά την αλλαγή του όγκου κάτω από αλλαγή συντεταγμένων, με παράγοντα την ορίζουσα του Ιακωβιανού πίνακα, και για $k=n$ εξειδικεύεται στην (113). Αυτή η ερμηνεία είναι πλήρως ανεξάρτητη από μετρική για τον πυρήνα της, αλλά η μετρική προσθέτει νέα γεωμετρική δομή που επιτρέπει την έκφραση των εννοιών σε όρους καθετότητας, ορθογώνιας προβολής και δυϊσμού Hodge. Στις οριακές περιπτώσεις $k=1$, $k=n-1$ και $k=n$, η γενική αυτή εικόνα εξειδικεύεται σε μέτρηση ορθογώνιων προβολών κατά μήκος μιας ευθείας, μέτρηση ροών μέσα από υπερεπιφάνειες με κάθετα διανύσματα, και μέτρηση όγκων με το κανονικό στοιχείο όγκου, αντίστοιχα.

\section{Δυϊσμός Ροής - Έργου}

Οι οριακές περιπτώσεις $k=1$ και $k=n-1$ της θεωρίας των απλών $k$-μορφών αποκαλύπτουν μια θεμελιώδη δυϊκότητα ανάμεσα στις έννοιες της \textbf{ροής} και του \textbf{έργου}. Είναι κρίσιμο να κατανοήσουμε εξαρχής ότι τόσο το έργο όσο και η ροή είναι, ως έννοιες, απολύτως ανεξάρτητες από την ύπαρξη μετρικής. Η μετρική δεν ορίζει καμία από τις δύο· αυτό που κάνει είναι να παρέχει μια γέφυρα — μια δυϊκότητα — που επιτρέπει τη μετάφραση της μίας στην άλλη. Αυτή η γέφυρα είναι ο τελεστής Hodge.

\subsection{Το έργο ως προ-μετρική έννοια: η περίπτωση $k=1$}

Μια 1-μορφή $\omega \in V^*$ και ένα διάνυσμα $\xi \in V$ είναι δυϊκά αντικείμενα εκ φύσεως: η 1-μορφή ανήκει στον δυϊκό χώρο $V^*$ και το διάνυσμα στον $V$, και η μεταξύ τους σύζευξη
\[
\omega(\xi)
\]
είναι η θεμελιώδης πράξη της γραμμικής άλγεβρας που δεν απαιτεί καμία πρόσθετη δομή. Στη φυσική, αυτή η σύζευξη ερμηνεύεται ως \textbf{έργο}: αν η $\omega$ είναι μια δύναμη και το $\xi$ είναι μια μετατόπιση, τότε $\omega(\xi)$ είναι το έργο που παράγεται. Σε τοπικές συντεταγμένες, με $\omega = \omega_j \, dx^j$ και $\xi = \xi^i \, \partial / \partial x^i$, το έργο είναι απλώς
\begin{equation}
\omega(\xi) = \omega_j \, \xi^j. \qquad (125)
\end{equation}
Καμία μετρική δεν εμφανίζεται σε αυτή την έκφραση. Η δύναμη είναι φυσικά μια 1-μορφή — η κλίση ενός δυναμικού — και η μετατόπιση είναι φυσικά ένα διάνυσμα — εφαπτόμενο σε μια τροχιά. Το έργο είναι η φυσική σύζευξη αυτών των δύο αντικειμένων.

Γεωμετρικά, σύμφωνα με τη θεωρία του Κεφαλαίου 16, ο υπόχωρος $\pi(\omega)$ είναι η μονοδιάστατη ευθεία που παράγεται από το δυϊκό διάνυσμα $\partial / \partial \omega$, και ο ριζικός υπόχωρος $\rho(\omega) = \ker \omega$ είναι το υπερεπίπεδο διάστασης $n-1$ που περιγράφεται από τη γραμμική εξίσωση $\omega_j x^j = 0$. Η $\omega$ μετρά την προβολή του $\xi$ πάνω στην ευθεία $\pi(\omega)$ — δηλαδή, το έργο εξαρτάται μόνο από τη συνιστώσα της μετατόπισης που είναι παράλληλη στη δύναμη. Κάθε μετατόπιση που ανήκει στον $\rho(\omega)$ παράγει μηδενικό έργο. Όλα αυτά ισχύουν χωρίς καμία αναφορά σε μετρική.

\subsection{Η ροή ως προ-μετρική έννοια: η περίπτωση $k=n-1$}

Μια $(n-1)$-μορφή $\eta$ τρώει $n-1$ διανύσματα $\xi_1, \dots, \xi_{n-1} \in V$ και επιστρέφει έναν αριθμό $\eta(\xi_1, \dots, \xi_{n-1})$. Και αυτή η σύζευξη είναι ανεξάρτητη από μετρική: είναι μια καθαρά πολυγραμμική αλγεβρική πράξη. Στη φυσική, αυτή η δράση ερμηνεύεται ως \textbf{ροή} μέσα από ένα $(n-1)$-διάστατο παραλληλεπίπεδο που ορίζουν τα $\xi_1, \dots, \xi_{n-1}$. Σε τοπικές συντεταγμένες, μια απλή $(n-1)$-μορφή $\eta$ γράφεται ως
\[
\eta = \eta^1 \wedge \dots \wedge \eta^{n-1},
\]
και η ροή είναι, σύμφωνα με την (111),
\[
\eta(\xi_1, \dots, \xi_{n-1}) = \det \left( \eta^a(\xi_i) \right)_{a,i=1}^{n-1}.
\]
Σε όρους γενικών συνιστωσών,
\[
\eta = \frac{1}{(n-1)!} \, \eta_{i_1 \dots i_{n-1}} \, dx^{i_1} \wedge \dots \wedge dx^{i_{n-1}},
\]
και η ροή είναι
\[
\eta(\xi_1, \dots, \xi_{n-1}) = \eta_{i_1 \dots i_{n-1}} \, \xi^{i_1}_1 \dots \xi^{i_{n-1}}_{n-1}.
\]
Γεωμετρικά, ο $\pi(\eta)$ είναι ένα υπερεπίπεδο διάστασης $n-1$, και ο $\rho(\eta)$ είναι η μονοδιάστατη ευθεία που είναι ο πυρήνας της προβολής στον $\pi(\eta)$. Η $\eta$ μετρά τον προσανατολισμένο όγκο της προβολής του παραλληλεπιπέδου στον $\pi(\eta)$, αγνοώντας τη συνιστώσα του κατά μήκος του $\rho(\eta)$. Τίποτα από αυτά δεν απαιτεί μετρική.

\subsection{Η μετρική ως γέφυρα: ο τελεστής Hodge και η δυϊκότητα}

Αν και το έργο και η ροή είναι ανεξάρτητες έννοιες, σε έναν χώρο εφοδιασμένο με μετρική $g$ αποκτούν μια βαθιά σχέση δυϊκότητας. Ο τελεστής Hodge $\star$, ο οποίος ορίζεται μόνο με τη βοήθεια της μετρικής, απεικονίζει 1-μορφές σε $(n-1)$-μορφές και αντιστρόφως. Συγκεκριμένα, η δράση του Hodge στις δύο αυτές τάξεις είναι:
\[
\star: \Omega^1(V) \to \Omega^{n-1}(V), \qquad \star: \Omega^{n-1}(V) \to \Omega^1(V).
\]
Αυτή η απεικόνιση είναι αμφιμονοσήμαντη (είναι ισομορφισμός) και ικανοποιεί την (124): εναλλάσσει τον υπόχωρο προβολής με τον ριζικό υπόχωρο.

Η δυϊκότητα ροής-έργου έγκειται ακριβώς σε αυτή την αντιστοιχία. Μια δύναμη $\omega$ (1-μορφή) μπορεί, μέσω του Hodge, να μεταφραστεί σε μια ροή $\star \omega$ ($(n-1)$-μορφή). Αντίστροφα, μια ροή $\eta$ ($(n-1)$-μορφή) μπορεί, μέσω του Hodge, να μεταφραστεί σε μια δύναμη $\star \eta$ (1-μορφή). Χωρίς μετρική, αυτή η μετάφραση είναι αδύνατη — το έργο και η ροή παραμένουν ανεξάρτητες, μη συνδεόμενες έννοιες. Η μετρική είναι η γέφυρα που επιτρέπει τη μετάβαση από τη μία στην άλλη.

Σε τοπικές συντεταγμένες, αυτή η δυϊκότητα εκφράζεται ως εξής. Για μια 1-μορφή $\omega = \omega_j \, dx^j$, η δυϊκή της Hodge είναι η $(n-1)$-μορφή
\[
\star \omega = \frac{1}{(n-1)!} \, \varepsilon_{i_1 \dots i_{n-1} j} \, \omega^j \, dx^{i_1} \wedge \dots \wedge dx^{i_{n-1}},
\]
όπου $\omega^j = g^{j\alpha} \omega_\alpha$ είναι οι ανυψωμένες συνιστώσες της $\omega$. Οι συνιστώσες της $\star \omega$ είναι
\begin{equation}
(\star \omega)_{i_1 \dots i_{n-1}} = \varepsilon_{i_1 \dots i_{n-1} j} \, \omega^j. \qquad (126)
\end{equation}
Παρατηρούμε ότι η μετρική εμφανίζεται ρητά μέσω της ανύψωσης του δείκτη στην $\omega^j$. Αντίστροφα, για μια $(n-1)$-μορφή $\eta$, η δυϊκή της Hodge είναι η 1-μορφή
\begin{equation}
(\star \eta)_j = \frac{1}{(n-1)!} \, \varepsilon_{i_1 \dots i_{n-1} j} \, \eta^{i_1 \dots i_{n-1}}, \qquad (127)
\end{equation}
όπου $\eta^{i_1 \dots i_{n-1}} = g^{i_1 \alpha_1} \dots g^{i_{n-1} \alpha_{n-1}} \eta_{\alpha_1 \dots \alpha_{n-1}}$. Και πάλι, η μετρική εμφανίζεται ρητά στην ανύψωση των δεικτών.

Η γεωμετρική σημασία αυτής της δυϊκότητας είναι η εξής. Η 1-μορφή $\omega$ έχει ριζικό υπόχωρο $\rho(\omega)$ διάστασης $n-1$ (υπερεπίπεδο) και υπόχωρο προβολής $\pi(\omega)$ διάστασης $1$ (ευθεία). Η δυϊκή της $\star \omega$ έχει, σύμφωνα με την (124), ριζικό υπόχωρο $\rho(\star \omega) = \pi(\omega)$ (την ευθεία) και υπόχωρο προβολής $\pi(\star \omega) = \rho(\omega)$ (το υπερεπίπεδο). Αυτό σημαίνει ότι η $\star \omega$ ``βλέπει'' ακριβώς ό,τι ``αγνοεί'' η $\omega$, και ``αγνοεί'' ακριβώς ό,τι ``βλέπει'' η $\omega$. Το έργο και η ροή είναι λοιπόν συμπληρωματικές όψεις της ίδιας γεωμετρικής οντότητας, και η μετρική είναι αυτή που επιτρέπει τη μετάβαση από τη μία όψη στην άλλη.

\subsection{Φυσικογεωμετρικές εφαρμογές}

\subsubsection{Ηλεκτροστατική: ο νόμος του Gauss}

Στην ηλεκτροστατική, το ηλεκτρικό πεδίο είναι φυσικά μια 1-μορφή $E = E_j \, dx^j$. Το έργο που παράγει σε ένα δοκιμαστικό φορτίο $q$ κατά μήκος μιας απειροστής μετατόπισης $\xi$ είναι, χωρίς καμία μετρική,
\[
W = q \, E(\xi) = q \, E_j \, \xi^j.
\]
Αυτή η έκφραση είναι απολύτως θεμελιώδης και ανεξάρτητη από μετρική. Το γεγονός ότι το ηλεκτροστατικό πεδίο είναι συντηρητικό εκφράζεται ως $\oint E = 0$ για κάθε κλειστή καμπύλη — και αυτό επίσης δεν απαιτεί μετρική.

Ωστόσο, στη φυσική θέλουμε επίσης να μιλήσουμε για τη ροή του ηλεκτρικού πεδίου μέσα από μια επιφάνεια — και εδώ ακριβώς εμφανίζεται η μετρική. Στον τρισδιάστατο Ευκλείδειο χώρο με $g_{ij} = \delta_{ij}$, ο τανυστής Levi-Civita είναι $\varepsilon_{ijk} = \epsilon_{ijk}$, και η δυϊκή Hodge του $E$ είναι η 2-μορφή
\[
\star E = \frac{1}{2!} \, \varepsilon_{ijk} \, E^k \, dx^i \wedge dx^j,
\]
όπου $E^k = \delta^{k\alpha} E_\alpha$ είναι οι ανυψωμένες συνιστώσες. Οι μη μηδενικές συνιστώσες της $\star E$ είναι
\[
(\star E)_{12} = E^3, \quad (\star E)_{23} = E^1, \quad (\star E)_{31} = E^2.
\]
Η ροή του ηλεκτρικού πεδίου μέσα από ένα στοιχείο επιφάνειας που ορίζεται από τα $\xi, \eta$ είναι
\[
\Phi_E = (\star E)(\xi, \eta) = \varepsilon_{ijk} E^k \xi^i \eta^j.
\]
Αυτή είναι ακριβώς η γνωστή έκφραση $\mathbf{E} \cdot (\boldsymbol{\xi} \times \boldsymbol{\eta})$. Ο νόμος του Gauss,
\[
\int_{\partial V} \star E = \int_V \rho \, dV,
\]
συνδέει τη ροή του $\star E$ με την πυκνότητα φορτίου $\rho$, και απαιτεί μετρική τόσο για τον ορισμό του $\star E$ όσο και για τον ορισμό του στοιχείου όγκου $dV$. Χωρίς μετρική, μπορούμε να μιλήσουμε για το έργο του ηλεκτρικού πεδίου, αλλά όχι για τη ροή του.

\subsubsection{Ρευστοδυναμική: κυκλοφορία και ροή μάζας}

Στη ρευστοδυναμική, το πεδίο ταχύτητας είναι φυσικά ένα διανυσματικό πεδίο $u$, αλλά η φυσική ποσότητα που μετρά το έργο κατά μήκος μιας καμπύλης είναι η 1-μορφή $u^\flat$, η οποία ορίζεται μόνο με τη βοήθεια μετρικής:
\[
u^\flat = u_j \, dx^j, \qquad u_j = g_{ij} u^i.
\]
Η κυκλοφορία της $u^\flat$ κατά μήκος μιας καμπύλης $\gamma$ είναι
\[
\Gamma = \int_\gamma u^\flat = \int_\gamma u_j \, dx^j,
\]
και απαιτεί μετρική για τον ορισμό της $u^\flat$. Αντίστοιχα, η ροή μάζας μέσα από μια επιφάνεια περιγράφεται από την $(n-1)$-μορφή $\star u^\flat$. Στον τρισδιάστατο χώρο, αυτή είναι η 2-μορφή
\[
\star u^\flat = \frac{1}{2!} \, \varepsilon_{ijk} \, u^k \, dx^i \wedge dx^j,
\]
και η ροή μάζας μέσα από ένα στοιχείο επιφάνειας είναι
\[
\dot{m} = \rho \, (\star u^\flat)(\xi, \eta) = \rho \, \varepsilon_{ijk} u^k \xi^i \eta^j = \rho \, \mathbf{u} \cdot (\boldsymbol{\xi} \times \boldsymbol{\eta}).
\]
Εδώ, η μετρική εμφανίζεται σε δύο σημεία: πρώτον, στη μετάφραση του διανύσματος $u$ στην 1-μορφή $u^\flat$ (μέσω της μετρικής), και δεύτερον, στη μετάφραση της $u^\flat$ στην $\star u^\flat$ (μέσω του Hodge). Χωρίς μετρική, η ταχύτητα ως διάνυσμα και η ροή μάζας ως $(n-1)$-μορφή είναι ανεξάρτητα αντικείμενα. Η μετρική είναι αυτή που επιτρέπει τη σύνδεσή τους.

\subsubsection{Μαγνητοστατική: ο νόμος του Amp\`ere}

Στη μαγνητοστατική, το μαγνητικό πεδίο είναι φυσικά μια 2-μορφή $B = \frac{1}{2} B_{ij} \, dx^i \wedge dx^j$, η οποία μετρά τη μαγνητική ροή μέσα από μια επιφάνεια χωρίς να απαιτεί μετρική. Ωστόσο, το έργο του μαγνητικού πεδίου κατά μήκος μιας καμπύλης — η κυκλοφορία — απαιτεί μετρική: είναι η 1-μορφή $\star B$, η οποία ορίζεται μέσω του Hodge. Στον τρισδιάστατο χώρο,
\[
(\star B)_j = \frac{1}{2!} \, \varepsilon_{ijk} \, B^{ik},
\]
και η κυκλοφορία κατά μήκος μιας κλειστής καμπύλης $\gamma$ είναι
\[
\oint_\gamma \star B.
\]
Ο νόμος του Amp\`ere συνδέει αυτή την κυκλοφορία με το ρεύμα $J$ που διέρχεται από την επιφάνεια που ορίζει η $\gamma$:
\[
\oint_\gamma \star B = \int_\Sigma J,
\]
όπου $J$ είναι μια 2-μορφή πυκνότητας ρεύματος. Εδώ, η μετρική εμφανίζεται στη μετάφραση της μαγνητικής ροής (2-μορφή $B$) στην κυκλοφορία (1-μορφή $\star B$). Χωρίς μετρική, η μαγνητική ροή και η κυκλοφορία παραμένουν ασύνδετες.

\subsubsection{Θερμοδυναμική: έργο και θερμότητα}

Στη θερμοδυναμική, το έργο και η θερμότητα είναι 1-μορφές στον χώρο των θερμοδυναμικών μεταβλητών, και η ύπαρξή τους δεν απαιτεί μετρική. Σε ένα σύστημα με συντεταγμένες $(S, V)$, το έργο είναι $\delta W = P \, dV$ και η θερμότητα είναι $\delta Q = T \, dS$. Αυτές είναι 1-μορφές, και η δράση τους σε ένα διάνυσμα που αντιπροσωπεύει μια απειροστή μεταβολή δίνει το αντίστοιχο ποσό έργου ή θερμότητας.

Ο πρώτος νόμος της θερμοδυναμικής,
\[
dU = \delta Q - \delta W = T \, dS - P \, dV,
\]
είναι μια σχέση ανάμεσα σε 1-μορφές που δεν απαιτεί μετρική. Σε αυτό το πλαίσιο, δεν υπάρχει φυσικός λόγος να εισαχθεί ο τελεστής Hodge — ο χώρος των θερμοδυναμικών μεταβλητών δεν είναι εφοδιασμένος με μετρική, και η δυϊκότητα ροής-έργου δεν παίζει ρόλο. Το έργο και η θερμότητα είναι απλώς δύο διαφορετικές 1-μορφές, και ο πρώτος νόμος εκφράζει τη σχέση τους.

\subsection{Σύνοψη}

Το έργο, ως δράση μιας 1-μορφής σε ένα διάνυσμα, και η ροή, ως δράση μιας $(n-1)$-μορφής σε $n-1$ διανύσματα, είναι έννοιες απολύτως ανεξάρτητες από μετρική. Η σύζευξη $\omega(\xi)$ και η ορίζουσα $\eta(\xi_1, \dots, \xi_{n-1})$ είναι θεμελιώδεις αλγεβρικές πράξεις που δεν απαιτούν καμία πρόσθετη γεωμετρική δομή.

Η μετρική εισέρχεται στο προσκήνιο όταν θελήσουμε να συνδέσουμε αυτές τις δύο ανεξάρτητες έννοιες. Ο τελεστής Hodge $\star$, ο οποίος ορίζεται μόνο με τη βοήθεια μετρικής, παρέχει έναν ισομορφισμό ανάμεσα στις 1-μορφές και τις $(n-1)$-μορφές, και συνεπώς μια δυϊκότητα ανάμεσα στο έργο και τη ροή. Οι σχέσεις (126) και (127) δίνουν τους τύπους αυτής της μετάφρασης σε τοπικές συντεταγμένες.

Στη φυσική, αυτή η δυϊκότητα εκδηλώνεται με συγκεκριμένους τρόπους. Στην ηλεκτροστατική, το έργο του ηλεκτρικού πεδίου $E(\xi)$ είναι ανεξάρτητο από μετρική, αλλά η ηλεκτρική ροή $\star E$ και ο νόμος του Gauss απαιτούν μετρική. Στη ρευστοδυναμική, τόσο η κυκλοφορία $u^\flat$ όσο και η ροή μάζας $\star u^\flat$ απαιτούν μετρική για να συνδεθούν με το διανυσματικό πεδίο ταχύτητας $u$. Στη μαγνητοστατική, η μαγνητική ροή $B$ είναι ανεξάρτητη από μετρική, αλλά η κυκλοφορία $\star B$ και ο νόμος του Amp\`ere την απαιτούν. Στη θερμοδυναμική, τέλος, η μετρική απουσιάζει εντελώς, και το έργο με τη θερμότητα παραμένουν ανεξάρτητες 1-μορφές χωρίς καμία δυϊκότητα Hodge να τις συνδέει.

Η κεντρική ιδέα είναι λοιπόν η εξής: η φύση μάς δίνει το έργο και τη ροή ως ανεξάρτητες γεωμετρικές οντότητες. Η μετρική είναι μια πρόσθετη δομή που επιλέγουμε να εφοδιάσουμε τον χώρο μας, και αυτή η επιλογή είναι που επιτρέπει τη μετάφραση ανάμεσα στις δύο — μια μετάφραση που, αν και βολική και διαφωτιστική, δεν είναι θεμελιωδώς απαραίτητη για την ύπαρξη καμίας από τις δύο έννοιες.

\section{Το Εξωτερικό Γινόμενο $n-1$ Διανυσμάτων σε $n$-Διάστατο Χώρο}

\subsection{Ορισμός και θεμελιώδεις ιδιότητες}

Θεωρούμε έναν $n$-διάστατο διανυσματικό χώρο $V$ εφοδιασμένο με μια Ρημάννεια μετρική $g$. Για $n-1$ διανύσματα $\xi_1, \dots, \xi_{n-1} \in V$, ορίζουμε το \textbf{εξωτερικό γινόμενο} (ή διανυσματικό γινόμενο) ως το μοναδικό διάνυσμα $\xi_1 \times \dots \times \xi_{n-1} \in V$ που ικανοποιεί
\begin{equation}
g(\xi_1 \times \dots \times \xi_{n-1}, \eta) = dV(\xi_1, \dots, \xi_{n-1}, \eta), \qquad \forall \eta \in V, \qquad (128)
\end{equation}
όπου $dV$ είναι το κανονικό στοιχείο όγκου που επάγεται από τη μετρική. Με άλλα λόγια, το εξωτερικό γινόμενο είναι το διάνυσμα που, μέσω της μετρικής, αντιστοιχεί στην $(n-1)$-μορφή που προκύπτει από την εισαγωγή των $\xi_1, \dots, \xi_{n-1}$ ως πρώτα ορίσματα στο $dV$. Το εξωτερικό γινόμενο είναι αντισυμμετρικό ως προς οποιαδήποτε εναλλαγή των ορισμάτων του:
\begin{equation}
\xi_{\sigma(1)} \times \dots \times \xi_{\sigma(n-1)} = \operatorname{sgn}(\sigma) \ \xi_1 \times \dots \times \xi_{n-1}, \qquad (129)
\end{equation}
για κάθε μετάθεση $\sigma$ των $n-1$ δεικτών. Είναι επίσης πολυγραμμικό: γραμμικό ως προς κάθε όρισμα χωριστά.

\subsection{Έκφραση σε τοπικές συντεταγμένες}

Σε ένα τοπικό σύστημα συντεταγμένων $x^1, \dots, x^n$, η μετρική αναπαρίσταται από τις συνιστώσες $g_{ij}$, και το στοιχείο όγκου είναι
\begin{equation}
dV = \sqrt{|\det(g_{ij})|} \ dx^1 \wedge \dots \wedge dx^n. \qquad (130)
\end{equation}
Για $n-1$ διανύσματα $\xi_a = \xi_a^i \, \partial / \partial x^i$, $a = 1, \dots, n-1$, το εξωτερικό γινόμενο έχει συνιστώσες που δίνονται από τη σχέση
\begin{equation}
(\xi_1 \times \dots \times \xi_{n-1})^i = \frac{1}{\sqrt{|\det(g_{ij})|}} \ \varepsilon^{i}_{\ j_1 \dots j_{n-1}} \ \xi_1^{j_1} \dots \xi_{n-1}^{j_{n-1}}, \qquad (131)
\end{equation}
όπου $\varepsilon_{i_1 \dots i_n} = \sqrt{|\det(g_{ij})|} \, \epsilon_{i_1 \dots i_n}$ είναι ο τανυστής Levi-Civita, και $\varepsilon^{i}_{\ j_1 \dots j_{n-1}} = g^{i\alpha} \varepsilon_{\alpha j_1 \dots j_{n-1}}$ είναι οι ανυψωμένες συνιστώσες του. Ισοδύναμα, μπορούμε να γράψουμε τη συνιστώσα $i$ ως την ορίζουσα
\begin{equation}
(\xi_1 \times \dots \times \xi_{n-1})^i = \frac{1}{\sqrt{|\det(g_{ij})|}} \ \det \left( \xi_a^{k} \right)_{k \neq i, \ a=1, \dots, n-1}, \qquad (132)
\end{equation}
όπου ο πίνακας έχει γραμμές τους δείκτες $k = 1, \dots, i-1, i+1, \dots, n$ και στήλες τα διανύσματα $a = 1, \dots, n-1$, και το πρόσημο καθορίζεται από τον κανόνα $(-1)^{i+1}$. Σε συμπεπτυγμένο συμβολισμό με το σύμβολο Levi-Civita,
\begin{equation}
(\xi_1 \times \dots \times \xi_{n-1})^i = \frac{1}{\sqrt{|\det(g_{ij})|}} \ \epsilon^{i}_{\ j_1 \dots j_{n-1}} \ \xi_1^{j_1} \dots \xi_{n-1}^{j_{n-1}}, \qquad (133)
\end{equation}
όπου $\epsilon_{i_1 \dots i_n}$ είναι το σύμβολο Levi-Civita με $\epsilon_{12\dots n} = 1$.

\subsection{Γεωμετρικές ιδιότητες}

Το εξωτερικό γινόμενο έχει τις ακόλουθες θεμελιώδεις γεωμετρικές ιδιότητες, οι οποίες το καθιστούν το φυσικό ανάλογο του διανυσματικού γινομένου σε αυθαίρετη διάσταση.

\textbf{Καθετότητα.} Το εξωτερικό γινόμενο $\xi_1 \times \dots \times \xi_{n-1}$ είναι ορθογώνιο σε καθένα από τα διανύσματα που το παράγουν. Πράγματι, για κάθε $a = 1, \dots, n-1$, θέτοντας $\eta = \xi_a$ στον ορισμό (128), έχουμε
\begin{equation}
g(\xi_1 \times \dots \times \xi_{n-1}, \xi_a) = dV(\xi_1, \dots, \xi_{n-1}, \xi_a) = 0, \qquad (134)
\end{equation}
αφού το $dV$ είναι πλήρως αντισυμμετρικό και δύο από τα ορίσματά του είναι ίσα. Συνεπώς, το εξωτερικό γινόμενο είναι κάθετο στο υπερεπίπεδο που παράγουν τα $\xi_1, \dots, \xi_{n-1}$.

\textbf{Μέτρο και όγκος.} Το μέτρο (μήκος) του εξωτερικού γινομένου δίνεται από
\begin{equation}
\|\xi_1 \times \dots \times \xi_{n-1}\|^2 = g(\xi_1 \times \dots \times \xi_{n-1}, \xi_1 \times \dots \times \xi_{n-1}) = dV(\xi_1, \dots, \xi_{n-1}, \xi_1 \times \dots \times \xi_{n-1}). \qquad (135)
\end{equation}
Από τον ορισμό, το $dV$ μετρά τον προσανατολισμένο $n$-διάστατο όγκο. Το τελευταίο όρισμα είναι ένα μοναδιαίο διάνυσμα κάθετο στο υπερεπίπεδο που παράγουν τα $\xi_1, \dots, \xi_{n-1}$, και μάλιστα στην κατεύθυνση που κάνει το σύστημα $(\xi_1, \dots, \xi_{n-1}, \xi_1 \times \dots \times \xi_{n-1})$ θετικά προσανατολισμένο ως προς το $dV$. Το αποτέλεσμα είναι ότι
\begin{equation}
\|\xi_1 \times \dots \times \xi_{n-1}\| = \text{Vol}_{n-1}(\xi_1, \dots, \xi_{n-1}), \qquad (136)
\end{equation}
όπου $\text{Vol}_{n-1}(\xi_1, \dots, \xi_{n-1})$ είναι ο $(n-1)$-διάστατος όγκος του παραλληλεπιπέδου που ορίζουν τα $\xi_1, \dots, \xi_{n-1}$ στο υπερεπίπεδό τους, υπολογισμένος με τη μετρική που επάγεται από την $g$ σε αυτό το υπερεπίπεδο. Με άλλα λόγια, το μήκος του εξωτερικού γινομένου ισούται με τον όγκο του παραλληλεπιπέδου που παράγουν τα διανύσματα.

\textbf{Προσανατολισμός.} Το εξωτερικό γινόμενο είναι προσανατολισμένο: το σύστημα $(\xi_1, \dots, \xi_{n-1}, \xi_1 \times \dots \times \xi_{n-1})$ είναι πάντα θετικά προσανατολισμένο ως προς το στοιχείο όγκου $dV$. Αυτό σημαίνει ότι το εξωτερικό γινόμενο ``δείχνει'' προς την κατεύθυνση που συμπληρώνει το υπερεπίπεδο των $\xi_a$ σε ένα θετικά προσανατολισμένο $n$-διάστατο σύστημα. Αν αλλάξουμε τον προσανατολισμό του $dV$ (π.χ., αλλάζοντας τη διάταξη των συντεταγμένων), το εξωτερικό γινόμενο αλλάζει πρόσημο.

\textbf{Μηδενισμός.} Το εξωτερικό γινόμενο μηδενίζεται αν και μόνο αν τα $\xi_1, \dots, \xi_{n-1}$ είναι γραμμικά εξαρτημένα. Αυτό είναι άμεση συνέπεια της σχέσης (136) με τον όγκο: ο $(n-1)$-διάστατος όγκος μηδενίζεται ακριβώς όταν τα διανύσματα είναι γραμμικά εξαρτημένα, δηλαδή όταν το παραλληλεπίπεδο που ορίζουν είναι εκφυλισμένο.

\subsection{Ορθοκανονικότητα και εξωτερικό γινόμενο}

Το εξωτερικό γινόμενο έχει μια αξιοσημείωτη σχέση με την έννοια της ορθοκανονικότητας, η οποία αναδεικνύει τον ρόλο του ως ``συμπληρωτή'' μιας ορθοκανονικής βάσης ενός υπερεπιπέδου σε μια ορθοκανονική βάση ολόκληρου του χώρου.

Θεωρούμε $n-1$ διανύσματα $\xi_1, \dots, \xi_{n-1}$ που είναι ορθοκανονικά:
\begin{equation}
g(\xi_a, \xi_b) = \delta_{ab}, \qquad a, b = 1, \dots, n-1. \qquad (137)
\end{equation}
Από την ιδιότητα της καθετότητας (134), το εξωτερικό γινόμενο είναι ήδη ορθογώνιο σε καθένα από τα $\xi_a$. Επιπλέον, ο $(n-1)$-διάστατος όγκος του παραλληλεπιπέδου που ορίζουν τα $\xi_a$ είναι, λόγω της ορθοκανονικότητας, ακριβώς $1$. Από την (136), το μέτρο του εξωτερικού γινομένου είναι
\begin{equation}
\|\xi_1 \times \dots \times \xi_{n-1}\| = 1. \qquad (138)
\end{equation}
Συνεπώς, το σύνολο
\[
\{\xi_1, \dots, \xi_{n-1}, \xi_1 \times \dots \times \xi_{n-1}\}
\]
είναι μια \textbf{ορθοκανονική βάση} ολόκληρου του $n$-διάστατου χώρου $V$. Το εξωτερικό γινόμενο συμπληρώνει μια ορθοκανονική βάση του υπερεπιπέδου σε μια ορθοκανονική βάση του χώρου, και μάλιστα με τον σωστό προσανατολισμό: το σύστημα είναι θετικά προσανατολισμένο ως προς το $dV$, όπως έχουμε ήδη δει. Αυτό σημαίνει ότι
\begin{equation}
dV(\xi_1, \dots, \xi_{n-1}, \xi_1 \times \dots \times \xi_{n-1}) = 1. \qquad (139)
\end{equation}
Αντίστροφα, αν το σύνολο $\{\xi_1, \dots, \xi_{n-1}, \xi_1 \times \dots \times \xi_{n-1}\}$ είναι ορθοκανονικό, τότε τα $\xi_1, \dots, \xi_{n-1}$ είναι αναγκαστικά ορθοκανονικά (αφού είναι υποσύνολο ορθοκανονικής βάσης), και το εξωτερικό γινόμενο έχει μήκος $1$. Από την (136), αυτό σημαίνει ότι ο $(n-1)$-διάστατος όγκος τους είναι $1$, που είναι ισοδύναμο με την ορθοκανονικότητά τους.

Σε όρους του τελεστή Hodge, αυτή η ιδιότητα εκφράζεται ως εξής. Αν τα $\xi_1, \dots, \xi_{n-1}$ είναι ορθοκανονικά, τότε η $(n-1)$-μορφή
\begin{equation}
\omega = (\xi_1)^\flat \wedge \dots \wedge (\xi_{n-1})^\flat \qquad (140)
\end{equation}
ικανοποιεί
\begin{equation}
\star \omega = (\xi_1 \times \dots \times \xi_{n-1})^\flat, \qquad (141)
\end{equation}
και η αντίστοιχη 1-μορφή έχει μήκος $1$ ως προς τη μετρική. Ο τελεστής Hodge απεικονίζει λοιπόν ορθοκανονικές $(n-1)$-μορφές σε ορθοκανονικές 1-μορφές, διατηρώντας τη δομή ορθοκανονικότητας. Σε τοπικές συντεταγμένες, αν τα $\xi_a$ είναι τα διανύσματα βάσης $\partial / \partial x^{i_a}$ για $n-1$ επιλεγμένους δείκτες, και η μετρική είναι η Ευκλείδεια, τότε το εξωτερικό γινόμενο είναι ακριβώς το εναπομείναν διάνυσμα βάσης $\partial / \partial x^{i_n}$ (με κατάλληλο πρόσημο), επιβεβαιώνοντας ότι η ορθοκανονική βάση $\{\partial / \partial x^1, \dots, \partial / \partial x^n\}$ είναι αυτο-δυϊκή ως προς το εξωτερικό γινόμενο.

\subsection{Το θεμελιώδες θεώρημα δυϊκότητας}

Το εξωτερικό γινόμενο και ο τελεστής Hodge συνδέονται μέσω μιας αξιοσημείωτης ταυτότητας, η οποία αποκαλύπτει ότι η μετρική, αν και αναγκαία για τον ορισμό και των δύο, τελικά απαλοίφεται από το αποτέλεσμα όταν αυτά συνδυάζονται με την κατάλληλη σειρά.

\textbf{Θεώρημα.} Έστω $\omega = \omega^1 \wedge \dots \wedge \omega^{n-1}$ μια απλή $(n-1)$-μορφή και $\xi_1, \dots, \xi_{n-1} \in V$ αυθαίρετα διανύσματα. Τότε
\begin{equation}
\omega(\xi_1, \dots, \xi_{n-1}) = (\star \omega)(\xi_1 \times \dots \times \xi_{n-1}). \qquad (142)
\end{equation}
Το αριστερό μέλος είναι η προ-μετρική ροή της $\omega$ μέσα από το παραλληλεπίπεδο των $\xi_a$. Το δεξί μέλος είναι το έργο της 1-μορφής $\star \omega$ κατά μήκος του εξωτερικού γινομένου — μια έκφραση που απαιτεί μετρική τόσο για τον Hodge όσο και για το $\times$. Και όμως, τα δύο μέλη είναι ίσα.

\textbf{Απόδειξη.} Το αριστερό μέλος είναι εξ ορισμού
\[
\omega(\xi_1, \dots, \xi_{n-1}) = \det \left( \omega^a(\xi_i) \right)_{a,i=1}^{n-1},
\]
μια έκφραση ανεξάρτητη από μετρική. Για το δεξί μέλος, χρησιμοποιούμε τον ορισμό (128) του εξωτερικού γινομένου και τον ορισμό του Hodge. Από την (128), για κάθε $\eta \in V$,
\[
g(\xi_1 \times \dots \times \xi_{n-1}, \eta) = dV(\xi_1, \dots, \xi_{n-1}, \eta).
\]
Από τον ορισμό του Hodge για 1-μορφές και $(n-1)$-μορφές, η $\star \omega$ είναι η 1-μορφή που ικανοποιεί
\[
\star \omega(\eta) = dV((\omega^1)^\sharp, \dots, (\omega^{n-1})^\sharp, \eta), \qquad \forall \eta \in V.
\]
Συγκρίνοντας με την προηγούμενη σχέση, βλέπουμε ότι
\[
\star \omega(\eta) = g((\omega^1)^\sharp \times \dots \times (\omega^{n-1})^\sharp, \eta), \qquad \forall \eta \in V,
\]
άρα $(\star \omega)^\sharp = (\omega^1)^\sharp \times \dots \times (\omega^{n-1})^\sharp$. Συνεπώς,
\[
(\star \omega)(\xi_1 \times \dots \times \xi_{n-1}) = g((\star \omega)^\sharp, \xi_1 \times \dots \times \xi_{n-1}) = g((\omega^1)^\sharp \times \dots \times (\omega^{n-1})^\sharp, \xi_1 \times \dots \times \xi_{n-1}).
\]
Από την ιδιότητα του εξωτερικού γινομένου και του Hodge, αυτό ισούται με τον προσανατολισμένο $(n-1)$-διάστατο όγκο του παραλληλεπιπέδου των $\xi_a$, που είναι ακριβώς το $\det(\omega^a(\xi_i))$. $\square$

Το θεώρημα αυτό έχει μια αξιοσημείωτη συνέπεια: η μετρική, αν και παρούσα σε κάθε βήμα της κατασκευής του δεξιού μέλους, εξαφανίζεται από το τελικό αποτέλεσμα. Αυτό συμβαίνει διότι η μετρική εισέρχεται με δύο ακριβώς αντίστροφους τρόπους — μέσω του Hodge και μέσω του εξωτερικού γινομένου — οι οποίοι αλληλοαναιρούνται. Ας το παρακολουθήσουμε σε τοπικές συντεταγμένες.

Γράφουμε $\omega^a = \omega^a_j \, dx^j$. Το αριστερό μέλος είναι
\[
\omega(\xi_1, \dots, \xi_{n-1}) = \det \left( \omega^a_j \, \xi_i^j \right)_{a,i=1}^{n-1},
\]
όπου καμία μετρική δεν εμφανίζεται. Για το δεξί μέλος, έχουμε από την (127)
\[
(\star \omega)_j = \frac{1}{(n-1)!} \, \varepsilon_{i_1 \dots i_{n-1} j} \, \omega^{i_1 \dots i_{n-1}},
\]
και από την (131)
\[
(\xi_1 \times \dots \times \xi_{n-1})^j = \frac{1}{\sqrt{|\det(g_{ij})|}} \ \varepsilon^{j}_{\ k_1 \dots k_{n-1}} \ \xi_1^{k_1} \dots \xi_{n-1}^{k_{n-1}}.
\]
Το γινόμενό τους είναι
\[
(\star \omega)_j (\xi_1 \times \dots \times \xi_{n-1})^j = \frac{1}{(n-1)!} \varepsilon_{i_1 \dots i_{n-1} j} \, \omega^{i_1 \dots i_{n-1}} \cdot \frac{1}{\sqrt{|\det(g_{ij})|}} \varepsilon^{j}_{\ k_1 \dots k_{n-1}} \xi_1^{k_1} \dots \xi_{n-1}^{k_{n-1}}.
\]
Χρησιμοποιώντας την ταυτότητα συναίρεσης
\[
\varepsilon_{i_1 \dots i_{n-1} j} \, \varepsilon^{j}_{\ k_1 \dots k_{n-1}} = (n-1)! \, \delta^{[i_1}_{k_1} \dots \delta^{i_{n-1}]}_{k_{n-1}},
\]
το δεξί μέλος γίνεται
\[
\frac{1}{\sqrt{|\det(g_{ij})|}} \ \omega^{i_1 \dots i_{n-1}} \, \xi_1^{[i_1} \dots \xi_{n-1}^{i_{n-1}]}.
\]
Αλλά $\omega^{i_1 \dots i_{n-1}} = g^{i_1 \alpha_1} \dots g^{i_{n-1} \alpha_{n-1}} \omega_{\alpha_1 \dots \alpha_{n-1}}$, και $\omega_{\alpha_1 \dots \alpha_{n-1}}$ περιέχει έναν παράγοντα $\sqrt{|\det(g_{ij})|}$ από τον τανυστή Levi-Civita (αφού $\omega_{\alpha_1 \dots \alpha_{n-1}} = \sqrt{|\det(g_{ij})|} \, \epsilon_{\alpha_1 \dots \alpha_{n-1} \beta} \, (\omega^1)^\beta \dots (\omega^{n-1})^\beta$). Ο παράγοντας αυτός ακυρώνεται με τον $1/\sqrt{|\det(g_{ij})|}$ του εξωτερικού γινομένου. Οι ανυψώσεις $g^{ij}$ αντισταθμίζονται από τις συστολές με τα $\delta$ που προκύπτουν από τη συναίρεση των $\varepsilon$. Το τελικό αποτέλεσμα είναι ακριβώς το $\det(\omega^a_j \, \xi_i^j)$, χωρίς κανένα ίχνος της μετρικής.

\subsection{Σχέση με τον τελεστή Hodge}

Το εξωτερικό γινόμενο συνδέεται στενά με τον τελεστή Hodge. Συγκεκριμένα, για $n-1$ διανύσματα $\xi_1, \dots, \xi_{n-1}$, η $(n-1)$-μορφή που αντιστοιχεί στο παραλληλεπίπεδό τους μέσω της μετρικής είναι
\begin{equation}
\omega = (\xi_1)^\flat \wedge \dots \wedge (\xi_{n-1})^\flat. \qquad (143)
\end{equation}
Η δυϊκή Hodge αυτής της $(n-1)$-μορφής είναι μια 1-μορφή, και το αντίστοιχο διάνυσμα (μέσω του $\sharp$) είναι ακριβώς το εξωτερικό γινόμενο:
\begin{equation}
(\star \omega)^\sharp = \xi_1 \times \dots \times \xi_{n-1}. \qquad (144)
\end{equation}
Ισοδύναμα, σε όρους του ριζικού υπόχωρου και του υπόχωρου προβολής, το εξωτερικό γινόμενο παράγει τον $\rho(\omega)$ — τον μονοδιάστατο υπόχωρο που είναι το ορθογώνιο συμπλήρωμα του υπερεπιπέδου που παράγουν τα $\xi_1, \dots, \xi_{n-1}$. Αυτό είναι απόλυτα συνεπές με τη γεωμετρική ερμηνεία του Κεφαλαίου 16: η $(n-1)$-μορφή $\omega$ έχει $\pi(\omega)$ το υπερεπίπεδο των $\xi_a$, και $\rho(\omega)$ την κάθετη ευθεία· το εξωτερικό γινόμενο είναι ένα διάνυσμα πάνω σε αυτή την κάθετη ευθεία, με μέτρο ίσο με τον όγκο του παραλληλεπιπέδου στον $\pi(\omega)$.

\subsection{Πρακτικοί τύποι σε τοπικές συντεταγμένες}

Σε πρακτικές εφαρμογές, είναι χρήσιμο να έχουμε μια συμπαγή έκφραση για το εξωτερικό γινόμενο. Σε ένα σύστημα συντεταγμένων, το εξωτερικό γινόμενο δίνεται από
\begin{equation}
\xi_1 \times \dots \times \xi_{n-1} = \frac{1}{\sqrt{|\det(g_{ij})|}} \ \varepsilon^{i}_{\ j_1 \dots j_{n-1}} \ \xi_1^{j_1} \dots \xi_{n-1}^{j_{n-1}} \ \frac{\partial}{\partial x^i}. \qquad (145)
\end{equation}
Το μέτρο του υπολογίζεται μέσω της μετρικής:
\begin{equation}
\|\xi_1 \times \dots \times \xi_{n-1}\|^2 = g_{ij} (\xi_1 \times \dots \times \xi_{n-1})^i (\xi_1 \times \dots \times \xi_{n-1})^j. \qquad (146)
\end{equation}
Ισοδύναμα, το τετράγωνο του μέτρου δίνεται από την ορίζουσα Gram:
\begin{equation}
\|\xi_1 \times \dots \times \xi_{n-1}\|^2 = \det \left( g(\xi_a, \xi_b) \right)_{a,b=1}^{n-1}. \qquad (147)
\end{equation}
Αυτή η έκφραση είναι ανεξάρτητη από την επιλογή συντεταγμένων και αναδεικνύει τη γεωμετρική σημασία του εξωτερικού γινομένου ως του διανύσματος του οποίου το μήκος ισούται με τον $(n-1)$-διάστατο όγκο του παραλληλεπιπέδου των $\xi_a$. Στην ορθοκανονική περίπτωση, η ορίζουσα Gram γίνεται ο ταυτοτικός πίνακας, και το μέτρο είναι $1$, επιβεβαιώνοντας την (138).

Για $n-1$ διανύσματα που είναι εφαπτόμενα σε μια υπερεπιφάνεια $\Sigma$ που ορίζεται από μια συνάρτηση $\phi(x) = 0$, το εξωτερικό γινόμενο είναι παράλληλο στο κάθετο διάνυσμα $\nabla \phi$, και μάλιστα
\begin{equation}
\xi_1 \times \dots \times \xi_{n-1} = \frac{dV(\xi_1, \dots, \xi_{n-1}, \cdot)^\sharp}{\|\nabla \phi\|} \ \frac{\nabla \phi}{\|\nabla \phi\|}, \qquad (148)
\end{equation}
όπου $\nabla \phi = (d\phi)^\sharp$ είναι η κλίση της $\phi$. Το πρόσημο καθορίζεται από τον προσανατολισμό της $\Sigma$ ως προς το $dV$. Αυτή η σχέση είναι θεμελιώδης για τον υπολογισμό ροών μέσα από υπερεπιφάνειες.

\subsection{Σύνοψη}

Το εξωτερικό γινόμενο $n-1$ διανυσμάτων σε έναν $n$-διάστατο Ρημάννειο χώρο είναι το μοναδικό διάνυσμα που είναι κάθετο στο υπερεπίπεδο που παράγουν τα διανύσματα, έχει μέτρο ίσο με τον $(n-1)$-διάστατο όγκο του παραλληλεπιπέδου τους, και είναι προσανατολισμένο έτσι ώστε το σύστημα $(\xi_1, \dots, \xi_{n-1}, \xi_1 \times \dots \times \xi_{n-1})$ να είναι θετικά προσανατολισμένο. Σε τοπικές συντεταγμένες, οι συνιστώσες του δίνονται από τη συστολή του τανυστή Levi-Civita με τις συνιστώσες των διανυσμάτων, όπως εκφράζεται στις (131)-(133). Το εξωτερικό γινόμενο συνδέεται άμεσα με τον τελεστή Hodge μέσω της (144): είναι το διάνυσμα που αντιστοιχεί στη δυϊκή Hodge της $(n-1)$-μορφής που σχηματίζουν οι 1-μορφές $(\xi_a)^\flat$. Αν τα $\xi_1, \dots, \xi_{n-1}$ είναι ορθοκανονικά, το εξωτερικό γινόμενο συμπληρώνει το σύστημα σε μια ορθοκανονική βάση ολόκληρου του χώρου. Το θεμελιώδες θεώρημα (142) αποκαλύπτει ότι ο συνδυασμός του Hodge και του εξωτερικού γινομένου απαλοίφει τη μετρική: η προ-μετρική ροή $\omega(\xi_1, \dots, \xi_{n-1})$ ισούται με το μετρικό έργο $(\star \omega)(\xi_1 \times \dots \times \xi_{n-1})$, και η ισότητα αυτή οφείλεται στο γεγονός ότι η μετρική εισέρχεται με δύο αντίστροφους τρόπους που αλληλοαναιρούνται. Οι ιδιότητες αυτές γενικεύουν το γνωστό διανυσματικό γινόμενο του τρισδιάστατου Ευκλείδειου χώρου σε αυθαίρετη διάσταση και μετρική.

\section{Εξωτερικό Γινόμενο και Δυϊσμός Ροής-Έργου}

Στα δύο προηγούμενα κεφάλαια αναπτύξαμε ανεξάρτητα τη θεωρία του δυϊσμού ροής-έργου (Κεφάλαιο 17) και τη θεωρία του εξωτερικού γινομένου $n-1$ διανυσμάτων (Κεφάλαιο 18). Σε αυτό το κεφάλαιο, συνθέτουμε τις δύο θεωρίες και αποκαλύπτουμε τη βαθιά σχέση που τις συνδέει: το εξωτερικό γινόμενο είναι ο γεωμετρικός μηχανισμός που υλοποιεί τον δυϊσμό ροής-έργου, και ο δυϊσμός ροής-έργου είναι η φυσική ερμηνεία του εξωτερικού γινομένου.

\subsection{Το θεμελιώδες θεώρημα}

\textbf{Θεώρημα.} Έστω $\omega = \omega^1 \wedge \dots \wedge \omega^{n-1}$ μια απλή $(n-1)$-μορφή σε έναν $n$-διάστατο διανυσματικό χώρο $V$ εφοδιασμένο με Ρημάννεια μετρική $g$, και έστω $\xi_1, \dots, \xi_{n-1} \in V$. Τότε
\begin{equation}
\omega(\xi_1, \dots, \xi_{n-1}) = (\star \omega)(\xi_1 \times \dots \times \xi_{n-1}). \qquad (149)
\end{equation}
Το αριστερό μέλος είναι η ροή της $\omega$ μέσα από το παραλληλεπίπεδο των $\xi_a$, μια ποσότητα απολύτως ανεξάρτητη από μετρική. Το δεξί μέλος είναι το έργο της 1-μορφής $\star \omega$ (της δυϊκής Hodge της $\omega$) κατά μήκος του εξωτερικού γινομένου $\xi_1 \times \dots \times \xi_{n-1}$, μια ποσότητα που απαιτεί μετρική τόσο για τον ορισμό του $\star$ όσο και για τον ορισμό του $\times$. Η ταυτότητα δηλώνει ότι αυτές οι δύο ποσότητες είναι ίσες.

\textbf{Απόδειξη.} Το αριστερό μέλος είναι, από την (111),
\[
\omega(\xi_1, \dots, \xi_{n-1}) = \det \left( \omega^a(\xi_i) \right)_{a,i=1}^{n-1},
\]
μια έκφραση που δεν περιέχει τη μετρική. Για το δεξί μέλος, χρησιμοποιούμε τον ορισμό (128) του εξωτερικού γινομένου:
\[
g(\xi_1 \times \dots \times \xi_{n-1}, \eta) = dV(\xi_1, \dots, \xi_{n-1}, \eta), \qquad \forall \eta \in V.
\]
Από τον ορισμό του τελεστή Hodge για 1-μορφές και $(n-1)$-μορφές, η $\star \omega$ είναι η 1-μορφή που ικανοποιεί
\[
\star \omega(\eta) = dV((\omega^1)^\sharp, \dots, (\omega^{n-1})^\sharp, \eta), \qquad \forall \eta \in V.
\]
Συγκρίνοντας με την προηγούμενη σχέση, έχουμε
\[
\star \omega(\eta) = g((\omega^1)^\sharp \times \dots \times (\omega^{n-1})^\sharp, \eta), \qquad \forall \eta \in V,
\]
άρα $(\star \omega)^\sharp = (\omega^1)^\sharp \times \dots \times (\omega^{n-1})^\sharp$. Συνεπώς,
\[
(\star \omega)(\xi_1 \times \dots \times \xi_{n-1}) = g((\star \omega)^\sharp, \xi_1 \times \dots \times \xi_{n-1}) = g((\omega^1)^\sharp \times \dots \times (\omega^{n-1})^\sharp, \xi_1 \times \dots \times \xi_{n-1}).
\]
Από τις ιδιότητες του εξωτερικού γινομένου και του Hodge, αυτό ισούται με τον προσανατολισμένο $(n-1)$-διάστατο όγκο του παραλληλεπιπέδου των $\xi_a$ ως προς τη μονάδα όγκου που ορίζουν οι $\omega^a$, που είναι ακριβώς το αριστερό μέλος. $\square$

\subsection{Γεωμετρική ερμηνεία της ταυτότητας}

Η ταυτότητα (149) έχει μια βαθιά γεωμετρική ερμηνεία. Ένα $(n-1)$-διάστατο παραλληλεπίπεδο μπορεί να περιγραφεί με δύο ισοδύναμους τρόπους:

\begin{itemize}
\item \textbf{Περιγραφή ροής:} από τις $n-1$ πλευρές του $\xi_1, \dots, \xi_{n-1}$. Η ροή μιας $(n-1)$-μορφής $\omega$ μέσα από αυτό είναι $\omega(\xi_1, \dots, \xi_{n-1})$. Αυτή η περιγραφή είναι προ-μετρική.
\item \textbf{Περιγραφή έργου:} από ένα μοναδικό κάθετο διάνυσμα, το εξωτερικό γινόμενο $\xi_1 \times \dots \times \xi_{n-1}$, του οποίου το μέτρο ισούται με τον $(n-1)$-διάστατο όγκο του παραλληλεπιπέδου και του οποίου η κατεύθυνση είναι κάθετη στο υπερεπίπεδο που το περιέχει. Το έργο της 1-μορφής $\star \omega$ κατά μήκος αυτού του διανύσματος είναι $(\star \omega)(\xi_1 \times \dots \times \xi_{n-1})$. Αυτή η περιγραφή απαιτεί μετρική.
\end{itemize}

Η ταυτότητα (149) δηλώνει ότι αυτές οι δύο περιγραφές είναι ισοδύναμες: η ροή μέσα από το παραλληλεπίπεδο ισούται με το έργο της δυϊκής 1-μορφής κατά μήκος του κάθετου διανύσματος. Το εξωτερικό γινόμενο είναι ο μηχανισμός που μεταφράζει την περιγραφή ροής σε περιγραφή έργου, και ο τελεστής Hodge είναι ο μηχανισμός που μεταφράζει τη $(n-1)$-μορφή (ροή) στην αντίστοιχη 1-μορφή (δύναμη). Οι δύο μηχανισμοί συνεργάζονται άψογα, και το αποτέλεσμα είναι ανεξάρτητο από την επιλογή της μετρικής.

\subsection{Απαλοιφή της μετρικής}

Το γεγονός ότι το αριστερό μέλος της (149) είναι ανεξάρτητο από μετρική, ενώ το δεξί κατασκευάζεται εξ ολοκλήρου από μετρικά αντικείμενα, εγείρει το ερώτημα: πώς γίνεται η μετρική να εξαφανίζεται από το τελικό αποτέλεσμα; Η απάντηση βρίσκεται στην αλληλοαναίρεση δύο αντίστροφων εξαρτήσεων.

Σε τοπικές συντεταγμένες, με $\omega^a = \omega^a_j \, dx^j$, το αριστερό μέλος είναι
\[
\omega(\xi_1, \dots, \xi_{n-1}) = \det \left( \omega^a_j \, \xi_i^j \right)_{a,i=1}^{n-1},
\]
όπου δεν εμφανίζεται πουθενά η μετρική. Το δεξί μέλος υπολογίζεται ως εξής. Από την (127),
\[
(\star \omega)_j = \frac{1}{(n-1)!} \, \varepsilon_{i_1 \dots i_{n-1} j} \, \omega^{i_1 \dots i_{n-1}},
\]
και από την (131),
\[
(\xi_1 \times \dots \times \xi_{n-1})^j = \frac{1}{\sqrt{|\det(g_{ij})|}} \ \varepsilon^{j}_{\ k_1 \dots k_{n-1}} \ \xi_1^{k_1} \dots \xi_{n-1}^{k_{n-1}}.
\]
Το γινόμενό τους είναι
\[
(\star \omega)_j (\xi_1 \times \dots \times \xi_{n-1})^j = \frac{1}{(n-1)!} \varepsilon_{i_1 \dots i_{n-1} j} \, \omega^{i_1 \dots i_{n-1}} \cdot \frac{1}{\sqrt{|\det(g_{ij})|}} \varepsilon^{j}_{\ k_1 \dots k_{n-1}} \xi_1^{k_1} \dots \xi_{n-1}^{k_{n-1}}.
\]
Χρησιμοποιώντας την ταυτότητα συναίρεσης
\[
\varepsilon_{i_1 \dots i_{n-1} j} \, \varepsilon^{j}_{\ k_1 \dots k_{n-1}} = (n-1)! \, \delta^{[i_1}_{k_1} \dots \delta^{i_{n-1}]}_{k_{n-1}},
\]
το δεξί μέλος γίνεται
\[
\frac{1}{\sqrt{|\det(g_{ij})|}} \ \omega^{i_1 \dots i_{n-1}} \, \xi_1^{[i_1} \dots \xi_{n-1}^{i_{n-1}]}.
\]
Αλλά $\omega^{i_1 \dots i_{n-1}} = g^{i_1 \alpha_1} \dots g^{i_{n-1} \alpha_{n-1}} \omega_{\alpha_1 \dots \alpha_{n-1}}$, και οι συνιστώσες $\omega_{\alpha_1 \dots \alpha_{n-1}}$ περιέχουν έναν παράγοντα $\sqrt{|\det(g_{ij})|}$ από τον τανυστή Levi-Civita. Αυτός ο παράγοντας ακυρώνεται ακριβώς με τον $1/\sqrt{|\det(g_{ij})|}$ του εξωτερικού γινομένου. Οι ανυψώσεις $g^{ij}$ αντισταθμίζονται από τις συστολές με τα $\delta$ που προκύπτουν από τη συναίρεση των $\varepsilon$. Το τελικό αποτέλεσμα είναι ακριβώς το $\det(\omega^a_j \, \xi_i^j)$, χωρίς κανένα ίχνος της μετρικής.

Η μετρική είναι λοιπόν ένας ενδιάμεσος — ένα βολικό εργαλείο που επιτρέπει τη μετάφραση ανάμεσα στις δύο περιγραφές. Αλλά το τελικό αποτέλεσμα είναι ανεξάρτητο από την επιλογή της: όποια μετρική και αν χρησιμοποιήσουμε για να ορίσουμε τον Hodge και το εξωτερικό γινόμενο, η σύνθεσή τους θα δώσει πάντα την ίδια, προ-μετρική ποσότητα.

\subsection{Φυσικογεωμετρικές ερμηνείες}

Η ταυτότητα (149) βρίσκει συγκεκριμένες εφαρμογές σε διάφορους κλάδους της φυσικής, όπου ο δυϊσμός ροής-έργου και το εξωτερικό γινόμενο συνδυάζονται για να δώσουν θεμελιώδεις νόμους.

\subsubsection{Ηλεκτροστατική: ροή του ηλεκτρικού πεδίου}

Στην ηλεκτροστατική, το ηλεκτρικό πεδίο περιγράφεται από την 1-μορφή $E = E_j \, dx^j$. Η ροή του ηλεκτρικού πεδίου μέσα από ένα στοιχείο επιφάνειας που ορίζεται από τα διανύσματα $\xi, \eta$ (στον τρισδιάστατο χώρο) είναι
\[
\Phi_E = (\star E)(\xi, \eta).
\]
Σύμφωνα με την (149), αυτή η ροή μπορεί να εκφραστεί ισοδύναμα ως
\[
\Phi_E = (\star E)(\xi, \eta) = E(\xi \times \eta),
\]
αφού $\star (\star E) = E$ (με κατάλληλο πρόσημο). Το δεξί μέλος είναι το έργο της 1-μορφής $E$ κατά μήκος του εξωτερικού γινομένου $\xi \times \eta$. Αυτό σημαίνει ότι η ηλεκτρική ροή μέσα από ένα στοιχείο επιφάνειας ισούται με το έργο του ηλεκτρικού πεδίου κατά μήκος ενός διανύσματος που είναι κάθετο στην επιφάνεια και έχει μέτρο ίσο με το εμβαδόν της. Σε διανυσματική γλώσσα, αυτό είναι ακριβώς η γνωστή σχέση
\[
\Phi_E = \mathbf{E} \cdot \mathbf{n} \, \Delta A,
\]
όπου $\mathbf{n}$ είναι το μοναδιαίο κάθετο διάνυσμα και $\Delta A$ είναι το εμβαδόν. Το εξωτερικό γινόμενο $\xi \times \eta$ είναι ακριβώς το διάνυσμα $\mathbf{n} \, \Delta A$. Ο νόμος του Gauss,
\[
\int_{\partial V} \star E = \int_V \rho \, dV,
\]
μπορεί λοιπόν να ερμηνευθεί ως η ισότητα ανάμεσα στη συνολική ροή (ολοκλήρωμα της $\star E$) και στο συνολικό φορτίο, ή ισοδύναμα, ως η ισότητα ανάμεσα στο έργο του $E$ κατά μήκος του εξωτερικού γινομένου των πλευρών κάθε στοιχείου επιφάνειας και στο φορτίο που περικλείεται.

\subsubsection{Ρευστοδυναμική: ροή μάζας και κυκλοφορία}

Στη ρευστοδυναμική, το πεδίο ταχύτητας ενός ρευστού είναι ένα διανυσματικό πεδίο $u$. Η αντίστοιχη 1-μορφή είναι η $u^\flat = u_j \, dx^j$, όπου $u_j = g_{ij} u^i$. Η ροή μάζας μέσα από ένα στοιχείο επιφάνειας που ορίζεται από τα $\xi, \eta$ (στον τρισδιάστατο χώρο) είναι
\[
\dot{m} = \rho \, (\star u^\flat)(\xi, \eta),
\]
όπου $\rho$ είναι η πυκνότητα. Χρησιμοποιώντας την (149) με $\omega = \star u^\flat$ (οπότε $\star \omega = u^\flat$), έχουμε
\[
\dot{m} = \rho \, u^\flat(\xi \times \eta) = \rho \, g(u, \xi \times \eta).
\]
Αυτό σημαίνει ότι η ροή μάζας μέσα από ένα στοιχείο επιφάνειας ισούται με την πυκνότητα επί το εσωτερικό γινόμενο της ταχύτητας με το κάθετο διάνυσμα $\xi \times \eta$. Σε διανυσματική γλώσσα, $\dot{m} = \rho \, \mathbf{u} \cdot \mathbf{n} \, \Delta A$. Η κυκλοφορία της $u^\flat$ κατά μήκος μιας κλειστής καμπύλης $\gamma$ συνδέεται με τη ροή της $\star u^\flat$ μέσα από την επιφάνεια που ορίζει η $\gamma$, και ο δυϊσμός αυτός εκφράζεται ακριβώς μέσω της (149).

\subsubsection{Μαγνητοστατική: μαγνητική ροή και έργο}

Στη μαγνητοστατική, το μαγνητικό πεδίο περιγράφεται φυσικά από μια 2-μορφή $B = \frac{1}{2} B_{ij} \, dx^i \wedge dx^j$. Η μαγνητική ροή μέσα από μια επιφάνεια που ορίζεται από τα $\xi, \eta$ είναι
\[
\Phi_B = B(\xi, \eta).
\]
Στον τρισδιάστατο χώρο, η δυϊκή Hodge του $B$ είναι η 1-μορφή $\star B$, και η (149) δίνει
\[
B(\xi, \eta) = (\star B)(\xi \times \eta).
\]
Το δεξί μέλος είναι το έργο της 1-μορφής $\star B$ κατά μήκος του εξωτερικού γινομένου $\xi \times \eta$. Σε διανυσματική γλώσσα, $\star B$ αντιστοιχεί στο γνωστό διάνυσμα του μαγνητικού πεδίου $\mathbf{B}$, και η παραπάνω σχέση γράφεται
\[
\Phi_B = \mathbf{B} \cdot \mathbf{n} \, \Delta A.
\]
Ο νόμος του Amp\`ere,
\[
\oint_\gamma \star B = \int_\Sigma J,
\]
συνδέει το έργο της $\star B$ κατά μήκος μιας κλειστής καμπύλης με τη ροή της 2-μορφής $J$ (πυκνότητα ρεύματος) μέσα από την επιφάνεια. Και πάλι, ο δυϊσμός εκφράζεται μέσω της (149).

\subsection{Σύνοψη}

Το εξωτερικό γινόμενο και ο δυϊσμός ροής-έργου είναι δύο όψεις του ίδιου γεωμετρικού νομίσματος, και η ταυτότητα (149) είναι η ακριβής μαθηματική έκφραση αυτής της ενότητας. Το αριστερό μέλος, $\omega(\xi_1, \dots, \xi_{n-1})$, είναι η προ-μετρική ροή μιας $(n-1)$-μορφής μέσα από ένα παραλληλεπίπεδο. Το δεξί μέλος, $(\star \omega)(\xi_1 \times \dots \times \xi_{n-1})$, είναι το μετρικό έργο της δυϊκής 1-μορφής κατά μήκος του εξωτερικού γινομένου. Η ισότητα των δύο μελών δείχνει ότι η μετρική, αν και απαραίτητη για τον ορισμό του Hodge και του εξωτερικού γινομένου, απαλοίφεται από το τελικό αποτέλεσμα — η σύνθεση $\star$ και $\times$ είναι, στην πραγματικότητα, μια προ-μετρική κατασκευή. Αυτό σημαίνει ότι ο δυϊσμός ροής-έργου δεν είναι μια ιδιότητα που επιβάλλεται από τη μετρική, αλλά μια θεμελιώδης γεωμετρική αλήθεια που η μετρική απλώς αποκαλύπτει. Στη φυσική, αυτή η ταυτότητα εκδηλώνεται σε θεμελιώδεις νόμους όπως ο νόμος του Gauss, ο νόμος του Amp\`ere και η σχέση ροής-κυκλοφορίας στη ρευστοδυναμική, επιβεβαιώνοντας ότι η ροή και το έργο είναι συμπληρωματικές όψεις της ίδιας φυσικής πραγματικότητας.

\section{Η Εξωτερική Παράγωγος}

\subsection{Ορισμός μέσω δράσης σε διανύσματα}

Στα προηγούμενα κεφάλαια, μελετήσαμε τις διαφορικές μορφές ως αντικείμενα που δρουν σε διανύσματα σε ένα μεμονωμένο σημείο — δηλαδή, ως τανυστές σε έναν διανυσματικό χώρο $V$. Όταν όμως θεωρούμε μια διαφορική μορφή $\omega$ ως πεδίο πάνω σε μια πολλαπλότητα $M$, δηλαδή ως μια αντιστοίχιση που σε κάθε σημείο $p \in M$ αποδίδει μια $k$-μορφή $\omega(p)$ στον εφαπτόμενο χώρο $T_pM$, τίθεται το ερώτημα: πώς μπορούμε να παραγωγίσουμε μια τέτοια μορφή με τρόπο που να σέβεται τη γεωμετρική της φύση; Η απάντηση είναι η \textbf{εξωτερική παράγωγος} $d$, η οποία απεικονίζει $k$-μορφές σε $(k+1)$-μορφές.

Ο αφηρημένος ορισμός της εξωτερικής παραγώγου βασίζεται στη δράση της πάνω σε διανύσματα. Για μια $k$-μορφή $\omega$, η $d\omega$ είναι η $(k+1)$-μορφή που ορίζεται από τον τύπο

\[
\begin{aligned}
d\omega(\xi_0, \xi_1, \dots, \xi_k) &= \sum_{i=0}^{k} (-1)^i \ \xi_i \left( \omega(\xi_0, \dots, \widehat{\xi_i}, \dots, \xi_k) \right) \\
&\quad + \sum_{0 \le i < j \le k} (-1)^{i+j} \ \omega([\xi_i, \xi_j], \xi_0, \dots, \widehat{\xi_i}, \dots, \widehat{\xi_j}, \dots, \xi_k),
\end{aligned} \qquad (150)
\]

όπου $\xi_0, \dots, \xi_k$ είναι διανυσματικά πεδία, η περισπωμένη $\widehat{\xi_i}$ δηλώνει ότι το όρισμα $\xi_i$ παραλείπεται, και $[\xi_i, \xi_j]$ είναι το γινόμενο Lie των δύο διανυσματικών πεδίων. Ο πρώτος όρος μετρά τη μεταβολή της $\omega$ κατά μήκος των διανυσματικών πεδίων (είναι ουσιαστικά μια κατευθυνόμενη παράγωγος), ενώ ο δεύτερος όρος διορθώνει για τη μη-αντιμεταθετικότητα των ροών, όπως αυτή εκφράζεται μέσω του γινομένου Lie.

Είναι κρίσιμο να παρατηρήσουμε ότι το γινόμενο Lie $[\xi_i, \xi_j]$ εμφανίζεται ακριβώς στον δεύτερο όρο του ορισμού. Αυτός ο όρος είναι υπεύθυνος για την κωδικοποίηση της πληροφορίας σχετικά με την ολοκληρωσιμότητα της κατανομής που ορίζουν τα διανύσματα. Όταν τα διανύσματα $\xi_i$ ανήκουν σε μια κατανομή, το γινόμενο Lie $[\xi_i, \xi_j]$ μετρά το "κενό" ανάμεσα στις ροές τους — αν το γινόμενο Lie παραμένει στην κατανομή, οι ροές είναι συμβατές και η κατανομή είναι ολοκληρώσιμη. Ο δεύτερος όρος της (150) αφαιρεί ακριβώς τη συνεισφορά αυτού του κενού, ώστε η $d\omega$ να μετρά μόνο την "εξωτερική" μεταβολή της $\omega$, ανεξάρτητα από τη μη-αντιμεταθετικότητα των ροών. Με άλλα λόγια, η $d\omega$ είναι κατασκευασμένη έτσι ώστε να "αγνοεί" τις εσωτερικές στρεβλώσεις της κατανομής και να μετρά μόνο την αποτυχία της να είναι ολοκληρώσιμη.

Για τις ανάγκες της παρούσας ανάπτυξης, θα επικεντρωθούμε στην περίπτωση όπου τα διανυσματικά πεδία είναι τα διανύσματα βάσης $\partial / \partial x^i$ ενός συστήματος συντεταγμένων. Σε αυτή την περίπτωση, τα γινόμενα Lie μηδενίζονται ($[\partial / \partial x^i, \partial / \partial x^j] = 0$), και ο παραπάνω τύπος απλοποιείται δραματικά. Για μια $k$-μορφή $\omega$ και για διανύσματα βάσης $\partial / \partial x^{i_0}, \dots, \partial / \partial x^{i_k}$, έχουμε

\[
d\omega \left( \frac{\partial}{\partial x^{i_0}}, \dots, \frac{\partial}{\partial x^{i_k}} \right) = \sum_{a=0}^{k} (-1)^a \ \frac{\partial}{\partial x^{i_a}} \left( \omega \left( \frac{\partial}{\partial x^{i_0}}, \dots, \widehat{\frac{\partial}{\partial x^{i_a}}}, \dots, \frac{\partial}{\partial x^{i_k}} \right) \right). \qquad (151)
\]

Αυτός ο τύπος είναι η βάση για την εξαγωγή της έκφρασης της $d\omega$ σε τοπικές συντεταγμένες.

\subsection{Έκφραση σε τοπικές συντεταγμένες}

Σε ένα τοπικό σύστημα συντεταγμένων $x^1, \dots, x^n$, μια $k$-μορφή $\omega$ αναπαρίσταται ως

\[
\omega = \frac{1}{k!} \ \omega_{i_1 \dots i_k} \ dx^{i_1} \wedge \dots \wedge dx^{i_k},
\]

όπου οι συνιστώσες $\omega_{i_1 \dots i_k}$ είναι ομαλές συναρτήσεις των συντεταγμένων και είναι πλήρως αντισυμμετρικές στους δείκτες. Η εξωτερική παράγωγος $d\omega$ είναι τότε η $(k+1)$-μορφή

\[
d\omega = \frac{1}{k!} \ \frac{\partial \omega_{i_1 \dots i_k}}{\partial x^j} \ dx^j \wedge dx^{i_1} \wedge \dots \wedge dx^{i_k}.
\]

Λόγω της αντισυμμετρίας του σφηνοειδούς γινομένου, μπορούμε να αντισυμμετρικοποιήσουμε πλήρως τους δείκτες. Το αποτέλεσμα είναι

\[
d\omega = \frac{1}{(k+1)!} \ (d\omega)_{i_0 i_1 \dots i_k} \ dx^{i_0} \wedge dx^{i_1} \wedge \dots \wedge dx^{i_k},
\]

όπου οι συνιστώσες της $d\omega$ δίνονται από τον τύπο

\[
(d\omega)_{i_0 i_1 \dots i_k} = \sum_{a=0}^{k} (-1)^a \ \frac{\partial}{\partial x^{i_a}} \ \omega_{i_0 \dots \widehat{i_a} \dots i_k}. \qquad (152)
\]

Αυτός είναι ο κλασικός τύπος για την εξωτερική παράγωγο σε τοπικές συντεταγμένες. Ισοδύναμα, μπορούμε να γράψουμε

\[
d\omega = \sum_{j=1}^{n} \frac{\partial \omega}{\partial x^j} \wedge dx^j,
\]

όπου η μερική παράγωγος $\partial \omega / \partial x^j$ νοείται ως η $k$-μορφή που προκύπτει από την παραγώγιση των συνιστωσών της $\omega$ ως προς $x^j$.

\subsection{Βασικές ταυτότητες}

Η εξωτερική παράγωγος ικανοποιεί μια σειρά από θεμελιώδεις ταυτότητες που την καθιστούν ένα από τα πιο χρήσιμα εργαλεία της διαφορικής γεωμετρίας.

\textbf{Μηδενισμός της διπλής παραγώγου.} Για οποιαδήποτε $k$-μορφή $\omega$, ισχύει

\[
d(d\omega) = 0. \qquad (153)
\]

Αυτή η ταυτότητα είναι συνέπεια της ισότητας των μεικτών μερικών παραγώγων, $\partial^2 / \partial x^i \partial x^j = \partial^2 / \partial x^j \partial x^i$, και της αντισυμμετρίας του σφηνοειδούς γινομένου. Συνοψίζεται στη φράση "το σύνορο ενός συνόρου είναι κενό".

\textbf{Κανόνας Leibniz.} Για μια $k$-μορφή $\omega$ και μια $l$-μορφή $\eta$, ισχύει

\[
d(\omega \wedge \eta) = d\omega \wedge \eta + (-1)^k \ \omega \wedge d\eta. \qquad (154)
\]

Αυτός ο κανόνας γενικεύει τον κανόνα του Leibniz για την παράγωγο γινομένου, με το πρόσημο $(-1)^k$ να οφείλεται στο γεγονός ότι η $d$ "περνά" μέσα από την $\omega$ για να δράσει στην $\eta$.

\textbf{Εξωτερική παράγωγος απλής μορφής.} Για μια απλή $k$-μορφή $\omega = \omega^1 \wedge \dots \wedge \omega^k$, η εξωτερική παράγωγος αναπτύσσεται ως

\[
d\omega = \sum_{a=1}^{k} (-1)^{a-1} \ \omega^1 \wedge \dots \wedge d\omega^a \wedge \dots \wedge \omega^k. \qquad (155)
\]

Αυτός ο τύπος προκύπτει άμεσα από επαναλαμβανόμενη εφαρμογή του κανόνα Leibniz (154), χρησιμοποιώντας το γεγονός ότι $d(\omega^a)$ είναι μια 2-μορφή και ότι το σφηνοειδές γινόμενο είναι αντισυμμετρικό. Το πρόσημο $(-1)^{a-1}$ οφείλεται στο γεγονός ότι η $d$ πρέπει να "περάσει" μέσα από $a-1$ μορφές $\omega^1, \dots, \omega^{a-1}$ για να φτάσει στην $\omega^a$.

\textbf{Δράση σε βαθμωτές συναρτήσεις.} Για μια βαθμωτή συνάρτηση $f$ (μια $0$-μορφή), η $df$ είναι η 1-μορφή

\[
df = \frac{\partial f}{\partial x^j} \ dx^j, \qquad (156)
\]

της οποίας η δράση σε ένα διάνυσμα $\xi$ είναι η κατευθυνόμενη παράγωγος $\xi(f) = \mathcal{L}_\xi f$.

\textbf{Αναλλοιότητα υπό αλλαγή συντεταγμένων.} Η εξωτερική παράγωγος είναι ανεξάρτητη από την επιλογή συντεταγμένων. Αν αλλάξουμε σύστημα συντεταγμένων, η $d\omega$ που υπολογίζεται στο νέο σύστημα ταυτίζεται με τη $d\omega$ που υπολογίζεται στο παλιό — ο τύπος μετασχηματισμού των συντεταγμένων εξασφαλίζει αυτή την αναλλοιότητα.

\subsection{Γεωμετρική ερμηνεία της εξωτερικής παραγώγου απλών μορφών}

Η πιο σημαντική γεωμετρική ερμηνεία της εξωτερικής παραγώγου προκύπτει όταν την εφαρμόζουμε σε απλές μορφές, και συνδέεται άμεσα με τη θεωρία του ριζικού υπόχωρου και του μηδενοχώρου που αναπτύξαμε στο Κεφάλαιο 16. Το κλειδί αυτής της σύνδεσης είναι ακριβώς το γινόμενο Lie που εμφανίζεται στον ορισμό (150): αυτό κωδικοποιεί την πληροφορία για τη μη-αντιμεταθετικότητα των ροών, η οποία με τη σειρά της καθορίζει αν ο ριζικός υπόχωρος είναι ολοκληρώσιμος.

Θεωρούμε μια απλή $k$-μορφή $\omega = \omega^1 \wedge \omega^2 \wedge \dots \wedge \omega^k$. Σε κάθε σημείο, η $\omega$ ορίζει έναν ριζικό υπόχωρο $\rho(\omega)$ διάστασης $n-k$ και έναν υπόχωρο προβολής $\pi(\omega)$ διάστασης $k$. Η $\omega$ μηδενίζεται ακριβώς σε εκείνα τα $k$-παραλληλεπίπεδα που δεν είναι πλήρως εγκάρσια στον $\rho(\omega)$ — δηλαδή, των οποίων η προβολή στον $\pi(\omega)$ έχει διάσταση μικρότερη του $k$.

Η εξωτερική παράγωγος $d\omega$ είναι μια $(k+1)$-μορφή. Από τον τύπο (155), η $d\omega$ είναι ένα άθροισμα όρων, καθένας από τους οποίους περιέχει την εξωτερική παράγωγο μιας από τις 1-μορφές $\omega^a$. Τι σημαίνει γεωμετρικά ο μηδενισμός ή μη της $d\omega$;

\subsubsection{Η σημασία του μηδενισμού $d\omega = 0$}

Αν $d\omega = 0$, τότε από την (155) έχουμε

\[
\sum_{a=1}^{k} (-1)^{a-1} \ \omega^1 \wedge \dots \wedge d\omega^a \wedge \dots \wedge \omega^k = 0.
\]

Αυτό σημαίνει ότι κάθε $d\omega^a$ μπορεί να γραφεί ως γραμμικός συνδυασμός σφηνοειδών γινομένων που περιέχουν τουλάχιστον μία από τις $\omega^b$. Με άλλα λόγια, η $d\omega^a$ δεν έχει "ανεξάρτητες" συνιστώσες — όλη η πληροφορία της περιέχεται στις $\omega^1, \dots, \omega^k$.

Για να καταλάβουμε τι σημαίνει αυτό γεωμετρικά, ας θυμηθούμε τον ορισμό (150) για $k=1$. Για μια 1-μορφή $\omega^a$ και δύο διανύσματα $\xi, \eta$ που ανήκουν στον $\rho(\omega) = \bigcap_{b=1}^k \ker \omega^b$, έχουμε $\omega^a(\xi) = \omega^a(\eta) = 0$, και ο ορισμός δίνει

\[
d\omega^a(\xi, \eta) = \xi(\omega^a(\eta)) - \eta(\omega^a(\xi)) - \omega^a([\xi, \eta]) = -\omega^a([\xi, \eta]).
\]

Εδώ βλέπουμε ρητά πώς το γινόμενο Lie $[\xi, \eta]$ υπεισέρχεται: η $d\omega^a(\xi, \eta)$ μετρά την τιμή της $\omega^a$ πάνω στο γινόμενο Lie των $\xi$ και $\eta$. Αν $d\omega = 0$, τότε $d\omega^a(\xi, \eta) = 0$ για κάθε $\xi, \eta \in \rho(\omega)$, άρα $\omega^a([\xi, \eta]) = 0$ για κάθε $a$. Αυτό σημαίνει ότι το γινόμενο Lie $[\xi, \eta]$ ανήκει επίσης στον $\rho(\omega)$. Συνεπώς, ο $\rho(\omega)$ είναι κλειστός ως προς το γινόμενο Lie — δηλαδή, είναι \textbf{ολοκληρώσιμος}.

Γεωμετρικά, η ολοκληρωσιμότητα του $\rho(\omega)$ σημαίνει ότι υπάρχουν $(n-k)$-διάστατες υποπολλαπλότητες που είναι παντού εφαπτόμενες στον $\rho(\omega)$. Σε κάθε σημείο μιας τέτοιας υποπολλαπλότητας, ο εφαπτόμενος χώρος της ταυτίζεται με τον $\rho(\omega)$ στο σημείο αυτό. Οι υπόχωροι $\rho(\omega)$ "κολλάνε" μεταξύ τους για να σχηματίσουν μια οικογένεια υποπολλαπλοτήτων που φυλλώνουν τον χώρο. Αυτές οι υποπολλαπλότητες είναι ακριβώς οι επιφάνειες πάνω στις οποίες η $\omega$ μηδενίζεται ταυτοτικά. Η $\omega$ "ζει" ουσιαστικά στον χώρο πηλίκο $V/\rho(\omega)$, μετρώντας προσανατολισμένους $k$-διάστατους όγκους στον $\pi(\omega)$.

\subsubsection{Η σημασία του μη μηδενισμού $d\omega \neq 0$}

Αν $d\omega \neq 0$, τότε τουλάχιστον μία από τις $d\omega^a$ έχει μια συνιστώσα που δεν μπορεί να εκφραστεί ως γραμμικός συνδυασμός των $\omega^b$. Από την παραπάνω ανάλυση, αυτό σημαίνει ότι υπάρχουν $\xi, \eta \in \rho(\omega)$ τέτοια ώστε $d\omega^a(\xi, \eta) \neq 0$, άρα $\omega^a([\xi, \eta]) \neq 0$. Το γινόμενο Lie $[\xi, \eta]$ δεν ανήκει πλέον στον $\rho(\omega)$ — έχει μια συνιστώσα έξω από αυτόν. Αυτό σημαίνει ότι ο $\rho(\omega)$ \textbf{δεν είναι ολοκληρώσιμος}: οι ροές δύο διανυσματικών πεδίων που ανήκουν στον $\rho(\omega)$ παράγουν ένα "κενό" που βγαίνει έξω από τον $\rho(\omega)$. Οι υπόχωροι $\rho(\omega)$ σε γειτονικά σημεία "στρίβουν" με τρόπο που δεν επιτρέπει τη συγκόλλησή τους σε μια υποπολλαπλότητα. Η $d\omega$ μετρά ακριβώς αυτή την αποτυχία ολοκλήρωσης, και το γινόμενο Lie είναι ο μηχανισμός μέσω του οποίου αυτή η αποτυχία κωδικοποιείται.

\subsubsection{Συγκεκριμένες περιπτώσεις}

\textbf{Περίπτωση $k=1$: η στροβιλότητα μιας δύναμης.} Μια 1-μορφή $\omega$ αντιπροσωπεύει μια δύναμη. Η $d\omega$ είναι μια 2-μορφή που μετρά τη "στροβιλότητα" της δύναμης. Ο ριζικός υπόχωρος $\rho(\omega) = \ker \omega$ είναι ένα υπερεπίπεδο διάστασης $n-1$ σε κάθε σημείο.

Αν $d\omega = 0$, τότε για κάθε $\xi, \eta \in \ker \omega$, έχουμε $\omega([\xi, \eta]) = 0$, άρα $[\xi, \eta] \in \ker \omega$. Ο $\ker \omega$ είναι ολοκληρώσιμος: υπάρχουν υπερεπιφάνειες (διάστασης $n-1$) παντού εφαπτόμενες στον $\ker \omega$. Αυτές είναι οι ισοδυναμικές επιφάνειες ενός δυναμικού: υπάρχει τοπικά μια βαθμωτή συνάρτηση $f$ τέτοια ώστε $\omega = df$. Το έργο της $\omega$ κατά μήκος μιας κλειστής καμπύλης είναι μηδέν — η δύναμη είναι συντηρητική.

Αν $d\omega \neq 0$, υπάρχουν $\xi, \eta \in \ker \omega$ με $\omega([\xi, \eta]) \neq 0$. Ο $\ker \omega$ δεν είναι ολοκληρώσιμος, δεν υπάρχει δυναμικό, και η δύναμη δεν είναι συντηρητική. Η $d\omega$ μετρά τη "στροβιλότητα" με την εξής έννοια: για δύο διανύσματα $\xi, \eta$, η ποσότητα $d\omega(\xi, \eta)$ είναι ακριβώς το έργο της $\omega$ γύρω από το απειροστό παραλληλόγραμμο που ορίζουν τα $\xi, \eta$. Αυτό το παραλληλόγραμμο "κλείνει" μόνο αν οι ροές των $\xi$ και $\eta$ αντιμετατίθενται — δηλαδή, αν $[\xi, \eta] = 0$. Το γινόμενο Lie μετρά το "κενό" στο παραλληλόγραμμο, και η $d\omega$ μετρά το έργο της $\omega$ γύρω από αυτό το κενό. Αν $d\omega \neq 0$, το έργο γύρω από έναν κλειστό βρόχο δεν είναι μηδέν, και αυτό οφείλεται ακριβώς στο γεγονός ότι ο $\ker \omega$ δεν είναι ολοκληρώσιμος — οι ισοδυναμικές επιφάνειες δεν υπάρχουν.

Η στροβιλότητα έχει λοιπόν μια βαθιά γεωμετρική ερμηνεία: είναι το μέτρο της αποτυχίας του πυρήνα της 1-μορφής να είναι ολοκληρώσιμος. Σε έναν τρισδιάστατο χώρο, η στροβιλότητα ενός διανυσματικού πεδίου $v$ είναι η 2-μορφή $d(v^\flat)$, και ο μηδενισμός της σημαίνει ότι τα υπερεπίπεδα που είναι κάθετα στο $v$ (τα οποία είναι ακριβώς ο $\ker(v^\flat)$) είναι ολοκληρώσιμα — δηλαδή, υπάρχουν επιφάνειες παντού κάθετες στο $v$. Αυτές είναι οι ισοδυναμικές επιφάνειες του δυναμικού από το οποίο προέρχεται το $v$. Αν η στροβιλότητα δεν είναι μηδέν, αυτές οι επιφάνειες δεν υπάρχουν, και το πεδίο "στροβιλίζεται".

\textbf{Περίπτωση $k=n-1$: η απόκλιση μιας ροής.} Μια $(n-1)$-μορφή $\eta$ αντιπροσωπεύει μια ροή. Ο $\rho(\eta)$ είναι μια μονοδιάστατη ευθεία σε κάθε σημείο — η κατεύθυνση ροής. Μια μονοδιάστατη κατανομή είναι πάντα ολοκληρώσιμη (οι ολοκληρωτικές καμπύλες υπάρχουν πάντα), οπότε η συνθήκη $d\eta = 0$ δεν σχετίζεται με την ολοκληρωσιμότητα αυτή καθεαυτή. Αντίθετα, η $d\eta$ είναι μια $n$-μορφή, δηλαδή μια πυκνότητα. Σε έναν χώρο με μετρική, η $d\eta$ μπορεί να εκφραστεί ως η απόκλιση του διανυσματικού πεδίου που αντιστοιχεί στην $\eta$ μέσω του Hodge. Ο μηδενισμός $d\eta = 0$ σημαίνει ότι η ροή είναι ασυμπίεστη: ο όγκος ενός στοιχείου ρευστού διατηρείται κατά μήκος της ροής.

Γεωμετρικά, η $d\eta$ μετρά τη "διαστολή" της ροής. Αν $d\eta = 0$, ο όγκος ενός παραλληλεπιπέδου που μεταφέρεται από τη ροή παραμένει σταθερός. Αν $d\eta \neq 0$, ο όγκος αλλάζει, και η $d\eta$ ποσοτικοποιεί αυτή την αλλαγή.

\subsection{Το θεώρημα Frobenius σε όρους διαφορικών μορφών}

Η γεωμετρική ερμηνεία που αναπτύξαμε στην προηγούμενη ενότητα δεν είναι απλώς μια διαισθητική εικόνα — είναι ένα αυστηρό θεώρημα, η δυϊκή διατύπωση του θεωρήματος Frobenius που μελετήσαμε στο Κεφάλαιο 15. Εκεί είχαμε διατυπώσει το θεώρημα σε όρους διανυσματικών πεδίων και γινομένου Lie, όπου η συνθήκη ολοκληρωσιμότητας ήταν ακριβώς $[\xi, \eta] \in \rho$ για κάθε $\xi, \eta \in \rho$. Τώρα είμαστε σε θέση να διατυπώσουμε και να αποδείξουμε την ισοδύναμη εκδοχή του σε όρους διαφορικών μορφών και εξωτερικής παραγώγου, όπου το γινόμενο Lie αντικαθίσταται από την $d$.

\textbf{Θεώρημα (Frobenius — εκδοχή διαφορικών μορφών).} Έστω $\omega^1, \dots, \omega^k$ γραμμικά ανεξάρτητες 1-μορφές σε μια πολλαπλότητα $M$ διάστασης $n$, που ορίζουν μια κατανομή $\rho = \bigcap_{a=1}^k \ker \omega^a$ διάστασης $n-k$. Τότε η κατανομή $\rho$ είναι ολοκληρώσιμη αν και μόνο αν

\[
d\omega^a = \sum_{b=1}^{k} \theta^a_b \wedge \omega^b, \qquad a = 1, \dots, k,
\]

για κάποιες 1-μορφές $\theta^a_b$. Ισοδύναμα, η απλή $k$-μορφή $\omega = \omega^1 \wedge \dots \wedge \omega^k$ ικανοποιεί

\[
d\omega = \theta \wedge \omega,
\]

για κάποια 1-μορφή $\theta$.

\textbf{Απόδειξη.} Η απόδειξη βασίζεται ακριβώς στη σύνδεση ανάμεσα στην εξωτερική παράγωγο και το γινόμενο Lie που είναι ενσωματωμένη στον ορισμό (150). Θεωρούμε δύο διανυσματικά πεδία $\xi, \eta$ που ανήκουν στην κατανομή $\rho$, δηλαδή $\omega^a(\xi) = \omega^a(\eta) = 0$ για κάθε $a = 1, \dots, k$. Από τον ορισμό (150) της $d\omega^a$ για $k=1$, έχουμε

\[
d\omega^a(\xi, \eta) = \xi(\omega^a(\eta)) - \eta(\omega^a(\xi)) - \omega^a([\xi, \eta]).
\]

Αφού $\omega^a(\xi) = \omega^a(\eta) = 0$, οι δύο πρώτοι όροι μηδενίζονται, και παίρνουμε

\[
d\omega^a(\xi, \eta) = -\omega^a([\xi, \eta]). \qquad (157)
\]

Αυτή η ταυτότητα είναι το κλειδί που συνδέει την εξωτερική παράγωγο με την ολοκληρωσιμότητα: η $d\omega^a$ πάνω σε δύο πεδία της κατανομής μετρά ακριβώς την τιμή της $\omega^a$ πάνω στο γινόμενο Lie τους.

Υποθέτουμε τώρα ότι η κατανομή $\rho$ είναι ολοκληρώσιμη. Από το θεώρημα Frobenius στη διατύπωση του Κεφαλαίου 15 (εξίσωση (105)), αυτό σημαίνει ότι για κάθε $\xi, \eta \in \rho$, έχουμε $[\xi, \eta] \in \rho$, δηλαδή $\omega^a([\xi, \eta]) = 0$ για κάθε $a$. Από την (157), αυτό συνεπάγεται

\[
d\omega^a(\xi, \eta) = 0 \qquad \forall \xi, \eta \in \rho.
\]

Αυτό σημαίνει ότι η 2-μορφή $d\omega^a$ μηδενίζεται πάνω στην κατανομή $\rho$. Με άλλα λόγια, η $d\omega^a$ ανήκει στο γραμμικό άνοιγμα των σφηνοειδών γινομένων που περιέχουν τις $\omega^b$. Αυτό ακριβώς σημαίνει ότι υπάρχουν 1-μορφές $\theta^a_b$ τέτοιες ώστε

\[
d\omega^a = \sum_{b=1}^{k} \theta^a_b \wedge \omega^b.
\]

Αντίστροφα, αν ισχύει αυτή η σχέση, τότε για κάθε $\xi, \eta \in \rho$, έχουμε $\omega^b(\xi) = \omega^b(\eta) = 0$ για κάθε $b$, και το δεξί μέλος μηδενίζεται όταν δρα στα $\xi, \eta$. Άρα $d\omega^a(\xi, \eta) = 0$. Από την (157), αυτό δίνει $\omega^a([\xi, \eta]) = 0$ για κάθε $a$, δηλαδή $[\xi, \eta] \in \rho$. Από το θεώρημα Frobenius στη διατύπωση του Κεφαλαίου 15, η $\rho$ είναι ολοκληρώσιμη.

Τέλος, για την απλή $k$-μορφή $\omega = \omega^1 \wedge \dots \wedge \omega^k$, χρησιμοποιούμε τον τύπο (155):

\[
d\omega = \sum_{a=1}^{k} (-1)^{a-1} \ \omega^1 \wedge \dots \wedge d\omega^a \wedge \dots \wedge \omega^k.
\]

Αντικαθιστώντας την $d\omega^a = \sum_b \theta^a_b \wedge \omega^b$, κάθε όρος στο άθροισμα περιέχει έναν παράγοντα $\omega^b \wedge \omega^a$ (ή κάτι ανάλογο), και λόγω της αντισυμμετρίας, οι όροι με $b \neq a$ δίνουν μηδέν (αφού η $\omega^b$ εμφανίζεται δύο φορές στο σφηνοειδές γινόμενο). Ο μόνος όρος που επιζεί είναι ο $b = a$, και το αποτέλεσμα είναι

\[
d\omega = \left( \sum_{a=1}^{k} \theta^a_a \right) \wedge \omega = \theta \wedge \omega,
\]

όπου $\theta = \sum_{a=1}^{k} \theta^a_a$. Αυτό ολοκληρώνει την απόδειξη. $\square$

Αυτή η εκδοχή του θεωρήματος Frobenius είναι ιδιαίτερα χρήσιμη σε εφαρμογές, καθώς η συνθήκη $d\omega^a = \sum_b \theta^a_b \wedge \omega^b$ μπορεί να ελεγχθεί απευθείας σε τοπικές συντεταγμένες χρησιμοποιώντας τον τύπο (152). Το γινόμενο Lie, που ήταν ο πρωταγωνιστής της διατύπωσης του Κεφαλαίου 15, έχει τώρα αντικατασταθεί από την εξωτερική παράγωγο — αλλά η σύνδεση παραμένει μέσω της ταυτότητας (157), η οποία δείχνει ότι η $d$ και το $[\cdot, \cdot]$ είναι δύο όψεις του ίδιου νομίσματος.

\subsection{Το Λήμμα του Poincaré}

\subsubsection{Εισαγωγή και διατύπωση}

Είδαμε ότι η εξωτερική παράγωγος ικανοποιεί τη θεμελιώδη ταυτότητα $d^2 = 0$. Αυτό σημαίνει ότι κάθε ακριβής μορφή (μια μορφή που γράφεται ως $\omega = d\eta$) είναι αυτομάτως κλειστή ($d\omega = 0$). Το ερώτημα που τίθεται φυσικά είναι το αντίστροφο: κάθε κλειστή μορφή είναι τοπικά ακριβής; Η απάντηση είναι καταφατική, και αυτό είναι το περιεχόμενο ενός από τα πιο θεμελιώδη αποτελέσματα της διαφορικής γεωμετρίας — του \textbf{λήμματος του Poincaré}.

\textbf{Θεώρημα (Λήμμα του Poincaré).} Έστω $\omega$ μια $k$-μορφή ($k \ge 1$) ορισμένη σε μια αστεροειδή περιοχή $U \subseteq \mathbb{R}^n$ (ως προς κάποιο σημείο). Αν η $\omega$ είναι κλειστή, δηλαδή $d\omega = 0$, τότε η $\omega$ είναι ακριβής: υπάρχει μια $(k-1)$-μορφή $\eta$ ορισμένη στο $U$ τέτοια ώστε $\omega = d\eta$.

Μια αστεροειδής περιοχή ως προς ένα σημείο $x_0$ είναι ένα σύνολο $U$ τέτοιο ώστε για κάθε $x \in U$, το ευθύγραμμο τμήμα που συνδέει το $x_0$ με το $x$ περιέχεται εξ ολοκλήρου στο $U$. Το $\mathbb{R}^n$ είναι αστεροειδές ως προς οποιοδήποτε σημείο του, και κάθε κυρτή περιοχή είναι αστεροειδής. Η απαίτηση αυτή είναι ουσιώδης: σε μια μη απλά συνεκτική περιοχή, μια κλειστή μορφή μπορεί να μην είναι ακριβής.

\subsubsection{Απόδειξη}

Η απόδειξη είναι κατασκευαστική: ορίζουμε έναν τελεστή ομοτοπίας $H$ που απεικονίζει $k$-μορφές σε $(k-1)$-μορφές, τέτοιον ώστε

\[
\omega = H(d\omega) + d(H\omega). \qquad (158)
\]

Αν $d\omega = 0$, τότε $\omega = d(H\omega)$, και η $\eta = H\omega$ είναι η ζητούμενη $(k-1)$-μορφή.

Χωρίς βλάβη της γενικότητας, θεωρούμε ότι η αστεροειδής περιοχή είναι το $\mathbb{R}^n$ και ότι το σημείο αναφοράς είναι η αρχή $0$. Ο τελεστής ομοτοπίας $H$ ορίζεται ως εξής. Για μια $k$-μορφή $\omega$, η $H\omega$ είναι η $(k-1)$-μορφή που δίνεται από τον τύπο

\[
(H\omega)_{i_1 \dots i_{k-1}}(x) = \sum_{a=1}^{k} (-1)^{a-1} \int_0^1 t^{k-1} x^{i_a} \omega_{i_1 \dots i_{a-1} i_a i_{a+1} \dots i_{k-1}}(tx) \, dt.
\]

Ισοδύναμα, σε πιο γεωμετρική γλώσσα, η δράση της $H\omega$ σε $k-1$ διανύσματα $\xi_1, \dots, \xi_{k-1}$ στο σημείο $x$ είναι

\[
(H\omega)_x(\xi_1, \dots, \xi_{k-1}) = \int_0^1 t^{k-1} \omega_{tx}(x, \xi_1, \dots, \xi_{k-1}) \, dt,
\]

όπου $\omega_{tx}$ είναι η τιμή της $\omega$ στο σημείο $tx$, και το διάνυσμα $x$ (το διάνυσμα θέσης) εισάγεται ως πρώτο όρισμα. Αυτός ο τελεστής έχει την ιδιότητα ότι για κάθε $k$-μορφή $\omega$, ισχύει η (158).

Η απόδειξη της (158) γίνεται με άμεσο υπολογισμό σε τοπικές συντεταγμένες. Ας την επαληθεύσουμε για την περίπτωση $k=1$, που είναι και η πιο σημαντική για τις εφαρμογές. Για μια 1-μορφή $\omega = \omega_j \, dx^j$, η $d\omega$ είναι η 2-μορφή με συνιστώσες

\[
(d\omega)_{ij} = \frac{\partial \omega_j}{\partial x^i} - \frac{\partial \omega_i}{\partial x^j}.
\]

Η $H\omega$ είναι μια 0-μορφή (βαθμωτή συνάρτηση) που δίνεται από

\[
(H\omega)(x) = \int_0^1 x^j \omega_j(tx) \, dt.
\]

Η $H(d\omega)$ είναι μια 1-μορφή με συνιστώσες

\[
(H(d\omega))_i(x) = \int_0^1 t x^j (d\omega)_{ji}(tx) \, dt.
\]

Η $d(H\omega)$ είναι η 1-μορφή με συνιστώσες

\[
d(H\omega)_i(x) = \frac{\partial}{\partial x^i} \int_0^1 x^j \omega_j(tx) \, dt.
\]

Παραγωγίζοντας κάτω από το ολοκλήρωμα και χρησιμοποιώντας τον κανόνα της αλυσίδας, βρίσκουμε

\[
d(H\omega)_i(x) = \int_0^1 \left( \omega_i(tx) + t x^j \frac{\partial \omega_j}{\partial x^i}(tx) \right) dt.
\]

Από την άλλη,

\[
(H(d\omega))_i(x) = \int_0^1 t x^j \left( \frac{\partial \omega_i}{\partial x^j}(tx) - \frac{\partial \omega_j}{\partial x^i}(tx) \right) dt.
\]

Προσθέτοντας τις δύο συνεισφορές, οι όροι με $t x^j \partial \omega_j / \partial x^i$ αλληλοαναιρούνται, και παίρνουμε

\[
(H(d\omega))_i(x) + d(H\omega)_i(x) = \int_0^1 \left( \omega_i(tx) + t x^j \frac{\partial \omega_i}{\partial x^j}(tx) \right) dt.
\]

Αναγνωρίζουμε την ολοκληρωτέα ποσότητα ως την παράγωγο του $t \omega_i(tx)$ ως προς $t$:

\[
\frac{d}{dt} \left( t \omega_i(tx) \right) = \omega_i(tx) + t x^j \frac{\partial \omega_i}{\partial x^j}(tx).
\]

Συνεπώς, το ολοκλήρωμα είναι ακριβώς $t \omega_i(tx) \big|_0^1 = \omega_i(x)$. Άρα

\[
\omega_i = (H(d\omega))_i + d(H\omega)_i,
\]

που είναι η (158) για $k=1$. Αν $d\omega = 0$, τότε $\omega_i = d(H\omega)_i$, δηλαδή $\omega = d(H\omega)$.

Η γενική απόδειξη για αυθαίρετο $k$ ακολουθεί την ίδια λογική και βασίζεται στην παραγώγιση κάτω από το ολοκλήρωμα και στις ιδιότητες συμμετρίας των μερικών παραγώγων. Το ουσιώδες σημείο είναι ότι ο τελεστής $H$ είναι καθολικά ορισμένος σε μια αστεροειδή περιοχή, και η ταυτότητα (158) ισχύει για κάθε $k$-μορφή. Αν $d\omega = 0$, τότε $\omega = d(H\omega)$, και η απόδειξη ολοκληρώνεται.

\subsubsection{Κλασικό παράδειγμα: η 1-μορφή στον $\mathbb{R}^2$}

Το κλασικό παράδειγμα που αναδεικνύει τη σημασία της αστεροειδούς συνθήκης είναι η 1-μορφή

\[
\omega = -\frac{y}{x^2 + y^2} \, dx + \frac{x}{x^2 + y^2} \, dy,
\]

ορισμένη στον $\mathbb{R}^2 \setminus \{0\}$. Μπορούμε εύκολα να ελέγξουμε ότι η $\omega$ είναι κλειστή:

\[
d\omega = \left( \frac{\partial}{\partial x} \left( \frac{x}{x^2 + y^2} \right) + \frac{\partial}{\partial y} \left( \frac{y}{x^2 + y^2} \right) \right) dx \wedge dy = 0.
\]

Ωστόσο, η $\omega$ δεν είναι ακριβής σε ολόκληρο τον $\mathbb{R}^2 \setminus \{0\}$. Πράγματι, αν ολοκληρώσουμε την $\omega$ κατά μήκος του μοναδιαίου κύκλου $x^2 + y^2 = 1$, παίρνουμε

\[
\oint_{S^1} \omega = \int_0^{2\pi} (-\sin\theta \cdot (-\sin\theta) + \cos\theta \cdot \cos\theta) \, d\theta = \int_0^{2\pi} d\theta = 2\pi \neq 0.
\]

Αν η $\omega$ ήταν ακριβής, $\omega = df$, το ολοκλήρωμα κατά μήκος οποιασδήποτε κλειστής καμπύλης θα ήταν μηδέν. Το γεγονός ότι δεν είναι μηδέν οφείλεται στο ότι ο $\mathbb{R}^2 \setminus \{0\}$ δεν είναι αστεροειδής — έχει μια "τρύπα" στο $0$, και η $\omega$ "αντιλαμβάνεται" αυτή την τρύπα μέσω της ολικής καμπυλότητας $2\pi$.

Αν περιοριστούμε σε μια αστεροειδή περιοχή, για παράδειγμα το δεξί ημιεπίπεδο $x > 0$, τότε η $\omega$ είναι ακριβής: $\omega = d\theta$, όπου $\theta = \arctan(y/x)$ είναι η γωνία. Το λήμμα του Poincaré εγγυάται ότι αυτό μπορεί να γίνει τοπικά, αλλά όχι καθολικά.

\subsubsection{Γεωμετρική ερμηνεία}

Το λήμμα του Poincaré έχει μια βαθιά γεωμετρική ερμηνεία που συνδέεται άμεσα με τη θεωρία του ριζικού υπόχωρου και της ολοκληρωσιμότητας που αναπτύξαμε στα προηγούμενα κεφάλαια.

Όπως είδαμε, για μια απλή $k$-μορφή $\omega$, η συνθήκη $d\omega = 0$ είναι ισοδύναμη με την ολοκληρωσιμότητα του ριζικού υπόχωρου $\rho(\omega)$. Το λήμμα του Poincaré μας λέει ότι, τοπικά, αυτή η ολοκληρωσιμότητα μπορεί να "ολοκληρωθεί" με την έννοια ότι υπάρχει μια $(k-1)$-μορφή $\eta$ της οποίας η εξωτερική παράγωγος είναι η $\omega$. Για $k=1$, αυτό είναι το γνωστό γεγονός ότι μια συντηρητική δύναμη προέρχεται από δυναμικό: $d\omega = 0$ συνεπάγεται τοπικά $\omega = df$.

Για $k > 1$, η γεωμετρική εικόνα είναι πιο πλούσια. Η $(k-1)$-μορφή $\eta$ που κατασκευάζεται από τον τελεστή $H$ έχει την ιδιότητα ότι ο ριζικός της υπόχωρος $\rho(\eta)$ περιέχει τον $\rho(\omega)$, και η $\omega$ "μετρά" τον προσανατολισμένο $k$-διάστατο όγκο των προβολών στον $\pi(\omega)$. Η ύπαρξη της $\eta$ σημαίνει ότι η $\omega$ μπορεί να εκφραστεί ως η "στροβιλότητα" ή η "απόκλιση" (ανάλογα με το $k$) μιας μορφής χαμηλότερης τάξης.

Σε όρους του θεωρήματος Frobenius, το λήμμα του Poincaré είναι το αντίστοιχο της ύπαρξης ολοκληρωμένων υποπολλαπλοτήτων: η ολοκληρωσιμότητα (κλειστότητα) εξασφαλίζει την ύπαρξη ενός "δυναμικού" (την $\eta$) από το οποίο προέρχεται η $\omega$. Αυτή η αντιστοιχία είναι πλήρης για $k=1$ (δυνάμεις και δυναμικά) και γενικεύεται σε ανώτερες τάξεις μέσω της εξωτερικής παραγώγου.

\subsubsection{Φυσική ερμηνεία με όρους δυναμικού}

Η πιο άμεση φυσική ερμηνεία του λήμματος του Poincaré προκύπτει για $k=1$, όπου συνδέεται με την έννοια του δυναμικού στη φυσική.

\textbf{Συντηρητικές δυνάμεις και δυναμικό.} Μια δύναμη αναπαρίσταται από μια 1-μορφή $\omega$. Η συνθήκη $d\omega = 0$ σημαίνει ότι η δύναμη είναι συντηρητική: το έργο της κατά μήκος μιας κλειστής καμπύλης είναι μηδέν. Το λήμμα του Poincaré εγγυάται ότι, τοπικά, υπάρχει μια βαθμωτή συνάρτηση $f$ (το δυναμικό) τέτοια ώστε $\omega = df$. Σε Καρτεσιανές συντεταγμένες, αυτό γράφεται

\[
\omega_j = \frac{\partial f}{\partial x^j},
\]

που είναι ακριβώς η σχέση δύναμης-δυναμικού: η δύναμη είναι η κλίση του δυναμικού. Ο τελεστής $H$ που κατασκευάσαμε στην απόδειξη δίνει έναν ρητό τύπο για το δυναμικό:

\[
f(x) = \int_0^1 x^j \omega_j(tx) \, dt.
\]

Αυτό είναι το γνωστό επικαμπύλιο ολοκλήρωμα της δύναμης κατά μήκος της ευθύγραμμης διαδρομής από την αρχή ως το $x$.

\textbf{Ηλεκτροστατική.} Στην ηλεκτροστατική, το ηλεκτρικό πεδίο $E$ είναι μια 1-μορφή. Σε μια περιοχή χωρίς φορτία, ο νόμος του Gauss δίνει $dE = 0$ (η στροβιλότητα του ηλεκτρικού πεδίου είναι μηδέν). Το λήμμα του Poincaré εγγυάται ότι τοπικά υπάρχει ένα βαθμωτό δυναμικό $\phi$ τέτοιο ώστε $E = -d\phi$. Το αρνητικό πρόσημο είναι συμβατικό και δηλώνει ότι το ηλεκτρικό πεδίο δείχνει προς την κατεύθυνση ελάττωσης του δυναμικού. Σε περιοχές με φορτία, $dE = \rho \, dV \neq 0$, και το δυναμικό δεν υπάρχει καθολικά — αλλά τοπικά, μακριά από τα φορτία, εξακολουθεί να υπάρχει.

\textbf{Μαγνητοστατική και το διανυσματικό δυναμικό.} Για $k=2$, το λήμμα του Poincaré αποκτά μια πιο πλούσια φυσική ερμηνεία. Μια 2-μορφή $B$ αναπαριστά ένα μαγνητικό πεδίο. Ο νόμος του Gauss για τον μαγνητισμό λέει ότι $dB = 0$ (δεν υπάρχουν μαγνητικά μονόπολα). Το λήμμα του Poincaré εγγυάται ότι τοπικά υπάρχει μια 1-μορφή $A$ (το διανυσματικό δυναμικό) τέτοια ώστε $B = dA$. Σε Καρτεσιανές συντεταγμένες, αυτό γράφεται

\[
B_{ij} = \frac{\partial A_j}{\partial x^i} - \frac{\partial A_i}{\partial x^j},
\]

που είναι ακριβώς η σχέση μαγνητικού πεδίου και διανυσματικού δυναμικού. Ο τελεστής $H$ δίνει έναν ρητό τύπο για το διανυσματικό δυναμικό (στην "ακτινωτή γεωμετρία"). Το γεγονός ότι το $A$ δεν είναι μοναδικό (μπορούμε να προσθέσουμε την κλίση μιας βαθμωτής συνάρτησης, $A \to A + df$, χωρίς να αλλάξει το $B$) είναι μια άμεση συνέπεια του γεγονότος ότι $d^2 = 0$: $d(A + df) = dA + d^2f = dA = B$. Αυτή η αβεβαιότητα στο $A$ είναι η περίφημη ελευθερία βαθμίδας της ηλεκτροδυναμικής.

\subsubsection{Σύνδεση με τη γεωμετρική ερμηνεία των απλών μορφών}

Το λήμμα του Poincaré εντάσσεται φυσικά στη γεωμετρική θεωρία των απλών μορφών που αναπτύξαμε στο Κεφάλαιο 16. Για μια απλή $k$-μορφή $\omega$, η συνθήκη $d\omega = 0$ σημαίνει ότι ο ριζικός υπόχωρος $\rho(\omega)$ είναι ολοκληρώσιμος. Το λήμμα του Poincaré μας λέει ότι αυτή η ολοκληρωσιμότητα μπορεί να "πραγματοποιηθεί" μέσω μιας $(k-1)$-μορφής $\eta$: η $\omega$ είναι η εξωτερική παράγωγος της $\eta$.

Γεωμετρικά, αυτό σημαίνει ότι οι ολοκληρωμένες υποπολλαπλότητες του $\rho(\omega)$ (που υπάρχουν λόγω της $d\omega = 0$) μπορούν να περιγραφούν ως οι επιφάνειες στάθμης μιας οικογένειας συναρτήσεων — των "δυναμικών" που κωδικοποιούνται στην $\eta$. Για $k=1$, η $\eta = f$ είναι μια βαθμωτή συνάρτηση, και οι ολοκληρωμένες υποπολλαπλότητες του $\ker \omega$ είναι οι επιφάνειες $f = \text{σταθερά}$. Για $k=2$, η $\eta$ είναι μια 1-μορφή, και η $\omega = d\eta$ μετρά τη "στροβιλότητα" της $\eta$ — δηλαδή, την αποτυχία του $\ker \eta$ να είναι ολοκληρώσιμος. Το γεγονός ότι $d\omega = d^2\eta = 0$ είναι ταυτολογικό, και το λήμμα του Poincaré είναι το αντίστροφο: αν η "στροβιλότητα της στροβιλότητας" είναι μηδέν, τότε η στροβιλότητα προέρχεται από κάτι.

Με άλλα λόγια, το λήμμα του Poincaré ολοκληρώνει την εικόνα: η εξωτερική παράγωγος $d$ μετρά την αποτυχία ολοκληρωσιμότητας, και ο τελεστής ομοτοπίας $H$ παρέχει την "ολοκλήρωση" — την κατασκευή ενός δυναμικού από μια κλειστή μορφή. Αυτή η αλληλεπίδραση ανάμεσα στην τοπική και την καθολική δομή, ανάμεσα στην ολοκληρωσιμότητα και την ύπαρξη δυναμικού, είναι ένα από τα κεντρικά θέματα της διαφορικής γεωμετρίας και των φυσικών της εφαρμογών.

\subsection{Σχέση εξωτερικής παραγώγου και παραγώγου Lie}

Η εξωτερική παράγωγος $d$ και η παράγωγος Lie $\mathcal{L}_\xi$ είναι δύο θεμελιώδεις τελεστές στη διαφορική γεωμετρία, και η μεταξύ τους σχέση είναι βαθιά και αποκαλυπτική. Σε αυτή την ενότητα, αποδεικνύουμε τρεις κεντρικές ταυτότητες που τις συνδέουν.

\subsubsection{Ο μαγικός τύπος του Cartan}

Ο \textbf{μαγικός τύπος του Cartan} συνδέει την παράγωγο Lie, την εξωτερική παράγωγο και τη συστολή (εσωτερικό γινόμενο) $\iota_\xi$. Για μια $k$-μορφή $\omega$ και ένα διανυσματικό πεδίο $\xi$, ισχύει

\[
\mathcal{L}_\xi \omega = \iota_\xi (d\omega) + d(\iota_\xi \omega). \qquad (159)
\]

\textbf{Απόδειξη.} Η απόδειξη γίνεται με επαγωγή στον βαθμό $k$ της μορφής. Για $k = 0$, η $\omega = f$ είναι μια βαθμωτή συνάρτηση. Τότε $\iota_\xi f = 0$ (η συστολή μιας 0-μορφής είναι μηδέν εξ ορισμού), οπότε το δεξί μέλος γίνεται $0 + df(\xi) = \xi(f) = \mathcal{L}_\xi f$, που είναι ακριβώς το αριστερό μέλος.

Υποθέτουμε ότι ο τύπος ισχύει για $k$-μορφές και τον αποδεικνύουμε για $(k+1)$-μορφές. Αρκεί να τον αποδείξουμε για μια $(k+1)$-μορφή της μορφής $\omega = f \, dx^{i_1} \wedge \dots \wedge dx^{i_{k+1}}$, αφού κάθε $(k+1)$-μορφή είναι γραμμικός συνδυασμός τέτοιων όρων. Γράφουμε $\omega = \eta \wedge dx^j$, όπου $\eta = f \, dx^{i_1} \wedge \dots \wedge dx^{i_k}$ είναι μια $k$-μορφή. Τότε

\[
\begin{aligned}
\mathcal{L}_\xi (\eta \wedge dx^j) &= (\mathcal{L}_\xi \eta) \wedge dx^j + \eta \wedge (\mathcal{L}_\xi dx^j) \\
&= (\iota_\xi d\eta + d\iota_\xi \eta) \wedge dx^j + \eta \wedge d(\xi^j),
\end{aligned}
\]

χρησιμοποιώντας την επαγωγική υπόθεση για την $\eta$ και το γεγονός ότι $\mathcal{L}_\xi dx^j = d(\mathcal{L}_\xi x^j) = d(\xi^j)$. Από την άλλη,

\[
\begin{aligned}
\iota_\xi d(\eta \wedge dx^j) + d\iota_\xi (\eta \wedge dx^j) &= \iota_\xi (d\eta \wedge dx^j) + d(\iota_\xi \eta \wedge dx^j + (-1)^k \eta \wedge \iota_\xi dx^j) \\
&= (\iota_\xi d\eta) \wedge dx^j + (-1)^k d\eta \wedge \iota_\xi dx^j + d\iota_\xi \eta \wedge dx^j \\
&\quad + (-1)^{k-1} \iota_\xi \eta \wedge d(dx^j) + (-1)^k d\eta \wedge \xi^j + (-1)^k \eta \wedge d(\xi^j).
\end{aligned}
\]

Χρησιμοποιώντας το γεγονός ότι $d(dx^j) = 0$, $\iota_\xi dx^j = \xi^j$, και ότι $(-1)^k d\eta \wedge \xi^j + (-1)^k d\eta \wedge \xi^j$ αλληλοαναιρούνται (το ένα προέρχεται από τον πρώτο όρο και το άλλο από τον τρίτο), βλέπουμε ότι το δεξί μέλος ταυτίζεται με το παραπάνω αποτέλεσμα. Αυτό ολοκληρώνει την επαγωγή. $\square$

Ο μαγικός τύπος του Cartan έχει μια βαθιά γεωμετρική ερμηνεία: η παράγωγος Lie μετρά τη μεταβολή μιας μορφής κατά μήκος της ροής ενός διανυσματικού πεδίου. Το δεξί μέλος την εκφράζει ως το άθροισμα δύο όρων: ο πρώτος, $\iota_\xi (d\omega)$, μετρά τη "στροβιλότητα" της $\omega$ στην κατεύθυνση του $\xi$, και ο δεύτερος, $d(\iota_\xi \omega)$, μετρά τη "διαφορική" μεταβολή της συστολής της $\omega$ με το $\xi$.

\subsubsection{Κανόνας γινομένου για την παράγωγο Lie}

Η παράγωγος Lie ικανοποιεί έναν κανόνα γινομένου ως προς το σφηνοειδές γινόμενο, ανάλογο με αυτόν της εξωτερικής παραγώγου. Για μια $k$-μορφή $\omega$ και μια $l$-μορφή $\eta$, ισχύει

\[
\mathcal{L}_\xi (\omega \wedge \eta) = (\mathcal{L}_\xi \omega) \wedge \eta + \omega \wedge (\mathcal{L}_\xi \eta). \qquad (160)
\]

\textbf{Απόδειξη.} Χρησιμοποιούμε τον μαγικό τύπο του Cartan (159) και τον κανόνα Leibniz για την εξωτερική παράγωγο (154). Έχουμε

\[
\begin{aligned}
\mathcal{L}_\xi (\omega \wedge \eta) &= \iota_\xi d(\omega \wedge \eta) + d\iota_\xi (\omega \wedge \eta) \\
&= \iota_\xi (d\omega \wedge \eta + (-1)^k \omega \wedge d\eta) + d(\iota_\xi \omega \wedge \eta + (-1)^k \omega \wedge \iota_\xi \eta).
\end{aligned}
\]

Για τη συστολή του πρώτου όρου, χρησιμοποιούμε την ιδιότητα $\iota_\xi (\alpha \wedge \beta) = (\iota_\xi \alpha) \wedge \beta + (-1)^{\deg \alpha} \alpha \wedge (\iota_\xi \beta)$. Έτσι,

\[
\begin{aligned}
\mathcal{L}_\xi (\omega \wedge \eta) &= (\iota_\xi d\omega) \wedge \eta + (-1)^{k-1} d\omega \wedge (\iota_\xi \eta) + (-1)^k (\iota_\xi \omega) \wedge d\eta + (-1)^{2k} \omega \wedge (\iota_\xi d\eta) \\
&\quad + (d\iota_\xi \omega) \wedge \eta + (-1)^{k-1} (\iota_\xi \omega) \wedge d\eta + (-1)^k d\omega \wedge (\iota_\xi \eta) + \omega \wedge (d\iota_\xi \eta).
\end{aligned}
\]

Παρατηρούμε ότι οι όροι $(-1)^{k-1} d\omega \wedge (\iota_\xi \eta)$ και $(-1)^k d\omega \wedge (\iota_\xi \eta)$ αλληλοαναιρούνται (αφού $(-1)^{k-1} = -(-1)^k$). Το ίδιο συμβαίνει με τους όρους $(-1)^k (\iota_\xi \omega) \wedge d\eta$ και $(-1)^{k-1} (\iota_\xi \omega) \wedge d\eta$. Αυτό που απομένει είναι

\[
\mathcal{L}_\xi (\omega \wedge \eta) = (\iota_\xi d\omega + d\iota_\xi \omega) \wedge \eta + \omega \wedge (\iota_\xi d\eta + d\iota_\xi \eta) = (\mathcal{L}_\xi \omega) \wedge \eta + \omega \wedge (\mathcal{L}_\xi \eta),
\]

χρησιμοποιώντας πάλι τον μαγικό τύπο του Cartan. $\square$

\subsubsection{Αντιμετάθεση παραγώγων}

Η εξωτερική παράγωγος και η παράγωγος Lie αντιμετατίθενται. Για κάθε $k$-μορφή $\omega$ και κάθε διανυσματικό πεδίο $\xi$, ισχύει

\[
d(\mathcal{L}_\xi \omega) = \mathcal{L}_\xi (d\omega). \qquad (161)
\]

\textbf{Απόδειξη.} Χρησιμοποιούμε τον μαγικό τύπο του Cartan (159) και την ταυτότητα $d^2 = 0$ (153). Έχουμε

\[
d(\mathcal{L}_\xi \omega) = d(\iota_\xi d\omega + d\iota_\xi \omega) = d\iota_\xi d\omega + d^2 \iota_\xi \omega = d\iota_\xi d\omega,
\]

αφού $d^2 = 0$. Από την άλλη,

\[
\mathcal{L}_\xi (d\omega) = \iota_\xi d(d\omega) + d\iota_\xi d\omega = 0 + d\iota_\xi d\omega = d\iota_\xi d\omega,
\]

αφού $d^2 = 0$. Άρα $d(\mathcal{L}_\xi \omega) = \mathcal{L}_\xi (d\omega)$. $\square$

Αυτή η ταυτότητα είναι εξαιρετικά σημαντική: δηλώνει ότι η εξωτερική παράγωγος και η παράγωγος Lie είναι συμβατές — η σειρά με την οποία τις εφαρμόζουμε δεν έχει σημασία. Γεωμετρικά, αυτό σημαίνει ότι η ροή ενός διανυσματικού πεδίου σέβεται την εξωτερική παράγωγο: αν αφήσουμε μια μορφή να εξελιχθεί κατά μήκος της ροής και μετά πάρουμε την εξωτερική παράγωγο, είναι το ίδιο με το να πάρουμε πρώτα την εξωτερική παράγωγο και μετά να την αφήσουμε να εξελιχθεί. Αυτή η συμβατότητα είναι θεμελιώδης για τη διατύπωση των νόμων διατήρησης σε γεωμετρική γλώσσα: αν μια μορφή είναι κλειστή ($d\omega = 0$), τότε η παράγωγος Lie της είναι επίσης κλειστή, και αν μια μορφή είναι ακριβής ($\omega = d\eta$), τότε η παράγωγος Lie της είναι επίσης ακριβής ($\mathcal{L}_\xi \omega = d(\mathcal{L}_\xi \eta)$).

\subsection{Εξωτερική παράγωγος και τελεστής Hodge υπό την παρουσία Ρημάννειας μετρικής}

Σε αυτή την ενότητα, εξετάζουμε πώς η παρουσία μιας Ρημάννειας μετρικής $g$ επιτρέπει τη σύνδεση της εξωτερικής παραγώγου με τον τελεστή Hodge $\star$, οδηγώντας σε θεμελιώδεις διαφορικούς τελεστές όπως η απόκλιση και ο Laplace-Beltrami. Αυτή η σύνδεση είναι που δίνει φυσικό νόημα στην εξωτερική παράγωγο σε εφαρμογές της φυσικής.

\subsubsection{Η συναλλοίωτη παράγωγος σε όρους μορφών: η απόκλιση}

Σε έναν $n$-διάστατο Ρημάννειο χώρο $(M, g)$, η μετρική επιτρέπει τον ορισμό της \textbf{συναλλοίωτης απόκλισης} ενός διανυσματικού πεδίου. Στη γλώσσα των διαφορικών μορφών, αυτή η έννοια εκφράζεται με κομψό τρόπο μέσω της εξωτερικής παραγώγου και του τελεστή Hodge.

Θεωρούμε ένα διανυσματικό πεδίο $\xi$ στην $M$. Μέσω του μουσικού ισομορφισμού $\flat$, αντιστοιχούμε στο $\xi$ την 1-μορφή $\xi^\flat$. Η δυϊκή Hodge αυτής της 1-μορφής είναι η $(n-1)$-μορφή $\star \xi^\flat$. Η εξωτερική παράγωγος αυτής της $(n-1)$-μορφής είναι μια $n$-μορφή, δηλαδή μια πυκνότητα. Εφαρμόζοντας ξανά τον τελεστή Hodge, παίρνουμε μια 0-μορφή (βαθμωτή συνάρτηση), η οποία ορίζεται ως η \textbf{απόκλιση} του $\xi$:

\[
\text{div} \, \xi = \star \, d \star \xi^\flat. \qquad (162)
\]

Σε τοπικές συντεταγμένες, αυτή η έκφραση παίρνει τη γνωστή μορφή

\[
\text{div} \, \xi = \frac{1}{\sqrt{|\det(g_{ij})|}} \frac{\partial}{\partial x^i} \left( \sqrt{|\det(g_{ij})|} \ \xi^i \right). \qquad (163)
\]

Για να το δούμε αυτό, θυμόμαστε ότι το κανονικό στοιχείο όγκου είναι $dV = \sqrt{|\det(g_{ij})|} \, dx^1 \wedge \dots \wedge dx^n$. Η $(n-1)$-μορφή $\star \xi^\flat$ έχει συνιστώσες

\[
(\star \xi^\flat)_{i_1 \dots i_{n-1}} = \varepsilon_{i_1 \dots i_{n-1} j} \, \xi^j,
\]

και η εξωτερική της παράγωγος είναι

\[
d(\star \xi^\flat) = \frac{\partial}{\partial x^i} \left( \sqrt{|\det(g_{ij})|} \ \xi^i \right) dx^1 \wedge \dots \wedge dx^n.
\]

Εφαρμόζοντας τον Hodge, η $n$-μορφή γίνεται βαθμωτή συνάρτηση, και ο παράγοντας $\sqrt{|\det(g_{ij})|}$ δίνει ακριβώς τον τύπο (163).

Η γεωμετρική σημασία της απόκλισης είναι ότι μετρά τη "διαστολή" ενός στοιχείου όγκου που μεταφέρεται από τη ροή του $\xi$. Αν $\text{div} \, \xi = 0$, η ροή είναι ασυμπίεστη: ο όγκος διατηρείται. Αν $\text{div} \, \xi > 0$, η ροή διαστέλλεται· αν $\text{div} \, \xi < 0$, η ροή συστέλλεται.

Η σύνδεση με τον δυϊσμό ροής-έργου είναι άμεση. Από το Κεφάλαιο 17, μια $(n-1)$-μορφή $\eta$ αντιπροσωπεύει μια ροή, και η δυϊκή της Hodge $\star \eta$ είναι μια 1-μορφή που αντιπροσωπεύει μια δύναμη. Η εξωτερική παράγωγος $d\eta$ είναι μια $n$-μορφή που μετρά την απόκλιση της ροής. Στην περίπτωση που $\eta = \star \xi^\flat$, έχουμε $d\eta = \text{div} \, \xi \, dV$, και η απόκλιση είναι ακριβώς η "πυκνότητα πηγής" της ροής.

\subsubsection{Ο τελεστής Laplace-Beltrami}

Ο \textbf{τελεστής Laplace-Beltrami} (ή απλώς Λαπλασιανή) είναι μια άμεση γενίκευση του συνηθισμένου τελεστή Laplace σε καμπύλους χώρους. Στη γλώσσα των διαφορικών μορφών, ορίζεται για μια βαθμωτή συνάρτηση $f$ (0-μορφή) ως

\[
\Delta f = \text{div} \, (\nabla f) = \star \, d \star df. \qquad (164)
\]

Σε τοπικές συντεταγμένες, χρησιμοποιώντας την (163) και το γεγονός ότι $(\nabla f)^i = g^{ij} \partial f / \partial x^j$, παίρνουμε

\[
\Delta f = \frac{1}{\sqrt{|\det(g_{ij})|}} \frac{\partial}{\partial x^i} \left( \sqrt{|\det(g_{ij})|} \ g^{ij} \frac{\partial f}{\partial x^j} \right). \qquad (165)
\]

Για μια γενική $k$-μορφή $\omega$, ο Laplace-Beltrami ορίζεται μέσω της εξωτερικής παραγώγου και του Hodge ως

\[
\Delta \omega = (d \star d \star + \star d \star d) \, \omega = (d + \star d \star)^2 \, \omega. \qquad (166)
\]

Αυτός ο ορισμός γενικεύει την κλασική Λαπλασιανή και σέβεται τη γεωμετρική δομή των μορφών. Ο τελεστής $\Delta$ είναι αυτοσυζυγής ως προς το εσωτερικό γινόμενο που επάγει η μετρική στις μορφές, και ικανοποιεί $\Delta = (d + d^*)^2$, όπου $d^* = \pm \star d \star$ είναι ο τυπικός συζυγής της $d$ (το πρόσημο εξαρτάται από τη διάσταση και τον βαθμό της μορφής).

\subsubsection{Γεωμετρική ερμηνεία και συσχέτιση με τον δυϊσμό ροής-έργου}

Η παρουσία της μετρικής επιτρέπει μια πλούσια γεωμετρική ερμηνεία των τελεστών που εμπλέκονται, συνδέοντας την εξωτερική παράγωγο με τον δυϊσμό ροής-έργου που αναπτύξαμε στο Κεφάλαιο 17.

Θεωρούμε μια 1-μορφή $\omega$, η οποία αντιπροσωπεύει μια δύναμη. Μέσω της μετρικής, αντιστοιχεί στο διάνυσμα $(\omega)^\sharp$. Η εξωτερική παράγωγος $d\omega$ είναι μια 2-μορφή που μετρά τη στροβιλότητα της δύναμης. Η δυϊκή Hodge $\star \omega$ είναι μια $(n-1)$-μορφή που αντιπροσωπεύει τη ροή που είναι δυϊκή στη δύναμη. Η εξωτερική παράγωγος $d \star \omega$ είναι μια $n$-μορφή που μετρά την απόκλιση αυτής της ροής — δηλαδή, την "πυκνότητα πηγής" της δύναμης. Τέλος, η $\star d \star \omega$ είναι μια 0-μορφή (βαθμωτή συνάρτηση) που δίνει την απόκλιση του διανύσματος $(\omega)^\sharp$.

Με άλλα λόγια, η σύνθεση $\star d \star$ μεταφράζει την εξωτερική παράγωγο μιας $(n-1)$-μορφής (ροής) σε μια βαθμωτή συνάρτηση (απόκλιση). Αυτή είναι ακριβώς η γεωμετρική έκφραση του γεγονότος ότι η απόκλιση μιας ροής μετρά την "πηγή" της. Αντίστροφα, η σύνθεση $d \star$ μεταφράζει μια 1-μορφή (δύναμη) σε μια $(n-1)$-μορφή (ροή), και η $d \star d$ δίνει την απόκλιση της κλίσης.

Ο τελεστής Laplace-Beltrami για 0-μορφές, $\Delta f = \star d \star df$, είναι λοιπόν η απόκλιση της κλίσης. Γεωμετρικά, μετρά τη "διαφορά" ανάμεσα στην τιμή της $f$ σε ένα σημείο και τη μέση τιμή της σε μια μικρή σφαίρα γύρω από αυτό. Στη φυσική, η εξίσωση $\Delta f = 0$ (εξίσωση Laplace) περιγράφει ένα βαθμωτό πεδίο χωρίς πηγές, ενώ η εξίσωση $\Delta f = \rho$ (εξίσωση Poisson) περιγράφει ένα πεδίο με πυκνότητα πηγής $\rho$.

\subsubsection{Εφαρμογές στη φυσική}

\textbf{Ηλεκτροστατική και βαρύτητα.} Στην ηλεκτροστατική, το ηλεκτρικό δυναμικό $\phi$ ικανοποιεί την εξίσωση Poisson

\[
\Delta \phi = -\rho,
\]

όπου $\rho$ είναι η πυκνότητα φορτίου. Σε όρους μορφών, το ηλεκτρικό πεδίο είναι η 1-μορφή $E = -d\phi$, και η εξωτερική παράγωγος της δυϊκής της Hodge είναι

\[
d \star E = \rho \, dV.
\]

Αυτή είναι ακριβώς η διαφορική μορφή του νόμου του Gauss: η ροή του ηλεκτρικού πεδίου μέσα από μια κλειστή επιφάνεια ισούται με το ολικό φορτίο που περικλείεται. Στη βαρύτητα, το Νευτώνειο δυναμικό $\Phi$ ικανοποιεί $\Delta \Phi = 4\pi G \rho$, όπου $\rho$ είναι η πυκνότητα μάζας, και η ίδια γεωμετρική δομή εφαρμόζεται.

\textbf{Διάχυση και θερμότητα.} Η εξίσωση της θερμότητας (ή διάχυσης) σε έναν καμπύλο χώρο γράφεται ως

\[
\frac{\partial u}{\partial t} = \Delta u,
\]

όπου $u$ είναι η θερμοκρασία (ή η συγκέντρωση ενός διαχεόμενου υλικού). Σε όρους μορφών, αυτή είναι η εξέλιξη μιας 0-μορφής κάτω από τη ροή της βαθμίδας του Laplace-Beltrami. Η γεωμετρική δομή του $\Delta$ εξασφαλίζει ότι η διάχυση σέβεται τη μετρική του χώρου: σε περιοχές όπου η μετρική "στενεύει", η διάχυση επιβραδύνεται· σε περιοχές όπου "πλαταίνει", επιταχύνεται.

\textbf{Κυματική εξίσωση.} Η κυματική εξίσωση για ένα βαθμωτό πεδίο $\phi$ σε καμπύλο χωροχρόνο γράφεται ως

\[
\square \phi = \Delta \phi - \frac{\partial^2 \phi}{\partial t^2} = 0,
\]

όπου $\square$ είναι ο d'Alembertian (η κυματική τελεστής). Σε όρους μορφών, αυτό γενικεύεται στην εξίσωση

\[
(d \star d \star + \star d \star d) \phi = 0,
\]

η οποία είναι ανεξάρτητη από την επιλογή συντεταγμένων και σέβεται την αιτιακή δομή του χωροχρόνου.

\subsection{Σύνοψη}

Η εξωτερική παράγωγος $d$ είναι ένας τελεστής που απεικονίζει $k$-μορφές σε $(k+1)$-μορφές, οριζόμενος μέσω της δράσης του σε διανυσματικά πεδία από τον τύπο (150). Το γινόμενο Lie, που εμφανίζεται ρητά στον ορισμό, είναι το κλειδί που κωδικοποιεί την πληροφορία για την ολοκληρωσιμότητα του ριζικού υπόχωρου. Σε τοπικές συντεταγμένες, οι συνιστώσες της $d\omega$ δίνονται από τον τύπο (152). Οι θεμελιώδεις ταυτότητες $d^2 = 0$ (153), ο κανόνας Leibniz (154) και ο τύπος για την εξωτερική παράγωγο απλής μορφής (155) την καθιστούν τον ακρογωνιαίο λίθο του εξωτερικού λογισμού. Για απλές $k$-μορφές, η $d\omega$ μετρά την αποτυχία του ριζικού υπόχωρου $\rho(\omega)$ να είναι ολοκληρώσιμος, και αυτή η αποτυχία εκφράζεται μέσω του γινομένου Lie των διανυσματικών πεδίων που ανήκουν στον $\rho(\omega)$. Το θεώρημα Frobenius στη διατύπωση των διαφορικών μορφών (ενότητα 20.5) καθιστά αυτή τη σύνδεση αυστηρή, δείχνοντας ότι η συνθήκη $d\omega^a = \sum_b \theta^a_b \wedge \omega^b$ είναι ισοδύναμη με την $[\xi, \eta] \in \rho$ μέσω της ταυτότητας-κλειδιού (157). Το λήμμα του Poincaré (ενότητα 20.6) ολοκληρώνει την εικόνα: σε μια αστεροειδή περιοχή, κάθε κλειστή μορφή είναι ακριβής, και ο τελεστής ομοτοπίας $H$ παρέχει μια ρητή κατασκευή του "δυναμικού". Στις ειδικές περιπτώσεις $k=1$ και $k=n-1$, η $d\omega$ ερμηνεύεται ως στροβιλότητα και απόκλιση αντίστοιχα, ενώ το λήμμα του Poincaré για $k=1$ δίνει το βαθμωτό δυναμικό μιας συντηρητικής δύναμης και για $k=2$ το διανυσματικό δυναμικό ενός μαγνητικού πεδίου. Η σχέση της εξωτερικής παραγώγου με την παράγωγο Lie εκφράζεται μέσω του μαγικού τύπου του Cartan (159), του κανόνα γινομένου (160) και της αντιμετάθεσης παραγώγων (161). Τέλος, υπό την παρουσία Ρημάννειας μετρικής, η εξωτερική παράγωγος και ο τελεστής Hodge συνδυάζονται για να δώσουν την απόκλιση (162)-(163) και τον τελεστή Laplace-Beltrami (164)-(166), οι οποίοι βρίσκουν θεμελιώδεις εφαρμογές στην ηλεκτροστατική, τη βαρύτητα, τη διάχυση και την κυματική εξίσωση. Η θεωρία αυτή αποτελεί το θεμέλιο για το θεώρημα του Stokes, τη θεωρία των χαρακτηριστικών επιφανειών και τη γεωμετρική διατύπωση των νόμων διατήρησης.

\section{Ολοκλήρωση Διαφορικών Μορφών}

\subsection{Εισαγωγή}

Στα προηγούμενα κεφάλαια, αναπτύξαμε την απειροστική θεωρία των διαφορικών μορφών: είδαμε πώς μια $k$-μορφή δρα σε $k$ διανύσματα μετρώντας προσανατολισμένους όγκους, και πώς η εξωτερική παράγωγος $d$ κωδικοποιεί πληροφορία για την ολοκληρωσιμότητα και τις πηγές. Η ολοκλήρωση είναι η γέφυρα που συνδέει αυτή την απειροστική εικόνα με καθολικές ποσότητες: μετατρέπει μια $k$-μορφή $\omega$ σε έναν αριθμό, το ολοκλήρωμά της πάνω σε μια προσανατολισμένη $k$-διάστατη περιοχή.

Η θεμελιώδης παρατήρηση είναι ότι ο νόμος μετασχηματισμού των $k$-μορφών υπό αλλαγή συντεταγμένων — ο οποίος εμπλέκει την ορίζουσα του Ιακωβιανού πίνακα — είναι ακριβώς ο ίδιος με τον νόμο αλλαγής μεταβλητών στο ολοκλήρωμα Riemann. Αυτή η σύμπτωση επιτρέπει να ορίσουμε το ολοκλήρωμα μιας $k$-μορφής με τρόπο ανεξάρτητο από την επιλογή παραμετρικοποίησης. Για $k=1$, παίρνουμε το επικαμπύλιο ολοκλήρωμα· για $k=n$, το ολοκλήρωμα όγκου. Η γενική περίπτωση ενοποιεί αυτές τις έννοιες.

\subsection{Τοπικός ορισμός και προσέγγιση Riemann}

Θεωρούμε μια $k$-μορφή $\omega$ ορισμένη σε μια πολλαπλότητα $M$, και μια $k$-διάστατη προσανατολισμένη υποπολλαπλότητα $\Sigma \subset M$. Υποθέτουμε αρχικά ότι η $\Sigma$ καλύπτεται από έναν μόνο χάρτη: υπάρχει μια αμφιδιαφορική απεικόνιση
\[
x: U \subseteq \mathbb{R}^k \to \Sigma,
\]
όπου $U$ είναι ένα ανοιχτό υποσύνολο του $\mathbb{R}^k$ με συντεταγμένες $u^1, \dots, u^k$. Η απεικόνιση $x$ αναπαρίσταται από τις συναρτήσεις $x^i(u^1, \dots, u^k)$, $i = 1, \dots, n$.

Αντί να γράψουμε την $\omega$ με τον παραγοντικό παράγοντα $1/k!$, τη γράφουμε ως άθροισμα πάνω σε διατεταγμένους δείκτες:
\[
\omega = \sum_{i_1 < \dots < i_k} \omega_{i_1 \dots i_k}(x) \ dx^{i_1} \wedge \dots \wedge dx^{i_k},
\]
ή, χρησιμοποιώντας τη σύμβαση άθροισης χωρίς περιορισμό στη διάταξη αλλά χωρίς τον παράγοντα $1/k!$, όπου οι συνιστώσες $\omega_{i_1 \dots i_k}$ είναι αντισυμμετρικές. Σε αυτή την περίπτωση, η $\omega$ δρα σε $k$ διανύσματα δίνοντας την ορίζουσα των τιμών των 1-μορφών βάσης πάνω στα διανύσματα, ακριβώς όπως στην (111).

Για να ολοκληρώσουμε την $\omega$ πάνω στη $\Sigma$, αντικαθιστούμε την παραμετρικοποίηση. Κάθε διαφορικό $dx^i$ γίνεται
\[
dx^i = \frac{\partial x^i}{\partial u^a} \ du^a.
\]
Αντικαθιστώντας στην $\omega$, παίρνουμε
\[
\omega = \sum_{i_1 < \dots < i_k} \omega_{i_1 \dots i_k}(x(u)) \ \frac{\partial x^{i_1}}{\partial u^{a_1}} \dots \frac{\partial x^{i_k}}{\partial u^{a_k}} \ du^{a_1} \wedge \dots \wedge du^{a_k},
\]
όπου η άθροιση πάνω στους $a_1, \dots, a_k$ είναι ελεύθερη (καθένας από $1$ ως $k$). Λόγω της αντισυμμετρίας του σφηνοειδούς γινομένου, το $du^{a_1} \wedge \dots \wedge du^{a_k}$ είναι μη μηδενικό μόνο όταν οι δείκτες $a_1, \dots, a_k$ είναι μια μετάθεση του $1, \dots, k$, και τότε ισούται με $\operatorname{sgn}(\sigma) \, du^1 \wedge \dots \wedge du^k$. Το άθροισμα πάνω σε όλους τους $a_1, \dots, a_k$ έχει $k!$ μη μηδενικούς όρους, καθένας από τους οποίους δίνει την ίδια απόλυτη συνεισφορά αλλά με πρόσημο που καθορίζεται από τη μετάθεση. Αυτό το άθροισμα είναι ακριβώς $k!$ φορές η ορίζουσα του Ιακωβιανού πίνακα. Το αποτέλεσμα είναι
\[
\omega = f(u) \ du^1 \wedge \dots \wedge du^k,
\]
όπου η συνάρτηση $f$ δίνεται από δύο ισοδύναμες εκφράσεις. Σε όρους μεταθέσεων,
\begin{equation}
f(u) = \sum_{i_1 < \dots < i_k} \omega_{i_1 \dots i_k}(x(u)) \sum_{\sigma \in S_k} \operatorname{sgn}(\sigma) \ \frac{\partial x^{i_1}}{\partial u^{\sigma(1)}} \dots \frac{\partial x^{i_k}}{\partial u^{\sigma(k)}}. \qquad (167a)
\end{equation}
Σε όρους του τανυστή Levi-Civita,
\begin{equation}
f(u) = \frac{1}{k!} \ \omega_{i_1 \dots i_k}(x(u)) \ \epsilon^{a_1 \dots a_k} \ \frac{\partial x^{i_1}}{\partial u^{a_1}} \dots \frac{\partial x^{i_k}}{\partial u^{a_k}}, \qquad (167b)
\end{equation}
όπου ο παράγοντας $1/k!$ εμφανίζεται τώρα φυσιολογικά στην έκφραση με τον τανυστή Levi-Civita, ακριβώς επειδή η ορίζουσα μέσω του $\epsilon^{a_1 \dots a_k}$ αθροίζει πάνω σε όλους τους $a_1, \dots, a_k$ (καθένας από $1$ ως $k$) και χρειάζεται κανονικοποίηση. Στην έκφραση με τις μεταθέσεις, η κανονικοποίηση είναι ενσωματωμένη στο γεγονός ότι αθροίζουμε πάνω σε $k!$ μεταθέσεις, που είναι ακριβώς το σωστό πλήθος.

Το \textbf{ολοκλήρωμα} της $\omega$ πάνω στη $\Sigma$ ορίζεται τότε ως το ολοκλήρωμα Riemann της $f$ πάνω στο $U$:
\begin{equation}
\int_\Sigma \omega = \int_U f(u) \, du^1 \dots du^k. \qquad (168)
\end{equation}

\subsection{Ανεξαρτησία από την παραμετρικοποίηση}

Για να είναι καλά ορισμένο το ολοκλήρωμα, πρέπει να δείξουμε ότι η τιμή του δεν εξαρτάται από την επιλογή της παραμετρικοποίησης $x$, αρκεί να διατηρείται ο προσανατολισμός. Θεωρούμε δύο παραμετρικοποιήσεις $x: U \to \Sigma$ και $y: V \to \Sigma$, με συντεταγμένες $u^a$ και $v^a$, και υποθέτουμε ότι η αλλαγή συντεταγμένων $v = \phi(u)$ έχει θετική Ιακωβιανή ορίζουσα.

Η σχέση $x(u) = y(\phi(u))$ και ο κανόνας της αλυσίδας δίνουν
\[
\frac{\partial x^i}{\partial u^a} = \frac{\partial y^i}{\partial v^b} \frac{\partial \phi^b}{\partial u^a}.
\]
Αντικαθιστώντας στην (167a), η $f(u)$ γράφεται
\[
f(u) = \sum_{i_1 < \dots < i_k} \omega_{i_1 \dots i_k}(y(\phi(u))) \sum_{\sigma \in S_k} \operatorname{sgn}(\sigma) \ \frac{\partial x^{i_1}}{\partial u^{\sigma(1)}} \dots \frac{\partial x^{i_k}}{\partial u^{\sigma(k)}}.
\]
Εισάγοντας τις παραγώγους της $y$ και της $\phi$ μέσω του κανόνα της αλυσίδας, κάθε όρος $\partial x^{i_a} / \partial u^{\sigma(a)}$ γίνεται
\[
\frac{\partial x^{i_a}}{\partial u^{\sigma(a)}} = \frac{\partial y^{i_a}}{\partial v^{b_a}} \frac{\partial \phi^{b_a}}{\partial u^{\sigma(a)}}.
\]
Το άθροισμα πάνω στις μεταθέσεις $\sigma$ παραγοντοποιείται: το κομμάτι που περιέχει τις παραγώγους των $y^{i_a}$ ως προς $v^{b_a}$ και το κομμάτι που περιέχει τις παραγώγους των $\phi^{b_a}$ ως προς $u^{\sigma(a)}$. Από τον ορισμό της ορίζουσας μέσω μεταθέσεων, το άθροισμα πάνω στις $\sigma$ του γινομένου των $\partial \phi^{b_a} / \partial u^{\sigma(a)}$ με το πρόσημο $\operatorname{sgn}(\sigma)$ είναι ακριβώς η ορίζουσα του Ιακωβιανού πίνακα της $\phi$. Το υπόλοιπο κομμάτι αναπαράγει την $g(\phi(u))$. Έτσι,
\[
f(u) = g(\phi(u)) \ \det \left( \frac{\partial \phi^a}{\partial u^b} \right),
\]
όπου $g(v)$ είναι η αντίστοιχη συνάρτηση για την παραμετρικοποίηση $y$. Το ολοκλήρωμα Riemann ικανοποιεί τον ίδιο νόμο μετασχηματισμού:
\[
\int_U f(u) \, du^1 \dots du^k = \int_V g(v) \, dv^1 \dots dv^k.
\]
Συνεπώς, το ολοκλήρωμα είναι ανεξάρτητο από την επιλογή παραμετρικοποίησης που σέβεται τον προσανατολισμό.

\subsection{Ύπαρξη του ολοκληρώματος}

Η ύπαρξη του ολοκληρώματος (168) ανάγεται στην ύπαρξη του ολοκληρώματος Riemann της $f$ στο $U$. Αυτό προϋποθέτει ότι η $f$ είναι Riemann-ολοκληρώσιμη — συνθήκη που ικανοποιείται όταν η $\omega$ είναι ομαλή και το $U$ είναι φραγμένο με σύνορο μηδενικού μέτρου. Για συμπαγή $\Sigma$, αυτό ισχύει πάντα.

\subsection{Ολική επέκταση σε υποπολλαπλότητα}

Όταν η $\Sigma$ δεν μπορεί να καλυφθεί από έναν μόνο χάρτη, χρησιμοποιούμε την εξής κατασκευή. Καλύπτουμε τη $\Sigma$ με μια πεπερασμένη συλλογή από χάρτες (παραμετρικοποιήσεις), που αλληλοεπικαλύπτονται. Σε κάθε χάρτη, η $\omega$ ανάγεται σε μια συνάρτηση $f_\alpha$ όπως στην (167a)-(167b), και το τοπικό ολοκλήρωμα ορίζεται από την (168).

Για να συνδυάσουμε τα τοπικά ολοκληρώματα χωρίς να μετρήσουμε πολλαπλά τις περιοχές επικάλυψης, εισάγουμε μια οικογένεια ομαλών συναρτήσεων $\rho_\alpha$, μία για κάθε χάρτη, με τις εξής ιδιότητες: κάθε $\rho_\alpha$ είναι μη αρνητική, μηδενίζεται έξω από τον αντίστοιχο χάρτη, και σε κάθε σημείο της $\Sigma$, το άθροισμα όλων των $\rho_\alpha$ ισούται με $1$. Αυτές οι συναρτήσεις ``μοιράζουν'' κάθε σημείο στους χάρτες που το καλύπτουν.

Το ολικό ολοκλήρωμα ορίζεται τότε ως
\begin{equation}
\int_\Sigma \omega = \sum_\alpha \int_{U_\alpha} \rho_\alpha(x_\alpha(u)) \ f_\alpha(u) \, du^1 \dots du^k. \qquad (169)
\end{equation}
Η ανεξαρτησία από την παραμετρικοποίηση που αποδείξαμε στην ενότητα 23.3, μαζί με την ιδιότητα $\sum_\alpha \rho_\alpha = 1$, εξασφαλίζει ότι το αποτέλεσμα δεν εξαρτάται ούτε από την επιλογή των χαρτών ούτε από την επιλογή των συναρτήσεων $\rho_\alpha$. Το αντικείμενο που προκύπτει εξαρτάται μόνο από τη γεωμετρία της $\Sigma$ και την $\omega$.

Οι συνθήκες υπό τις οποίες υπάρχει το ολοκλήρωμα είναι η ομαλότητα της $\omega$ και της $\Sigma$, η συμπάγεια της $\Sigma$ (ή, αν η $\Sigma$ δεν είναι συμπαγής, η επαρκής φθίση της $\omega$ στο άπειρο ώστε το άθροισμα στην (169) να συγκλίνει), και η προσανατολισιμότητα της $\Sigma$ με όλους τους χάρτες συμβατούς με τον επιλεγμένο προσανατολισμό. Για το θεώρημα του Stokes, απαιτείται επιπλέον η $\Sigma$ να είναι συμπαγής πολλαπλότητα με ομαλό σύνορο διάστασης $k-1$.

\subsection{Γεωμετρική ερμηνεία}

Το ολοκλήρωμα μιας $k$-μορφής $\omega$ πάνω σε μια $k$-διάστατη επιφάνεια $\Sigma$ έχει άμεση γεωμετρική ερμηνεία σε σχέση με τη θεωρία των απλών μορφών του Κεφαλαίου 16. Αν η $\omega$ είναι μια απλή $k$-μορφή, σε κάθε σημείο μετρά τον προσανατολισμένο $k$-διάστατο όγκο ενός $k$-παραλληλεπιπέδου ως προς τις μονάδες που ορίζουν οι 1-μορφές που την απαρτίζουν. Το ολοκλήρωμα $\int_\Sigma \omega$ είναι η συνολική ποσότητα αυτού του προσανατολισμένου όγκου πάνω σε ολόκληρη την $\Sigma$. Για $k=1$, αυτό είναι το επικαμπύλιο ολοκλήρωμα — το συνολικό έργο μιας δύναμης κατά μήκος μιας καμπύλης. Για $k=n$, είναι το ολοκλήρωμα όγκου — ο συνολικός όγκος, ή η συνολική μάζα, ή το συνολικό φορτίο μιας περιοχής.

Το θεώρημα του Stokes, που ακολουθεί στο επόμενο κεφάλαιο, συνδέει το ολοκλήρωμα της $d\omega$ πάνω σε μια περιοχή με το ολοκλήρωμα της $\omega$ πάνω στο σύνορό της. Αυτή η σύνδεση είναι το κεντρικό αποτέλεσμα του εξωτερικού λογισμού και το θεμέλιο για τη γεωμετρική διατύπωση των νόμων διατήρησης στη φυσική.

\section{Θεώρημα Stokes}

\subsection{Εισαγωγή}

Στα δύο προηγούμενα κεφάλαια, αναπτύξαμε ανεξάρτητα την εξωτερική παράγωγο $d$ και την ολοκλήρωση διαφορικών μορφών. Το θεώρημα του Stokes είναι η γέφυρα που συνδέει αυτές τις δύο πράξεις: δηλώνει ότι το ολοκλήρωμα της εξωτερικής παραγώγου μιας μορφής πάνω σε μια προσανατολισμένη περιοχή ισούται με το ολοκλήρωμα της ίδιας της μορφής πάνω στο σύνορο της περιοχής. Είναι η γενίκευση του θεμελιώδους θεωρήματος του απειροστικού λογισμού και ενοποιεί τα κλασικά θεωρήματα της διανυσματικής ανάλυσης.

\subsection{Διατύπωση}

\textbf{Θεώρημα (Stokes).} Έστω $\Sigma$ μια συμπαγής, προσανατολισμένη $k$-διάστατη ομαλή πολλαπλότητα με σύνορο $\partial \Sigma$ (το οποίο είναι μια $(k-1)$-διάστατη ομαλή πολλαπλότητα, προσανατολισμένη με τον επαγόμενο προσανατολισμό). Έστω $\omega$ μια ομαλή $(k-1)$-μορφή ορισμένη σε μια περιοχή της $\Sigma$. Τότε
\begin{equation}
\int_\Sigma d\omega = \int_{\partial \Sigma} \omega. \qquad (170)
\end{equation}
Το αριστερό μέλος είναι το ολοκλήρωμα της $k$-μορφής $d\omega$ πάνω στη $\Sigma$. Το δεξί μέλος είναι το ολοκλήρωμα της $(k-1)$-μορφής $\omega$ πάνω στο σύνορο $\partial \Sigma$. Ο επαγόμενος προσανατολισμός του $\partial \Sigma$ είναι εκείνος για τον οποίο ένα θετικά προσανατολισμένο εφαπτόμενο πλαίσιο του $\partial \Sigma$, ακολουθούμενο από ένα εξωτερικά κατευθυνόμενο κάθετο διάνυσμα, δίνει ένα θετικά προσανατολισμένο πλαίσιο της $\Sigma$.

\subsection{Απόδειξη}

Η απόδειξη δομείται σε τρία βήματα: αναγωγή σε τοπική κατάσταση μέσω διαμέρισης της μονάδας, απόδειξη για την περίπτωση ορθογωνίου στον $\mathbb{R}^k$, και μεταφορά στη γενική περίπτωση μέσω παραμετρικοποίησης.

\textbf{Βήμα 1: Αναγωγή σε τοπική κατάσταση.} Καλύπτουμε τη $\Sigma$ με έναν πεπερασμένο αριθμό από χάρτες $(U_\alpha, x_\alpha)$, όπου κάθε $x_\alpha: U_\alpha \subseteq \mathbb{R}^k \to \Sigma$ είναι μια παραμετρικοποίηση. Εισάγουμε μια διαμέριση της μονάδας $\{\rho_\alpha\}$ υποταγμένη σε αυτή την κάλυψη, με $\sum_\alpha \rho_\alpha = 1$ και κάθε $\rho_\alpha$ ομαλή με φορέα εντός του αντίστοιχου χάρτη. Τότε
\[
\omega = \sum_\alpha \rho_\alpha \omega, \qquad d\omega = \sum_\alpha d(\rho_\alpha \omega),
\]
αφού η $d$ είναι γραμμική. Το ολοκλήρωμα γίνεται
\[
\int_\Sigma d\omega = \sum_\alpha \int_\Sigma d(\rho_\alpha \omega), \qquad \int_{\partial \Sigma} \omega = \sum_\alpha \int_{\partial \Sigma} \rho_\alpha \omega.
\]
Αρκεί λοιπόν να αποδείξουμε την (170) για κάθε όρο $\rho_\alpha \omega$ χωριστά. Κάθε τέτοιος όρος έχει φορέα μέσα σε έναν μόνο χάρτη. Μέσω της παραμετρικοποίησης $x_\alpha$, το ολοκλήρωμα ανάγεται σε ένα ολοκλήρωμα πάνω στο $U_\alpha \subseteq \mathbb{R}^k$. Συνεπώς, αρκεί να αποδείξουμε το θεώρημα για μια $(k-1)$-μορφή $\omega$ ορισμένη σε ένα ανοιχτό υποσύνολο του $\mathbb{R}^k$, με φορέα σε ένα συμπαγές υποσύνολο.

\textbf{Βήμα 2: Απόδειξη στον $\mathbb{R}^k$.} Θεωρούμε μια $(k-1)$-μορφή $\omega$ στον $\mathbb{R}^k$ με συμπαγή φορέα. Λόγω γραμμικότητας, αρκεί να αποδείξουμε το θεώρημα για μια μορφή της μορφής
\[
\omega = f(x) \ dx^1 \wedge \dots \wedge \widehat{dx^j} \wedge \dots \wedge dx^k,
\]
όπου η περισπωμένη δηλώνει ότι το $dx^j$ παραλείπεται, και $f$ είναι ομαλή συνάρτηση με συμπαγή φορέα. Πράγματι, κάθε $(k-1)$-μορφή είναι πεπερασμένο άθροισμα τέτοιων όρων.

Υπολογίζουμε την εξωτερική παράγωγο της $\omega$. Από τον ορισμό της $d$, έχουμε
\[
d\omega = \frac{\partial f}{\partial x^j} \ dx^j \wedge dx^1 \wedge \dots \wedge \widehat{dx^j} \wedge \dots \wedge dx^k.
\]
Μετακινώντας το $dx^j$ στην πρώτη θέση, κάθε εναλλαγή με ένα από τα $dx^1, \dots, dx^{j-1}$ επιφέρει ένα πρόσημο $-1$. Μετά από $j-1$ εναλλαγές, παίρνουμε
\[
d\omega = (-1)^{j-1} \ \frac{\partial f}{\partial x^j} \ dx^1 \wedge \dots \wedge dx^j \wedge \dots \wedge dx^k.
\]

Θεωρούμε τώρα τη $\Sigma$ ως ένα ορθογώνιο $R = [a^1, b^1] \times \dots \times [a^k, b^k]$ στον $\mathbb{R}^k$, επιλεγμένο έτσι ώστε ο φορέας της $f$ να περιέχεται στο εσωτερικό του $R$, εκτός ενδεχομένως από τις έδρες όπου η $f$ μπορεί να μη μηδενίζεται. Το σύνορο $\partial R$ αποτελείται από $2k$ έδρες· κάθε έδρα είναι ένα $(k-1)$-διάστατο ορθογώνιο όπου μία συντεταγμένη είναι σταθερή (ίση με $a^j$ ή $b^j$) και οι υπόλοιπες μεταβάλλονται.

Το ολοκλήρωμα της $d\omega$ πάνω στο $R$ είναι
\[
\int_R d\omega = \int_R (-1)^{j-1} \ \frac{\partial f}{\partial x^j} \ dx^1 \dots dx^k.
\]
Από το θεώρημα Fubini, ολοκληρώνουμε πρώτα ως προς $x^j$. Έστω $R'$ το ορθογώνιο των υπολοίπων $k-1$ συντεταγμένων. Τότε
\[
\int_R d\omega = (-1)^{j-1} \int_{R'} \left( \int_{a^j}^{b^j} \frac{\partial f}{\partial x^j} \ dx^j \right) dx^1 \dots \widehat{dx^j} \dots dx^k.
\]
Το εσωτερικό ολοκλήρωμα υπολογίζεται από το θεμελιώδες θεώρημα του απειροστικού λογισμού:
\[
\int_{a^j}^{b^j} \frac{\partial f}{\partial x^j} \ dx^j = f(x^1, \dots, b^j, \dots, x^k) - f(x^1, \dots, a^j, \dots, x^k).
\]
Συνεπώς,
\[
\begin{aligned}
\int_R d\omega &= (-1)^{j-1} \int_{R'} f(x^1, \dots, b^j, \dots, x^k) \ dx^1 \dots \widehat{dx^j} \dots dx^k \\
&\quad - (-1)^{j-1} \int_{R'} f(x^1, \dots, a^j, \dots, x^k) \ dx^1 \dots \widehat{dx^j} \dots dx^k.
\end{aligned}
\]

Από την άλλη, το ολοκλήρωμα της $\omega$ πάνω στο $\partial R$ είναι το άθροισμα των ολοκληρωμάτων πάνω στις $2k$ έδρες. Για μια έδρα όπου $x^j = \text{σταθερά}$, ο περιορισμός της $\omega$ δίνεται από την ίδια έκφραση, αλλά το $dx^j$ μηδενίζεται (αφού η $x^j$ είναι σταθερή στην έδρα). Οι υπόλοιπες έδρες, όπου κάποια άλλη συντεταγμένη $x^i$ ($i \neq j$) είναι σταθερή, δεν συνεισφέρουν, διότι η $\omega$ περιέχει τον όρο $dx^i$ ο οποίος μηδενίζεται πάνω στην αντίστοιχη έδρα. Έτσι, μόνο οι δύο έδρες με $x^j = b^j$ και $x^j = a^j$ συνεισφέρουν.

Ο προσανατολισμός του $\partial R$ είναι ο επαγόμενος. Για την έδρα $x^j = b^j$, το εξωτερικά κατευθυνόμενο κάθετο διάνυσμα είναι το $\partial / \partial x^j$, και ο προσανατολισμός της έδρας είναι $(-1)^{j-1}$ φορές ο φυσικός προσανατολισμός των $x^1, \dots, \widehat{x^j}, \dots, x^k$. Για την έδρα $x^j = a^j$, το εξωτερικά κατευθυνόμενο κάθετο διάνυσμα είναι το $-\partial / \partial x^j$, και ο προσανατολισμός είναι ο αντίθετος. Συνεπώς,
\[
\begin{aligned}
\int_{\partial R} \omega &= (-1)^{j-1} \int_{R'} f(x^1, \dots, b^j, \dots, x^k) \ dx^1 \dots \widehat{dx^j} \dots dx^k \\
&\quad - (-1)^{j-1} \int_{R'} f(x^1, \dots, a^j, \dots, x^k) \ dx^1 \dots \widehat{dx^j} \dots dx^k.
\end{aligned}
\]
Αυτό είναι ακριβώς το ίδιο με το $\int_R d\omega$ που υπολογίσαμε παραπάνω. Άρα
\[
\int_R d\omega = \int_{\partial R} \omega,
\]
που είναι η (170) για την περίπτωση του ορθογωνίου.

\textbf{Βήμα 3: Μεταφορά στη γενική περίπτωση.} Για μια γενική $\Sigma$, κάθε όρος $\rho_\alpha \omega$ έχει φορέα σε έναν χάρτη, και μέσω της παραμετρικοποίησης $x_\alpha$ ανάγεται στην περίπτωση του $\mathbb{R}^k$ που μόλις αποδείξαμε. Το ολοκλήρωμα είναι ανεξάρτητο από την παραμετρικοποίηση (ενότητα 23.3), και η διαμέριση της μονάδας εξασφαλίζει ότι το άθροισμα των τοπικών συνεισφορών δίνει το ολικό ολοκλήρωμα. Αυτό ολοκληρώνει την απόδειξη. $\square$

\subsection{Άμεσα πορίσματα}

\subsubsection{Πόρισμα 1: Ολοκλήρωμα ακριβούς μορφής πάνω σε κλειστή πολλαπλότητα}

Αν η $\Sigma$ είναι μια κλειστή πολλαπλότητα (συμπαγής χωρίς σύνορο, $\partial \Sigma = \emptyset$), τότε για οποιαδήποτε ακριβή $k$-μορφή $\omega = d\eta$,
\begin{equation}
\int_\Sigma \omega = \int_\Sigma d\eta = \int_{\partial \Sigma} \eta = 0. \qquad (171)
\end{equation}
Ισοδύναμα, αν μια κλειστή μορφή ($d\omega = 0$) έχει μη μηδενικό ολοκλήρωμα πάνω σε μια κλειστή πολλαπλότητα, δεν μπορεί να είναι ακριβής.

Γεωμετρικά, το ολοκλήρωμα μιας κλειστής μορφής πάνω σε έναν κλειστό $k$-κύκλο είναι ένα τοπολογικό εμπόδιο για την ακρίβεια: μετρά την "καθολική μη-ακρίβεια" της μορφής. Το κλασικό παράδειγμα είναι η 1-μορφή $\omega = (x \, dy - y \, dx)/(x^2 + y^2)$ στον $\mathbb{R}^2 \setminus \{0\}$, η οποία είναι κλειστή αλλά το ολοκλήρωμά της πάνω στον κύκλο $S^1$ είναι $2\pi \neq 0$, άρα δεν είναι ακριβής.

Φυσικά, αυτό εκφράζει την αρχή ότι το έργο μιας συντηρητικής δύναμης κατά μήκος κλειστής διαδρομής είναι μηδέν. Αν $\omega$ είναι μια 1-μορφή δύναμης, τότε $\omega = df$ (ύπαρξη δυναμικού) συνεπάγεται $\oint \omega = 0$ για κάθε κλειστή καμπύλη. Το αντίστροφο ισχύει τοπικά (Poincaré) αλλά όχι καθολικά αν ο χώρος έχει "τρύπες".

\subsubsection{Πόρισμα 2: Το λήμμα του Poincaré μέσω του Stokes}

Σε μια αστεροειδή περιοχή $U$, κάθε κλειστή $k$-διάστατη επιφάνεια $C$ είναι το σύνορο μιας $(k+1)$-διάστατης περιοχής $\Sigma \subseteq U$. Αυτό οφείλεται στη δυνατότητα κατασκευής ενός "κώνου" με κορυφή το κέντρο της αστεροειδούς περιοχής και βάση την $C$. Αν $\omega$ είναι μια κλειστή $k$-μορφή ($d\omega = 0$), τότε για κάθε τέτοια κλειστή επιφάνεια $C = \partial \Sigma$,
\begin{equation}
\int_C \omega = \int_{\partial \Sigma} \omega = \int_\Sigma d\omega = 0. \qquad (172)
\end{equation}
Το γεγονός ότι το ολοκλήρωμα μιας κλειστής μορφής πάνω σε οποιαδήποτε κλειστή επιφάνεια μηδενίζεται είναι η συνθήκη που εξασφαλίζει την καθολική ακρίβεια στην αστεροειδή περιοχή. Ο τελεστής ομοτοπίας $H$ της ενότητας 20.6 είναι ακριβώς η ολοκλήρωση κατά μήκος των ακτινικών γραμμών αυτού του κώνου.

Γεωμετρικά, το πόρισμα δείχνει ότι σε μια αστεροειδή περιοχή, η τοπική ακρίβεια (Poincaré) και η καθολική ακρίβεια (Stokes) ταυτίζονται: η απουσία "τρυπών" εξασφαλίζει ότι κάθε κλειστή επιφάνεια είναι σύνορο, και άρα κάθε κλειστή μορφή είναι ακριβής.

Φυσικά, σε μια περιοχή χωρίς τρύπες, κάθε διανυσματικό πεδίο με μηδενική στροβιλότητα είναι συντηρητικό. Το θεώρημα του Stokes εγγυάται ότι η συνθήκη $d\omega = 0$ είναι επαρκής για την ύπαρξη βαθμωτού δυναμικού. Σε μια περιοχή με τρύπες (π.χ., $\mathbb{R}^2 \setminus \{0\}$), αυτό δεν ισχύει, και η φυσική εκδήλωση είναι φαινόμενα όπως το Aharonov-Bohm, όπου ένα μαγνητικό πεδίο περιορισμένο σε απροσπέλαστη περιοχή επηρεάζει την κβαντική φάση ενός σωματιδίου.

\section{Θεωρία του διανυσματικού πεδίου}

\subsection{Θεώρημα Reynolds}

\subsection{Απόκλιση}

\subsubsection{Θεώρημα Gauss-Ostrogradsky}

\subsection{Στροβιλισμός στον 3-διάστατο χώρο}

\subsubsection{Θεώρημα Kelvin-Stokes}

\subsection{Θεώρημα Helmholtz}

\section{Κατανομές και Φυλώσσεις}


\section{Συμπλεκτική γεωμετρία}

\end{document}